\documentclass[a4paper]{article}

% --- PACKAGES (kept as you had them) ---
\usepackage{fancyhdr}
\usepackage{lipsum}
\usepackage{graphicx}
\usepackage{tabularx}
\usepackage{booktabs}
\usepackage{longtable}
\usepackage{tabularx, longtable}
\setlength{\extrarowheight}{2pt}
\usepackage{array}
\usepackage{pdfpages}
\usepackage[table]{xcolor}
\usepackage{longtable}
\usepackage[table]{xcolor}
\usepackage{fix-cm}
\usepackage{tocloft}
\newcommand{\listofexamplesname}{\normalsize Examples of High Marking Responses}
\newlistof{examples}{exp}{\listofexamplesname}

\renewcommand{\cftexamplesfont}{\normalsize}
\renewcommand{\cftexamplespagefont}{\normalsize}
\usepackage{hyperref}
\hypersetup{
    colorlinks,
    citecolor=black,
    filecolor=black,
    linkcolor=black,
    urlcolor=black
}
\pagestyle{fancy}

% %\usepackage{draftwatermark}
% %\SetWatermarkText{DRAFT}
% %\SetWatermarkScale{1}
\usepackage{scholax}
\usepackage{attachfile2}
\usepackage{embedfile}
\usepackage{tocloft}
\usepackage{tocloft}
\setlength{\cftsecnumwidth}{3em}   % section number width
\setlength{\cftsubsecnumwidth}{3em}% subsection number width

% --- HEADER / FOOTER (no overlap; LEFT = section/subsection, RIGHT = part) ---
\setlength{\headheight}{16pt}
\fancyhf{} % clear defaults

% Marks:
% - LEFT header uses \rightmark  (section/subsection)
% - RIGHT header uses \leftmark  (part)
\renewcommand{\sectionmark}[1]{\markright{\thesection\ #1}}
\renewcommand{\subsectionmark}[1]{\markright{\thesubsection\ #1}} % comment out if you prefer section over subsection

% Define \partmark and make \part set it automatically (article doesn't do this)
\makeatletter
\providecommand{\partmark}[1]{\markboth{\partname~\thepart\ #1}{\partname~\thepart\ #1}}
\let\old@part\part
\renewcommand{\part}[1]{%
  \old@part{#1}%
  \partmark{#1}%
}
\makeatother

% Header boxes that wrap (prevent collisions)
\fancyhead[LO]{\parbox[t]{0.70\headwidth}{\raggedright\itshape\nouppercase{\rightmark}}}
\fancyhead[LE]{\parbox[t]{0.70\headwidth}{\raggedright\itshape\nouppercase{\rightmark}}}
\fancyhead[RO]{\parbox[t]{0.27\headwidth}{\raggedleft\bfseries\nouppercase{\leftmark}}}
\fancyhead[RE]{\parbox[t]{0.27\headwidth}{\raggedleft\bfseries\nouppercase{\leftmark}}}

\fancyfoot[C]{\thepage}
\renewcommand{\headrulewidth}{0.4pt}

% Optional: add a little extra space between header and text if needed
% \setlength{\headsep}{15pt}


%%%%%%%%%%%%%%%%%%%%%%%%%%%
\begin{document}


\begin{titlepage}
    \begin{center}
            
        \Huge
        \textbf{South Oakleigh College}
            
        \vspace{0.5cm}
        \LARGE
                 
        \vspace{0.5cm}
         
        
\includegraphics[width=0.75\textwidth]{SOC LOGO_VERTICAL.png}   

\vspace{1.2cm}
\centerline{\hbox to 13cm{\hrulefill}}
\vspace{0.3cm}
\centerline{\sc \Large \uppercase{\textbf{Units 3/4 English}}}
\vspace{0.3cm}
\centerline{\sc \Large \uppercase{\textbf{Exam Preparation Guide}}}
\centerline{\hbox to 13cm{\hrulefill}}

            
        
     \end{center}  
     
\begin{center}
    
\Large
    Compiled by
    \vspace{0.125cm}\\
\Large
    Mr. Morlin.
  \end{center}


\begin{center}
    \vspace{1cm}
        For the VCE English Cohort of 2026\\
\vspace{1cm}
Version 3.1
\end{center}
    
\end{titlepage}

\clearpage
\thispagestyle{empty}
\hfill
\clearpage

\part{Preface}

This booklet is the result of many thoughtful hours spent in careful collaboration by the English teaching team, created with a single guiding aim: to equip you with the tools, knowledge, and inspiration to approach your final English examination with clarity, confidence, and purpose.\\

Far more than a simple collection of notes, this guide seeks to be your companion in these final weeks, an intellectual compass to help you navigate the intricate terrain of literary analysis, argument construction, and creative exploration. Within these pages you will find a carefully curated array of resources: from exemplar essays and annotated responses, to thematic quote banks and character breakdowns; from breakdowns of past assessments to guidance drawn from VCAA feedback, every element has been selected with your success in mind.\\

This is not a booklet to be read in one sitting, nor is it intended to be passively consumed from cover to cover. Rather, it is a reference to be dipped into with purpose, a treasury of ideas and strategies to consult in moments of uncertainty or reflection. Like a well-thumbed dictionary or a trusted atlas, it offers direction precisely when you need it most.\\

The road to the English examination is often a challenging one, strewn with doubts, decisions, and the pressure of time. Yet it is also a journey rich in discovery, a final opportunity to let your voice be heard, your thinking be sharpened, and your understanding of the texts you’ve studied take flight. Whether you are refining your essay structure, revising key quotes, or contemplating the deeper implications of authorial intent, may this guide be a source of insight and encouragement.\\

And remember that your teachers are here to support you every step of the way. They know your strengths, understand your learning, and believe in your potential. Should you require further guidance or simply a reassuring word, do not hesitate to reach out.\\

On behalf of the English team, we extend to you our best wishes. May your hard work bear fruit, your confidence grow, and your skills shine clearly and proudly in the exam room.\\


 - The VCE English Team\\
 
Mr. Morlin, Ms. Levin, Ms. Hoernel and Ms. Zhiv.



\cleardoublepage

\clearpage
\thispagestyle{empty}
\hfill
\clearpage

\tableofcontents
\listoffigures
%%%%%%%%%%%%%%%%%%%%%%%%%%
\listoftables
\listofexamples

\cleardoublepage

\clearpage
\thispagestyle{empty}
\hfill
\clearpage


\part{\underline{Sunset Boulevard}}
\vspace{8cm}
\small{
\textit{'I hate that word. It's a return, a return to the millions of people who have never forgiven me for deserting the screen.'}
}

\clearpage
\clearpage
\thispagestyle{empty}
\hfill
\clearpage

\section{Practice Analytical Text Response Topics}
\addcontentsline{toc}{section}{Practice Analytical Text Response Topics}

1. 'Maybe that’s why it’s good, it’s true” \underline{Sunset Boulevard} shows that the truth, however ordinary or painful, is what matters. Do you agree?\\

2. In Billy Wilder’s \underline{Sunset Boulevard} things are never what they seem. Discuss.\\

3. Wilder’s \underline{Sunset Boulevard} opposes exploitation in all its forms. Discuss.\\

4. \underline{Sunset Boulevard} explores different kinds of madness. Discuss.\\

5. 'There’s nothing tragic about being 50, not unless you try to be 25”. \underline{Sunset Boulevard} is about the need to face reality. Do you agree?\\

6. In \underline{Sunset Boulevard} fantasy prevents the truth from being confronted. To what extent do you agree?\\

7. 'Sagittarius! I like Sagittarians - you can trust them” \underline{Sunset Boulevard} warns against trusting others. Discuss.\\

8. 'I felt trapped like the cigarette in the contraption on her finger”. \underline{Sunset Boulevard} is concerned with entrapment. Do you agree?\\

9. \underline{Sunset Boulevard} explores the intermingling of public and private realms. Discuss.\\

10. Billy Wilder’s \underline{Sunset Boulevard} explores the intersection of the written and the real. Discuss.\\

11. \underline{Sunset Boulevard} touches on the true horrors of Hollywood. Do you agree?\\

12. The structure of \underline{Sunset Boulevard} is crucial to its message. Discuss.\\

13. \underline{Sunset Boulevard} is not just a sad tale of deception, greed, and jealousy. Discuss.\\

14. 'I am big, it's the pictures that got small' \underline{Sunset Boulevard} shows that mistaken beliefs are responsible for suffering. Do you agree?\\

15. \underline{Sunset Boulevard} argues that we have little control over our lives. To what extent do you agree?\\

16. 'But before you hear it all distorted and blown out of proportion, before those Hollywood columnists get their hands on it, maybe you'd like to hear the facts, the whole truth”. In \underline{Sunset Boulevard} truth is the real victim. Discuss.\\

17. \underline{Sunset Boulevard} reveals Hollywood as a world of predation and exploitation. Do you agree?\\

18. \underline{Sunset Boulevard} explores the consequences of people being in the wrong place or time. Discuss.\\

19. Billy Wilder’s \underline{Sunset Boulevard} is a contest between selfishness and self-sacrifice. Discuss.\\

20. Norma and Joe are very similar people. To what extend Do you agree?\\

21. How does Wilder create suspense for the audience in \underline{Sunset Boulevard}?\\

22. “I remember waving to you. I was dancing on a table.” \underline{Sunset Boulevard} shows the lengths people will go to in order to keep the past alive. Discuss.\\

23. “Look sweetie, be practical.” Joe’s pursuit of fame is a symptom of the selfish society he inhabits. To what extent do you agree?\\

24. How does Wilder’s choice of a noir thriller position viewers to feel uncomfortable in \underline{Sunset Boulevard}?\\

25. “That’s the trouble with you readers, you know all the plots.” ‘In the world of \underline{Sunset Boulevard}, it is impossible to experience anything true.’ Do you agree?\\

26. ‘In \underline{Sunset Boulevard}, the magic of the movies carries also the potential for disenchantment.’ Discuss.\\

27. ‘Joe’s death is the result of his failure to see the tragic nature of his relationship with Norma.’ Do you agree?\\

28. ‘\underline{Sunset Boulevard} represents a condemnation of Hollywood’s idea of itself as the dream factory.’ To what extent do you agree?\\

29. ‘\underline{Sunset Boulevard} exposes the less than benevolent treatment of creative talent in the pursuit of profit.’ Discuss.\\

30. Wilder’s film is a scathing critique of the Hollywood film industry. To what extent do you agree?\\

31. How does Wilder use Franz Waxman’s score to intensify the events of the film?\\

32. Joe Gillis was complicit in all the events of \underline{Sunset Boulevard} and deserves no sympathy. Discuss.\\

33. Norma Desmond is a tragic victim of a cruel industry. Discuss.\\

34. How does Wilder’s use of a voiceover contribute to the menacing nature of the film?\\

35. “I am big. It’s the pictures that got small.” \underline{Sunset Boulevard} is an examination of the dangers of pride. To what extent do you agree?\\

36. Norma Desmond’s delusions are fuelled and supported by men hoping to profit from her fame and fortune. Discuss.\\

37. In Wilder’s \underline{Sunset Boulevard}, every character is selfish and opportunistic. To what extent do you agree?\\

38. How does Wilder capture the power of delusion in his film \underline{Sunset Boulevard}?\\

39. “DeMille didn’t have the heart to tell you. None of us had the heart.” How does deception contribute to the tragic events of \underline{Sunset Boulevard}?\\



40. What film techniques has Billy Wilder used to further his themes and main messages?\\

41. How does Billy Wilder use film techniques to further his main messages in Sunset Boulevard?\\

42. What role does self-interested thinking play in Sunset Boulevard?\\

43. How does Billy Wilder use film techniques to further his message that isolation can be a dangerous thing?\\

44. What narrative devices does Wilder use to further his main messages?\\

45. 'One of Wilder's significant themes is that living in one's own created world can have its advantages.' To what extent do you agree?\\

46. 'Wilder uses his film to criticise a society that is in love with fame.' Discuss.\\

47. How does Wilder use his film to criticise modern social foibles?\\

48. 'The screenwriter and director of Sunset Boulevard, Billy Wilder, leaves his audience feeling disappointed, yet hopeful.' To what extent do you agree?\\

49. How does Wilder use his film text to explore the problem of excessive self-interest?\\

50. 'In the film Sunset Boulevard, the crucial conflict is between reality and ambition.' Discuss.\\

51. How does the structure of this film influence the reader's response to Joe Gillis?\\

52. 'In Norma Desmond, Billy Wilder has created a character who doesn't need anyone?' To what extent do you agree?\\

53. To what extent does Wilder use his film text to criticise modern society?\\

54. How does Wilder use his film to explore the problems caused by excessive pride?\\

55. 'From her Sunset Boulevard mansion, Norma Desmond imposes her view of the world upon others.' To what extent do you agree?\\

56. How does Wilder construct his film to further his significant themes?\\

57. What lessons does the director of Sunset Boulevard have about the way fame should be regarded?\\

58. 'In his film, Wilder points an accusing finger at the weaknesses of human behaviour.' Do you agree?\\

59. How does Wilder's directing style influence the audience's response?\\

60. 'Joe Gillis's view of the world is changed by his interaction with Norma Desmond.' Discuss.\\

61. How does Wilder represent Norma Desmond as a way to further his main messages?\\

62. Discuss the role of morality in Sunset Boulevard.\\

63. How does Wilder represent Joe Gillis as a way to further his theme of the value of autonomy?\\

64. 'The life choices of every character in Sunset Boulevard are flawed.' Discuss.\\

65. How does Wilder utilise setting to portray the world as an imperfect realm?\\

66. 'Despite appearances, in Sunset Boulevard nothing is easy.' Discuss.\\

67. Who is the main victim in Sunset Boulevard?\\

68. How does Wilder represent minor characters as a way to further his main messages?\\

69. 'A central message of Sunset Boulevard is to do with the stories we tell.' To what extent do you agree?\\

70. 'The most significant motif of Wilder's film is to do with the importance of connection.' Discuss.\\

71. How does Wilder represent minor characters as a way to further his theme of the dangers of deceit?\\

72. To what extent does Wilder suggest that there are benefits in believing in yourself?\\

73. How does Wilder attempt to convince his audience of the evils of self-adulation?\\

74. In Wilder's film, Sunset Boulevard, what is considered to be essential?\\

75. How does Wilder use symbolism to further his main messages?\\

76. 'Sunset Boulevard shows the importance of trying to see the world more clearly, even though this can come at great personal cost.' Discuss.\\

77. 'The greatest crime depicted in this film is to not pay enough attention.' Discuss.\\


\subsection{Topics from Previous Exams (VCAA and others)}
1. ‘In Sunset Boulevard it is the capacity for self-delusion, rather than fame and wealth, that is shown to be destructive.’\\ Discuss.\\

2. ‘All the characters in Sunset Boulevard allow themselves to be entrapped.’\\ To what extent do you agree?\\

3. Sunset Boulevard shows the corrupt and corrupting side of Hollywood.’\\ Do you agree?\\

4. ‘Sunset Boulevard is as funny and absurd as it is tragic.’\\ Discuss.\\

5. 'We can pity the characters in Sunset Boulevard, but we cannot like or admire them.'\\ Do you agree?\\

6. 'Joe and other characters do see a real if faded glamour in Norma Desmond.'\\ Discuss.\\

7. ‘Sunset Boulevard highlights the importance of adapting to new circumstances.’\\ Discuss.\\

8. “Great stars have great pride.”\\ To what extent does Norma cause her own downfall?\\

9. ‘Sunset Boulevard exposes Hollywood as a cruel and exploitative environment.’\\ Discuss.\\

10. “I made her a star and I cannot let her be destroyed.”‘It is Max who bears the most responsibility for Norma’s descent into madness.’\\ Do you agree?\\

11. ‘In this film, the rife presence of discontent underscores the predatory and vicious nature of Hollywood.’\\ Discuss.\\

12. "All right, Mr. DeMille, I'm ready for my close-up.”‘By the end of Sunset Boulevard, the audience is left feeling nothing but sympathy for Norma.’\\ To what extent do you agree?\\

13. How is death narratively and symbolically integral to Sunset Boulevard?\\

14. ‘In this film, the distinctions between fantasy and reality are blurred.’\\ Discuss.\\

15. ‘Max’s actions are to blame for Norma’s dangerous detachment from reality.’\\ Do you agree?\\

16. ‘Deception affects the exterior and interior lives of the characters.’\\ Discuss.\\

17. ‘Norma’s obsession with her own celebrity status is detrimental to herself and others.’\\ Discuss.\\

18. ‘Relationships in Sunset Boulevard are not completely devoid of love.’\\ To what extent do you agree?\\\

19. ‘The film is a cautionary warning about the deception of Hollywood.’\\ Discuss.\\

20. ‘There is nothing more to Norma than her superficial appearance.’\\ To what extent do you agree?\\




\newpage
\section{Background and Context}
\subsection{Overview}
Understanding the historical and cinematic context of \underline{Sunset Boulevard} is essential to appreciating its form, themes, and tone. Billy Wilder’s 1950 film emerges from a period of immense transformation in Hollywood — from the decline of the silent era to the rise of the studio system and the eventual disillusionment of post-war America. The film reflects both the glamour and the decay of the entertainment industry, satirising its obsession with fame, morality, and illusion.\\

\subsection{The Silent Film Era}
From the early 1910s until the late 1920s, cinema existed largely without synchronised sound. The silent era represented the height of early film artistry, when narrative was conveyed through title cards, exaggerated gestures, and musical accompaniment. Although these films demonstrated remarkable creativity, the arrival of sound — the so-called “talkies” — revolutionised the industry.\\

Many silent stars found their careers abruptly curtailed as they struggled to adapt to dialogue-driven storytelling. This shift from visual expressiveness to vocal performance exposed the fragility of celebrity, a theme later echoed in Wilder’s portrayal of Norma Desmond. Her inability to transition from silent stardom to the spoken screen mirrors the real-life decline of many former icons.\\

\subsubsection{Analysis}
The silent era symbolises a world of illusion sustained by artifice and image — a notion central to \underline{Sunset Boulevard}. Wilder’s film draws nostalgia and melancholy from this lost cinematic world, exposing the cost of technological progress and the impermanence of fame.\\

\subsection{The Golden Age of Hollywood}
From the 1930s through to the early 1960s, Hollywood entered what became known as its “Golden Age”. During this time, the major studios — MGM, Paramount, Warner Bros., and others — dominated the global film industry. By the mid-1940s, these studios produced hundreds of films annually, shaping popular imagination through idealised visions of love, success, and moral order.\\

Films of this period were characterised by clear narratives, linear plots, and archetypal heroes and heroines. Their moral tone was often conservative, reinforcing social hierarchies and traditional gender roles. However, by the late 1950s, the system began to weaken under the pressures of television, shifting audiences, and artistic fatigue.\\

However, the flim industry in the United States was going through a rough patch:
\begin{itemize}
\item Television was emerging as the dominant entertainment medium in the U.S.
\item Weekly cinema attendance was declining.
\item The Paramount anti-trust consent decree was implemented, requiring separate ownership for production companies and theatre chains.
\item There was a population shift to the suburbs, impacting downtown movie theatres.
\item There were a set of rules for moral guidelines that filmmakers had to follow under The Hays Code. 
\end{itemize}

\subsubsection{Analysis}
Wilder’s critique of Hollywood in \underline{Sunset Boulevard} coincides with the decline of this era. The film’s depiction of faded glamour and moral corruption exposes the darker underbelly of an industry once devoted to producing dreams. Through satire, Wilder dismantles the myth of Hollywood’s innocence.\\

\subsection{The Production Code}
Between 1930 and 1968, the film industry was governed by the Production Code Administration (PCA), a body that imposed strict moral guidelines on cinematic content. Initially voluntary, these regulations became mandatory by the mid-1930s and dictated that all films must uphold public decency and moral virtue.\\

Depictions of nudity, profanity, excessive violence, or sexual behaviour were prohibited. Filmmakers were compelled to ensure that wrongdoing was punished and virtue rewarded. Wilder and his collaborator Charles Brackett often tested these boundaries, embedding subversive commentary within seemingly conventional narratives.\\

In \underline{Sunset Boulevard}, the tragic fates of Joe Gillis and Norma Desmond adhere to the code’s demand for moral consequence — the transgressors must pay for their sins. Yet Wilder’s irony turns this moralism on its head, inviting sympathy for both victim and perpetrator.\\

\subsubsection{Analysis}
The Production Code’s influence gives the film a paradoxical tone — simultaneously moralistic and cynical. While the narrative satisfies censorship standards through punishment, it also exposes the hypocrisy of Hollywood’s moral pretensions. Wilder uses the very language of the Code to critique the system that enforced it.\\

\subsection{Film Noir}
\underline{Sunset Boulevard} belongs to the film noir tradition — a style originating in 1940s post-war America and named by French critics to describe the “dark films” emerging from Hollywood. Noir was marked by chiaroscuro lighting, urban settings, complex psychology, and an atmosphere of mistrust and fatalism.\\

The genre’s visual style relied on unbalanced compositions, low-key lighting, and stark black-and-white contrasts, evoking a world of corruption and moral ambiguity. In tone, noir films explored alienation and moral decay, often narrated by disillusioned protagonists trapped by circumstance.\\

\subsubsection{Analysis}
Wilder fuses noir aesthetics with Hollywood satire, using shadow and perspective to visualise moral collapse. Joe’s cynical narration, the claustrophobic mansion, and the interplay of light and darkness all reinforce themes of delusion, decay, and entrapment. The film thus becomes both a noir thriller and a self-reflexive critique of the industry that produced it.\\

\subsection{America in the 1950s}
Post-war America entered a period of unparalleled prosperity and conservatism. The 1950s were shaped by the Cold War, the rise of consumer culture, and rigid gender expectations. Prosperity bred conformity, and political paranoia — intensified by McCarthyism and the House Un-American Activities Committee — fostered fear of dissent.\\

Socially, women were expected to embody domestic virtue while men resumed positions of economic dominance. Yet beneath this façade of order lay widespread disillusionment and anxiety about identity and purpose. \underline{Sunset Boulevard} channels these tensions through its portrayal of isolation, ambition, and moral corruption within an ostensibly glamorous world.\\

\subsubsection{Analysis}
The film captures the contradictions of 1950s America: wealth and decay, progress and repression, fame and emptiness. Norma’s madness and Joe’s moral exhaustion personify the cultural malaise of a nation seduced by success yet haunted by its own discontents. Wilder’s critique of Hollywood thus becomes a mirror for post-war America’s soul.\\
\newpage
\section{Values in \underline{Sunset Boulevard}}

\subsection{Social Values}
\subsubsection{Traditional Gender Roles}
The era upheld traditional gender roles, with men as breadwinners and women primarily responsible for homemaking and child-raising. The suburban housewife became an iconic image of the time. Norma’s character highlights the fading role of ageing actresses in Hollywood, mirroring the limited opportunities for women beyond traditional roles as wives and mothers.\\

\subsubsection{Consumerism}
The 1950s were characterised by a post-war economic boom, and there was a strong emphasis on consumerism. Owning a home, car, and various consumer goods became aspirational goals. The opulent mansion on Sunset Boulevard and Norma’s extravagant lifestyle represent a distorted version of the consumerist dream, where wealth and possessions mask deeper issues.\\

\subsection{Cultural Values}
\subsubsection{Hollywood's Golden Age}
Hollywood’s Golden Age is a term used to describe a period in the history of the American film industry that roughly spans from the late 1920s to the early 1960s. It was a time of unprecedented creativity, innovation, and commercial success in Hollywood.
The Studio System – a way of making movies where big companies, like Warner Bros. and Paramount, had total control over everything. They decided which actors to use, what movies to make, and even how the movies were shown in theatres.\\

\subsubsection{Technological Advancements}
This era witnessed significant advancements in filmmaking technology, including the transition from silent films to ‘talkies’, the development of colour cinematography, and the use of widescreen formats.\\

\subsubsection{The Star System}
The method of creating, promoting and exploiting stars in Hollywood films. Movie studios would select promising young actors and glamourise them through various promotion tactics.\\






\newpage
\section{Thematic Quote Bank}

\subsection{Fame and Obsession}
\begin{itemize}
  \item Norma Desmond: 'I am big. It's the pictures that got small.'
  \item Norma Desmond: 'They still write me fan letters every day. Why do they beg me for my photographs? Because they want to see me, me, me! Norma Desmond.'
  \item Gillis' Voice: 'She was still sleepwalking along the giddy heights of a lost career — plain crazy when it came to that one subject: her celluloid self, the great Norma Desmond.'
  \item Norma Desmond: 'Salome — what a woman! What a part! The Princess in love with a Holy man. She dances the Dance of the Seven Veils. He rejects her, so she demands his head on a golden tray, kissing his cold, dead lips.'
  \item Gillis' Voice: 'I could sense her eyes on me from behind those dark glasses, defying me not to like what I read, or maybe begging me in her own proud way to like it. It meant so much to her.'
  \item Norma Desmond: 'I wrote that with my heart.'
  \item Norma Desmond: 'I hate that word. It is a return. A return to the millions of people who have never forgiven me for deserting the screen.'\\
\end{itemize}

\subsection{Isolation and Decay}
\begin{itemize}
  \item Gillis' Voice: 'A neglected house gets an unhappy look. This one had it in spades. It was like that old woman in Great Expectations — Miss Havisham in her rotting wedding dress.'
  \item Gillis' Voice: 'There was a tennis court, or rather the ghost of a tennis court, with faded markings and a sagging net.'
  \item Gillis' Voice: 'There was something else going on below: the last rites for that hairy old chimp, performed with the utmost seriousness — as if she were laying to rest an only child.'
  \item Gillis' Voice: 'The whole place seemed to have been stricken with a kind of creeping paralysis, out of beat with the rest of the world, crumbling apart in slow motion.'
  \item Gillis' Voice: 'And of course she had a pool. Who didn't then? It was empty now... or was it?'
  \item Gillis' Voice: 'That night I'd had a mixed-up dream. In it was an organ grinder. I couldn't see his face, but the organ was all draped in black, and a chimp was dancing for pennies.'
  \item Gillis' Voice: 'She'd sit very close to me, and she'd smell of tuberoses, which is not my favorite perfume, not by a long shot.'\\
\end{itemize}

\subsection{Reality vs. Illusion}
\begin{itemize}
  \item Norma Desmond: 'I hate that word. It is a return. A return to the millions of people who have never forgiven me for deserting the screen.'
  \item Norma Desmond: 'We didn't need dialogue. We had faces. There just aren't any faces like that any more.'
  \item Gillis' Voice: 'You don't yell at a sleepwalker — he may fall and break his neck. That's it — she was still sleepwalking along the giddy heights of a lost career.'
  \item Gillis' Voice: 'She sure could say a lot of things with those pale eyes of hers. They’d been her trademark. They’d made her the Number One Vamp of another era.'
  \item Norma Desmond: 'Still wonderful, isn't it? And no dialogue. We had faces.'
  \item Norma Desmond: 'The plain fact was that she was afraid of that world outside. Afraid it would remind her that time had passed.'
  \item Gillis' Voice: 'She sat coiled up like a watch spring, her cigarette clamped in a curious holder... begging me in her own proud way to like it.'\\
\end{itemize}

\subsection{Power Dynamics}
\begin{itemize}
  \item Norma Desmond: 'You proud boy, why didn't you tell me you were having difficulties. It's all taken care of. It's all paid for.'
  \item Max: 'I made your bed this afternoon.' (Indicating that Norma has already decided Joe will stay)
  \item Norma Desmond: 'Suppose you make up your mind. Do you want this job or don't you?'
  \item Gillis: 'I let him unpack my things. I wanted the dough and I wanted to get out of there as quickly as possible.'
  \item Gillis: 'I’m used to paying my own bills.'
  \item Norma Desmond: 'Now, now, don't let's be small about such matters. We won't keep books.'
  \item Gillis' Voice: 'So I let him unpack my things. I wanted the dough, and I wanted to get out of there as quickly as possible.'\\
\end{itemize}

\subsection{The Dark Side of Hollywood}
\begin{itemize}
  \item Gillis: 'Well, I needed about two hundred and ninety dollars real quick, or I'd lose my car... I was way ahead of the finance company.'
  \item Norma Desmond: 'Look at them in the front offices — the masterminds! They took the idols and smashed them.'
  \item Betty Schaefer: 'It's from hunger... Just a rehash of something that wasn't very good to begin with.'
  \item Gillis: 'Just sold an original for a hundred grand. The Life of the Warner Brothers. Starring the Ritz Brothers.' (sarcastically to Skolsky, showing how phony Hollywood is)
  \item Sheldrake: 'Who wants to see a Civil War picture?' (regarding Gone with the Wind)
  \item Gillis' Voice: 'Actors and stock girls and waiters. Kind of a combination office, Kaffee-Klatsch and waiting room. Waiting, waiting for the gravy train.' (about Schwab’s Drugstore)
  \item Gillis' Voice: 'The whole place seemed to have been stricken with a kind of creeping paralysis, out of beat with the rest of the world, crumbling apart in slow motion.'\\
\end{itemize}

\subsection{Gender Roles}
\begin{itemize}
\item \textbf{JOE} A woman almost twice my age ...
\item \textbf{JOE} A very simple set-up: An older woman who is well-to-do. A younger man who is not doing too well ... Can you figure it out yourself?
\item \textbf{JOE} Norma, grow up. You're a woman of fifty. There's nothing tragic about being fifty — not unless you try to be twenty-five.
\item \textbf{GILLIS} I don't need any clothes, and I certainly don't want you buying them for --
\item \textbf{NORMA} I'll buy you a boat and we'll sail to Hawaii.
\item \textbf{NORMA} I own three blocks downtown. I have oil in Bakersfield — pumping, pumping, pumping. What's it for but to buy us anything we want.
\item \textbf{NORMA} Just follow me.
\item \textbf{GILLIS} Pull up at the drugstore, will you, Max. (To Norma) I'll get you some.
\end{itemize}


\subsection{Love, Infatuation and Narcissism}
\begin{itemize}
\item \textbf{NORMA} Salome -- what a woman! What a part! The Princess in love with a Holy man. She dances the Dance.
\item \textbf{GILLIS} They'll love it in Pomona.
\item \textbf{NORMA} They will love it every place. Read it. Read the scene just before she has him killed!
\item \textbf{NORMA} Perfect. Wonderful shoulders. And I love that line.
\item \textbf{NORMA} Is it funny that I'm in love with you?
\item \textbf{NORMA} I'm in love with you. Don't you know that? I've been in love with you all along.
\item \textbf{NORMA} What you're trying to say is that you don't want me to love you. Is that it?
\item \textbf{ARTIE} Wait a minute. This is the woman I love. What's going on? Who was loaded?
\item \textbf{NORMA} To fall in love with you -- that was the idiotic thing.
\item \textbf{BETTY} Of course I love him. I always will. I'm just not in love with him any more.
\item \textbf{JOE} And she was in love with me. Me! She was a fool not to sense that there was something phoney in my set-up.
\item \textbf{NORMA} (In a whimpering monotone) I love you, Joe. I love you, Joe. I love you, Joe. I love you, Joe.
\item \textbf{BETTY} If you love me, Joe.
\item \textbf{NORMA} You're here, Joe ... When did you come home? Where were you? Is it a woman? I know it's a woman ... Who is she? Oh Joe, why can't I ask you? I must know, I must!
\end{itemize}

\newpage

\section{Character Quote Bank}


\subsection{Joe Gillis}
\subsubsection{Fame}
\begin{itemize}
\item That was last year. This year I'm trying to earn a living.
\item I know your face. You're Norma Desmond. You used to be in pictures. You used to be big.
\item I don't think so. It's your script. It's your show. Good luck.
\item If you get a hundred thousand for it, you buy me a box of chocolate creams. If you get an Oscar, I get the left foot.
\item He must be quite a guy.
\item Come on where? Back to a one-room apartment that I can't pay for? Back to a story that may sell and very possibly will not?
\item And I'd take it in a second, Norma — only it's a little too dressy for sitting behind the copy desk in Dayton, Ohio.
\item Norma, you'd be throwing it away. I don't qualify for the job, not any more.
\item Norma, grow up. You're a woman of fifty. There's nothing tragic about being fifty — not unless you try to be twenty-five.
\end{itemize}

\subsubsection{Reality and Illusion}

\begin{itemize}
\item Oh, you're one of the message kids. Just a story won't do. You'd have turned down \textit{Gone With the Wind}.
\item Forget Romanoff's. It's the car I'm talking about. If I lose my car it's like having my legs cut off.
\item You know, I'm pretty expensive. I get five hundred a week.
\item Why? Her career? She got enough out of it. She's not forgotten. She still gets those fan letters.
\item Yes, I was playing hooky every evening along in there.
\item Oh, that. It's from a friend of mine. A middle-aged lady, very foolish and very generous.
\item Sure you are. And we're not helping her any, feeding her lies and more lies. Getting herself ready for a picture ... What happens when she finds out?
\item That's between you and yourself, Norma.
\item Wake up, Norma. You'd be killing yourself to an empty house. The audience left twenty years ago. Now face it.
\item De Mille didn't have the heart to tell you. None of us has had the heart.
\item Madame is the greatest star of them all... I will take Mr. Gillis' bags.
\end{itemize}



\subsubsection{Hollywood Power and Transitions}
\begin{itemize}
\item 'I knew a big shot over there who’d always liked me”
\item 'They're pretty hot about it over at Twentieth, but I think Zanuck's all wet. Can you see Ty Power as a shortstop?”
\item 'Any kind of assignment. Additional dialogue.”
\item 'Sure you did. That's what makes it great.”
\item 'What it needs is a little more dialogue.”
\item 'Honestly, it's a little old hat. They don't want that any more.”
\item 'Maybe deMille is shooting.”
\item 'Psychopaths sell like hotcakes.”
\item 'It's that old familiar story, you help a timid little soul across a crowded street. She turns out to be a multimillionaire and leaves you all her money.”
\item 'Ever been in one of these old Hollywood palazzos? That's from when they were making eighteen thousand a week, and no taxes.”
\item 'Oh, sure — if it would make a good scene.”
\item 'He was trying to spare your feelings. The studio wanted to rent your car.”
\end{itemize}


\subsubsection{Gender Roles}

\begin{itemize}
\item 'And she’d mentioned something to drink”
\item 'It's all taken care of. It's all paid for.”
\item 'The ashtray, Joe dear! Can we have the ashtray?”
\item 'I don't need any clothes, and I certainly don't want you buying them for —”
\item 'A camel's hair will do.”
\item 'Stop it. You aren't going to buy me anything more.”
\item 'What right do you have to take me for granted?”
\item 'I'm all wrong for you. You want a Valentino — somebody with polo ponies — a big shot –”
\item 'Pull up at the drugstore, will you, Max. (To Norma) I'll get you some.”
\item 'All my things? All the eighteen suits, all the custom-made shoes and the eighteen dozen shirts, and the cuff-links and the platinum key-chains, and the cigarette cases”
\end{itemize}

\subsubsection{Love, Infatuation and Narcissism}

\begin{itemize}
\item 'Sagittarius. I like Sagittarians. You can trust them.”
\item 'I wouldn't worry about money. I'll make it worth your while.”
\item 'Her eyes are like two moonlit waterholes, where strange animals come to drink.”
\item 'I didn’t argue with her. You don't yell at a sleepwalker.”
\item 'She’d clutch my arm or my hand.”
\item 'Is it funny that I'm in love with you?”
\item 'It sure would have made attractive headlines: Great Star Kills Herself for Unknown Writer.”
\item 'All right, I drove to the beach. Norma, you don't want me to feel I'm locked up in this house?”
\item 'I haven't done anything, Norma.”
\item 'Constantly. No, in all seriousness, it's really good. It's fun writing again. I'm happy here, honest I am.”
\item 'Three cheers for Betty Schaefer! I will now kiss that nose of yours.”
\item 'Of course I love him. I always will. I'm just not in love with him any more.”
\item 'All right, baby. This way out.”
\end{itemize}


\subsection{Norma Desmond}
\subsubsection{Fame/Opulence}
\begin{itemize}
\item 'I don't care anyway. I want the coffin to be white. And I want it specially lined with satin. White, or deep pink.”
\item 'How much will it be? I warn you — don't give me a fancy price just because I'm rich.”
\item 'Now, now, don't let's be small about such matters. We won't keep books.”
\item \textbf{MAX:} 'She was the greatest. You wouldn't know. You are too young. In one week she got seventeen thousand fan letters.”
\item 'An Isotta-Fraschini. Have you ever heard of Isotta-Fraschinis? All hand-made. Cost me twenty-eight thousand dollars.”
\item 'I'll open my house in Malibu, and you can have the whole ocean. Or I'll buy you a boat and we'll sail to Hawaii.”
\item 'I own three blocks downtown. I have oil in Bakersfield — pumping, pumping, pumping. What's it for but to buy us anything we want.”
\item \textbf{JONESY:} 'Miss Desmond can. Come on.”
\item 'Last time I saw you was someplace very gay. I remember waving to you. I was dancing on a table.”
\item 'But remember, darling — I don't work before ten in the morning, and never after 4:30 in the afternoon.”
\item 'I'm not worried. Everything will be fine. The old team together. Nothing can stop us.”
\end{itemize}



\subsubsection{Reality and Illusion}
\begin{itemize}
\item 'I hate that word. It is a return. A return to the millions of people who have never forgiven me for deserting the screen.”
\item 'They will love it every place.”
\item 'What for? I can say anything I want with my eyes.”
\item \textbf{JOE:} 'Say, she's quite a character, that Norma Desmond.”
\item 'Because they want to see me, me, me! Norma Desmond.”
\item \textbf{JOE:} 'More Norma Desmond and still more Norma Desmond.”
\item 'I'll show them. I'll be up there again. So help me!”
\item 'I'll do it again, I'll do it again, I'll do it again!”
\item 'The chart from my astrologer. She read deMille's horoscope. She read mine.”
\item 'I never looked better in my life. Do you know why? Because I've never been as happy in my life.”
\item 'I know that trick! He wants to belittle me. He's trying to get my price down. I've waited twenty years for this call. Now Mr. deMille can wait till I'm good and ready.”
\item 'Did you see them? Did you see how they came?”
\item 'I don't care about the money. I just want to work again. You don't know what it means to know that you want me.”
\item 'You don't care, do you? But hundreds of thousands of people will care!”
\item 'That's a lie! They still want me! I'm the greatest star of them all.”
\item 'No one leaves a star. That makes one a star. Stars are ageless, aren't they?”
\item 'Oh, yes, yes. They're below, waiting for the Princess ... I'm ready.”
\item 'All right, Mr. deMille, I'm ready for my close-up.”
\end{itemize}


\subsubsection{Hollywood Power and Transition}
\begin{itemize}
\item I am big. It's the pictures that got small.
\item They're dead. They're finished. There was a time when this business had the eyes of the whole wide world. But that wasn't good enough. Oh, nol They wanted the ears of the world, too. So they opened their big mouths, and out came talk, talk, talk...
\item 'They took the idols and smashed them.”
\item 'You've made a rope of words and strangled this business.”
\item 'Still wonderful, isn't it? And no dialogue. We didn't need dialogue. We had faces.”
\item 'I'm not just selling the script. I'm selling me.”
\item 'deMille always said I was his greatest star.”
\item 'Now, now do you believe me? I told you deMille would jump at it.”
\item 'Say that I'm busy, and hang up.”
\item \textbf{deMille:} 'Listen, thirty million fans have given her the brush. Isn't that enough?”
\item \textbf{deMille:} 'You didn't know Norma Desmond as a plucky little girl of seventeen, with more courage and wit and heart than ever came together in one youngster.”
\item \textbf{deMille:} 'Only toward the end. A dozen press agents working overtime can do terrible things to the human spirit.”
\item 'We'll be working again, won't we, Chief? We'll make our greatest picture.”
\item 'I've lost half a pound since Tuesday.”
\item 'I can't go on with the scene. I'm too happy.”
\item 'Do you mind, Mr. deMille, if I say a few words? Thank you. I just want to tell you how happy I am to be back in the studio making a picture again.”
\end{itemize}


\subsubsection{Gender Roles}
\begin{itemize}
\item 'I wouldn't worry about money. I'll make it worth your while.”
\item 'It's all taken care of. It's all paid for.”
\item 'You proud boy, why didn't you tell me you were having difficulties.”
\item 'The ashtray, Joe dear! Can we have the ashtray?”
\item 'I don't need any clothes, and I certainly don't want you buying them for —”
\item 'As long as the lady is paying for it, why not take the Vicuna?”
\item 'She takes a dollar bill from her purse and gives it to him.”
\item 'What is it you want — money?”
\end{itemize}


\subsubsection{Love, Infatuation and Narcissism}
\begin{itemize}
\item 'Who? I'd have to have somebody I can trust. When were you born? I mean, what sign of the zodiac?”
\item 'Are you married, Mr. -- ?”
\item 'Perfect. Wonderful shoulders. And I love that line.”
\item 'Here's to us.”
\item 'There are no other guests. We don't want to share this night with other people. This is for you and me.”
\item 'Hold me tighter.”
\item 'Is it funny that I'm in love with you?”
\item 'I'm in love with you. Don't you know that? I've been in love with you all along.”
\item 'Mad about the boy.”
\item 'What you're trying to say is that you don't want me to love you. Is that it?”
\item 'To fall in love with you — that was the idiotic thing.”
\item 'Go away. Go to that girl of yours.”
\item 'Happy New Year, darling.”
\item 'I just came to say good night. I don't want you to see me — I'm not very attractive.”
\item 'It's just that I don't want to be left alone.”
\item 'Not now, while I'm under this terrible strain. My nerves are being torn apart. All I ask is for you to be a little patient and a little kind.”
\item 'Is it a woman? I know it's a woman...”
\item 'Don't hate me, Joe. I did it because I need you.”
\item 'I love you, Joe. I love you, Joe.”
\item 'May I come in? I've stopped crying. I'm all right again.”
\item 'What are you doing, Joe? What are you doing? You're not leaving me?”
\item 'You can't do this! Max! Max! I can't face life without you, and I'm not afraid to die, you know.”
\item 'You're not leaving me!”
\end{itemize}



\subsection{Max Von Mayerling}
\begin{itemize}
\item Madame is waiting.
\item Max was the only other person in that grim Sunset castle.
\item She was the greatest. You wouldn't know. You are too young. In one week she got seventeen thousand fan letters.
\item I wouldn't look too closely at the Postmarks.
\item Be careful. Do not race up the stairs. The musicians must not know what has happened.
\item Nobody here can give you any information. You will please not call again.
\item You must be very careful as you cross the patio. Madame may be watching.
\item She never will. That is my job. It has been for a long time. You must understand I discovered her when she was eighteen. I made her a star. I cannot let her be destroyed.
\item I directed all her early pictures. There were three young directors who showed promise in those days: D.W. Griffith, C.B. deMille, and Max von Mayerling.
\item It was I who asked to come back, humiliating as it may seem. I could have gone on with my career, only I found everything unendurable after she divorced me. You see, I was her first husband.
\item Madame is the greatest star of them all... I will take Mr. Gillis' bags.
\item The cameras have arrived, Madame.
\item Please ...
\item Is everything set up, gentlemen? Are the lights ready?
\item Quiet, everybody! Lights! Are you ready, Norma?
\item This is the staircase of the palace.
\item All right. (To cameramen) Camera!
\item (To Norma) Action!
\end{itemize}

\subsection{Betty Schaefer}
\begin{itemize}
\item 'It’s just a rehash of something that wasn’t very good to begin with”
\item 'I wish I could crawl in a hole and pull it in after me”
\item 'I just think a picture should say a little something”
\item 'One of those message kids”
\item 'I’d always heard that you'd had some talent”
\item 'There's one called windows”
\item 'You’ve got a flashback there”
\item 'Maybe that's why it's good, it's true”
\item 'With a wildly beating heart”
\item No, Phillip, no. We must be strong. You're still wearing the uniform of the Coldstream Guards!
\item Furthermore, you can have the phone now.
\item Is this Crestview 5-1733? ... I'm sorry to bother you again, but I've confirmed the number. I must speak to Mr. Gillis.
\item Sheldrake likes that angle about the teacher.
\item It's not your career -- it's mine. I kind of hoped to get in on this deal. I don't want to be a reader all my life. I want to write.
\item Couldn't we work in the evenings?
\item Six o'clock in the morning? This next month I'm completely at your disposal. Artie is out of town.
\item I was born just two blocks from this studio. Right on Lemon Grove Avenue. Father was head electrician here till he died. Mother still works in Wardrobe.
\item Not once. What's wrong with being on the other side of the cameras? It's really more fun.
\item Of course I love him. I always will. I'm just not in love with him any more.
\item If you love me, Joe.
\item Joe, I can't look at you any more.
\item Good luck to you, Betty. You can finish that story on the way to Arizona. When you and Artie get back, if the two of you ever feel like a swim, here's the pool ...
\end{itemize}


\subsection{Cecil B. deMille}
\begin{itemize}
\item 'Norma Desmond”
\item 'I hate to think where that puts me”
\item 'What can I tell her, what can I say”
\item '30 million fans have given her the brush, isn’t that enough”
\item 'You didn’t know Norma Desmond as a lovely little girl of 17, more courage and wit and heart than ever came together in a youngster.”
\item 'A dozen press agents working overtime can do terrible things to the human spirit”
\item 'Hello young fellow”
\item 'It’s good to see you”
\item 'A lot of people were, Lindbergh had just landed in Paris”
\item 'Norma I must apologise for not calling you”
\item 'I’m terribly busy”
\item 'I’m in the middle of a rehearsal”
\item 'Sit up there in my chair and make yourself comfortable”
\item 'Hog-guy, turn that light back where it belongs”
\item 'You know, some crazy things happen in this business Norma”
\item 'We've missed you too dear”
\item 'It’d be a very expensive picture”
\item 'Nothing would please me more Norma, if it were possible”
\item 'Pictures have changed quite a bit”
\item 'Goodbye Norma, we’ll see what we can do”
\item 'The old team, yeah”
\item 'Goodbye dear”
\item 'Tell him to forget about our car ... I'll buy him five old cars if necessary”
\end{itemize}

\subsection{Sheldrake}
\begin{itemize}
\item 'I knew a bigshot over there who always liked me.”
\item 'He was a smart producer”
\item 'Alright Gillis, you’ve got five minutes, what's your story about?”
\item 'No, that was me, I said who wants to see a civil war picture”
\item 'You sound like a bunch of New York critics”
\item 'If we made it a girls' softball team, put in a few numbers, that’d make a cute musical – it happened in the bullpen, the story of a woman”
\item 'This year I had to mortgage the ranch, so that I could keep my life insurance, so I could borrow on my insurance”
\end{itemize}

\subsection{Artie Green}
\begin{itemize}
\item 'An awful nice guy, an assistant director, he could let me have 20, but twenty wouldn't do”
\item 'Well what do you know, Joe Gillis”
\item 'Where you been keeping that gorgeous face of yours”
\item 'Come on, give me your coat”
\item 'Who did you borrow that from”
\item 'It just so happens we have a vacancy on the couch”
\item 'Wait a minute, this is the woman I love, who was loaded”
\item 'Hey Joe, I said you could have my couch, I didn’t say you could have my girl”
\end{itemize}
\newpage


\section{Character Relationships}

\subsection{Joe Gillis and Norma Desmond}
Joe and Norma's relationship is built on dependence, illusion, and emotional imbalance. Norma clings to Joe as a last connection to youth and relevance, while Joe allows himself to be financially supported, though emotionally trapped. Their dynamic is parasitic and destructive, culminating in Joe’s death and Norma’s complete psychological collapse.

\subsubsection*{Key Quotes}
\begin{itemize}
    \item \textit{Joe:} 'Norma, you'd be throwing it away. I don't qualify for the job, not any more.'
    \item \textit{Norma:} 'Is it funny that I'm in love with you?'
    \item \textit{Joe:} 'Wake up, Norma. You’d be killing yourself to an empty house. The audience left twenty years ago.'
    \item \textit{Norma:} 'You're not leaving me!'
\end{itemize}

\subsubsection*{Key Ideas}
\begin{itemize}
    \item Dependency
    \item Reality vs. Illusion
    \item Power imbalance
    \item Gender roles
\end{itemize}

\subsection{Joe Gillis and Betty Schaefer}
Joe and Betty’s relationship symbolises hope, creativity, and authentic emotional connection. While Joe is initially cynical, Betty helps him reconnect with his talent and ideals. However, Joe’s entrapment in Norma’s world ultimately sabotages their bond.

\subsubsection*{Key Quotes}
\begin{itemize}
    \item \textit{Betty:} 'If you love me, Joe.'
    \item \textit{Joe:} 'Three cheers for Betty Schaefer! I will now kiss that nose of yours.'
    \item \textit{Betty:} 'Joe, I can't look at you any more.'
    \item \textit{Joe:} 'Good luck to you, Betty... When you and Artie get back... here’s the pool.'
\end{itemize}

\subsubsection*{Key Ideas}
\begin{itemize}
    \item Authenticity
    \item Missed opportunities
    \item Love vs. obligation
    \item Hope and disappointment
\end{itemize}

\subsection{Norma Desmond and Max Von Mayerling}
Max and Norma share a hauntingly devoted and co-dependent relationship. Once her director and husband, Max now lives as her butler, fuelling her illusions to protect her fragile psyche. His loyalty borders on tragic obsession.

\subsubsection*{Max as a Pygmalion Archetype}
Max can be read as a Pygmalion archetype: like the mythical sculptor who fell in love with the statue he carved, Max created the star 'Norma Desmond' \textemdash{} he discovered her at sixteen, directed her early pictures, and 'made her a star' \textemdash{} and then fell in love with his own creation. Unlike Pygmalion, however, Max cannot bring his creation to life in the present; instead he devotes himself to preserving the illusion he built, writing the fan letters that sustain it. This reading deepens the tragedy of their relationship: Norma's delusion is not merely her own weakness but a construction Max continues to maintain, making him both her most loyal protector and the quiet architect of her downfall.

\subsubsection*{Key Quotes}
\begin{itemize}
    \item \textit{Max:} 'I made her a star. I cannot let her be destroyed.'
    \item \textit{Max:} 'It was I who asked to come back... I was her first husband.'
    \item \textit{Max:} 'The cameras have arrived, Madame.'
    \item \textit{Norma:} 'All right, Mr. deMille, I'm ready for my close-up.'
\end{itemize}

\subsubsection*{Key Ideas}
\begin{itemize}
    \item Loyalty
    \item Delusion
    \item Power reversal
    \item Sacrifice
    \item The Pygmalion archetype \textemdash{} the creator in service of his creation
\end{itemize}

\subsection{Norma Desmond and Cecil B. deMille}
Norma views deMille as the key to her return to fame, though he treats her with nostalgia and pity. Their relationship highlights the **Hollywood power shift** from silent films to sound, and the industry's neglect of ageing stars.

\subsubsection*{Key Quotes}
\begin{itemize}
    \item \textit{Norma:} 'Now Mr. deMille can wait till I'm good and ready.'
    \item \textit{deMille:} 'Thirty million fans have given her the brush. Isn’t that enough?'
    \item \textit{Norma:} 'We’ll be working again, won’t we, Chief?'
    \item \textit{deMille:} 'Tell him to forget about our car. I’ll buy him five old cars if necessary.'
\end{itemize}

\subsubsection*{Key Ideas}
\begin{itemize}
    \item Fame and decay
    \item Nostalgia vs. reality
    \item Hollywood’s cruelty
    \item Gendered ageism
\end{itemize}

\subsection{Joe Gillis and Max Von Mayerling}
Though they barely trust one another, Joe and Max are both caught in Norma’s gravitational pull. Max warns Joe of Norma’s delusions but ultimately enables them, silently observing Joe’s downfall.

\subsubsection*{Key Quotes}
\begin{itemize}
    \item \textit{Max:} 'You must understand, I discovered her when she was eighteen.'
    \item \textit{Max:} 'Madame is the greatest star of them all... I will take Mr. Gillis’ bags.'
    \item \textit{Joe:} 'Oh, that. It’s from a friend of mine. A middle-aged lady, very foolish and very generous.'
\end{itemize}

\subsubsection*{Key Ideas}
\begin{itemize}
    \item Powerlessness
    \item Enabling behaviour
    \item Masculinity
    \item Pity and rivalry
\end{itemize}

\subsection{Joe Gillis and Artie Green}
Joe and Artie are casual friends, connected through Betty. Artie is warm, generous, and unsuspecting—offering Joe a place to stay and unknowingly introducing him to Betty, who will fall in love with Joe. Their relationship represents Hollywood's working-class camaraderie, but it is ultimately undermined by Joe’s betrayal.

\subsubsection*{Key Quotes}
\begin{itemize}
    \item \textit{Artie:} 'Hey Joe, I said you could have my couch, I didn’t say you could have my girl.'
    \item \textit{Artie:} 'Where you been keeping that gorgeous face of yours?'
    \item \textit{Artie:} 'Well what do you know, Joe Gillis.'
\end{itemize}

\subsubsection*{Key Ideas}
\begin{itemize}
    \item Friendship
    \item Jealousy and betrayal
    \item Trust
\end{itemize}

\subsection{Betty Schaefer and Artie Green}
Betty and Artie are engaged, but her growing admiration for Joe shifts the dynamic. Artie symbolises comfort and stability, while Joe represents ambition and excitement. Their relationship quietly dissolves as Betty’s emotional attention shifts.

\subsubsection*{Key Quotes}
\begin{itemize}
    \item \textit{Artie:} 'Wait a minute, this is the woman I love, who was loaded.'
    \item \textit{Betty:} 'Of course I love him. I always will. I'm just not in love with him any more.'
\end{itemize}

\subsubsection*{Key Ideas}
\begin{itemize}
    \item Love vs. stability
    \item Unspoken heartbreak
    \item Emotional conflict
\end{itemize}

\subsection{Joe Gillis and Sheldrake}
Sheldrake represents the industry gatekeepers—pragmatic and dismissive. Joe’s interactions with him underscore how impersonal and transactional Hollywood has become. Their relationship exposes the disillusionment of struggling writers.

\subsubsection*{Key Quotes}
\begin{itemize}
    \item \textit{Sheldrake:} 'Alright Gillis, you’ve got five minutes, what's your story about?'
    \item \textit{Joe:} 'You’d have turned down \textit{Gone With the Wind}.'
    \item \textit{Sheldrake:} 'Put in a few numbers, that’d make a cute musical.'
\end{itemize}

\subsubsection*{Key Ideas}
\begin{itemize}
    \item Hollywood as business
    \item Artistic compromise
    \item Creative frustration
\end{itemize}

\subsection{Norma Desmond and Betty Schaefer}
Although they never interact directly, Norma and Betty represent **opposing ideals of womanhood**. Betty is young, modern, and practical; Norma is older, delusional, and clinging to a past era. Their symbolic rivalry culminates in Norma's jealousy-fuelled murder of Joe.

\subsubsection*{Key Quotes (Implied Tension)}
\begin{itemize}
    \item \textit{Norma (to Joe):} 'Is it a woman? I know it's a woman...'
    \item \textit{Joe:} 'All right, baby. This way out.'
    \item \textit{Betty:} 'If you love me, Joe.'
\end{itemize}

\subsubsection*{Key Ideas}
\begin{itemize}
    \item Rivalry
    \item Gendered insecurity
    \item Age and beauty
\end{itemize}

\subsection{Norma Desmond and Cecil B. deMille}
Norma idolises deMille as her old mentor, believing he will revive her career. deMille pities Norma, indulges her briefly, but ultimately distances himself. Their relationship exemplifies Hollywood’s betrayal of its past stars.

\subsubsection*{Key Quotes}
\begin{itemize}
    \item \textit{Norma:} 'Now Mr. deMille can wait till I'm good and ready.'
    \item \textit{deMille:} 'Thirty million fans have given her the brush. Isn’t that enough?'
    \item \textit{Norma:} 'We'll be working again, won't we, Chief?'
    \item \textit{deMille:} 'A dozen press agents working overtime can do terrible things to the human spirit.'
\end{itemize}

\subsubsection*{Key Ideas}
\begin{itemize}
    \item Nostalgia
    \item Fame and decline
    \item Compassion vs. cruelty
\end{itemize}

\subsection{Max Von Mayerling and Betty Schaefer}
Though they never meet, Max and Betty represent two opposite relationships with Joe: one built on silent loyalty and loss of self (Max), the other on mutual respect and potential for growth (Betty). Max is the ghost of what Joe could become if he fully submits to Norma.

\subsubsection*{Key Quotes (Implied Contrast)}
\begin{itemize}
    \item \textit{Max:} 'I directed all her early pictures. I made her a star.'
    \item \textit{Betty:} 'I want to write. I don’t want to be a reader all my life.'
\end{itemize}

\subsubsection*{Key Ideas}
\begin{itemize}
    \item Self-sacrifice vs. ambition
    \item Emotional entrapment vs. liberation
    \item Generational contrast
\end{itemize}

\subsection{Sheldrake and Betty Schaefer}
Betty and Sheldrake represent opposite ends of the studio power structure. Sheldrake controls opportunities and trivialises creative work, while Betty strives to have a voice as a writer. Their professional relationship reflects the dismissiveness of studio executives toward aspiring creatives.

\subsubsection*{Key Quotes}
\begin{itemize}
    \item \textit{Sheldrake:} 'Make it a girls’ softball team... it’d make a cute musical.'
    \item \textit{Betty:} 'I kind of hoped to get in on this deal. I don’t want to be a reader all my life.'
\end{itemize}

\subsubsection*{Key Ideas}
\begin{itemize}
    \item Industry hierarchy
    \item Gender and creative voice
    \item Power and dismissiveness
\end{itemize}

subsection{Joe Gillis and Cecil B. deMille}
Joe’s view of deMille is shaped by disillusionment; deMille represents the **industry's indifference** toward people like Norma and himself. Their interaction is limited, but deMille’s treatment of Norma is indirectly relayed to Joe—highlighting Hollywood’s tendency to sentimentalise the past while refusing to truly engage with it.

\subsubsection*{Key Quotes}
\begin{itemize}
    \item \textit{deMille:} 'Thirty million fans have given her the brush. Isn’t that enough?'
    \item \textit{Joe (on deMille):} 'He didn’t have the heart to tell you. None of us has had the heart.'
\end{itemize}

\subsubsection*{Key Ideas}
\begin{itemize}
    \item Fame and abandonment
    \item Disillusionment with power
    \item Truth vs. kindness
\end{itemize}

\subsection{Betty Schaefer and Norma Desmond}
Though Betty and Norma never speak directly, their relationship exists as a powerful symbolic contrast—two women representing vastly different eras, values, and roles in Hollywood. Betty is young, ambitious, and grounded in reality; Norma is ageing, delusional, and obsessed with regaining lost fame. They become rivals not only for Joe’s affection but also as emblematic opposites: one looking forward, the other trapped in the past. Norma’s jealousy toward Betty culminates in Joe’s death—making their indirect relationship one of the most consequential in the film.

\subsubsection*{Key Quotes}
\begin{itemize}
    \item \textit{Betty:} 'It’s not your career — it’s mine. I kind of hoped to get in on this deal. I don’t want to be a reader all my life. I want to write.'
    \item \textit{Norma:} 'Is it a woman? I know it's a woman...'
    \item \textit{Norma:} 'You're not leaving me!'
    \item \textit{Joe (to Betty):} 'Good luck to you, Betty... here's the pool.'
\end{itemize}

\subsubsection*{Key Ideas}
\begin{itemize}
    \item Generational conflict
    \item Female ambition vs. faded fame
    \item Jealousy and romantic rivalry
    \item Hope vs. delusion
\end{itemize}

\subsection*{Relationship with Ambition}
Ambition is a desire and determination to achieve success. It is a motivational force that compels individuals to set and pursue goals, often with a willingness to work hard, overcome obstacles, and make sacrifices to attain those goals.\\

\subsubsection{Norma Desmond}
\begin{itemize}
\item Norma’s ambition is rooted in her desire to regain her past Hollywood stardom and in her fear of becoming irrelevant.
\item It is also sustained by her delusional belief in her continued desirability as an actress, despite Hollywood moving on. 
\item Norma’s isolation and the presence of enablers who validate her fantasies further fuel her ambition and detachment from reality.
\item In response to the word comeback: ‘I hate that word. It’s a return, a return to the millions of people who have never forgiven me for deserting the screen.’
\end{itemize}

\subsubsection{Joe Gillis}
\begin{itemize}
\item Joe’s ambition arises from financial desperation and a desire for Hollywood success.
\item His pursuit of financial security and career advancement drives him into a complex relationship with Norma.
\item ‘The poor dope - he always wanted a pool. Well, in the end, he got himself a pool.’
\end{itemize}


\subsubsection{Betty Schaefer}
\begin{itemize}
\item Betty is driven by her genuine passion for screenwriting and a strong desire to succeed in the film industry.
\item ‘What’s wrong with being on the other side of the cameras? It’s really more fun.’
\end{itemize}
\subsubsection{Max von Mayerling}
\begin{itemize}
\item Max’s ambition is solely devoted to serving and preserving Norma Desmond’s illusion of fame.
\item He is driven by his unwavering dedication to Norma, having been her director and husband during her stardom.
\item ‘Madame is the greatest star of them all.’
\end{itemize}
%%%%%%%%%%%%%%%%%%%%%%%%
\subsection*{Relationship with Delusion}
\subsubsection{Joe’s agent, Morino}
Wilder characterises Joe’s agent Morino as a representative of the cutthroat nature of Hollywood. Morino says he could give Gillis money… but he won’t. Instead, Gillis should use his desperate state as inspiration – it’ll force him to write because he’ll need to do it.

\subsubsection{Betty Schaefer}
Betty is portrayed as a down-to-earth, sincere, and practical character. She represents a departure from the delusional world of Hollywood, focusing on her genuine passion for screenwriting and her desire for authentic, meaningful connections.

\subsubsection{Norma Desmond}
\begin{itemize}
\item ‘All right, Mr. DeMille, I’m ready for my close-up.’
\item In the film’s iconic closing line, Norma's delusional state of mind reaches its peak, marking her descent into madness. While she retained some connection to reality in earlier scenes, this quote signifies the complete erasure of any remaining self-awareness. Norma regresses to her past glory, even mistaking Max for director Cecil DeMille, vividly highlighting the profound transformation of her fantasy world.
\end{itemize}

\subsubsection{Joe Gillis}
\begin{itemize}
\item ‘You don’t yell at a sleepwalker – he may fall and break his neck. That’s it: She was still sleepwalking along the giddy heights of a lost career.’
\item Joe compares Norma to a sleepwalker, suggesting she is unaware of her actions and the reality of her situation, emphasising the delicate balance between preserving her illusions and confronting the truth of her faded career.
\end{itemize}

\subsubsection{Max von Mayerling}
\begin{itemize}
\item ‘That is my job. It has been for a long time. You must understand I discovered her when she was eighteen. I made her a star. I cannot let her be destroyed.’
\item By nourishing her illusions and crafting fan mail, Max sees himself as Norma’s guardian, dedicated to preserving her faded fame. Yet, as he fervently conceals the truth to protect her from despair, his deceptions inadvertently propel her further into madness.
\end{itemize}

\subsection*{Relationship with Greed}
\begin{figure}[h!]
    \centering
    \includegraphics[width=1\linewidth]{gtable.png}
    \caption{Who Decieves Who? in \underline{Sunset Boulevard}}
\end{figure}


\newpage
\section{Scene Summaries and Analyses}
\subsection{Scene 1 : The Corpse (0:00:0 - 0:02:38)}
Summary : Opening Credits; a body is discovered in the pool of a Hollywood Mansion.\\
1. The opening music and the long shot of the road in the film contribute to creating a mood of dramatic, lonely, mystery, tension, apprehension, terrified and shocked.\\
2. When the voice-over explains that the 'poor dope' always wanted a pool, the price 'turned out to be too high', the purpose of this voice-over to refer to the main\\
character Joe Gillis, the outcome of his desires and decision to the dark side of Hollywood. Other answer : The voice meant that the man always wanted a pod. But ironically he ended up dead in it, the price he paid was life.\\
3. Joe Gillis voice-over is a key storytelling tool. It not only leads us through the film's scenes but also gives us a deeper focus into Joe's character, his challenging, his struggling and the tough realities he encounters in the competitive Hollywood environment.' Other answer : We learn that the narrator is a young man who is nobody important, just a movie script-writer.\\

\subsection{Scene 2 : The Corpse (0:02:38 - 0:09:55)}
Summary : The finance company threatens to repossess Joe’s car. Joe tries unsuccessfully
to seek the job and borrow the money from his friend to save his car from the finance
company.\\
1. My first impressions of Joe Gillis are of a character who is disillusioned and
somewhat cynical about Hollywood. He appears to be a struggling screenwriter who
feels overlooked and underappreciated by the industry.\\
2. The entrance gates to Paramount Studios illustrate the attraction and exclusivity of
Hollywood. They represent the gateway to fame, success, and the charming world of
the movie industry. (represent wealth and power)\\
3. Schwab's Pharmacy is a drugstore, serving as a multifaceted symbol of Hollywood's
allure, desire, and the stark contrast between illusion and reality within the film
construction.\\
4. Betty Schaefer's appraisal of Joe's script 'Bases Loaded' has a particular impact on
Joe both personally and professionally and Joe is furious with Betty because the way
she said to Joe, it can affect him to seek his job. (being resentful, uncomfortable,
frustrated and saying she was just a message kid.)\\


\subsection{Scene 3 : A White Elephant of a House (0:09:55 - 0:12:55)}
Summary : Joe is pursued by the finance company who wants to obtain his car, he escapes
by landing himself in the driveway of some big mansion and storing his car into the huge
garage at that mansion.\\
1. The music score during the car chase supplies a powerful tool for conveying Joe's
emotions, as well as amplifying the suspense, drama, and darkly comedic elements
of the film. After he is landing himself in the driveway in a big mansion, the music
slows down which shows Joe’s emotion of peace. In addition, the panicky energy of
the car chase, the surrounding area in the driveway may provide the risk of serious
car accident.\\
2. Joe narrates about the house that ‘a neglected house gets an unhappy look. This
one had it in spades’, it evoked a sense of disquiet and desolation which makes the
audience assume that the people who live in this mansion might be struggling with
personal or emotional challenges.\\
3. Joe's plan to stash the car displays his resourcefulness, street smarts and
cautiousness, all of which are essential aspects of his character as he navigates the
complexities of life in Hollywood.\\

\subsection{Scene 4 : ‘I am big. It’s the picture that got small’ (0:12:55 - 0:17:04)}
Summary : Joe is mistaken for walking into this old huge mansion of his curiosity. He meets
the isolated former silent film star and recognises her as Norma Desmond.\\
1. We first see Norma from a low angle shot looking up at her in the window watching
Joe, which impresses with power, mystery, weakness, and greatness.\\
2. Norma's bedroom is a luxury designed space that serves as a reflection of her past
fame, present circumstances, and complex inner world.\\
3. Wilder productively presents Norma Desmond's belief that she is still ' big ' in
Hollywood through a combination of visual storytelling, dialogue, character
interactions, and narrative development.\\
4. Despite their contrasting personalities and conditions, Joe and Norma share common
themes and experiences related to their desires, struggles, and illusions in the world
of Hollywood.\\

\subsection{Scene 5 : The Salome Script (0:17:04 - 0:23:01)}
1. Norma comments about her comeback, ‘A return to the millions of people who’ve
never forgiven me for destroying the screen’, she means her complicated
relationship with fame, legacy, possessiveness and self-perception.\\
2. After Joe reads through Norma’s script, he views it as outdated, melodramatic, and
unsuitable for modern audiences.\\
3. The close-up on Norma's face when she suggests Joe stay the night reveals her
intricate emotions of loneliness, dependency, desperation and optimism about
returning to stardom again.\\
4. The mise-en-scène in Norma's salon is accurately crafted to reflect her personality,
status, history, mindset and the overall atmosphere of her character's world.\\

\subsection{Scene 6 : An Invitation to Stay (0:23:01 - 0:28:04)}
Summary : Joe spends the night in Norma's garage due to the need to modify Norma’s
script, Max takes Joe’s belongings from his apartment as what Norma desires.\\
1. The burial ceremony for the chimpanzee provides insight into Norma Desmond's life,
highlighting her extravagance, eccentricity, and underlying feelings of loneliness.\\
2. Norma's decision to pay Joe's unpaid rent is driven by a combination of infatuation,
manipulation, fear, and a desire for control.\\
3. Billy Wilder illustrates Norma as the one in control through her artful behaviour,
domineering presence, financial influence and emotional manipulation.\\

\subsection{Scene 7 : The Face of a Star (0:28:04 - 0:31:47)}
Summary : Joe works on Norma’s script. Joe and Norma watch a movie in the living room
together about Norma’s silent movie film in which she is a main actor.\\
1. When Joe attempts to change her script, Norma reacts irritated, terrified and shocked
due to Joe removing some scenes which relate to her acting. In addition, Norma
wants to have her in all scenes of the movie because she believes that her fans
obsess to see her every episode.\\
2. The panning of the photos of Norma in her mansion serves to illustrate her obsession
with her past stardom and her desperate attempt to hold onto her fading fame.\\
3. The irony in Norma showing the 'idiot producers' ' what a star looks like underlines
the themes of delusion, self-confidence, and tragedy. She stopped working at the
Paramount studio and she didn’t have a beautiful face anymore like before.\\
4. Norma's clutching on Joe's arm conveys a beginning of relationship emotions,
including desire, fear, insecurity, control, and longing.\\

\subsection{Scene 8 : Bridge with the Wax Works (0:31:47 - 0:36:10)}
Summary : Joe loses his car because the finance company comes to Norma’s house and
informs Max that they will repossess Joe’s car due to his 3 payments behind. Norma plays a
little bridge game with Norma’s actor friends. She takes Joe to a men's clothing shop to buy
the shirt for Joe as he’s wearing an old suit everyday.\\
1. The 'waxworks” were former silent film stars who, like Norma, have fading their
careers, delusional, and isolated.\\
2. The quote that shows Joe is becoming a kept man is 'He was the kept boy of a now
somewhat lady-like but very rich lady, about fifteen years older than he was.' which
Joe said to reflect his own situation. Well, as long as the lady's paying for it. Why not
take the vicuna.\\

\subsection{Scene 9 : The Husbands’ Bedroom (0:36:10 - 0:39:31)}
Summary : Joe moves into the main house to live in Norma’s ex-husbands room.\\
1. Joe's move into 'the room of the husband' which indicate various possibilities for the
track of their relationship and Joe is increasingly becoming a part of Norma's destiny.\\
2. This scene makes us believe that Norma is living in a fantasy world and not facing
reality based on the fact that she has moments of sadness which in her mansion
doesn’t have locks, sleeping pills, razor blades and shut off the gas in her room.\\
3. The mise en scene in Norma's bedroom is carefully crafted to convey themes of
longing, decay, and fantasy. In addition, the bedroom represents a place that reflects
Norma's emotion. Satin, drapes, bedding, ruffles, bed like a gilded rowboat (covered in
gold), old fashioned, old furniture. Norma’s living in the past. Keep her lavish
lifestyle.\\

\subsection{Scene 10 : New Year’s Eve (0:39:31 - 0:53:50)}
Summary : Norma declares that there are only two of them in this New Year’s eve party and
she confesses to Joe that she fell in love with him. Joe escapes Norma after they have an
argument between each other about their relationship and Joe returns because when he
calls Max to back up his belongings, someone will pick them up, but Max tells him that the
doctor has come to Norma’s mansion because Norma wants to kill herself.\\
1. The gold engraved cigarette case symbolises Norma’s present to Joe for New Year’s
eve and about her emotion of love with Joe.\\
2. New Year’s Eve from Max’s perspective in which he remembers his old love
memories with Norma when they were husband and wife.\\
3. The turning point for Joe is when he returns to the mansion, after Max calls him to tell
about Norma wants to murder herself because Joe rejects her love for him, then Joe
decides to become Norma’s partner from that night.\\
NOTE: Norma slit her wrists.\\

\subsection{Scene 11 : A Kept Man (0:53:50 - 0:58:38)}
Summary : Betty is trying to persuade and encourage Joe to collaborate on a screenplay
with her, Norma sends Max to ask Joe to come back to the car and Betty wants Joe to work
with her in a new story.\\
1. The interactions between Norma and Joe are characterised by power struggles,
complicated emotions and dependency.\\
2. In this scene, we learn about Betty's hopefulness, working hard and she wants to be
a screenwriter rather than read scripture. In addition, she persuades Joe to work with
her in a new movie that she creates by herself.\\
3. When Norma gives money to Joe to buy cigarettes, the emphasis of Joe doesn’t
have enough money to buy cigarettes. Moreover, Joe is dependent on Norma’s
wealth and Norma has the power to control Joe.\\

\subsection{Scene 12 : The Norma Desmond Follies (0:58:38 - 1:01:24)}
Summary : Norma shows a Charlie Chaplin comedy performance to Joe when she suspects
that Joe’s getting bored with her and Norma receives a telephone call from the Paramount
studio.\\
1. Norma Desmond entertains Joe Gillis as Charlie Chaplin because she doesn’t want
Joe to get bored with her, so she tries to show every comedy's famous movie to
make Joe enjoy living with her and she is scared that Joe will leave her one day.\\
2. The contrast between Joe when he first came to Hollywood and now is when Joe first
arrives at Hollywood, he is hopeful and believes that he can find a great job and be a
successful screenwriter. However, until now he becomes disillusioned with Hollywood
and lives a life which is dependent on Norma’s wealth.\\

\subsection{Scene 13 : Parading to Paramount 9 (1:01:24 - 1:04:42)}
Summary : Norma visits Paramount studio after urgent calls from Paramount.\\
1. The emphasis of the difference between past and present in this scene is the
homesickness and romanticism of Hollywood and the harsh realities of ageing and
unimportance faced by its former stars. For example, after 20 years, when Norma
goes back to Paramount Studio the new generation no one knows her except old
people who work with her in that time and also no one in that studio wants her
anymore due to her age and obsession.\\
2. Norma Desmond and Isotta Fraschini share similarities in their representation of
faded their famous, delusions of beauty and they both serve as powerful symbols in
the past.\\

\subsection{Scene 14 : A Meeting with Mr. Demille (1:04:42 - 1:12:48)}
Summary : Norma speaks to Mr. Demille about her script and someone calling her, Joe talks
with Betty.\\
1. The high angle shot of Stage 18 at Paramount Studios provides viewers with a
perspective of the studio that emphasises its greatness and glory. When the camera
is looking down on Norma Desmond it signifies her fading of fame and
powerlessness.\\
2. When Norma says to Demille, ' And remember, darling. I don’t work before 10 in the
morning and never after 4:30 in the afternoon”, she is failing to realise the extent of
her own delusion, unreasonableness and the reality of her situation in the modern
film industry.\\
3. In this scene, we learn that Max is loyal and intelligent because he found out the
reason for all those phone calls from Paramount, they’re not want Norma but they
want Norma’s car and he didn’t tell Norma about it because he is scared that Norma
will be furious and disappointed. Moreover, he is also a former employee who worked
in Paramount studio as a director.\\
4. The close-ups of Joe and Betty talking, allow Wilder to express the depth and
beginning of their relationship, complicated for them to communicate to each other
and the emotional difficulty of their interactions.\\

\subsection{Scene 15 : Joe Sneak Out (1:12:48 - 1:17:56)}
Summary : Norma is going through a merciless series of treatment. Joe meets Betty
Schafer at night to write a screenplay together.\\
1. Norma's efforts to prepare for her 'return' to the screen are characterised by
hopefulness, obsession and she thinks that there are a lot of fans who are waiting for
her.\\
2. Wilder shows the power Norma has over Joe though she doesn’t allow Joe to go
anywhere what he wants, when Joe goes somewhere he needs to tell Norma and
Joe doesn’t have freedom. She said ' You went out last night, didn’t you Joe ?”.\\
3. Joe desires to get free from his dependence on Norma and control over his own life
which is the symbol for it being taken away represents his demise.\\
4. The irony in Joe’s statement 'That’s the trouble with you reader, you know all the
plots.” Joe's statement carries irony in its criticism of Betty's understanding of plots
while at the same time reflecting his own fault as a writer and his final fall within the
imprison of the Hollywood industry.\\

\subsection{Scene 16 : Untitled Love Story (1:17:56 - 1:26:19)}
Summary : Joe and Betty become closer friends. Max reveals to Joe that he was the first
husband of Norma. Norma discovers that Joe and Betty declare they have feelings for each
other.\\
1. Wilder uses the collaborative process of writing a screenplay at night to illustrate the
slow growth of Joe and Betty's relationship, from writers' partners to close friends
and, eventually, romantic interests.\\
2. Joe and Max have a confrontation about Joe going out every night to write a
screenplay with Betty which can make Norma furious and jealous about it. Max, a
former husband who still loves her and protects Norma by writing fan letters to
Norma and lying to Norma about her comeback to the studio. In addition, when
Norma was young, he was one of the directors who made Norma become a famous
silent film star and he didn't want to destroy her.\\
3. Wilder uses camera angles and framing in the last scene between Joe and Betty to
represent their true love and deep connection between each other but they can’t be
together due to a lot of complicated drama. Ultimately, Joe and Betty's relationship
can’t continue anymore after Betty knows the truth about Joe living with Norma
Desmond.\\

\subsection{Scene 17 : Betty’s Exit (1:26:19 - 1:34:35)}
Summary : Norma tells the truth about Joe and her living arrangement to Betty. Betty comes
to Norma’s mansion to find out about it and Joe introduces/leaves/ Norma’s house to Betty.\\
1. Joe's decision to invite Betty Schaefer to Norma Desmond's mansion shows his
growing sense of honesty, responsibility, acceptance of reality and he tells Betty
about his situation that he has a good deal in this house and he doesn’t want to go
to one room apartment which he can’t pay for and even the screenplay which they
write together.\\
2. The techniques that Wilder uses to show a shift in power and build up of tension is
cinematography, lighting, sound design, dialogue, set design, editing and creating a
scene rich with dramatic intensity, emotional depth and low angels light camera
between Joe and Norma which makes him become more powerful on Norma.\\
3. Betty's feelings of curiosity, anxiety, concern for Joe,support him through whatever
challenges they may face and she deals with that by always staying on Joe’s side
whatever happens to him.\\
4. Norma’s actions and feelings in this scene is when she calls Betty to tell her the facts
about Joe and her relationship and their living arrangement because she is
increasingly jealous and scared that Joe will leave her to live with another girl one
day who is younger than her. Norma Desmond says 'men of his sort', she is likely
talking to men who are selfish, cowardly and take advantage of others for their own
benefits .\\
5. It is ironic that Joe turns the lights on in the pool because he wants to tell Betty about
the danger, deception and darkness of his life with Norma in her mansion.

\subsection{Scene 18 : Noone Ever Leaves A Star (1:34:35-1:39:13)}
Summary : Joe tells Norma that he will leave the house and go back to his hometown, she
tells Joe that she will kill herself if he leaves her.\\
1. The close-up on Norma's face as she looks in the mirror represents a max
complication of emotions, including hopeful and satisfied Joe deciding to be
desperate from Betty because he gets a lot of benefits from this house.\\
2. The context of the scene when Joe says to Norma 'The audience left twenty years
ago” is he wants to inform Norma to wake up from the past, self-obsession,
daydreaming and delusion because nowaday no one cares about her anymore.\\
3. Norma persuades and warns Joe not to abandate her, if he goes away from her, she
will shoot him but Joe doesn’t listen to Norma’s saying, then Norma decides to shoot
Joe on the way out the mansion and he falls into the swimming pool where Joe’s
end.\\

\subsection{Scene 19 : ‘All right Mr deMille, I’m ready for my close up’
(1:34:35-1:39:13)}
Summary : The house is surrounded by the policies, the reporter neighbours and passer-by.
Joe’s body is floating in the swimming pool. Norma has a mental shock after she kills Joe.\\
1. The narrative style adds depth to the storytelling, allowing the audience to empathise
with Joe's experiences while also providing important scenes through which to view
the themes of fame, obsession, examining the dark side of Hollywood and
disillusionment within the film.\\
2. Black humour is evident when the voice over says ' Here was an item everybody
could have fun with”, which is referring to Joe’s body floating in the swimming pool
who wishes to go into Hollywood, as a result end up his life in the water.\\
3. Norma won’t come down the stairs until she hears the word 'camera” Illustrate her
deep-emotions of delusions and obsession with her past fame as a silent film star.
Norma has become isolated and disconnected from reality, living in a huge mansion
filled with memories of her fame days in Hollywood.\\
4. Norma's feelings as she walks down the stairs are a mix of pride, hopefulness,
insecurity, loneliness, and delusion because she thinks she has come back to her
place which she waited for 20 years.\\
5. Yes, I do. Wilder conveys a sense of pity, shallowness and compassion for Norma
Doesmond through inviting the audience to empathise with her challenges and
struggles.\\

\newpage
\section{Film Techniques}
\subsection{List of Techniques}
\renewcommand{\arraystretch}{1.3}
\begin{longtable}{|p{5cm}|p{\dimexpr\linewidth-5cm-2\tabcolsep-2\arrayrulewidth}|}
\caption{Film Techniques and Metalanguage} \label{tab:filmtechmeta} \\
\hline
\rowcolor[gray]{0.9} \textbf{Film technique / metalanguage term} & \textbf{Definition / Effect} \\
\hline
\endfirsthead

\hline
\rowcolor[gray]{0.9} \textbf{Film technique / metalanguage term} & \textbf{Definition / Effect} \\
\hline
\endhead

% your rows
Close-up shot & Tightly frames a person or a subject. Enables the audience to see the emotions of a character or details on a subject. \\
\hline
Extreme close-up & A more intense close-up that might focus on a specific feature. \\
\hline
Mid shot & Usually shows an actor from the waist up. Focus is on the characters but also shows the environment they are in. \\
\hline
Long shot & Used to show a subject from a distance, establishing them within a setting. \\
\hline
Establishing shot & Usually the first shot of a scene to establish the location and environment. \\
\hline
Panning shot & When the camera moves horizontally from its fixed position and mirrors the motion of looking from side to side. \\
\hline
Point of view shot & Mimics what a character in a scene is seeing, putting the audience into their perspective. \\
\hline
Framing & The composition of characters or subjects within a single shot. \\
\hline
Two shot & A medium shot showing two actors in a single frame. \\
\hline
Mise en scène & All elements contained in a single frame (lighting, costuming, positioning, set, props). \\
\hline
Low angle shot & Subject is shot from below eye level to make them appear powerful, dominant, or morally superior. \\
\hline
High angle shot & Subject is filmed above eye level to make them appear vulnerable or weak. \\
\hline
Aerial shot & A shot with the camera positioned directly above a subject. \\
\hline
Over the shoulder shot & A medium close-up shot from behind another character’s shoulder, typically used in conversations. \\
\hline
Diegetic sound & Sound within the world of the film — characters can hear it. \\
\hline
Non-diegetic sound & Sound added in post-production — characters cannot hear it (e.g. soundtrack, narration). \\
\hline
Setting & The environment in which the film is set. Can shift over time. \\
\hline
Motifs & A dominant or recurring prop, image, or symbol (e.g. mirrors, frames, candles, chandeliers). \\
\hline
Zoom in / Zoom out & A change in focal length during a shot to magnify or reduce the subject. \\
\hline
Film noir lighting & Low-key lighting used to emphasise shadows and create strong contrasts. \\
\hline
Three-point lighting & Technique using key light, back light, and fill light to evenly illuminate a subject. \\
\hline

\end{longtable}
\newpage
\subsection{Analysis of Particular Scenes}
\underline{Sunset Boulevard} is perhaps the most famous film about film. A darkly funny yet disturbing noir, it follows washed-up screenwriter Joe Gillis being pulled into the murky world of even-more-washed-up former silent film star Norma Desmond, disingenuously helping with her screenplay. Critical commentary on the film industry is obviously included here, but Billy Wilder’s 1950 film digs deeper to explore the blurred line between fantasy and reality, as well as power, authenticity and self-delusion. Crucially, these themes are often shown in the film’s construction, via the cinematic techniques implemented by Wilder in each scene.\\

This subsection will explore the most important examples of these cinematic techniques. Remember, VCE examiners are on the lookout for students who can offer a close reading of the text they are discussing, giving specific examples of how its creator has constructed it to support their arguments. Just look at the difference between an essay that says:

\begin{quote}
    'Through the final shot of the film, Wilder shows Norma completely succumbing to her fantasy.’ 
\end{quote}

Compared to one that argues:

\begin{quote}
    'Through his utilisation of an increasingly glossy and distorted filter in the ominous final shot, Wilder depicts Norma being completely overtaken by her romanticised fantasy of ‘Old Hollywood’.'
\end{quote}

\subsubsection{Camera Techniques: Shot Types and Angles}
Camera techniques are arguably the primary way that a director will intentionally direct the eye of the audience, directly framing how they view a film. The two most basic ways in which the camera is used for this are through the distance between the subject (what the scene is about) and the camera, or the ‘shot type’ and the ‘camera angle’ at which the subject is being filmed.\\

\begin{figure} [h!]
    \centering
    
\includegraphics[width=0.75\linewidth]{ShotType1.jpeg}
    \caption{First Appearance of Norma Desmond}
    \label{fig:enter-label}
\end{figure}

Our first look at Norma Desmond is within the wide shot above, just as Joe Gillis has entered her dishevelled mansion early in the film. As a rule, the introductory shot of a character is always worth closely analysing, as the director typically establishes their characteristics and place within the film’s wider world. \\

Shown above, this distant first look at Norma establishes her distance, both physical and mental, from the world around her. Removing herself from an industry that has long since moved on from her, she is severely out of touch with the reality of the world outside her home. Crucially, as this same shot is from Joe’s perspective, Wilder also foreshadows the more specific character ‘distance’ that will emerge between the two. Here, the audience sees the space Joe will similarly leave between himself and Norma, disingenuously humouring her poor-quality scripts and romantic advances and, therefore, always keeping her ‘at a distance’.

\begin{figure} [h!]
    \centering
    \includegraphics[width=0.75\linewidth]{Shot2.jpeg}
    \caption{The Final Goodbye}
    \label{fig:enter-label}
\end{figure}Another shot conveying crucial information about character relationships is shown when Joe officially ‘loses’ Betty towards the end of the film, refusing to give up his ‘long-term contract’ with Norma. Here, Wilder consciously frames the scene’s subject (Betty) at a distance with a medium shot. Supported by her refusal to make eye contact with Joe and her literal statement that she ‘can't look at [him]’ we again see physical distance between the camera and the subject translating to emotional distance between two characters. The impact of them no longer ‘seeing eye to eye’ is additionally heightened by the clear chemistry they previously demonstrated across the film. \\

\subsubsection{Key Examples of Camera Angles}
\begin{figure} [h!]
    \centering
    \includegraphics[width=0.75\linewidth]{Ang1.jpeg}
    \caption{Opening Title}
    \label{fig:enter-label}
\end{figure}

Just like the introductory shot of a character is worth digging into, the opening shot of a film is also incredibly important to unpack. Sunset Boulevard’s seemingly straightforward opening shot simply includes the film’s title, by showing the real-life Hollywood street. However, notice that we are not seeing a ‘Sunset Boulevard’ street sign (the more obvious choice), but instead a dirty and stained curbside. Further, Wilder shoots this curb from a high angle. Therefore, the film’s opening shot establishes maybe the most central aim of Wilder’s film; offering a critical look at the superficiality and flawed nature of Hollywood. As such, we are literally looking down on the film industry in the first moment of the film, represented by this dirty and unflattering visual symbol of Hollywood. This, therefore, is setting the stage for the satire and critical commentary that will follow. \\

\begin{figure}[h!]
    \centering
    
\includegraphics[width=0.75\linewidth]{ang2.jpeg}
    \caption{After Losing Betty}
    \label{fig:enter-label}
\end{figure}
Wilder’s careful use of camera angles is further shown at the end of the film after Betty abandons Joe at the gate of Norma’s mansion. Crucially, this all happened due to the desperate exertion of power by Norma, who called Betty and revealed the details of her relationship with Joe. As such, Wilder shoots Norma at a low angle, as Joe looks up at her haughty gaze. The level of power that Norma has exerted over Joe may seem minimal within the moment, but when we consider what happens next, this shot becomes much more important. On the brink of descending completely into madness and taking Joe’s life, Wilder uses this shot to establish that Joe should be looking up in fear at Norma, and his dismissive and pitiful opinion of her will soon lead to his death.\\


\subsubsection{Mise-en-Scène}
Mise-en-scène is perhaps the most deceptively simple cinematic technique. It involves analysing what appears within a frame and where it has been placed by the director. This includes elements such as the actor’s costumes, the props and the design of the set. Often, mise-en-scène is used to reinforce something we are being told about a character already through the film’s dialogue and acting.  \\

\begin{figure} [h!]
    \centering
    \includegraphics[width=0.75\linewidth]{mes1.jpeg}
    \caption{Joe's Apartment I}
    \label{fig:enter-label}
\end{figure}

\begin{figure} [h!]
    \centering
    \includegraphics[width=0.75\linewidth]{mes2.jpeg}
    \caption{Joe's Apartment II}
    \label{fig:enter-label}
\end{figure}

We can see a key example of characterisation through mise-en-scène early in the film, where the audience’s introduction to Joe Gillis visually communicates his unconcerned and detached attitude, as well as his tendency to settle for something convenient despite its inauthenticity. His being dressed in a bathrobe with the blazing sun outside (and his debt collectors clearly up and doing their jobs) speaks to his slovenliness and uninvested approach to life. The set design within this scene further characterises Joe, with the script directly describing the ‘reproductions of characterless paintings’ that cover his walls. Here, the set arguably provides a visual metaphor for the profit-driven ‘Bases Loaded’ script he is writing at that very moment, later described by Betty as having come ‘from hunger.’ 


\begin{figure} [h!]
    \centering
    \includegraphics[width=0.75\linewidth]{nb1.jpeg}
    \caption{Norma's Bedroom}
    \label{fig:enter-label}
\end{figure}

\begin{figure} [h!]
    \centering
    \includegraphics[width=0.75\linewidth]{nb2.jpeg}
    \caption{Norma's Living Room}
    \label{fig:enter-label}
\end{figure}

Equally, our introduction to the home of Norma Desmond helps establish the key elements of her character. The house is, as Joe describes, ‘crowded with Norma Desmonds’, in the form of countless framed photos of her from her silent film era. These self-portraits constantly looking out onto Norma symbolise the deluded fantasy world she has placed herself in. They both show how this world is based around her still being a youthful and famous actress, and that this delusion is maintained through Norma only communicating inwardly, refusing to face the reality of the outside world.



ZZPDFEMBED ShotAnalysis.pdf ZZEND






\paragraph{Further Mise-en-Scène Analysis}
\subparagraph{Set Design and Locations}
\begin{itemize}
    \item Norma Desmond’s Mansion: The opulent yet decaying mansion is central to the film. It represents Norma’s past glory and her current isolation. The grand, dark interiors are filled with outdated, Gothic furnishings, creating a contrast with the modern Hollywood world outside. This setting emphasises the disconnect betweenN Norma’s delusions and reality.
    \item Joe Gillis’ Apartment: Joe’s modest apartment contrasts sharply with Norma’s mansion. Its simplicity reflects his struggle and the reality of his life as a struggling screenwriter. The clutter and starkness of the apartment underscore Joe's financial and emotional instability.
    \item Paramount Studios: The scenes at Paramount Studios, vibrant and bustling,symbolise the thriving, contemporary Hollywood that has left Norma and her era behind. The contrast between the lively studio and Norma’s decaying home highlights the industry's constant evolution and neglect of past stars.
\end{itemize}

\subparagraph{Lighting}
\begin{itemize}
    \item Low-Key Lighting: especially in Norma’s mansion. This creates deep shadows,enhancing the film’s noir atmosphere. The darkness of the mansion’s interiors mirrors Norma’s mental state and her descent into madness.
    \item High-Key Lighting: In contrast, scenes set in the modern parts of Hollywood, such as the Paramount Studios, use high-key lighting. This bright and even lighting reflects the vibrancy and energy of the current film industry and emphasises the disconnect between the past and present
\end{itemize}

\subparagraph{Costume Design}
\begin{itemize}
    \item Norma Desmond’s Costumes: Norma’s costumes are elaborate and dramatic,reflecting her former star status and her attempt to maintain an image of grandeur.Her wardrobe includes vintage, glamorous outfits that emphasise her clinging to a bygone era. The costumes also highlight her delusions of grandeur and her disconnect from reality.
    \item Joe Gillis’ Costumes: Joe’s wardrobe is more casual and practical, reflecting his status as a struggling screenwriter. His clothing is less flashy, indicating his lower position in the Hollywood hierarchy and his struggle to make a name for himself.
\end{itemize}

\subparagraph{Camera Angles and Framing}
\begin{itemize}
    \item Close-Ups: The fiilm frequently uses close-ups, particularly of Norma’s face. These shots reveal her emotional state and add to the film’s psychological intensity. Thec lose-ups emphasise her vulnerability, madness, and the contrast between her internal and external worlds.
    \item Wide Shots: Wide shots of Norma’s mansion often include extensive background details, emphasising the grandeur and the sense of decay. These shots also highlight the mansion’s isolation and Norma’s detachment from the world.
\end{itemize}

\subparagraph{Composition}
\begin{itemize}
    \item Symmetrical and Asymmetrical Composition: The film uses both symmetrical and asymmetrical compositions to reflect character relationships and themes.Symmetrical compositions often highlight Norma’s desire for control and order,while asymmetrical compositions reflect the chaos and imbalance in her life.
    \item Depth of Field: The use of depth of field helps separate characters from their environment, particularly in scenes where Norma’s fantasies clash with reality. This technique underscores the distance between her delusions and the real world.
\end{itemize}



\subsubsection{Symbols in Scenes}
\begin{figure} [h!]
    \centering
    \includegraphics[width=0.75\linewidth]{chimp.jpeg}
    \caption{Chimp}
    \label{fig:enter-label}
\end{figure}
One of the more seemingly inexplicable parts of Wilder's film actually contains one of its most important symbols, with Norma’s pet monkey playing a key foreshadowing role from beyond the grave. The chimp, a pet owned and trained by Norma to amuse her, leaves a vacant role that Joe will gradually fill after having unknowingly interrupted its funeral. From this point in the film, Joe is manipulated, or ‘trained’, by Norma to entertain and provide companionship to her. Naturally, Joe also ends up dead within the bounds of Norma’s estate, with this symbol, therefore, foreshadowing the full trajectory of his character. All of this is directly alluded to through Joe’s description of the ‘mixed-up dream’ he has the night of the funeral, imagining ‘an organ [player]’ and the ‘chimp…dancing for pennies’ that he will soon become.\\

\begin{figure} [h!]
    \centering
    
\includegraphics[width=0.75\linewidth]{organ.jpeg}
    \caption{Organ}
    \label{fig:enter-label}
\end{figure}
This naturally brings us to the organ itself, which serves as a physical reminder of the unflattering parts of the new role Joe must play. Included after Joe wakes from his ‘mixed-up dream’, the shot above frames Max’s organ-playing hands as massive and overpowering, as the much-smaller Joe storms in demanding to know why his ‘clothes and things’ were moved to Norma’s house without his say-so. Crucially, Norma then reveals that she ordered this action and that Joe's apartment debts are ‘all taken care of’, hand-waving his attempt at grasping back some control and dignity by proposing it be ‘deduct[ed]...from [his] salary’. This scene reveals the symbolic role the organ plays within Sunset Boulevard, reminding Joe of the shameful and powerless role of the ‘pet monkey’ that he now fills, as well as what he will be ‘dancing’ for.\\

\subsubsection{Allusions}
Finally, we come to allusions, one of the techniques that \underline{Sunset Boulevard} is most famous for. Allusions refer to anytime something from outside the world of the text is referenced, including other texts and real-world people, places, events, etc. Biblical and mythological allusions are commonly found in fiction, but references to something closer to our world can often bring a degree of realism to certain texts, working to strengthen their social commentary. \\

\paragraph{Cinematic Allusions}

\begin{figure} [h!]
    \centering
    \includegraphics[width=0.75\linewidth]{para.jpeg}
    \caption{Paramount Studios}
    \label{fig:enter-label}
\end{figure}

\begin{figure} [h!]
    \centering
    \includegraphics[width=0.75\linewidth]{para 2.jpeg}
    \caption{Sheldrake's Office}
    \label{fig:enter-label}
\end{figure}

Being a film about film, \underline{Sunset Boulevard} naturally contains many allusions to other films. However, Wilder does not shy away from adding an extra level of realism to his references to the film industry. Central to this is the use of the real (and still functional) Paramount Pictures studio to which Joe attempts to sell his clichéd baseball script. Notably, this is the studio that actually released Sunset Boulevard, all of which adds a self-deprecating edge to the satire of the film industry these scenes contain. The scene where the cigar-chomping Paramount executive, Mr Sheldrake, cynically suggests that changing Joe’s film concept to a ‘girls' softball team’ might ‘put in a few numbers’, packs an extra punch due to the use of the real film studio, therefore, showing the effect of this allusion in strengthening the film’s satire.\\

Allusions to specific films are additionally used for humorous purposes and character development. For instance, take Joe’s dry observation that the extravagance of the funeral for Norma’s pet means that he ‘must have been a very important chimp’, perhaps the ‘great-grandson of King Kong’. Here, Joe’s sardonic and witty character is revealed to the audience. Additionally, these kinds of references further place the film firmly in the world of real Hollywood, again working to strengthen the satire it offers of this industry. \\


\paragraph{Literary Allusions}
Similarly, allusions to the world of literature flesh out both the characters and the world of Sunset Boulevard. The most stand-out example of this is the allusion to Charles Dickens’ classic novel Great Expectations. Here, Joe muses that the ‘unhappy look’ of Norma’s house reminds him of ‘Miss Havisham’ from this text. This is a character, who, after being abandoned by her fiance, refuses to change her clothing and lives secluded in a ‘rotting wedding dress’. Havisham directly parallels Norma, being a tragic figure immovably stuck in the past, with Norma's excessive placement of young self-portraits being reminiscent of Havishman’s insistence on keeping her house’s clocks at the exact time she received her letter of marital rejection. Therefore, this comparison to the Dickens character, who engages in a more exaggerated version of Norma’s behaviour, seeks to highlight just how detached Norma is from reality through her attempts to live in the past, implying that what she is doing is just as deluded as refusing to remove a rotting wedding dress. Further, the eventual fate of Miss Havisham within Great Expectations, with her wedding dress catching fire and leaving her as an invalid, foreshadows Norma’s similar descent to invalidity through her madness.

\newpage
\section{Symbols}
\begin{itemize}
    \item \textbf{The Mansion}  
    \begin{itemize}
        \item Symbolises Norma's decaying fame and psychological isolation  
        \item A crumbling structure, 'out of beat with the rest of the world'  
        \item Reflects her detachment from reality and the passage of time  
    \end{itemize}

    \item \textbf{Black-and-White Cinematography}  
    \begin{itemize}
        \item Enhances contrast between Norma's world (dark, stagnant) and Betty's world (bright, full of life)  
        \item Reinforces themes of nostalgia and decay  
    \end{itemize}

    \item \textbf{Spotlights}  
    \begin{itemize}
        \item Symbolise attention, former fame, and Norma’s obsession with being seen  
        \item Used dramatically when Norma sits in deMille’s chair – a ghost-like figure in the light  
    \end{itemize}

    \item \textbf{Gates and Barriers}  
    \begin{itemize}
        \item Represent isolation, exclusion, and entrapment  
        \item Norma uses her home as a prison to keep herself—and Joe—separated from the real world  
    \end{itemize}

    \item \textbf{The Monkey’s Funeral}  
    \begin{itemize}
        \item Symbolises Norma’s lost career and her inability to let go of the past  
        \item Gothic mis-en-scène highlights the absurdity and tragedy of her delusion  
    \end{itemize}

    \item \textbf{Cars}  
    \begin{itemize}
        \item Represent control, freedom, and social mobility  
        \item Joe’s loss of his car mirrors his loss of independence when he becomes trapped in Norma’s world  
    \end{itemize}

    \item \textbf{Close-Ups of Norma’s Face}  
    \begin{itemize}
        \item Reflect her silent film origins and the over-expressiveness required in that era  
        \item Symbolise her inability to adapt and her clinging to outdated forms of identity and performance  
    \end{itemize}

\item \textbf{The Salome Script}  
    \begin{itemize}
        \item Symbolises Norma’s desperate attempt to reclaim her stardom and control her narrative  
        \item Reflects her delusion and obsession with youth, power, and seduction  
        \item A vehicle through which she imagines her triumphant return, though the industry has no real interest in it  
    \end{itemize}

\item \textbf{Paramount Studios}  
    \begin{itemize}
        \item Symbolises the modern, ever-evolving world of Hollywood that has left Norma behind  
        \item Represents both nostalgia and rejection—once her kingdom, now a place where she is merely tolerated  
        \item A cold contrast to the warmth of Norma’s imagined fame  
    \end{itemize}

\item \textbf{Max’s Role as Butler}  
    \begin{itemize}
        \item Max, who was once a respected director and is now reduced to being Norma’s butler,symbolises the destructive nature of Hollywood.
        \item His loyalty to Norma and his willingness to play into her delusions highlight how far people can fall in the industry and how they cling to past identities.
    \end{itemize}

\item \textbf{Max’s Gloves}  
    \begin{itemize}
        \item Symbol of Max’s servitude and control—he handles every aspect of Norma’s life with careful precision  
        \item Suggests a barrier between Max and reality, as he preserves the illusion with delicacy and ritual  
        \item Reinforces his role as both protector and enabler of Norma’s fantasy  
    \end{itemize}

\item \textbf{Norma’s Clothes}  
    \begin{itemize}
        \item Luxurious and extravagant, but styled from a bygone era—symbolising her refusal to move forward  
        \item Represent how Norma costumes herself in a fantasy, masking the reality of her decline  
        \item Reinforce the image of a star who is 'big', but only in her own memory  
    \end{itemize}

\item \textbf{The Pool}  
    \begin{itemize}
        \item The pool, where Joe’s body is found, symbolises the false allure of Hollywood’s glamour.
        \item Initially a symbol of success and luxury, it ultimately becomes a symbol of death and disillusionment.
        \item Joe’s lifeless body floating in the pool illustrates the deadly consequences of chasing hollow dreams.  
    \end{itemize}

\item \textbf{The Film Projector}  
    \begin{itemize}
        \item The film projector represents Norma’s obsession with the past.
        \item As she watches her old films, the projector becomes a symbol of her inability to move on from her glory days, trapping her in a loop of nostalgia and delusion. 
    \end{itemize}

\item \textbf{The Final Scene}  
    \begin{itemize}
        \item The final scene, where Norma believes she is on set filming a movie, symbolises her complete disconnection from reality.
        \item It also serves as a critique of Hollywood, where the lines between reality and illusion often blur, leading to tragic consequences.
    \end{itemize}




    
    
\end{itemize}


\section{Written Exemplars}

\subsection{Sample 1 - High Scoring \underline{Sunset Boulevard} Essay I with Annotations}


ZZEXP Sunset Boulevard Essay I with Annotations ZZEXPEND


Topic:
\textbf{\underline{Sunset Boulevard} reveals the dark side of celebrity and fame. Discuss.}\\

\textit{Introduction}

Exposing the powerlessness of naïve ‘dream[s]’ of fulfilment against the toxicity of celebrity which has ‘seized [their] heart[s]’, (1) Billy Wilder’s quintessentially American (2) film \underline{Sunset Boulevard} (1950) poignantly unveils the dark side of fame that is inextricable from Hollywood’s allure (3). Compounded by the skilful overlaying of a gothic lens, this noir mystery (4) performs as a deep meditation on the brutality of this business, and a sardonic, yet subtly sympathetic, depiction of the victims of its endless cycle of sadism (5). Through his intelligent blend of fiction and meta-cinema (6), Wilder’s iconoclastic opus emerges as a powerful vivisection of the dangerous illusions of stardom.\\

\textit{Annotations:}

(1) Beginning the piece with an embedded quote adds a degree of complexity and immediately highlights to markers that you have an understanding of the text. It also improves the poetic and sophisticated nature of the text, almost as if it is a literature essay, thus separating your piece from other purely analytical responses.

(2) Recognising the context from the outset also showcases one’s understanding of the wider societal issues and themes.

(3) An immediate reference to the topic is necessary. If the topic isn’t explicitly or implicitly addressed throughout the response, it cannot receive a high mark.\\


\textit{Body Paragraphs}

Personifying this complex world as a dystopian nightmare, Wilder satirically contextualises the hidden horrors of Hollywood (1). Immediately referencing a classic trope of film noir, (2) a disturbing amalgamation of fantasy and sobering realism is established through the eerie posthumous narration of this ‘poor dope’, with this non-linear structure exposing this inevitable tragedy from the outset. Heightening this strange sense of unease, (3) the foreboding non-diegetic string music of Franz Waxman’s pertinent score (4) warns the audience of this ensuing horror. Indeed, \underline{Sunset Boulevard} itself can no longer be hopefully associated with fame and luxury; rather, its disconcerting introduction in a gutter with ‘dead leaves’ and ‘scraps of paper’ alludes to its burgeoning corruption, as literally emphasised by the titular finiteness of ‘sunset’. Ominously towering over this decrepit street, the silhouetted palm trees embody personified judges (5) of the ‘distorted’ and ‘blown out’ stories that ‘Hollywood columnists’ weave, as fortified by the intense chiaroscuro lighting (6). Similarly, (7) although swimming pools were an iconic image of glamour, the floating corpse ironically exposes that this perpetual chase of fame is a deadly enterprise. As a poignant reflection of this danger, the toxic relationship between Joe Gillis – our cynical antihero – and Norma Desmond – the stereotypical antagonist of a femme fatale – (8) is inherently doomed. This Gothic synthesis of old and new Hollywood thematically foreshadows their demise, as both are dangerously delusional: one is confined by their obsolescence, the other ensnared in the shallowness of the present. Symbolising the subordination that defines their affair, the loss of Joe’s car is a ‘matter of life and death’ that resolutely deforms him into an item of Norma’s possession, as fame is too often synonymous with prostitution. Subverting the traditional gender stereotypes of their time, Edith Head’s pertinent costumes of shared leopard print – a predominant motif of Norma’s wardrobe, which is also reminiscent of her silent-era glamour – illuminates the oppression that this pursuit of celebrity has manufactured. (9) In a house with no locks, Joe is already irreparably caught in the predatory grip of Norma’s atrophic talons, and thus, by extension, trapped by the noxious allure of Hollywood. Functioning as a provocateur to brazenly ‘bite the hand that feeds him’ (Mayer) (10) Wilder controversially deplores the deceptive façade of fame through his pioneering direction.\\

\textit{Annotations:}

(1) Open with dependent clauses elevates the piece by shifting the pacing. Adverbs are a great opportunity to briefly note tone. Every topic sentence should reference the topic.

(2) Link back to the genre to reflect your wider understanding; most students can only analyse what they see, without considering how this fits into the genre as a whole

(3) It is sophisticated to examine how this makes the audience feel (i.e. unease).

(4) Music is often forgotten in analysis. Films provide opportunities for different types of analysis. Don’t forget the soundtrack/camera angles/acting/outfits.

(5) Develop unique interpretations of symbols. Most students would ignore or neglect the palm trees, simply viewing them as a background item. When you provide unique views and analyses, markers are always impressed.

(6) As mentioned before, having an inventory of metalanguage relating to lighting is greatly sophisticated.

(7) Using linking words like ‘similarly’ underlines to markers that your evidence has been specifically chosen to flow cohesively. Listing dissimilar evidence can weaken the sophistication of your writing

(8) Showing an understanding of character tropes is sophisticated.

(9) Referencing outfits and costume choices can elevate your writing above other pieces, especially when you recognise their purpose as motifs and their significance at the time, such as in relation to gender stereotypes.

(10) The inclusion of one or two external quotes throughout your piece shows a wider contextual understanding and advances your piece to an academic or literary text.\\

Challenging his fervent condemnation of Hollywood’s evil, Wilder manufactures a hopeful alternative that appears to exist untarnished by celebrity (1). Joe and Betty’s relationship functions as a tender antithesis to the toxically phantasmal alliance of Joe and Norma: this mystery is also an ‘Untitled love story’, saving it from merely being a lurid horror (2). Momentarily, Wilder permits Joe to escape his ‘peculiar prison’ of empty fame, returning to the corporeal world of Artie’s New Year’s gathering, an event imbued with a jollity that is crosscut against the solemn falseness of Norma’s decrepit mansion. From her entrance to the film, Betty Schafer embodies the forgotten passion for truth, ‘itching with ambition’ and presenting a virtuous alternative to Hollywood’s amorality. Illuminating her integrity, Schafer’s intelligence protects her from being deterred by the toxicity that is fundamental to this rapacious industry, as revealed when she is repeatedly misnamed by the sexist producer. Establishing a meta critique of the film industry, Wilder’s script sees this ‘sweet kid’ affirming that ‘pictures should say a little something’ rather than capitalising on being ‘glossy’ and slick’. Pertinently, whilst Norma is a victimised by-product of fame’s lethal illusion, Betty is an astute examiner who ‘know[s] all the plots’, acknowledging and embracing the deception at Hollywood’s core: she likes this ‘cardboard’ and ‘hollow’ town ‘better than any street in the world’. Further emphasising her innocence and purity, Betty’s practical attire of simple skirts and sweaters oppose Norma’s characteristic wardrobe of glamour to highlight the dichotomy between these women. Betty refuses to accept that amorality is inextricable from fame; rather, she empowers to write her own story, both literally and figuratively, reminding us that ‘life can be beautiful’. Hence, she is the perfect counterpart to Joe’s cynicism, embodying both his possible redemption as a man and Hollywood’s rehabilitation as an industry. Her ‘cubby hole’ of an office holds more promise and passion than the entire ‘white elephant’ of Norma’s mansion, highlighting the importance of truth. Their parallel storyline offers an escape for Joe, dislocating him from Norma’s persecution and empowering him to pursue his creative aspirations for motives nobler than shallow renown. Indeed, their witty dialogue – ‘Don’t you sometimes hate yourself?’ – and romantic chemistry symbolise this newfound sense of ambition, concurrently contrasting the sombre superficiality of Norma’s exchanges. Unveiling the wavering silver lining of celebrity, their classic Hollywood kiss offers a reprieve of romance. Emerging as perhaps the most admirable character of this satire, Betty’s genuineness empowers the audience to believe in an industry that prioritises creativity and love over dangerous materialism.\\

(1) The structure of your essay is ultimately your decision, but having a paragraph that grapples/disagrees with the topic sentence can add an interesting layer of analysis. It provides an opportunity to explore other paths, characters and ideas, and differentiates your work from other responses.

(2) As mentioned previously, it is important to recognise that the film is not one-dimensional: it can be multiple conflicting genres at once.

Notice how this paragraph is based on comparisons and contrasts (e.g. Artie’s party vs Norma’s mansion, and Betty vs Norma). This is far more effective in conveying understanding and nuance than simply analysing characters and settings individually. It is important to consider their wider role and implications within the film.\\

 

By the catastrophic climax of Wilder’s horrifying creation, it is overwhelmingly transparent that the dark side of fame will ultimately eclipse its idealistic counterpart, as this hope cannot survive unadulterated in the poisonous atmosphere of Hollywood. (1) Although Joe and Betty’s bond strives to pervert this amorality, even their romance eventually succumbs to the deception of this industry, for they cannot ‘get away with it’. Evidencing the moral dilapidation that is endemic to Hollywood, Betty’s chastity is defiled by the contagiousness of Joe’s shallow illusions; although she ignorantly maintains that she has ‘never heard any of this’, her newfound awareness of the toxicity of celebrity is palpable. As ominously revealed in the opening scenes, our protagonist’s flawed ‘long-term contract’ catalyses his spiritual and physical death, for the amoral motives of fame cannot be without consequence. Foreboded by the dramatic orchestral music, William Holden’s (2) purposeful characterisation of Joe’s premortem trance mirrors the enduring ignorance of Hollywood’s victims. Indeed, his final reach to his briefcase is analogous of the inherent desire of these hopefuls to continue to nourish fame’s deceptive image, even when that fatal bullet has been lodged deep in their unexpecting back. Rather than shattering Norma’s agonising delusions, this murder gruesomely precipitates her resolute descent into madness. Deprived of her moment of anagnorisis, as the role of Salome consumes her, Wilder’s use of mirrors as a motif conveys Norma’s dislocation of real and imaginary. Adorned with royal and theatrical music, Wilder’s staging of Norma’s slow-motion descent of her staircase in her bedazzled dress disturbingly typifies an embodiment of Hollywood’s ghastly underbelly erupting to the surface. With each step on her terminal pursuit of fame, this deranged ghost exits her corporeal sphere, her hyperreality finally becoming inescapable, as ‘there’s nothing else’. Surrounded by spectators yet entirely alone, as the camera closes in on her manic face – with her rigid claws and psychotic eyes blaming the culpable audience – she blurs into this phantasmal purgatory, destroying the alleged integrity of this business. As this ‘dream she had clung to so desperately’ finally ‘enfolded her’, Wilder’s writing expresses that celebrity, or the lack thereof, possesses the omnipotence to corrode a famed icon into an insane killer. Ultimately, \underline{Sunset Boulevard} serves as a dire caution to the horrors of Hollywood, as the ‘price’ of this stardom will always be ‘too high’. (3)\\

(1) Ensure that, even if your middle paragraph does provide an alternate view, the final paragraph expresses your overarching conclusion. Complexity and equivocation is important, but so is developing strong opinions.

(2) Referencing actors’ names proves that you have done extensive research.

(3) Referencing actors’ names proves that you have done extensive research.

\newpage

\subsection{Sample 2 - Example Contention and Body Paragraph}


ZZEXP Sunset Boulevard Contention and Body Paragraph ZZEXPEND



Topic
\textbf{'The only villain in \underline{Sunset Boulevard} is Hollywood.' To what extent do you agree?}


\textit{Contention}

Wilder critiques Hollywood’s self-interest, its obsession with youth and its ruthless pursuit for profit over their writers’ artistic intentions, yet he also acknowledges that succumbing to these demands, can be highly immoral and destructive.\\


\textit{Body Paragraph}

Wilder’s critique of Hollywood’s self-interest is exemplified through the decaying mansion, a visual metaphor for the discarded silent era stars. Norma Desmond, ‘the waxworks’, and Max personify how 1940’s Hollywood has relegated these former stars and directors to what resembles a gothic museum. Norma’s boudoir, draped in heavy fabrics and her sculpturally carved ‘gondola’ style bed, is a visual manifestation of her psychological stagnation, as a boat can’t move without water, just like Norma can’t move on from her former status. This is reinforced by Seitz’s camerawork of Norma who is often framed behind curtains, doors, or within mirrored reflections, signifying the prison like barriers she faces from the current era. Her Isotta-Fraschini, also a symbol of her former status, further reinforces stagnation as the car is immobile and a relic of the past. It is ironic that this old car is the only reason she returns to 1940s Hollywood, after receiving a call from Paramount studio, not from deMille whom Norma believed should have called, but Cole who wants her car as a prop ‘for a Bing Crosby film’, rather than Norma Desmond’s delusional role as Salome to be directed by deMille. Living in her delusion of a ‘return’ as ‘Salome’ in her self-written script, where she will play the lead role of a teenage girl at the age of fifty, Norma is exploited by Joe for some quick cash to pay off his car debt, demonstrating how people in the industry of Hollywood succumb to immoral actions in order to keep their dreams alive. In turn, Norma exploits Joe for his script editing and youthful addition to her life, and as a replacement for her dead chimp pet. The use chiaroscuro lighting demonstrates the interplay of light and darkness that underscores this exploitation that occurs. Thus, Wilder signals how both Hollywood of the 1940s, and former stars like Norma of the1920s, who were led to believe that stars never die because 'that’s what makes them a star,” can lead good people to desperate immoral acts, and severe mental illness.\\

\newpage

\subsection{Sample 3 - High Scoring \underline{Sunset Boulevard} Essay II}


ZZEXP Sunset Boulevard Essay II ZZEXPEND



\textbf{'\underline{Sunset Boulevard} shows that the pursuit of fame is a perilous journey. Discuss.”}\\

Billy Wilder’s 1950 noir film, \underline{Sunset Boulevard}, scrutinises the complexity of the Hollywood machine, and through the lens of the grim, sinister underbelly of the famous and fortunate, shatters the glamorous façade fed to the people. By drawing on the personal aberrations of characters and the delusional artefacts which arise with the pursuit of fame, Wilder comments on the fragile nature of stardom, and creates a sense of distrust between the film industry and the ever-elusive entertainment realm. The film reinforces an atmosphere of entrapment silent film actors as those the wisps of her eminence, a B-list script writer from Ohio who aims to keep his distance from the flashy and aspiring master script writer who falls into the Hollywood cookie cutter. By juxtaposing characters helps Wilder reflect the paradox of human nature and portray the foreboding dangers associated with unattainable dreams in the film industry. Through his characters and film genre, Wilder asserts that notoriety is perilous and far less rewarding that what is believed.\\

\underline{Sunset Boulevard}’s Joe Gillis, played by renowned actor William Holden acts as the perfect noir protagonist; cynical, aggrieved and pursued by a mob of sorts. However, Joe is far from a treacherous, undercover cop, black market infiltrating protagonist that Wilder offers other film offer. Unlike most films of the time, Wilder employs meta elements, having famous directors play directors, and legendary, but ageing silent actors such as Buster Keaton, to play themselves. This allows for him to abstractly comment on the brutality of Hollywood, the moral corruption kept so hidden from the public, and sets a sense of realism when inserting Joe Gillis in scenes with these characters. In order to boldly mimic the public of the personalities of fame, Wilder ties Joe’s story line based on a tragic tale as old as time; Icarus. Icarus is a tale of tragedy, a young boy who tried to escape Crete, grew arrogant, flew too close to the sun, and died in the Icarian Sea. Wilder presents Joe Gillis as Icarus, a tale of peril and hubris. By mirroring Icarus’ tragedy through Joe’s story line, with him escaping the sluggish debt collectors 'when things were tough”, becoming Norma Desmond’s gigolo, growing arrogant and eventually a 'has-been” with his lifeless body floating face down in a pool. Wilder showcases that danger associated with catching a 'break” in Hollywood, with tragedy being the likelier outcome. With '[projects] that turned out to be night”, and the misfortunes associated with Joe’s big break into Hollywood, Wilder comments on the complex, manipulative nature of Hollywood, reinforcing the ghastly horrifying truth about his industry for individuals who aren’t old accomplished Caucasian men.\\

While Billy Schader is by no means at fault for allegorical character, Wilder intentionally leaves her story as two dimensional and unfinished, favouring to emphasis the way that he represents most other women in the entertainment industry who do not succeed in these careers. Because of Betty’s duality in representing the unrepeatable aspect of women in Hollywood, and a child of nepotism, who spent 'ten years” on 'dramatic lessons, dictation and dancing” to ascend to stardom. Wilder portrays the frustration of women who want to become screenwriters, directors, and producers. Wilder, through Betty’s technique of Joe to work with her on his script 'bases loaded”, shares his own frustration on making a break in the perilous industry of Hollywood. He embodies his own frustrations and fears through betty, with many parallels between them including the fear of both Betty and Wilder are also shunned from acting due to their looks. Their connections, modern mannerisms, and Betty’s open and vocalised mind is crafted to portray to the audience of the life that may be had should she choose to proceed without acting and instead to fill her life with purpose. Wilder shows Betty’s line to proceed with a fulfilling life not through acting but also leading a fulfilling life, unlike other individuals in the film.\\

Wilder continues to explore this in the hazardous stardom hierarchy that many an actress is victim to, and which does not reflect the innocence and pleasantry of 1950’s woman. Norma Desmond, played by Gloria Swanson, encapsulates this idea. The narrator sets her up with golden names with silent films with 'just her” [face!]. Through Norma’s status as the 'greatest star” who 'made” millions at 'the lovely age of 17”, Wilder throws us in with Hollywood’s pessimistic, backward relationship in Hollywood, placing women stars far more than a role or a 'vehicle”, he dehumanises and contextualised Norma as a silent film actress who was unable to 'get to the talkies”. His account is laced with '100,000 dollar fan letters” which chases her youth, willing as far to kill and obtain Joe after his 'unused” past self, yet self-deluded as she keeps the paper headlines on clips of camera technologies and lighting, including a scene in which an old crewmate of Norma turns the spotlight onto her, before being told to 'return that to where it belongs”. The spotlight shinning down on Norma, paired with a high angle to make her look small, helps symbolise the beginning of the end. It shows Norma succumbing to years of pressure and pressure that was once enveloped in the warmth of her former years, but now has shattered the Hollywood machine. In the end, Wilder chalks up Norma as Narcissus. He shows a woman damaged by an industry that once enveloped her, a woman so afflicted by vanity, egotism, and self-delusion that she is a shell of her former self. Through a deluded Norma, playing an unlike character, and a dramatic story line, with an open end, Wilder caveats his audience of the precariousness and toxicity of the film industry.\\

Although \underline{Sunset Boulevard} is a film centred in the noir and macabre, highlighting the underworld of Hollywood, the ominous and disturbing truths of the famous and fortunate, Wilder portrays contrasting characters, with hopeful endings for some, and deserving, melancholic endings for others. These disparities allow for wilder to abstractly comment on his industry, portraying the notion that no true member of Hollywood is ever as innocent, and despite the dubious actions of his characters, draws a sense of sympathy for those who exploit others. Above all, Wilder allows for conflicting morals to take centre stage, often prompting his audience to takes sides, to showcase the complexity and multidimensional paradox of human nature. Through tragedy, and familiar tales, Wilder rewards courage, resilience and individuality, while disapproving of those embroidered into the Hollywood clothe, leading them to untimely outcomes, and others to fulfilling lives.\\
\newpage

\newpage
\subsection{Sample 4 - High-Scoring \underline{Sunset Boulevard} Essay III}


ZZEXP Sunset Boulevard Essay III ZZEXPEND


\textbf{Topic}\\
How does Sunset Boulevard explore the corrupting effects of Hollywood and the desperation it causes?\\

\textbf{Response}\\
Although a grimly cynical film, Billy Wilder’s \underline{Sunset Boulevard} presents a critical portrayal of the effects selfishness and desperation can have on characters. Through its backstudio narrative, the film explores the corrosive nature of a heartless system that mistreats its subjects, and the desperate means necessary to survive in such an industry. By the end of the film, Wilder begins to envision an alternative but concludes by exposing the entrapping nature of a corrupt system.\\

The Hollywood industry is constructed as a corrupt and heartless environment that acts as a catalyst for characters’ desperation. In the opening shot of the film a stark juxtaposition is formed between the glamour and corruption of Hollywood, the name connoting glamour and wealth, while the placement of the sign in the gutter indicating the opposite. Wilder presents Hollywood as corrupt because it promises fame and fortune but cannot deliver to the vast majority of people. From the outset of the film, the struggling screenwriter Joe experiences the cutthroat nature of Hollywood. SitÝng on his bed in a robe, Joe is ‘grinding out’ scripts in a desperate attempt to recover his career. Wilder highlights the dire nature ofhis situation, and how the promises Hollywood once made to the previously successful script writer are no longer being delivered upon. When debt collectors come to repossess his car, Joe’s perilous attempt to cling on to it serves as a parallel to his instability and financial insecurity. Joe's statement that ‘[he needs] about two hundred and ninety dollars and [he needs] it real quick’ works to develop this idea further, highlighting his desperation due to the importance of his car that is symbolic of success and wealth. Wilder exposes Hollywood’s corruption by exploring the false promises of fame and fortune that are never delivered upon.\\

In order to cope with the mistreatment of Hollywood and the ramifications of rejection,
characters seek to survive through desperate means and selfish behaviours. Prior to his
arrival at Norma’s mansion, Joe has been corrupted by Hollywood, making him utterly
desperate and vulnerable. Being rejected by the industry forces him into desperate and
unethical behaviours. He begins ‘rehashing’ old scripts, abandoning artistic integrity in his
desperation to figure out how to succeed in the industry. Wilder examines the effects of
desperation on relationships when Joe begins to exploit his contacts, using them as ‘assets’.
Gillis' deceptive nature is revealed through his exploitation of others, which is analogous to
Hollywood's commodification of its participants. Joe refers to the fact that Sheldrake had
‘always liked [him]’, displaying his intent to leverage personal relationships for his
advantage. Wilder utilises Norma’s house as a reflection of Norma herself, the place being a
‘white elephant’, and a decrepit shrine to the past. Norma’s construction of this shrine is
her way of dealing with the rejection of an industry that previously gave her validation and
identity. Joe and Norma both create a narrative that suits their own personal interests,
exploiting and manipulating one another. Norma is able to exploit Joe’s vulnerability, his
need for money. She buys Joe’s companionship and loyalty, buying him a ‘vicuna’ coat, and
seeing him as a young man that can make her feel young and attractive again. Joe exploits
Norma’s vanity and instability to get what he believes will be easy money to pay off his
debts and reestablish himself in the industry. Gillis is driven to deceive Norma by the
corrupt Hollywood system that causes his desperation for financial aid. Both characters are led to manipulate the other by their desire for personal gain, acting in self-interest due to
their desperation. This behaviour is analogous to the shallow values of Hollywood, living in
an alternate reality and manipulating each other for their own gain.\\

Whilst Wilder demonstrates an alternative approach to selfishly delusional living, the tragic
ending of the film ultimately reveals that there is no real escape from the corrupt and
Hollywood system for those that are too deeply entrenched. Wilder pictures Betty as a
counterpart to Norma, she is an earlier version of Norma, one that is ‘like all us writers
when we first hit Hollywood’ and has not yet been corrupted by the system. Betty
represents authenticity and artistic integrity, adamant in her belief that ‘film’s should say
something’. She is aware of the hollowness of Hollywood referring to it as ‘all cardboard’.
Wilder presents her as the antithesis to Norma, a beacon of hope that is realistic where
Norma is delusional. Wilder offers hope through Betty, an alternative to living in an
alternate reality. Furthermore, her working relationship with Joe is more collaborative and
transparent. She wants to ‘do it together’ and does not attempt to manipulate and leverage
his talent for her own profit. Wilder contrasts Betty’s behaviour to the other characters, as
she is willing to work with Joe rather than exploit him. However, although her sincerity
connects with Gillis, Wilder’s portrayal of Joe’s subsequent behaviours undermines any
optimism proffered by his relationship with Chafer. While Joe reconciles with reality and
attempts to leave Hollywood, he is ultimately unable to extricate himself from Norma’s
grasp, a representation of Hollywood’s hold over him. For Norma, Joe’s departure would
mean the end of her alternate reality. Joe’s death as he leaves Norma’s entrapment
underscores all characters’ inability to escape the corrupt Hollywood system. Even Joe, the
character that finally realised the effect Hollywood has had on him, is not allowed to escape,
with this realisation instead leading to his eventual demise. In line with the prevailing
cynicism throughout the movie, Joe meets his demise just as he attempts to break free.
Following Joe's death, Norma's descent into madness is complete. She retreats further into
her fantasy world, believing that she is starring in a movie scene rather than facing the grim
reality of her actions. The final words of Norma and of Hollywood itself are ‘no one leaves a
star’. Her last words are a testament to the way the system entraps and manipulates its
participants.\\

\underline{Sunset Boulevard} is a scathing criticism of the cruelty of Hollywood and the corrupting
effects of a mistreating system that causes desperation and selfishness. Relationships based
on manipulation and self-interest in addition to characters exploiting others emphasises the
film’s intention the expose the deceitful nature of Hollywood. Furthermore, while some
hope is generated, the corrupting nature of Hollywood overshadows any positive features,
overall revealing the corrosive effects that selfishness and desperation have on characters
within the industry.\\

\newpage

\subsection{Sample 5 - High Scoring \underline{Sunset Boulevard} Essay IV}


ZZEXP Sunset Boulevard Essay IV (VCAA 2023 Topic) ZZEXPEND


\begin{quote}
\textit{'Sunset Boulevard shows that the pursuit of fame is a perilous journey.' Discuss. (VCAA Exam 2023)}
\end{quote}

Released towards the end of 1950, Sunset Boulevard, Billy Wilder’s incisive expose of the Hollywood industry and American society more broadly was met with scathing reviews by critics who felt Wilder had maligned the all-American pursuit of fame. Indeed, Wilder criticises in his film how the pursuit of fame often creates vulnerability to the exploitation of powerful figures, and more concerningly, the loss of oneself to discontent and even madness. Despite this, Wilder endorses those in the industry who pursue meaning over fame, and illustrates how such a path can allow one to avoid the perils of the industry. Ultimately, Wilder suggests that while Hollywood is a dangerous place, it can be navigated by those who seek to produce good work.

In the movie, the pursuit of fame leaves characters open to exploitation and even endangers their lives. Wilder reminds his audience that predatory characters wait around every corner to profit off the naivety of ‘the talent’. Joe’s callous agent tells him he should be grateful to be short on money, since “the greatest things were written on an empty stomach.” Despite professing to be Joe’s “friend”, Morino represents the most exploitative and profit-driven parts of deep Hollywood, always lurking in wait for the next “10 percent” cut. Wilder symbolically illustrates the dirtiness of the industry right from the opening of the film, when he depicts the titular street name ‘Sunset Boulevard’ written not atop a street sign, but instead in the gutter. This placement represents the lowlife aspects of Hollywood, with the fallen leaves in particular subtly hinting at the industry’s fallen and dead stars, while the high camera angle looking down on it tellingly indicates Wilder’s disgust towards the industry. The potentially fatal outcomes of chasing fame are reiterated in the next scene when the audience is presented with Joe’s dead body floating face down in a pool, while photographers crowd trying to get a good picture. Speaking from beyond the grave, Joe muses with dry empathy that “the poor dope…always wanted a pool.” As the cameras flash, he adds that he “got himself a pool – only the price turned out to be a little high.” This quote epitomises from the outset the ultimate peril of seeking acclaim: death. Indeed, it was only through his murder that Joe achieved not only the “pool”, but also the success and fame this represents in Hollywood. Joe’s haunting narration keeps the thread of death running through every scene, reminding the audience that a reckless pursuit of fame can have potentially fatal consequences.

Perhaps even more sadly, Wilder disappointedly comments on the ability for the Hollywood industry to produce cynicism, discontent, and even madness. In his search for recognition, to “get [his name] up there”, Joe loses his desire to do anything other than chase cash. Although there is an indication Joe has “some talent”, Joe insists “that was last year; this year [he’s] trying to make a living”. As much as he feels hard-done-by in the industry, Joe here reveals that he himself has chosen to give up writing what is “true” in the name of trying to make money. Ironically, this leads Joe to a cycle of producing work that was either “too original” or “not original enough”, but in any case “it didn’t sell”. Through his cynical quips through the film and the lugubrious, jaded tone of Joe’s voiceover, Wilder conveys how for artists who approach Hollywood with the aim of making it big, only burnout and disappointment await. Likewise, Norma loses herself in the pursuit of public adoration. Despite beginning “as a lovely little girl of 17” full of “courage and wit and heart”, Norma became a “terror to work with” after years of “press agents working overtime” on her. Norma’s inability to separate her stage self from her real self is established early on through the thick makeup accentuating her lips and eyes, reminiscent of the silent film style. This makeup and her exaggerated facial expressions become only more extreme as the film continues, symbolising Norma’s mental decline because of her inability to adapt. The wish for fame and “a return” is a manifestation of Norma’s profound desire for attention, represented by attempts to entertain Joe with the “Norma Desmond Follies”, painfully incongruent with her ageing self. Watching Desmond’s cosmetic procedures and her increasing insanity is like watching a trainwreck in slow motion, and by the end of the film when Norma famously cries she’s “ready for [her] close-up”, the audience is left with a profound sense of horror for how the desire for the adoration of “those wonderful people out there in the dark” can lead actors to totally lose touch with reality. The parallel arcs of Norma and Joe both are symbolic of Wilder’s broader message that a desire for fame leads people to lose their sanity and meaning.

Although Wilder’s depiction of the dangers of chasing fame is unforgiving, the director also suggests that those who seek meaning in Hollywood are able to withstand the pressures those only seeking material reward buckle under. Even when faced with Joe’s cynicism, asking “who wants true [stories]? Who wants moving?” Betty maintains she “[wants] to write” rather than be a reader “all [her] life”. Betty’s ambition is not to achieve fame or wealth, but to develop herself, create a career, and ultimately make movies which “say a little something”. Betty is also able to bounce back from failure. As insistent as Joe is that Betty must “weep for those lost close-ups” after she was rejected as an actress, she responds that the whole ordeal “taught [her] a little sense”. Because she does not seek publicity or wealth, Betty is able to escape the disappointments Joe and Norma can not. The extras at Artie’s party, too, are positive, despite being “writers without a job, composers without a publisher”. When Joe enters their room on New Year’s, they are singing happily. It is not that their experience of Hollywood is different: indeed, they sing about how “Hollywood for [them] ain’t been so good”, and Artie jokingly breaks the fourth wall by commenting he can only afford “three drinks per extra”, a reminder of the financial limitations of Hollywood, particular for those on its lowest rung. However, the extras maintain a positive attitude, using humour to cope, and view this hardship more as a liveable reality of Hollywood rather than as a personal affront. Even DeMille’s extras, despite being old and conventionally unattractive, outlast Norma in Hollywood because they are willing to work in the background, creating the scene, rather than starring in it. Thus, Wilder illustrates that while Hollywood is rife with danger, this peril is avoidable for those who choose to work on movies for the love of it, rather than to achieve money or fame.

In Sunset Boulevard, Wilder pulls back the curtain on the multitude of personalities who seek to work in the movie business. Those who seek fame ultimately lose either their physical lives, or their mental stability. However, those who choose meaning and passion over flimsy materialism are able to survive the perils of Hollywood. Thus, though the original critics may not have seen it, Wilder ultimately leaves his audience with the American Dream: that if you do what you love, one day you might just make it.

\newpage

\section{Sunset Boulevard Script}
The full script is available online:

\href{https://www.dailyscript.com/scripts/sunset_bld_3_21_49.html}{Read the Sunset Boulevard script at Daily Script}



\newpage

\cleardoublepage

\clearpage
\thispagestyle{empty}
\hfill
\clearpage


\part{\underline{Rainbow's End}}
\vspace{8cm}
\newpage

\section{Practice Analytical Text Response Topics}

\newpage
\addcontentsline{toc}{section}{Practice Analytical Text Response Topics}

1. '\underline{Rainbow's End} is a bitter condemnation of colonisation.' Discuss.\\

2. 'A real home is where there are people looking out for each other' How does \underline{Rainbow's End}
explore the importance of home?\\

3. "An empty place for Papa Dear." The men, in Harrison's play are marginalised and have little impact on the play's events. To what extent do you agree?\\

4. The main hardships characters face in the play are caused by gender inequality. To what extent do you agree?\\

5. 'The dream sequences in \underline{Rainbow's End} are the characters' only escape from the depressing
reality of life.' Discuss.\\

6. 'Gladys and Dolly are torn between the world of white Australians and their Indigenous Australian heritage.' Discuss.\\

7. How does Harrison use setting to reflect characters' experiences in \underline{Rainbow's End}?\\

8. “Your mum chops wood! With an axe and all!” How is strength displayed throughout the play?\\

9. ‘Hope is a powerful tool for change.’ To what extent does this apply to \underline{Rainbow’s End}?\\

10. “Que sera sera.” \underline{Rainbow’s End} argues that people should accept their fates and not strive for better futures. To what extent do you agree?\\

11. “I don’t want my mother to be served last in the butcher’s.” How does \underline{Rainbow’s End} explore how individuals can survive and overcome their experiences of discrimination?\\

12. ‘\underline{Rainbow’s End} suggests that family relationships can be destructive.’ To what extent do you agree?\\

13. ‘\underline{Rainbow’s End} tells a tragic tale about the loss of culture and identity.’ To what extent do you agree?\\

14. ‘\underline{Rainbow’s End} is a bitter condemnation of colonisation.’ Discuss.\\

15. \underline{Rainbow’s End} highlights the futility of hoping for a better future in a prejudiced world. To what extent do you agree?\\

16. ‘The dream sequences in \underline{Rainbow’s End} are the characters’ only escape from the depressing reality of life.’ Discuss.\\

17. “A real home is where there are people looking out for each other.” How does \underline{Rainbow’s End} explore the importance of home?\\

18. ‘Characters in \underline{Rainbow’s End} are torn between the world of white Australians and their Indigenous Australian heritage.’ Discuss.\\

19. “I want what any mother, black, white or brindle, wants for her daughter.” Harrison presents mother–daughter relationships as complex. Discuss.\\

20. “An empty place for Papa Dear.” The men in Harrison’s play are marginalised and have little impact on the play’s events. To what extent do you agree?\\

21. The main hardships characters face in the play are caused by gender inequality. To what extent do you agree?\\

22. Ultimately, Harrison presents an uplifting message about individuals’ capacity to overcome adversity. Discuss.\\

23. How does Harrison employ dramatic techniques to highlight the struggle and despair faced by Indigenous Australians?\\

24. How does Harrison utilise setting to demonstrate the hardships that Indigenous Australians encounter?\\

25. To what extent does \underline{Rainbow’s End} show that optimism is justified in times of injustice?\\

26. “We demand the right to control our own destiny.” \underline{Rainbow’s End} suggests that vulnerable people can control their own destinies. Do you agree?\\

27. The characters in \underline{Rainbow’s End} share similar experiences but do not see things the same way. To what extent do you agree?\\

28. In \underline{Rainbow’s End} the women discover the power of having a voice. Discuss.\\

29. Relationships can be a catalyst for personal transformation and empowerment. How does Harrison explore the influence of relationships in her play?\\

30. “Are you saying you’d rather live in a humpy by the river? When I’m promising you the world?” In \underline{Rainbow’s End}, attempts to be helpful are often unsuccessful. Discuss.\\

31. Oppression and other negative human experiences can have lasting impacts on an individual’s life and outlook. How is this idea explored in \underline{Rainbow’s End}?\\

32. How does \underline{Rainbow’s End} explore ideas about the ‘other’?\\

33. Ultimately, Jane Harrison’s play \underline{Rainbow’s End} is a study in what it means to be human. Discuss.\\

34. ‘Cultural connection is the greatest resistance in \underline{Rainbow’s End}.’ Discuss.\\

35. In what ways does Harrison explore how tensions between groups shape individual and collective experiences?\\

36. How does Harrison use dramatic techniques to inspire audiences to sympathise with the characters in \underline{Rainbow’s End}?\\

37. How does Harrison’s play explore the way ethnocentric Australia creates social and psychological experiences for individuals?\\

38. “It’ll be all right.” How does Harrison explore the resilience of Indigenous Australians in \underline{Rainbow’s End}?\\

39. Explore the effectiveness of the ending of the play in positioning the audience to respond strongly to Harrison’s themes and ideas.\\

40. Harrison reveals that cultural values are significant in elucidating Aboriginal identity. How does she do this in \underline{Rainbow’s End}?\\

41. “A real home is where there are people looking out for each other.” How does \underline{Rainbow’s End} reveal the importance of home?\\

42. Social ostracisation is presented by Harrison as a catalyst for power and change. To what extent do you agree with this statement?\\

43. Errol’s attempt to be “a knight in shining armour” highlights the complexity of being a man in a patriarchal society. Discuss.\\

44. Compare and contrast the character arcs of all the main characters in \underline{Rainbow’s End}, highlighting their individual growth and the themes of resilience and self-discovery throughout the play.\\

45. Explore how voice within the play reflects the values and attitudes of the context within which it is set.\\

46. Discuss the stylistic choices made by Harrison to construct her play and to present its key themes.\\

47. To what extent does \underline{Rainbow’s End} reveal both the emotional and intellectual responses provoked by discovery?\\

48. Harrison uses voice to encourage her audience to empathise with others outside their own context. Do you agree?\\

49. \underline{Rainbow’s End} tells a tragic tale about the loss of culture and identity. Do you agree?\\

50. Examine the complexities of Dolly’s character, exploring her struggles with identity and the impact of her cultural heritage on her relationships.\\

51. The struggle for justice is worth it, even if it takes a lifetime. Discuss the themes of justice and equality in \underline{Rainbow’s End} and how the characters strive to overcome the barriers they face.\\

52. \underline{Rainbow’s End} explores the instinctive human need to belong through relationships and social status. Discuss.\\

53. How does Harrison encourage readers to challenge their assumptions of the world through her depiction of individual human experiences?\\

54. In what ways does Harrison present the women living on The Flats as unsung heroes?\\

55. Gladys and Dolly are torn between the world of white Australians and their Indigenous Australian heritage. Discuss.\\

56. Despite the damaging history and continuing trying circumstances in \underline{Rainbow’s End}, there is still hope. Discuss.\\

57. \underline{Rainbow’s End} is pivotal in helping audiences come to new insights about belonging. To what extent do you agree?\\

58. “They forced us to leave … Our home.” \underline{Rainbow’s End} explores the cruelty that humans can inflict on each other, no matter the time or place.\\

59. \underline{Rainbow’s End} is a tribute to the enduring power of love and patience. How is this revealed in the play?\\

60. Harrison demonstrates that, although it is difficult, individual independence and familial duty can be balanced. Discuss in relation to \underline{Rainbow’s End}.\\

61. \underline{Rainbow’s End} explores discovery as a transformative journey for individuals. Discuss.\\

62. It is through storytelling that our ideas about the individual and collective human experience are challenged. Discuss.\\

63. Belonging is as much about choice as it is about circumstances that an individual cannot control. How does Harrison’s play explore this idea?\\

64. \underline{Rainbow’s End} explores the life-altering discoveries of self, others, and the world. Discuss.\\

65. Discuss the stylistic choices made by Harrison to construct her play and to present its key themes.\\

66. Explore how voice within the play reflects the values and attitudes of the context within which it is set.\\

67. \underline{Rainbow’s End} reveals that the values of individuals are shaped by their circumstances. Discuss.\\

68. Explore the significance of language and communication in the play, considering how it empowers or restricts the characters’ experiences and relationships.\\

69. Harrison demonstrates that individual growth cannot occur if one does not challenge their pre-existing knowledge and understanding of their world. Do you agree?\\

70. Identify how Harrison uses language features to invite empathy for particular individuals in her play.\\

71. Analyse the connection between land, heritage, and cultural identity in \underline{Rainbow’s End}, focusing on how the characters’ relationships with the land shape their sense of self.\\

72. \underline{Rainbow’s End} is a bitter condemnation of colonisation. Discuss.\\

73. How does Harrison lead audiences to question the assumptions and values within Australian society?\\

74. Analyse how language patterns are used to create meaning in the play.\\

75. Examine the concept of family in \underline{Rainbow’s End}, exploring how it is defined and redefined by the characters throughout the play.\\

76. Change can be both scary and liberating. Discuss the theme of change in the play, considering the characters’ fears, desires, and the potential for personal growth.\\

77. Discuss the development of Gladys in the play, exploring her transformation from a submissive housewife to a confident and independent woman.\\

78. Explore the role of dreams and aspirations in the play and how they impact the characters’ choices and relationships.\\

79. The play suggests that characters have little control over their lives. Discuss.\\

80. It is the connection between storytelling and human experience that truly shapes an individual. To what extent does the play support this statement?\\

81. Society and perspectives may change over time, but the human condition remains constant. To what extent does \underline{Rainbow’s End} explore this statement?\\

82. How does Harrison use language or structural innovations to unsettle her audience?\\

83. To what extent are Harrison’s characters affected by injustices of the past?\\

84. How does the play \underline{Rainbow’s End} explore the experiences of displacement and assimilation for Aboriginal Australians?\\

85. “Just not going to happen. Not in my lifetime.” To what extent does Harrison demonstrate that change is possible?\\

86. \underline{Rainbow’s End} demonstrates the importance of family and community in the lives of Indigenous Australians. Discuss.\\

87. Discuss the role of family in \underline{Rainbow’s End}. How does family help the characters navigate their challenges, and how does it contribute to their sense of identity?\\

88. Analyse how different attitudes or perspectives are communicated through the construction of voices in \underline{Rainbow’s End}.\\

89. Discuss the significance of the play’s setting in shaping the experiences of the characters. How does the physical environment reflect their emotional states and their sense of belonging?\\

90. \underline{Rainbow’s End} illustrates the racism and discrimination faced by generations of Aboriginal people in their own Australian society. Discuss.\\

91. “Why do we have to prove we can live like whitefellas, before we get the same opportunities?” How does \underline{Rainbow’s End} explore injustice?\\

92. How does the play use language and symbolism to convey its themes? What techniques does the author use to engage the audience and create a sense of empathy for the characters?\\

93. Analyse the ending of the play. How does it resolve the conflicts that have been established, and what message does it leave with the audience?\\

94. Compare and contrast the experiences of different characters in \underline{Rainbow’s End}. How do their backgrounds and personalities affect the way they respond to the challenges they face?\\

95. How does \underline{Rainbow’s End} explore the theme of cultural heritage and the importance of maintaining cultural identity?\\

96. Cultural identity and heritage are key themes in \underline{Rainbow’s End}. Discuss their significance for the characters.\\

97. In her play, Jane Harrison exposes the effects of historical events on Indigenous Australians. Discuss.\\

98. \underline{Rainbow’s End} is essentially a study of assimilation versus cultural preservation. To what extent do you agree with this evaluation of the play?\\

99. Analyse gender roles and expectations as depicted in the play, and explore what views and values Harrison endorses through the attitudes and actions of her characters.\\

100. Harrison suggests that dreams and aspirations are especially important in circumstances when individuals are confronted with obstacles beyond their control. Discuss.\\

101. Explore the ways in which \underline{Rainbow’s End} is a story about displacement and the search for belonging.\\

102. Harrison expounds the significance of education and knowledge in her play. How does she do this, and what are their effects on the characters?\\

103. Jane Harrison’s play illustrates that there are many catalysts which lead to an individual’s shift in perceptions. Discuss.\\

104. Harrison uses voice to encourage her audience to empathise with others outside their own context. Do you agree?\\

105. How does Harrison lead audiences to question the assumptions and values within Australian society?\\

106. Discuss the effectiveness of the ending of the play to position the audience to respond strongly to the themes or ideas presented by Harrison.\\

107. Jane Harrison’s play \underline{Rainbow’s End} accentuates the correlation between individual and collective experiences. Discuss.\\

108. Analyse how language patterns are used to create meaning in the play.\\

109. “I’m not an interloper – I belong here – this is my land!” Analyse the tensions exposed by Jane Harrison in relation to the cultural conflicts experienced by various characters in \underline{Rainbow’s End}.\\

110. Explore the shifting power dynamics among the characters. How do characters come to understand and appreciate more about both themselves and others, and what is the effect of these understandings?\\

111. ‘The past is always present.’ Discuss how the past influences the characters’ lives and decisions in \underline{Rainbow’s End}.\\

112. Analyse the importance of friendship and cross-cultural relationships in \underline{Rainbow’s End}, exploring how these connections challenge stereotypes and foster understanding.\\

113. Explore how Harrison has crafted the voices of her characters to reveal their inner or hidden conflicts.\\

114. ‘The dream sequences in \underline{Rainbow’s End} are the characters’ only escape from the depressing reality of life.’ Discuss.\\

115. To what extent does \underline{Rainbow’s End} show that the characters’ optimism is justified?\\

116. “We demand the right to control our own destiny.” \underline{Rainbow’s End} suggests that vulnerable people can control their own destinies. Do you agree?\\

117. How does Harrison represent the emotions arising from human experiences through the features of drama?\\




118. Harrison suggests that for many white Australians, indigenous lives should be covered up, 'like a band-aid over a sore'. Discuss.\\

119. 'The play opens with the lines of the song Que Sera, Sera: "whatever will be, will be", as a foreshadowing of Harrison's main themes and ideas.' Discuss.\\

120. 'Rainbow's End exposes the racist underbelly of white Australia.' To what extent do you agree?\\

121. Harrison's strongest message is to do with the assumption that 'your world is better'. Do you agree?\\

122. 'In Rainbow's End, the characters receive an education as to the possibility of being able to make a breakthrough.' Discuss.\\

123. To what extent does Harrison use her text to criticise modern society?\\

124. How do the characters in Rainbow's End learn what they need to know?\\

125. What narrative devices does Harrison use to further her messages about the social evil of racism?\\

126. 'The great unfairness of modern society for indigenous Australians is to do with choices.' To what extent to do you agree?\\

127. How does Harrison use dramatic techniques to further her main messages?\\

128. Discuss the significance of prejudice in Rainbow's End.\\

129. 'Harrison points an accusing finger at the weaknesses of mainstream Australia.' Do you agree?\\

130. Errol's big problem was that 'he thought he was some kind of knight in shining armour'. Why wasn't he a hero for Dolly?\\

131. 'Harrison highlights that increased progress in society is built upon more effective communication'. Discuss.\\

132. Jane Harrison's play mourns the generations of disadvantage that Aboriginal Australians have endured. Do you agree?\\

133. Rainbow's End celebrates the small victories. Discuss\\

134. Rainbow's End works as a history, a voice for Indigenous Australians and as an educational guide for others. Discuss.\\

135. How does Jane Harrison create a sense of two worlds in her play?\\

136. 'We demand the right to control our own destiny.' In reality, the characters in Rainbow's End have no control over their lives.' Discuss.\\

137. Gladys lives in both worlds more than any other character at the beginning of the play. To what extent is this true by the conclusion of the play?\\

138. Initially, Errol sells encyclopedias but doesn't fully understand how the world works. To what extent do you agree?\\

139. Harrison's message is that voice and knowledge are the keys to fighting oppression. To what extent do you agree?\\

140. To what extent were white Australians aware of the living conditions of Aboriginals?\\

141. In what way do we see the strong sense of connection within the Aboriginal community?\\

142. The play is filled with tragedy but is overwhelmingly positive and hopeful. Discuss.\\

143. How does Rainbow's End depict family dynamics and their influence on individuality?\\

144. Rainbow's End depicts the impact of racism on Indigenous characters' self-perception and identity formation. Discuss\\

145. How does Rainbow's End portray the influence of culture on our sense of belonging?\\

146. The dreams the characters hold create difficulties in the characters' lives. Discuss.\\

147. Discuss the importance of education in Rainbow's End and its impact on characters' lives.\\

148. The characters in Rainbow's End are driven by a need for power. Discuss.\\

149. How does the staging of Rainbow's End explore the relationships between the characters and their environment?\\

150. "A real home is where the people are looking after each other." How does the play explore this idea?\\

151. "There is more hope than despair for the residents of The Flats." Do you agree?\\

152. "All the women in Rainbow's End are leaders in their own right." To what extent do you agree?\\

153. To what extent is Rainbow's End about the impact of racism?\\

154. Rainbow's End demonstrates the difficulty of changing your circumstances. Discuss.\\

155. How does Harrison use symbols to explore the changing dynamics at Rainbow's End?\\

156. In Rainbow's End, family is more important than anything. Do you agree?\\

157. The most remarkable thing about the Dear family is their lack of resentment and bitterness. Discuss.\\

158. Rainbow's End is based on the problem of Aboriginal housing, but the play shows relationships are more important than circumstances. Do you agree?\\

159. While the Dear women are put in humiliating situations, they are not crushed or ashamed. Do you agree?\\

160. How does Rainbow's End present a critical view of authority figures?\\

161. "You're white. I'm Aboriginal. Or haven't you noticed." To what extent does the play suggest differences can be overcome?\\

162. 'Discrimination and poverty are the driving forces in this play.' Discuss\\

163. Although audiences see few white characters, the impact of colonialism is pervasive in Rainbow's End. Discuss.\\

164. Despite the tragedies and hardships that befall the characters, Rainbow's End is ultimately more optimistic than it is pessimistic about the futures of the characters. Discuss.\\

165. How does Harrison explore the idea of family in Rainbow's End?\\

166. To what extent does Rainbow's End offer audiences a sense of hope and reconciliation?\\

167. In Rainbow's End discrimination has profound effects on individuals and cultural identities. Discuss.\\

168. Rainbow's End suggests that there are many types of knowledge. Discuss.\\

169. Rainbow's End shows the importance of family and where you belong to the people of the river flats. Do you agree?\\

170. Although Rainbow's End presents great loss and suffering, it offers an ultimately positive message. Do you agree?\\

171. The lives of the three generations of the Dear family show the enduring effects of discrimination. Discuss\\

172. Aboriginal activism and struggle shapes the story of Rainbow's End. Discuss\\

173. 'You get knocked down when you get too clever.' The characters in Rainbow's End are determined to improve their lives. Discuss.\\

174. 'Sometimes you have to move on. Leave things behind... Even things you love.' How does Rainbow's End suggest that change, although difficult, is inevitable?\\

175. Rainbow's End shows that it is impossible to overcome adversity. Do you agree?\\

176. 'We demand to be the equal of anyone.' How does Rainbow's End explore resistance to oppression?\\

177. The characters in Rainbow's End respond to trauma in different ways. Discuss\\

178. Rainbow's End suggests that it is ultimately not possible for the Indigenous characters to get what they want out of life. Discuss.\\


\subsection{Topics from Previous Exams}
\begin{itemize}
\item To what extent does Rainbow’s End show that the characters’ optimism is justified? 
\item “We demand the right to control our own destiny.”\\
‘Rainbow’s End suggests that vulnerable people can control their own destinies.’\\
Do you agree?
\item 'Least here we do things our way-no one breathin' down our necks.'\\
Home is central to identify in \underline{Rainbow's End}.\\
To what extend do you agree?
\item In \underline{Rainbow's End}, attempts to be helpful are often uncessful.\\ Discuss.


\end{itemize}
\newpage

\section{Introduction and Context}

\subsection{Author Background}
Jane Harrison, a descendant of the Muruwari people, is an acclaimed playwright, novelist, and critic. As a child, she “always loved books” but was disheartened by the absence of Aboriginal representation in Australian literature. In writing \underline{Rainbow's End}, Harrison sought to give voice to those experiences that had long been silenced or marginalised. Drawing upon oral histories from Aboriginal Elders, she captures the resilience and humour of families who lived on “The Flats” and later at Rumbalara—the First Nations housing estate in Mooroopna.\\

Harrison’s work strongly condemns the substandard housing and systemic neglect endured by First Nations families during the 1950s. Through her writing, she positions herself as a conduit for truth-telling, giving dramatic expression to both the hardships and triumphs of Aboriginal life. Her aim, as she has stated, is to use fiction as “an important way for First Nations people to have a voice.”\\

\subsubsection{Analysis}
Harrison’s authorial perspective intertwines personal heritage and political conviction. By dramatizing the everyday realities of Aboriginal families, she reclaims history from the margins and rehumanises voices often reduced to statistics or stereotypes. Her compassionate yet incisive storytelling challenges audiences to confront both the past and its lingering presence in modern Australia.\\


\subsection{Historical Context}
The colonisation of Australia, beginning with the arrival of the First Fleet in 1788, unleashed devastating consequences for First Nations peoples. The ensuing Frontier Wars—violent conflicts between Aboriginal people and British settlers—resulted in the massacre and displacement of countless Indigenous communities. Throughout much of the twentieth century, government policy continued to pursue assimilation and cultural eradication.\\

One of the most tragic outcomes of these policies was the forced removal of Aboriginal children from their families. These children, later known as the Stolen Generations, were stripped of their language, names, and cultural identity. Many were placed in institutions or adopted into white families, where they were exploited as unpaid labour and suffered emotional and physical abuse.\\

In \underline{Rainbow's End}, Harrison threads this history through the lives of Nan and Gladys, whose trauma reflects the generational wounds inflicted by colonial intervention. Gladys’s experience of being taken from her family and made to work as a domestic servant mirrors the broader policy of forced assimilation. Her obsession with Dolly’s education becomes a personal act of restitution—an attempt to reclaim what was taken from her by systemic injustice.\\

\subsubsection{Analysis}
By embedding the historical trauma of dispossession within intimate family narratives, Harrison compels her audience to recognise the enduring emotional and social consequences of colonisation. Her portrayal of the Stolen Generations not only condemns past wrongs but also foregrounds the resilience and agency of those who survived them.\\


\subsection{The Flats and Cummeragunja}
\textit{Cummeragunja}, founded in 1888 on Yorta Yorta/Bangerang Country, represented both refuge and resistance for Aboriginal people. Initially a self-sufficient farming community, it was later undermined by the New South Wales Protection Board, which seized control in 1915. Aboriginal workers were paid in meagre rations, and funds generated by their labour were appropriated by the Board. By the 1930s, disease and poverty had decimated the community.\\

In 1939, the residents of Cummeragunja staged a historic protest—the first mass strike by Aboriginal Australians. Many left the mission and settled on “The Flats,” a floodplain near Mooroopna and Shepparton. Despite harsh living conditions, including humpies made from “flattened kero tins” and exposure to extreme weather, the community retained its cultural vitality and kinship ties. By the early 1950s, around 300 people lived there.\\

\subsubsection{Analysis}
The Flats serve as both a literal and symbolic setting in \underline{Rainbow's End}. They embody dispossession yet also community, poverty yet perseverance. Through this setting, Harrison contrasts the instability imposed by colonial structures with the cultural rootedness of Aboriginal families. The Flats become a testament to the capacity of a people to adapt, resist, and endure despite systemic neglect.\\

\section{Analysis of the Historical and Cultural Context in \textit{Rainbow's End}}

\subsection{Overview}
Jane Harrison’s \underline{Rainbow's End} emerges from a landscape shaped by colonisation, segregation, and resilience. Set in 1950s Australia, the play reflects both the harsh realities and quiet triumphs of Aboriginal families living under systemic discrimination. The historical context is indispensable to understanding the characters’ identities, the injustices they face, and the endurance that defines their humanity. Harrison’s work gives voice to the unsung heroes—especially women—who sustained their communities despite persistent marginalisation.\\

\subsection{Acknowledgement and Purpose}
The play’s thematic foundation begins with an acknowledgment of Aboriginal and Torres Strait Islander peoples as the traditional custodians of the land. This recognition situates \underline{Rainbow's End} within an ethical framework that honours history and memory. Harrison’s intent is not merely to tell a story but to preserve a cultural legacy and challenge erasure. Commissioned to highlight Aboriginal women’s experiences, the play transforms domestic life into political narrative, demonstrating that resistance often begins within the home.\\

\subsubsection{Analysis}
By embedding acknowledgment into the narrative fabric, Harrison bridges past and present. The act of remembering becomes both art and activism, inviting audiences to reflect on how storytelling can restore dignity to those history overlooked. Her focus on Aboriginal women redefines heroism as endurance and nurturing in the face of dispossession.\\

\subsection{Historical Background: Colonisation and Segregation}
The origins of the play’s conflict trace back to the legacy of colonisation and government policies such as the Protection Acts, which sanctioned segregation and the forced removal of children. Aboriginal families were confined to reserves like Kamarra Ganja, where strict management and limited access to resources entrenched poverty. These lived realities form the backbone of Harrison’s narrative world, revealing how state control sought to diminish autonomy while families fought to preserve it.\\

\subsubsection{Analysis}
Harrison transforms historical record into lived drama. The reserve system becomes a stage for both oppression and resilience. By drawing from documented injustices, she crafts a narrative that demands empathy and accountability. The audience confronts how legislative cruelty translated into generational trauma.\\

\subsection{Life on the Reserves and the Flats}
After the walk-off from Kamarra Ganja, many Aboriginal families relocated to makeshift settlements known as “the flats.” These spaces, though impoverished, became hubs of culture, laughter, and communal strength. The families’ persistence in maintaining cultural traditions symbolises defiance against systemic neglect. In this environment, children learned resilience as a form of inheritance, and hope became both a necessity and an act of protest.\\

\subsubsection{Analysis}
The depiction of life on the flats captures a paradox — material deprivation coexists with emotional and cultural richness. Harrison celebrates community cohesion as a radical response to marginalisation. The domestic setting transforms into a microcosm of resistance where identity is preserved through humour, storytelling, and maternal leadership.\\

\subsection{Government Policy and the Metaphor of the Rainbow}
The government housing schemes of the 1950s, such as those in Rumbalara, promised integration but perpetuated segregation through inadequate living conditions. The metaphor of the rainbow recurs as a powerful symbol throughout the play. To the characters, the “end of the rainbow” represents dreams of equality and acceptance — visions that remain just out of reach. This unattainable horizon mirrors the historical struggle for recognition endured by Aboriginal Australians.\\

\subsubsection{Analysis}
The rainbow encapsulates the duality of hope and despair. It is both a dreamscape and a mirage — a reflection of the distance between aspiration and reality. Harrison employs this image to expose the cruel optimism embedded within promises of assimilation. The characters’ pursuit of the rainbow becomes a metaphor for the broader national pursuit of reconciliation still incomplete.\\

\subsection{The 1950s Social Landscape}
Post-war Australia was marked by conservative values, economic expansion, and racial inequality. For Aboriginal peoples, this period intensified the contradictions between visibility and exclusion. Women, in particular, bore the burden of navigating these intersecting oppressions. Through Gladys, Dolly, and Nan Dear, Harrison portrays how Aboriginal women negotiated identity, motherhood, and dignity within a society determined to silence them.\\

\subsubsection{Analysis}
Harrison’s representation of 1950s Aboriginal womanhood dismantles stereotypes of passivity. Her characters inhabit a world of limited choices yet wield moral and emotional power. Their courage lies not in grand rebellion but in persistence — in cleaning, teaching, loving, and speaking truths that white Australia refused to hear.\\

\subsection{Conclusion}
Understanding the historical and cultural context of \underline{Rainbow's End} enriches every reading of the play. Harrison’s depiction of segregation, displacement, and resilience transforms history into human experience. The struggles of the Banks and Dear families are not isolated but emblematic of a broader narrative of survival. The play stands as both testimony and challenge: to remember, to acknowledge, and to act. In recognising the truths of the past, Harrison invites audiences to see the enduring strength of Aboriginal identity — not as tragedy, but as triumph.\\



\subsection{Structure of the Play}
\underline{Rainbow's End} opens \textit{in medias res}, plunging the audience directly into the aftermath of a flood. This sudden beginning creates an atmosphere of crisis and displacement, immediately reflecting the instability that defines the characters’ world. The play’s structure is unconventional—beginning with a prologue, followed by scenes divided into parts A and B, reflecting fragmentation and lack of control in the characters’ lives.\\

Harrison integrates dream sequences that blur the boundary between fantasy and reality. These imaginative interludes, such as Gladys’s visions of meeting the Queen or Dolly’s aspirations for education, express both hope and futility. The juxtaposition of dreams and deprivation highlights the gap between aspiration and circumstance.\\

Two narrative threads underpin the play: Gladys’s political awakening and Dolly and Errol’s romance. Each culminates in an epiphany—Gladys’s courageous confrontation with the Shire Council and Dolly’s recognition of Errol’s unconscious racial assumptions. The play concludes with two significant symbols: Nan’s dream of Dolly’s wedding and the song “Que Sera, Sera.” The latter, appearing at both the beginning and end, transforms from a cynical reflection of government indifference into a subtle anthem of hope and change.\\

Harrison’s structural choices mirror her thematic concerns. The fragmented narrative echoes the disrupted lives of Aboriginal families, while the interweaving of dreams and realism underscores the endurance of hope amid systemic oppression. The dual climaxes—Gladys’s activism and Dolly’s assertion of independence—affirm the possibility of transformation both personal and political. The cyclical use of “Que Sera, Sera” redefines fatalism as faith: a belief that the future, though uncertain, can still hold renewal.\\

\section{Structure and Themes in \textit{Rainbow's End}}

\subsection{Overview}
Jane Harrison’s \underline{Rainbow's End} is constructed as a two-act play divided into short, interlinked scenes that illuminate the resilience, struggles, and dreams of First Nations families in 1950s Australia. The fragmented structure mirrors the fractured social and political realities of Aboriginal life, where moments of humour and hope coexist with hardship and prejudice. Harrison’s manipulation of stage direction, song, and dialogue builds dramatic tension and emotional intimacy, guiding the audience through a tapestry of generational experience.\\

\subsection{Act Structure and Thematic Design}
Unlike traditional plays divided by chapters, \underline{Rainbow's End} unfolds through acts and scenes that offer glimpses into domestic and emotional moments. This episodic format reflects the uncertainty that governs the Dear and Banks families’ lives. Scenes often end abruptly or transition during moments of crisis, symbolising the instability caused by poverty and discrimination. The two major narrative threads — Gladys’s political awakening and Dolly’s coming-of-age romance — develop in parallel, converging at the climax where both women redefine their strength and identity.\\

\subsubsection{Analysis}
Harrison’s two-act structure underscores her social commentary. The brevity and rhythm of scenes intensify audience empathy while exposing the invisible pressures imposed by systemic inequality. By structuring the play around fragmented lives, Harrison creates a dramatic form that mirrors cultural and emotional dislocation — yet also reveals endurance and connection within the Aboriginal matriarchal family.\\

\subsection{In Media Res and Dramatic Openings}
The play begins \textit{in media res} amid the chaos of a flood, a moment both literal and symbolic. Accompanied by the recurring song “Que Sera, Sera,” the opening sequence immediately juxtaposes hope with helplessness. This song returns at the play’s conclusion, binding the narrative through a motif that evokes the tension between aspiration and acceptance.\\

\subsubsection{Analysis}
Beginning in the middle of action thrusts the audience into the characters’ reality without exposition, reflecting the immediacy of their struggle for survival. The flood motif encapsulates emotional and social inundation, while the musical refrain serves as both comfort and irony — a lullaby masking pain.\\

\subsection{Dream Sequences and Juxtaposition}
Dream sequences punctuate the play, revealing the characters’ inner lives and unfulfilled desires. Gladys’s fantasy of meeting the Queen becomes a recurring image of colonial aspiration — an impossible vision shaped by internalised subservience and social exclusion. These dreams contrast sharply with the family’s reality in a riverside shack, where dignity must be defended daily.\\

\subsubsection{Analysis}
Through dream imagery, Harrison critiques the illusion of social mobility within a racist society. The juxtaposition of fantasy and deprivation exposes the psychological toll of colonisation, while simultaneously affirming the imagination as a site of survival and self-worth.\\

\subsection{Character Growth and Narrative Resolution}
Both major storylines reach their culmination in acts of self-realisation. Dolly challenges her mother’s fixation on external validation, forcing Gladys to confront the limits of her own deference. This confrontation propels Gladys toward activism, culminating in her powerful declaration at the Shire meeting: “We demand proper schooling … The white people too — they need to be educated about us” (p.197). Though the play ends with personal growth, systemic barriers remain unresolved, reflecting the persistence of structural inequality.\\

\subsubsection{Analysis}
The dual transformation of Gladys and Dolly demonstrates Harrison’s belief in intergenerational empowerment. Their evolution symbolises a shift from private endurance to public assertion — a necessary transition for community progress. The incomplete resolution emphasises realism over idealism, reinforcing that the struggle for equality continues beyond the stage.\\

\subsection{Dramatic Intensity and Unseen Forces}
Harrison’s use of short, tightly constructed scenes sustains dramatic tension and emphasises the relentlessness of external pressures. Many of the most significant events occur offstage, their impact conveyed through dialogue and reaction. The small cast heightens intimacy while unseen antagonistic figures — landlords, officials, and employers — personify the faceless nature of institutional discrimination.\\

\subsubsection{Analysis}
By keeping oppression unseen yet omnipresent, Harrison amplifies its psychological weight. The absence of visible villains shifts focus toward resilience, love, and humour as tools of survival. The minimal staging and interwoven dialogue allow audiences to inhabit the emotional pulse of the characters’ world, where silence often speaks louder than confrontation.\\

\subsection{Circular Structure and `Que Sera, Sera'}
The play opens and closes with `Que Sera, Sera', and the framing is more argument than decoration. On its surface the song counsels fatalism -- whatever will be, will be -- yet everything between the two renditions shows the women refusing passivity: Gladys finds her public voice, Dolly chooses her own future, Nan lets new truth revise old fear. The circle closes with an echo rather than a repetition: Gladys's early reassurance ``It'll be all right'' returns in the final scene in Dolly's mouth. Spoken after flood, humiliation and assault, the line is no longer naive comfort but a tested position -- hope as something inherited, examined and deliberately kept. Structure itself carries the play's meaning: what recurs is not circumstance but resolve.\\

\subsection{Conclusion}
The play’s structure functions as both narrative device and political metaphor. Through acts, scenes, dreams, and songs, Harrison crafts a world where fragmented experiences reveal collective truth. The rhythm of the play mirrors the cycles of struggle and renewal within Aboriginal communities, inviting readers and audiences to engage with the ongoing dialogue between past injustices and present resilience. Understanding the form of \underline{Rainbow's End} thus deepens appreciation for its portrayal of identity, endurance, and hope in mid-twentieth-century Australia.\\


\newpage
\section{Audience and Purpose in \textit{Rainbow's End}}

\subsection{Overview}
Jane Harrison’s \underline{Rainbow's End} (2005) was written to educate and engage both Indigenous and non-Indigenous audiences through its portrayal of three generations of Aboriginal women — Nan Dear, Gladys, and Dolly — living in 1950s Victoria. Harrison’s purpose extends beyond recounting hardship; she celebrates resilience, dignity, and familial love amidst oppressive social structures. The play invites audiences to confront the consequences of colonisation and assimilation while recognising the strength of community and cultural continuity.\\

\subsection{Purpose of the Play}
Harrison’s central purpose is to foreground the untold stories of Aboriginal women and their everyday heroism. Through the intimate domestic setting of the Dear and Banks family, she challenges the invisibility imposed upon First Nations peoples by historical narratives. Her play functions as both cultural preservation and political statement — illustrating that Indigenous survival is not defined by victimhood but by persistence and pride.\\

Harrison’s writing embodies the principle that art can both memorialise and mobilise. By humanising historical trauma, she transforms policy and prejudice into personal experience, allowing audiences to empathise rather than merely observe. The play’s tone oscillates between warmth and sorrow, comedy and critique, mirroring the emotional complexity of survival.\\

\subsubsection{Analysis}
Through \underline{Rainbow's End}, Harrison asserts that storytelling is a mode of reclamation. Each scene restores voices historically silenced by government policy and social discrimination. Her purpose is not only to depict resilience but to invite truth-telling — to reposition Aboriginal women as custodians of memory, agency, and identity.\\

\subsection{Intended Audience}
\underline{Rainbow's End} addresses a dual audience. For non-Indigenous viewers, it acts as an awakening — revealing the lived experiences obscured by whitewashed Australian history. For Indigenous audiences, it affirms continuity and pride, validating intergenerational strength and the endurance of culture.\\

By positioning humour and tenderness alongside hardship, Harrison ensures accessibility for classroom and theatre audiences alike. The play encourages educational engagement, prompting critical reflection among students while offering Aboriginal spectators the comfort of recognition. The inclusion of intergenerational female perspectives fosters empathy across age groups and cultural divides.\\

\subsubsection{Analysis}
Harrison’s nuanced approach ensures that the play operates on multiple emotional and intellectual levels. Non-Indigenous audiences are confronted with complicity in structural injustice, while Indigenous audiences are offered representation that transcends trauma. This balance of challenge and affirmation defines the play’s enduring impact within Australian theatre and education.\\

\subsection{Historical Awareness and Cultural Responsibility}
The play’s authenticity stems from Harrison’s meticulous research and personal connection to her Murrawari heritage. Drawing upon oral histories, newspaper archives, and community interviews, she constructs a historically grounded yet emotionally resonant narrative. Her inclusion of real figures, such as Pastor Doug Nichols, situates the fictional world within an authentic cultural framework.\\

This fidelity to truth underscores Harrison’s responsibility as both playwright and cultural custodian. Her depiction of the 1950s assimilation era reflects an ethical commitment to remembrance — to ensuring that audiences understand the policies that fractured families and communities, and the resilience that rebuilt them.\\

\subsubsection{Analysis}
Harrison’s creative purpose merges documentation with dramatisation. She transforms history into a living dialogue between past and present, ensuring that Aboriginal experience is neither mythologised nor forgotten. By embedding historical authenticity within intimate storytelling, she teaches audiences that understanding cultural identity begins with acknowledging historical reality.\\

\subsection{The Role of Representation in Theatre}
For Harrison, representation is activism. The play reclaims space within the Australian canon for Aboriginal voices, particularly those of women. Through humour, vernacular speech, and emotional candour, she dismantles stereotypes and restores humanity to communities too often reduced to statistics. \underline{Rainbow's End} thereby becomes a vessel for truth-telling, reconciliation, and education.\\

\subsubsection{Analysis}
Representation in \underline{Rainbow's End} is not symbolic but lived — an act of cultural assertion. Harrison’s purpose is fulfilled when audiences see Aboriginal women not as passive subjects of history but as authors of their own narratives. The play’s enduring presence in school curricula ensures that these voices continue to educate, inspire, and transform generations of Australians.\\

\subsection{Conclusion}
Harrison’s \underline{Rainbow's End} unites purpose and audience in a dialogue of empathy, education, and recognition. Her portrayal of three Aboriginal women transcends their fictional identities to stand as symbols of cultural endurance. By crafting a story grounded in truth and delivered with compassion, Harrison reshapes the Australian stage into a site of reconciliation — where understanding begins through listening, and hope endures through storytelling.\\



\newpage
\section{Language in \underline{Rainbow's End}}

\subsection{Overview}
Language in Jane Harrison’s \underline{Rainbow's End} is a powerful vehicle for expressing identity, class, and cultural conflict. The characters’ speech patterns reveal their personal histories and social positions, while Harrison’s authorial voice is embedded within their dialogue, reflecting her own commentary on racial inequality, resilience, and generational change. Language, therefore, functions not only as a mode of communication but also as an indicator of power, prejudice, and pride within the 1950s Australian context.\\

\subsection{Nan Dear}
Nan’s language is colloquial, humorous, and often sharply ironic. When she informs Gladys that she “won’t be needing any of them encyclops” (p.136), she dismisses education’s relevance to her family’s daily struggles, revealing both her pragmatic worldview and the limited opportunities afforded to her generation. Her unpolished speech reflects her lower-class status and lack of formal education—an outcome of systemic exclusion faced by Aboriginal Australians.\\

Yet Nan’s language also injects comedy into the play. Her dry humour and literal phrasing create moments of levity that highlight resilience in hardship. When mocking Gladys’s excitement about meeting the Queen, Nan’s irony underscores the absurdity of colonial class structures. Through Nan’s speech, Harrison reveals both the pain and pride of a woman grounded in survival, whose humour becomes an act of quiet rebellion against social inequality.\\

\subsubsection{Analysis}
Nan’s dialect situates her within a marginalised social class while asserting authenticity and self-respect. Her vernacular speech becomes an emblem of identity, opposing assimilationist pressures that sought to erase Aboriginal voice. Harrison uses Nan’s language to celebrate oral tradition and defiant humour, challenging notions that education alone defines intelligence.\\


\subsection{Dolly Banks}
Dolly’s language reflects her education and confidence but also her awareness of the social hierarchies constraining her. She is direct and assertive, particularly when confronting racism. Her rebuke of Errol’s remark about “you people” (p.141) exposes his unconscious bias and the normalisation of racial division. Likewise, her challenge to her mother — “You’re always telling me to stick up for myself, but when do you, eh?” (p.162) — forces Gladys to examine her internalised inferiority.\\

Dolly also code-switches, adopting Indigenous vernacular when she threatens to “kick Ester’s boys up to moon” (p.137), using humour and familiarity to assert belonging within her community. Her fluency, both in educated English and colloquial Aboriginal dialect, embodies her dual identity as a bridge between two worlds. Her closing assertion to Errol — that she would never trade her family or heritage for a “better life” (p.172) — epitomises pride and defiance.\\

\subsubsection{Analysis}
Dolly’s language signifies empowerment through articulation. Harrison crafts her speech to demonstrate that mastery of language is not assimilation but assertion. Dolly’s ability to speak for herself—and against injustice—makes her a vehicle for Harrison’s advocacy of Indigenous self-determination. Her linguistic versatility becomes symbolic of cultural strength rather than concession.\\



\subsection{Errol Fisher}
Errol’s halting and awkward speech reflects both his nervousness and his ignorance. As an encyclopedia salesman who depends on persuasion, his incompetence is ironically comic. When he awkwardly notes that Dolly’s glass is “made out of a can” (p.140), he unwittingly insults her. His language often reveals class and racial privilege cloaked in politeness — exemplified in his use of “you people” (p.140) to generalise Aboriginal Australians.\\

His verbal missteps and tone-deaf expressions highlight his sheltered upbringing and lack of cultural awareness. Over time, however, Errol’s language softens and becomes more considerate, mirroring his emotional growth and respect for Dolly’s world.\\

\subsubsection{Analysis}
Harrison uses Errol’s speech to critique the well-meaning but ignorant paternalism of mid-century white Australia. His progression from patronising to empathetic reflects the broader social journey toward understanding and equality. His linguistic awkwardness thus serves as a tool of both comedy and commentary.\\



\subsection{Gladys Banks}
Gladys’s voice develops from timid dreamer to assertive activist. During the Rodney Shire meeting, her passionate declaration — “We demand proper schooling … The white people too — they need to be educated about us, and our ways” (p.197) — becomes a turning point where her speech echoes Harrison’s own political message.\\

Her earlier language is often infused with fantasy and aspiration: daydreams of the Queen’s attention or Dolly’s graduation soften her struggles. Yet by the play’s end, Gladys has found a political and moral voice that embodies the collective frustrations of Indigenous Australians denied dignity. Her transition from domestic to public speech represents empowerment through language.\\

\subsubsection{Analysis}
Through Gladys, Harrison demonstrates the transformative power of speech as activism. The act of “speaking up” replaces passive dreaming with agency. Gladys’s voice becomes that of her people — articulate, righteous, and unafraid — asserting the necessity for visibility and reform.\\


\subsection{The Royals}
The Queen’s language is formal, elevated, and exclusionary. Her reference to Australia as “the birthplace of the nation” (p.124) erases 65,000 years of Aboriginal civilisation, while her statement about standing “at last among Australians” implicitly refers only to white citizens. This polished rhetoric embodies colonial arrogance.\\

By contrast, the Duke’s language is more spontaneous and unguarded. His offhand question — “Where’s Father?” (p.127) — humanises him but also underscores the royal couple’s detachment from the lived realities of those they visit.\\

\subsubsection{Analysis}
Harrison employs the Royals’ diction as social critique. Their words — ceremonious, detached, and self-congratulatory — expose the gulf between imperial symbolism and the Indigenous suffering it conceals. The Queen’s non-inclusive language mirrors the nation’s non-inclusive identity at mid-century.\\


\subsection{The Radio Announcer}
The radio announcer’s language is superficially cheerful yet profoundly condescending. In describing the royal visit, he mentions a demonstration of “how to throw a boomerang” and marvels that it “really does come back” (p.128). His tone trivialises a complex cultural artefact, reducing Indigenous skill to entertainment for the royal couple.\\

\subsubsection{Analysis}
The announcer’s commentary encapsulates institutional racism disguised as enthusiasm. Harrison’s use of irony highlights how mainstream media perpetuated ignorance, framing Aboriginal culture as novelty rather than heritage. His words parody the way national pride excluded First Nations people.\\


\subsection{Structural and Symbolic Language}
Beyond dialogue, Harrison uses language in her stage directions and symbolism to reinforce meaning. Descriptions of the Dear family’s humpy — “waterlogged” and “mud-spattered” (pp.123, 176) — visually communicate poverty and marginalisation. The title itself, \underline{Rainbow's End}, draws on the myth of a pot of gold, suggesting the illusion of prosperity and the unattainability of a promised better life.\\

The reference to “Somewhere Over the Rainbow” from \textit{The Wizard of Oz} introduces an ironic hopefulness. Its lyrics — “the dreams that you dare to dream really do come true” — echo the family’s optimism but also question the realism of those dreams in a society unwilling to grant equality.\\

\subsubsection{Analysis}
Harrison’s manipulation of both literal and figurative language constructs a layered critique of Australia’s racial landscape. The interplay between hope and futility in her metaphors underscores the central tension of the play — between aspiration and limitation, imagination and reality. Language, whether spoken, written, or symbolic, thus becomes both a weapon and a wound.\\




\newpage

\section{Analysing Stage Directions}
\addcontentsline{toc}{section}{Analysing Stage Directions}

In a play, the author speaks twice: once through the dialogue, and once through everything the audience is shown. Harrison's stage directions -- lighting changes, props, physical business, silences -- are deliberate authorial constructions, which makes them premium evidence in a Section A essay: quoting one instantly signals that you are treating the text as something \textit{made}. The rule, however, is absolute: you are never writing \textit{about} stage directions. You are writing about character, relationships, ideas, concerns and values -- and using stage directions as evidence to prove your interpretation. Identifying a direction and stopping is the mid-band trap; the marks live in what it proves.\\

\subsection{The Evidence Ladder}
Before a stage direction goes anywhere near your essay, climb it through five rungs: What is the direction? What does the audience literally see? What does that reveal? What idea, concern or value does it support? And finally -- what does Harrison want the audience to understand? Consider the difference in these three sentences about the moment Dolly and Errol look at each other ``longingly'': saying ``Harrison uses the stage direction `longingly''' merely identifies; saying it ``shows that Dolly and Errol like each other'' merely explains; but writing that ``Harrison's direction that the pair look at each other `longingly' allows their developing affection to be communicated visually, positioning their relationship as a source of emotional possibility'' analyses -- it moves from choice to meaning to purpose. Reach for precise analytical verbs as you climb: constructs, establishes, reveals, emphasises, reinforces, contrasts, symbolises, positions, foregrounds, intensifies, challenges.\\

\subsection{A Worked Ladder: The Blue Jewels}
Take the direction from Gladys's dream in Act 1, Scene 7: ``A tap appears from nowhere and from it flows blue jewels in an approximation of water.'' What we see: water presented as treasure. What it reveals: the intensity of Gladys's longing for something most of the audience has never once thought about. What idea it supports: dignity, housing and inequality -- the politics of the play made visible in a single image. What Harrison wants understood: that the ordinary becomes precious only when it has been withheld, so the audience is made to feel the family's deprivation rather than merely be told about it. Built into a paragraph, the ladder becomes: argument (Gladys's hopes are shaped by material deprivation), evidence (the ``blue jewels''), analysis (the theatrical image renders running water magical), significance (the depth of her longing for basic security), link (what many audiences regard as ordinary becomes a powerful aspiration for those living without it).\\

\subsection{Directions Worth Knowing Closely}
Some of the play's directions carry entire arguments within them. The Prologue's opening -- ``gloomy outside'', everything ``sodden and mud-splattered'' -- makes the setting perform the family's circumstances before a word is spoken, while the women's calm, practised rebuilding tells the audience this catastrophe is routine: Harrison stages normalised crisis, not singular disaster. The Inspection's ``nervous, anxious energy'' as Nan and Gladys straighten everything in sight physicalises the power the welfare state holds over the household -- fear expressed through housework. At the bank, the manager ``pours a cup for himself, but doesn't offer one to GLADYS'': four seconds of stage business that dramatise an entire social hierarchy without one hostile word. At the ball, ``DOLLY is frozen in shame'' -- stillness as violence, the visual opposite of her ``swirling and floating'' entrance moments earlier; the two directions analysed together map her humiliation more precisely than any dialogue. When ``their humpy disappears'' to the sound of bulldozers, Harrison erases the family home as an act of theatre, and the replacement housing rises ``concrete, small, white and featureless'' -- the adjective ``white'' doing quiet, double work. And in Act 2, Errol finds Gladys ``in a similar scene to the first time he saw her, chopping wood'' -- a structural echo that measures how far the relationship has travelled by holding the setting still. Small kindnesses are staged too: Errol ``notices GLADYS is sitting on a kero tin so he does the same'' -- respect performed as posture, the visual antithesis of the bank manager's tea.\\

\subsection{Stage Direction Evidence Bank}
\renewcommand{\arraystretch}{1}
\begin{longtable}{|p{3.4cm}|p{9.6cm}|}
\hline
\textbf{Act / Scene} & \textbf{Direction (verbatim)} \\
\hline
\endfirsthead
\hline
\textbf{Act / Scene} & \textbf{Direction (verbatim)} \\
\hline
\endhead
Prologue: Aftermath & ``It's late spring, late afternoon and gloomy outside. Inside their humpy NAN DEAR and GLADYS are rebuilding after a flood has devastated their home.'' \\
\hline
Prologue: Aftermath & ``GLADYS mops, wrings out and removes things that are destroyed. NAN finishes hanging a piece of hessian to replace a ruined piece that lined the interior walls.'' \\
\hline
Act 1, Scene 1 (A) & ``Humpy interior. Morning. GLADYS is getting dressed up and humming to herself. DOLLY has her head down over her schoolbooks.'' \\
\hline
Act 1, Scene 1 (A) & ``GLADYS flounces out to the back room.'' \\
\hline
Act 1, Scene 2 (B) & ``The lights come up on DOLLY pushing the pram, while a Brylcreemed lad, ERROL, wobbles up on a bicycle from the opposite direction. He nearly falls off his bike at the sight of her.'' \\
\hline
Act 1, Scene 2 (B) & ``ERROL dismounts and studies his map. He looks puzzled, but puts the map away and pushes on towards GLADYS, who is still chopping wood.'' \\
\hline
Act 1, Scene 3: Lino & ``DOLLY is rummaging at the town tip. She looks at, and discards, a few items. The lights change to a dream sequence. A well-groomed SALESMAN appears.'' \\
\hline
Act 1, Scene 3: Lino & ``DOLLY hoists the lino roll over her shoulder as the SALESMAN fades away. DOLLY walks past the cork trees. She sees someone in the shadows.'' \\
\hline
Act 1, Scene 4: House of Biba & ``GLADYS serves out the stew while DOLLY sets the table.'' \\
\hline
Act 1, Scene 5: The Delivery & ``ERROL notices GLADYS is sitting on a kero tin so he does the same.'' \\
\hline
Act 1, Scene 5: The Delivery & ``The lights change for a dream sequence. GLADYS sees DOLLY in a robe and clapboard hat.'' \\
\hline
Act 1, Scene 6: The Inspection & ``NAN and GLADYS are inside the humpy, straightening everything in sight with a nervous, anxious energy.'' \\
\hline
Act 1, Scene 7: The Turn & ``DOLLY, out of sight, is singing to the radio. NAN sits in the only chair, cleaning a pair of slingback shoes with white shoe cleaner.'' \\
\hline
Act 1, Scene 7: The Turn & ``The lights change for GLADYS's dream sequence. A tap appears from nowhere and from it flows blue jewels in an approximation of water.'' \\
\hline
Act 1, Scene 9: Home Sweet Home & ``As the lights come up, GLADYS is going around the humpy emptying the mouse traps of dead mice and setting new traps, while DOLLY does her homework.'' \\
\hline
Act 1, Scene 10: The Bank & ``On the BANK MANAGER's desk is a tray with two china cups of tea, milk jug and sugar bowl. He pours a cup for himself, but doesn't offer one to GLADYS.'' \\
\hline
Act 1, Scene 11: The Ball & ``When the lights come up, DOLLY is swirling and floating around the dance floor in her full-skirted dress, with ERROL as her partner. She is wearing an orchid corsage.'' \\
\hline
Act 1, Scene 11: The Ball & ``DOLLY is frozen in shame.'' \\
\hline
Act 1, Scene 14: The Flood & ``Rain, thunder, darkness. Time passes. The waters rise.'' \\
\hline
Act 2, Scene 1: After the Flood & ``Dawn finds GLADYS outside the humpy. The water has drained away, but the devastation has been wrought. Everything is saturated and muddy.'' \\
\hline
Act 2, Scene 2: Rumbalara & ``NAN and GLADYS hold their humble possessions as the sound of the bulldozers is heard. Their humpy disappears.'' \\
\hline
Act 2, Scene 2: Rumbalara & ``The lights go down, then come up on the new housing. It's concrete, small, white and featureless.'' \\
\hline
Act 2, Scene 3: The Broadcast & ``The lights come up on DOLLY who is peeling big dirty potatoes.'' \\
\hline
Act 2, Scene 4: The Contract & ``ERROL hurries up the track and is relieved to see GLADYS -- in a similar scene to the first time he saw her, chopping wood -- but this time outside the new Rumbalara housing.'' \\
\hline
Act 2, Scene 4: The Contract & ``So she nods and they both sit down on the kero tins. Resting a piece of paper on ERROL's book, they start to compose a letter together, ERROL writing it down.'' \\
\hline
Act 2, Scene 5: Pay the Rent & ``As the lights come up, it is early morning, but already blindingly hot. NAN is cooking outside, just like she used to.'' \\
\hline
Act 2, Scene 6: Errol Spills the Beans & ``The lights come up on GLADYS and ERROL, sitting on a park bench, with the Inspector's report.'' \\
\hline
\end{longtable}

\section{Quote Banks}

\subsection{Character Quote Bank}

\subsubsection{Gladys}
\begin{itemize}
    \item ``I know how you feel about Royalty'' (p.~124)
    \item ``Where’s my white gloves'' (p.~124)
    \item ``Gladys get off your high horse. Least we do things our way - no one breathin down our necks'' (p.~128)
    \item ``It should have been where they built us housing, being the highest land around here'' (p.~142)
    \item ``Oh for heaven’s sake, spit it out, will ya, son?!'' (p.~148)
    \item ``I’m sorry about last time, Errol. I wasn’t so smart then. I treated you unfairly'' (p.~183)
    \item ``Well… it achieved nothing…'' (reality) (p.~189)
    \item ``If you want me to read any part to you… I can. Thank you.'' (p.~198)
\end{itemize}

\subsubsection{Nan}
\begin{itemize}
    \item ``What’s he doing here? He’s not the Welfare?'' (p.~134)
    \item ``Ask him what it says in that there encyclops about the Aborigines, eh?'' (p.~134)
    \item ``Think I don’t know a snake-oil salesman when I see one!'' (p.~136)
    \item ``Now Gladys you won’t be needing any of them encyclops'' (p.~136)
    \item ``They may be drinkers but they are still our people … Show some respect girl, those lads have had it hard'' (p.~139)
    \item ``The bruised ones that otherwise go to the pigs'' (p.~146)
    \item ``Encyclops boy and he still knows nothing'' (p.~147)
    \item ``She’s with a boy. That’ll make her silly, not sensible'' (p.~173)
    \item ``You’ve got cheek. Showing your face'' (After Errol shows up after flood) (p.~177)
    \item ``Not quite the ‘new deal’'' (p.~178)
    \item ``And hospitals is where they take our babies away'' (p.~185)
    \item ``You are your visiting hours'' (p.~186)
    \item ``No matter how they come into the world, you still love ’em the same'' (p.~186)
    \item ``I can recognise a good man when I see one'' (p.~198)
\end{itemize}

\subsubsection{Dolly}
\begin{itemize}
    \item ``Mum reckons they think is too good for blackfellas'' (p.~135)
    \item ``Nan wouldn’t like that. City’s are full of sin'' (p.~141)
    \item ``I’m trying, Mum, I’m really trying'' (p.~144)
    \item ``Nan why do you always treat me like a child?'' (p.~153)
    \item ``I want to know. I’m not a child. I’m a woman, Nan!'' (p.~158)
    \item ``You think books and school, are the answer to everything'' (p.~159)
    \item ``It’s an interview for a job I’m not going to get. And that you want, not me. Why doesn’t anyone ask what I want'' (p.~178)
    \item ``Look at you, Mum. You go on about the things I should be doing… Why don’t you fix your own house? If you know what I mean…'' (p.~179)
    \item ``But I’m worried I won’t love it. Because of'' (p.~186)
    \item ``Keeping secrets, you mean'' (p.~197)
\end{itemize}

\subsubsection{Errol}
\begin{itemize}
    \item ``Yes. But all’s well that end well'' (p.~137)
    \item ``Why do you live out here!'' (p.~147)
    \item ``I like hearing out your life. You family… it’s different from mine'' (p.~170)
    \item ``My dad expects me to call him ‘sir’'' (p.~170)
    \item ``It’s none of my business, Mrs Banks, but you’re a smart lady'' (p.~183)
    \item ``My family? Same as always I guess'' (p.~190)
\end{itemize}

\subsubsection{Inspector}
\begin{itemize}
    \item ``Really that’s the way! And what does the future hold for you, Dolores?'' (p.~151)
    \item ``Yes… As a result of my report, things will change… Things must change!'' (p.~152)
\end{itemize}

\subsubsection{Dear}
\begin{itemize}
    \item ``You don’t ever get too clever, my girl'' (p.~155)
    \item ``You get knocked down when you get too cleaver'' (p.~155)
    \item ``That cyclops boy? That Grubby Fella'' (p.~155)
\end{itemize}

\subsubsection{Cousin}
\begin{itemize}
    \item ``No? Think you’re too good for us.'' (p.~167)
    \item ``This who you step out with?'' (p.~167)
    \item ``She’s no lady, she’s just a little'' (p.~167)
\end{itemize}

\subsection{Thematic Quote Bank}
\subsubsection{Female}
\begin{itemize}
    \item ``Gladys is chopping wood with an axe'' (p.~129)
\end{itemize}

\subsubsection{Racism}
\begin{itemize}
    \item ``A black contestant? I’d like to see that'' (p.~130)
    \item ``I don’t think so, Mister. They’re not for the likes of us.'' (p.~133)
    \item ``A girl from the flats? I don’t even see the town Aboriginal working in stores'' (p.~144)
    \item ``Do you have any idea what’ll happen if you walk into that dance with me’’ … “You’re white, I’m aboriginal’’ … “I’m from the flats. Not even one of those townie types of cross-over aboriginals’’ … “What matters is you. Not your address’’ … “That’s sweet’’ (pp.~148--149)
    \item ``Are you aware of the concept of assimilation, Mrs Banks?’’  
    ``The aborigine needs to be absorbed into the community… But how can he be absorbed until he learns to live like us?’’ (p.~152)
    \item ``They won’t take her.. She’s seventeen, they’d make her work for someone like they did you’’ (p.~154)
    \item ``When your last to be served, you hear lots of things’’ (p.~158)
    \item ``Nan why don’t we have a normal life’’ … ``But why? It’s like we’re fighting nature all the time’’ (p.~159)
    \item ``It’s me that get’s stones thrown at her when I walk down the street. It’s me that get’s snide remarks’’ … ``You think I don’t get my fair share? Or Nan? Even Papa Dear… you have to learn to let them not shame you’’ … ``Have you Mum? Have you learnt not to be shamed by them? You’re always telling me to stick up for myself but when do you, eh?’’ (p.~162)
    \item ``Feel like getting coshed every time we step out together? Can’t walk down the street holding hands for being called names? And what would your dad say if you took me home to meet him?’’ (p.~171)
    \item ``Us Aboriginals are not welcome in the township… as if we choose to live on a floodplain—not realising we need water too—to cook and to clean’’ … ``Maybe you don’t think we do wash’’ (p.~181)
    \item ``We’ve been humble too long’’ … ``so we don’t need privacy, not like regular folk, is that it?’’ … ``Why do we have to prove we can live like whitefellas’’ … ``No, we’re not even citizens’’ (pp.~196--197)
\end{itemize}

\subsubsection{Dreams}
\begin{itemize}
    \item ``Least not till you’ve finished your studies. And get a good job. In town!’’ (p.~131)
    \item ``My girl, a graduate’’ (p.~133)
    \item ``A girl like you… You could do anything you want’’ (p.~141)
    \item ``But it sounds like a good job… a good opportunity’’ (p.~143)
    \item ``At seventeen you have these silly dreams’’ (p.~187)
\end{itemize}

\subsubsection{Hope}
\begin{itemize}
    \item ``Set them up for life, you reckon?’’ (p.~134)
    \item ``You need to rally yourselves… Form a delegation’’ (p.~152)
    \item ``And what bloody good would that do? Daydreams!’’ (p.~154)
    \item ``Do you think they’re wonderful though? You’re expanding your possibilities’’ (p.~161)
    \item ``Look, it’s real important, For my daughter, see. Not for me, but for my daughter. Please’’ (p.~174)
    \item ``Every one, Gladys? Every single one, ruined’’ (p.~176)
\end{itemize}

\subsubsection{Reality}
\begin{itemize}
    \item ``Hardly anything in that tin’’ (p.~145)
    \item ``How frequently does it flood?’’ (p.~151)
    \item ``Sorry rapport means’’ … ``You put your pennies in here’’ (p.~164)
    \item ``When it was our sunroom curtains. But, you know, I thought we took them to the tip.’’ (p.~166)
    \item ``I knew that’d happen’’ (p.~168)
    \item ``Nice? You heard them. It’ll never change… some things never do’’ (p.~168)
    \item ``The future’s different now’’ (p.~193)
\end{itemize}

\subsubsection{Distrust}
\begin{itemize}
    \item ``Hard to tell a good man from a bad’’ (p.~136)
\end{itemize}

\subsubsection{Dreams/Reality}
\begin{itemize}
    \item ``At seventeen you have these silly dreams’’ (p.~187)
\end{itemize}

\subsubsection{Racism/Stereotype}
\begin{itemize}
    \item ``Are you aware of the concept of assimilation, Mrs Banks?’’ (p.~152)
\end{itemize}

\subsubsection{Racism/Inferiority}
\begin{itemize}
    \item ``When your last to be served, you hear lots of things’’ (p.~158)
\end{itemize}

\subsubsection{Trauma}
\begin{itemize}
    \item ``At least here we sink or swim on our own. Not like the Cummeragunja days, always at the mercy of the manager’’ (p.~159)
    \item ``You ask me what I want? Well I want any mother’’ (p.~159)
\end{itemize}

\subsubsection{Mother/fights}
\begin{itemize}
    \item ``And was it her place to give you permission… Well it’s not right I’m her mother’’ (p.~161)
\end{itemize}

\subsubsection{Mother/conflict}
\begin{itemize}
    \item ``I’ll make the decisions regarding Dolores, thank you’’ (p.~162)
\end{itemize}

\subsubsection{Contension}
\begin{itemize}
    \item ``Sorry rapport means’’ (p.~164)
\end{itemize}

\subsubsection{Oblivious}
\begin{itemize}
    \item ``I can’t imagine how many pressies you have under the tree’’ (p.~169)
\end{itemize}

\subsubsection{Family}
\begin{itemize}
    \item ``I’d have to leave the river’’ … ``I’d have to leave my family’’ … ``A real home’’ (pp.~171--172)
    \item ``This is my place. I’m staying right here with my mum and Nan’’ (p.~172)
    \item ``He didn’t understand that it could never leave you’’ (p.~187)
    \item ``Because, where you belong, and your family, is important. To you, and to me’’ (p.~193)
\end{itemize}

\subsubsection{Family/Place}
\begin{itemize}
    \item ``A real home is where there are people looking out for each other. Do they do that in your home, in your family, Errol?’’ … ``No Errol. Our life isn’t perfect, but like Nan says, it’s ours. And you can’t respect that. I’m sorry.’’ (pp.~171--172)
    \item ``The Murray – that’s your place, isn’t it?’’ (p.~194)
\end{itemize}

\subsubsection{Abandon}
\begin{itemize}
    \item ``Well, Papa Dear is not here to hear is he’’ (p.~186)
\end{itemize}

\subsubsection{Gender}
\begin{itemize}
    \item ``What kind of man has soft hands’’ (p.~195)
\end{itemize}

\subsubsection{Strength/Family}
\begin{itemize}
    \item ``It’ll be all right’’ (p.~199)
\end{itemize}

\subsection{Stage Devices Quote Bank}
\begin{itemize}
    \item ``She toes the old, ruined lino’’ (p.~123)
    \item ``But Nan doesn’t want to talk about that business’’ (p.~126)
    \item ``Instead of being formal, the queen pulls her into a hug’’ (p.~126)
    \item ``Instant spark of attraction’’ (p.~132)
    \item ``Celestial music is heard’’ (Around Encyclopedia) (p.~133)
    \item ``Nan springs into action’’ (p.~135)
    \item ``[hissing] Call her back right now’’ (p.~135)
    \item ``Dolly is humiliated but Errol doesn’t notice’’ (p.~140)
    \item ``Kisses Gladys on the forehead’’ (p.~143)
    \item ``She shakes her head and that’s that’’ (p.~144)
    \item ``She is openly hostile to Errol’’ (p.~144)
    \item ``Dolly knows their relationship cannot work’’ (p.~149)
    \item ``Nan appears and the music stops suddenly’’ (p.~150)
    \item ``perhaps *Nan snorts*’’ (p.~153)
    \item ``Nan shakes her head and vanishes outside to hide the anxiety she is feeling’’ (p.~153)
    \item ``Nan hides the white shoes’’ (p.~156)
    \item ``Nan looks relieved’’ (p.~157)
    \item ``Who’s all dressed up’’ (p.~163)
    \item ``But doesn’t offer one to Gladys’’ (p.~163)
    \item ``They look at each other in despair, barely able to hold it together (after realising what happened to Dolly)’’ (p.~177)
    \item ``The two women move to stand by Dolly’s side, in unity’’ (p.~177)
    \item ``Nan and Gladys hold their humble possession’’ (After radio talks about attempts to solve Aboriginal housing) (p.~178)
    \item ``Who looks at disdain at her body’’ (Inspector to Dolly pregnant) (p.~185)
    \item ``He runs. Nan chucks eggs, one after another, at his departing form’’ (p.~186)
    \item ``Nan looks away to hide her upset’’ (after mention of Papa Dear not being here) (p.~186)
    \item ``She’s finally a woman in Nan’s eyes’’ (p.~186)
\end{itemize}

\subsection{Literacy Devices Quote Bank}
\begin{itemize}
    \item ``Like a band-aid over a sore’’ (p.~127)
    \item \textit{Pathetic Fallacy} – ``Reply is like the rising wind’’ (p.~168)
    \item ``Vignettes’’ (no page given)
\end{itemize}

\subsection{Foreshadowing Quote Bank}
\begin{itemize}
    \item ``If I catch you going past them cork trees’’ (p.~139)
    \item ``There’s these gnarled trees down the track… All twisty and rough’’ (p.~149)
    \item ``He’s Errol Fisher? A Fisher?’’ (p.~156)
\end{itemize}

\subsection{Symbols Quote Bank}
\begin{itemize}
    \item Hessian (p.~123)
    \item Cork trees (p.~126)
    \item Encyclopedia (p.~133)
    \item Fountain pen (p.~135)
    \item Business card (p.~138)
    \item Lino (p.~140)
    \item Rainbows end – ``end of a rainbow’’ Rumbalara (p.~155)
    \item ``I’ll be home by Pat Boone’’ (p.~160)
    \item Catch a falling star (p.~163)
    \item Storm – (Storm brewing as life gets worse for the family, leading up to Dolly’s rape) (p.~167)
    \item Thunder, Lighting – (Leading up to rape) (p.~172)
    \item Dress (Made out of Nancy’s curtains from the tip) – ``It’s a stupid dress’’ (p.~166, p.~169)
    \item ``I asked my sweetheart, what lies ahead’’ (p.~175)
    \item Kero Tin (p.~175)
    \item Rent collector (p.~175)
    \item Nancy (p.~175)
    \item Chopping wood (p.~175)
    \item Catching a Falling Star (p.~183)
    \item Somewhere over the rainbow (p.~184)
    \item Hospital (p.~184)
    \item Inspector (p.~184)
\end{itemize}
\newpage



ZZPDFEMBED QuoteTableSummary.pdf ZZEND



\newpage
\section{Theme Analysis}

\subsection{Love and Belonging}
Harrison explores love and belonging in rainbows end, with Nan and Gladys creating a warm, loving home for Dolly. When Errol offers Dolly a “better life” in the city with a “real home,” Dolly is taken aback by his naivety to her culture. Even though the humpy had its downside, Dolly acknowledged that it's the fierce love and loyalty that made living there easier. 

Nan. Sometimes, you have to move on. Leave things behind. Even things you love. (P.187) \\
Dolly. No matter how they come into the word, you still love'em the same. (P.186) \\
Errol. What matters is [her] / not [her] address (p. 149). \\
Nan. Oh, for heaven's sake, Dolly marry that boy, before someone else does. (P.198) \\

\subsection{Leadership}
Leadership is explored through Gladys’ speech where she took Papa Dears place at the petition.
Harrison displays not only the leadership in the family dynamic but also the government taking leadership over the indigenous Australians and believing that they know what is best for them and displaying leadership in a very different way. This is displayed by the Inspector as being the character showing the government in ‘Rainbow’s End’.

Inspector: “but someone must take leadership” \\
Gladys: “I will present the petition” \\
Gladys: “She needs to know the world is bigger than just this.” (pg 131) \\
Gladys: “I’m not an interloper- I belong here –this is my land!” \\
Gladys: “We’ve been humble for too long.” \\
Gladys: “I make the decisions regarding Dolores, thank you.” \\
Nan Dear: “Grab the flour, Glad, and the tea.” \\
Gladys: “Oi, Re the so-called ‘housing problem’, it is only a problem because of us Aboriginals-” \\
Nan Dear: “I don’t care what you think! You and your visiting hours. Your rules. No singalongs after dark. Your spying. You, mister, can go to blazes! I’ll give you ‘one’...” \\
Gladys: “We demand to be heard.” \\

\subsection{Race and Racism}
- The love plot in the book is purposefully between Dolly and Errol, Harrsion aims to defy the otherwise segregated nature of the society the play is set in. \\
- Whiteness is a constant theme throughout the play, symbolising the belief that white is somehow ‘cleaner’. \\
- All white characters in the play besides Errol are a source of authority; The councilmen, bank manager and inspector all represent different amenities needed by Aboriginal people that have been controlled by the white population. \\

Nan Dear: “They forced us to leave. Forced us to leave Cummeragunja. Our home” the repetition of “forced” and the use of the words “our home” suggests that Harison is trying to show that the white people think they know better than the indigenous people in the play. p126. \\
Nan Dear: “and hospitals is where they take our baby’s away” \\
Nan Dear: “A girl from The Flats? I don’t even see the town Aboriginals working in stores.” \\
P127 “the hessian that I couldn’t get through and couldn’t even peek over.” \\
P127... “stop the likes of her seeing our humpies.” \\

P152 “I don’t know how you do it...With river water, no less!” \\
P164 “- sorry, report means -” bank manager assumes she doesn’t know. \\
P140 “You people” - Errol – oblivious \\
P163 “What I'm wondering is how she’d fit in” \\
“and hospitals is where they take our babies away...” (Nan pg129) \\
“and your whites are so white!” (Inspector pg152) \\

\subsection{The Intersection of Class and Race}
- The intersection of race and class in Rainbow’s End is exemplified by the dynamic between the white population and the Indigenous population. This relationship is explored in scenarios such as Dolly’s search for job, in which it is seen as unlikely she will be able to receive employment. This creates the inescapable cycle that is forged when race meets class, that to escape a lower class, it’s required to achieve a job that could be considered higher status to society. Race hinders this progress due to the prejudices and barriers held against Indigenous people. \\
- The inspector's view in regard to the flood, as he seems to have no awareness of the effect this will take upon first nations families, particularly the Dear family. \\
- Gladys having to save every penny to afford the encyclopedias just to invest in the cultural capital. \\
- Profit from proximity to whiteness (town Aboriginals) \\
- Government policies \\

Inspector: “property can be replaced” \\
“I'd kind of like to be a nurse”- Dolly, which she believes to be silly. \\
Dolly- “What tree? It would be nice to have a tree and presents.” Pg 169 \\
Inspector “I’m not familiar with the department stores in this area” – pg. 152 (in reference to Daisch’s the town tip” \\
“Your whites are so white! With river water no less!” Pg 152, the inspector. Subtle reference to white is believed better. \\
“When it was our sunroom curtains. But you know i thought we took them to the tip” (in reference to dolly’s ball gown) pg.166 \\
Gladys- “Not on a picker’s wage. Gawd we’d be living on Johnny cakes” pg. 174 (in reference to when Jungi says that property can be replaced.) \\
“The old, ruined lino” pg. 123 The dears have to make dear with what they can get access to unlike the whites who it is privileged. \\
“Love your dress, Dolly. Love the fabric. My mother quite liked it too. When it was our sunroom curtains. But, you know, I thought we took them to the tips.” Pg 166 Nancy embarrasses dolly about her dress when Dolly doesn’t have access to the same privileges that Nancy does. \\
Bank manager: “But everyone needs savings” when Gladys says she’s never been in a bank, the white population has the privilege of using banks while the aborigines make do with tins to store their money. Pg 164 \\
“Rapport with our customers is important –sorry”. Pg. 164 \\
“Isn’t there any way?” Errol, when Gladys says she no longer needs the encyclopedia; The dears unlike whites would have trouble with paying for the encyclopedias and the metre box. Errol doesn’t understand this and there for asks the quote. \\
“I’d think a hospital would be a more suitable place-“when Dolly goes into labour. Unlike white people Aboriginals babies are taken away at the hospitals and the dears are unlikely to be able to afford to go to the hospital. \\
Gladys: -read and write” Errol starts to say that he’s never met anyone who cannot read or write and Gladys cuts him off. Gladys unlike Errol did not have the privilege of being taught yo read and write. Gladys had been sent to work for a family by a terrible mission manager and then had children and then there was the war, so she never had the opportunity Errol did. \\
“Homemade glue” pg. 123 \\
Errol: “What matters is you, not your address...” \\

\subsection{Connection to Land and the Importance of Home}
- When the floods hit Nan, Dolly and Gladys, their house is destroyed and they are forced to move out of their home, the importance of home comes from Gladys when she is determined to “Fix my own house.” Although Gladys does not have the money to fix her home she is sharing the importance of the connection to the land by not leaving and letting the white population win. Dolly also expresses the importance of home when Errol and Dolly argue about leaving the flats to live in the city, Dolly tells Errol that her world is about family rather than leaving them for what Errol calls a real home in the city. \\

“I'm gunna go back and die there”- pg. 125 nan \\
Gladys: “I’m not an interloper – I belong here – this is my land! \\
“a real home is where there are people looking out for each other”- pg. 172 dolly \\
Nan Dear: “Least here we do things our way.” \\
The Rumbalara housing is “white and featureless” and “Anything but loveable.” (Stage directions) \\
Radio: “the most vigorous attempt yet to solve Aboriginal housing” \\
Gladys: ‘[I’m going] to fix my own house.” \\
Errol: “It’s like your family’s from another country or something” \\
Dolly: “Our life isn’t perfect, but like Nan says, it’s ours.” \\
Gladys: “People is what matters.” \\
Gladys: “Mrs Windsor, would you live at Rumbalara? Then why is it good enough for us?” \\

\subsection{Types of Knowledge and Ways of Knowing}
Much of the family’s energy is expended on searching for access to the institutions of the white majority. Gladys never had time for learning, to busy to earn a living to have time to make up for the education that she was denied when she was removed from her home. Gladys has knowledge that far eclipses that of the ‘encyclops boy’. Harrison addresses the sparce of the colonial model of education and the aboriginal knowledge and ways of learning have faced since colonisation. \\



\newpage

\section{Character Analysis}

\subsection{Dolly Banks}
Dolly is the teenage daughter of Gladys and
granddaughter of Nan Dear. Throughout the play, she serves as
a receptacle for her mother’s dreams and her grandmother’s
fears. Gladys fantasies often about what a better life might
look like, but she doesn’t believe that such a life is possible for
her. Instead, she dreams on Dolly’s behalf, and she tries to
encourage Dolly to pursue goals that seem impossible, like
getting a job at a bank that does not want to hire Aboriginal
employees. Nan Dear, meanwhile, views Dolly as vulnerable to
the same discrimination and violence that Nan Dear and
Gladys have experienced. She is especially wary when Dolly
begins seeing Errol, a White boy. Dolly’s relationship with Errol
is complicated by Errol’s unconscious biases, and she
temporarily breaks things off with him when he is unable to put
aside his assumption that White middle-class life is superior to
the more community-based life Dolly wants. Her relationship with Errol also enrages Dolly’s cousin Leon, who sexually
assaults her and leaves her pregnant. Despite the violent origin
of the pregnancy, Dolly loves her baby and names it after Papa
Dear, cementing baby Regina’s status as a beloved member of
the family included in the traditions and legacy it carries.\\

\subsubsection{Analysis}
Dolly Banks is a crucial character in Jane Harrison's play "Rainbow's End," capturing the 
experiences of a young Aboriginal woman during the 1950s, a time rife with changes and challenges 
for her community. The daughter of Gladys Banks and the granddaughter of Nan Dear, Dolly is a 
character who embodies the bridge between the prior generation's traditional views and the new 
generation's aspirations.\\

Throughout the play, Dolly asserts her desires and individuality in a society that would typically 
dismiss them. In one scene, she interacts with Errol, a character symbolising the broader society's 
views, displaying her self-assured yet personable character. When Errol calls her "Miss Banks," she 
promptly corrects him, saying, "Dolly. We’re not too fussed about fancy titles," indicating both 
her lack of pretense. Errol complements her name, but once he learns her full name is Dolores, 
Dolly momentarily moves away, reflecting a subtle consciousness of boundaries and self-respect as 
she navigates her interactions with a non-Aboriginal man.\\

Dolly's aspirations and accomplishments are a point of pride for her mother, Gladys, who notes that 
Dolly is the first in the family to complete her Leaving Certificate, a testament to her 
intelligence and drive. The circumstances in which this achievement is presented, however, bring to 
light the stark realities of racial prejudices they confront. Gladys holds a photo of Dolly, which 
prompts the bank manager to acknowledge Dolly's physical appearance but does not immediately 
understand the reason for Gladys's visit. The scene showcases

Dolly through her mother's eyes—as capable and promising— and reveals the obstacles that she would 
face in achieving her potential in a society that does not readily afford opportunities to 
Aboriginal people.\\

Dolly's character arc within the play reveals the nuanced complexities of identity, as she must 
balance the expectations of her cultural heritage with her own personal ambitions in a society that 
often sees her and her community as lesser. Her actions and decisions throughout the play 
demonstrate the younger generation's determination to both honour their heritage and pursue a 
future they define for themselves.
Through Dolly, Harrison presents themes of cultural identity,
generational change, and the possibilities that lie ahead for
Aboriginal women in a changing world.



Dolly's most decisive act is a refusal. When Errol offers her a life in Melbourne -- a flat, a Kelvinator, the department stores -- he imagines himself as a rescuer, but the offer quietly asks her to give up her family, her community and her Country in exchange for comfort on his terms. Her reply, ``Your world. And you're just assuming that your world is better\ldots This is my place. I'm staying right here with my mum and my nan,'' overturns the assumption at the heart of assimilation: that a white life is self-evidently a better life. In choosing the humpy over the flat, Dolly redefines ``a real home'' as belonging rather than comfort, and asserts a self-determination that the play treats as strength, not sacrifice.\\

Even her scavenging is complicated. At Daish's Paddock tip she adopts a mock-genteel accent -- ``Why yars, I'm after new linoleum'' -- performing the department-store shopper she is denied the chance to be. The joke is self-aware rather than deluded, and the linoleum she carries home genuinely improves the humpy. Harrison lets imagination and dignity work together: the fantasy does not replace real action, it fuels it.\\


\subsection{Gladys Banks}
Gladys is Nan Dear’s middle-aged daughter
and Dolly’s mother. As a child, Gladys was taken away from her
family by the government and put to work as a maid for a White
family. As such, she never learned to read, and the only way she
could make a life for herself was to follow in her mother’s
footsteps as a manual labourer. Despite the hardship she has
faced, Gladys is relentlessly optimistic. She frequently enters
into fantasy sequences where the world is just, her family is
happy, and Dolly has endless possibility before her. Even
outside those fantasy sequences, Gladys believes that Dolly
can live a better life than the generations of women before her.
She is steadfast in this belief, even when Dolly herself doubts it.
Gladys’s optimism proves a vital tool for her eventual activism.
After Dolly accuses Gladys of projecting her own dreams onto
Dolly, Gladys commits herself to working toward making the
world a better place for the sake of herself and her family. She
confronts the city council on the radio about Aboriginal
housing, and she takes over Papa Dear’s speech at a public
presentation of a petition when he fails to arrive on time.
Gladys’s optimism lets her imagine and believe in a better
future, which in turn empowers her to fight for that future.\\

\subsubsection{Analysis}
Gladys Banks represents the aspirations and challenges of an
Aboriginal woman in 1950s Australia. This multi-dimensional
character is caught between her cultural heritage and the
desire for progress and recognition within the broader, often
unfriendly, society.\\

As the daughter of Nan Dear and mother to Dolly, Gladys
mediates between the old ways and new possibilities that her
daughter may have. A scene in the play that captures her
ambitions and protective instincts as a mother is her
confrontation with the bank manager where she presents her
daughter Dolly for a teller's position. Proudly asserting
"She's just
completed her Leaving Certificate—the first in the family—with
real good grades… 'N' top of her class in algebra,"
Gladys is met
not with recognition but indifference, with the bank manager
pouring himself tea but pointedly ignoring her. The scene
captures the societal barriers that Gladys and her daughter
face, as well as Gladys's determination to push against them.\\

Gladys's agency is further highlighted in her firm stance on her
daughter's
autonomy
and social life. Despite Nan Dear's
apprehensions, Gladys declares decisively,
"I’ll make the
decisions regarding Dolores thank you. She’s going to the ball.
And Errol Fisher is walking her home."
This reveals her resolve
and the assertion of her role as a parent, as well as the
modernising viewpoint she holds compared to Nan Dear, who is
more rooted in traditionalist concerns. Moreover, Gladys's
readiness to confront and challenge Nan Dear's objections
demonstrates her strong will as she endeavours to protect her
children's rights to forge their own paths.\\

Gladys's complexity extends into her views on housing and
respect. When discussing the Queen's visit and the humiliation of having their poor living conditions concealed behind hessian curtains, Gladys expresses upset not just over the cosmetic cover-ups but the deeper issue of decent housing and respect for her community. This highlights her awareness of the broader socio-political context and her aspiration for improved conditions for her people. Yet, she also values the autonomy they possess in their current living situation, preferring the freedom they enjoy at the shanty on the river to the restrictions experienced in more controlled environments like Cummeragunja, an Aboriginal settlement she and Nan Dear remember with claustrophobia.\\

Gladys Banks embodies the transition between tradition and progress, acting as both a nurturer and a pioneer. Her interactions with others showcase her resolve to advance her family's status despite the systemic challenges imposed by aw orld that often disregards their aspirations and dignity.



Gladys's quiet arc from shame to voice binds the whole play together. Her illiteracy -- a legacy of the mission schooling she was denied -- is first a secret managed with jam tins and pride, then something she confronts directly, learning to read so that no institution can again decide things about her family over her head. The arc completes at the Melbourne meeting: when Papa Dear fails to arrive, Gladys steps up unplanned, begins haltingly with the petition's borrowed words, then abandons them for her own -- ``Why humbly? We've been humble too long\ldots we demand to be heard.'' The movement from reading someone else's sentence to speaking her own is Harrison's clearest image of empowerment through voice.\\

Tellingly, the crowd's tribute -- that she is ``Papa Dear's daughter all right'' -- hands her courage back to a patriarch, showing how even a triumphant woman is read through the men around her. And when her earlier interjection at the council is carried into the humpy by radio, Harrison literalises the point: a woman may be shut out of the room, but her voice can travel beyond it.\\

\subsection{Nan Dear}
Nan Dear is the matriarch of the family at the
centre of the play. She is Gladys’s mother and Dolly’s
grandmother, and she takes her duty of care to them seriously.
Nan Dear is a bastion of tradition, teaching Dolly about her
family history and aspects of their culture that she doesn’t
learn at school. She has suffered under institutions of racism
and ideologies of White supremacy, especially at the Aboriginal
reserve in Cummeragunja, where Nan Dear was abused and
her daughter Gladys was taken away from her. As such, Nan
Dear has little faith that the world can improve, and she doesn’t
understand Gladys and Dolly’s attempts to improve it. She
especially fears for Dolly when Dolly begins a romantic
relationship with Errol, a White boy with the same last name as
a White man who raped Nan Dear years ago. Nan Dear works
to separate Dolly and Errol, but at the end of the play she
realises that Errol is a considerate and respectful young man.
When she finds out for certain that Errol is not related to her
attacker, she gives him her blessing and encourages Dolly to
marry him, demonstrating that even someone as traditional as
Nan Dear can grow and change.

\subsubsection{Analysis}
Born by the Murray River as her birth certificate states, Nan Dear
is steeped in her Aboriginal heritage and connections to her
traditional lands. Her identity is closely tied to the Murray River
and the land from which she was
displaced
, specifically from
Cummeragunja, marking a bitter and forcible removal from her
home. This
dislocation
is a source of deep-seated resentment
and shapes her worldview. Despite the trials she has faced, Nan
Dear is characterised by a deep sense of duty and
pragmatism
.
She is concerned with everyday survival and the practical needs
of her family—working to put food on the table and keeping her
family safe.\\

As the matriarch, Nan Dear's relationships with her daughter,Gladys, and granddaughter, Dolly, are a mix of affection and stern guidance. She is quick to offer wisdom to Dolly, such as reminding her of the importance of knowing family connections to avoid unwitting relationships with close relatives. Nan Dear's knowledge of family history is implied to be extensive, though she may feign ignorance when it suits her, suggesting a complicated relationship with the past.\\

Her cultural beliefs surface subtly through her reluctance to embrace the visiting monarchy, mocking the idea of needing white gloves for a royal visit and declaring, "She's not my queen". This resistance to the monarch's visit and the detachment she expresses towards Gladys' interest in royalty reveals an underlying rejection of colonial authority and ap reference for the Aboriginal ways of life.\\

Nan Dear's resilience is not just a personal trait but serves as symbolic of her community's endurance. Statements like "Born[by the Murray River] and by crikey I'm gunna go back and die there" demonstrate her strong will to return to her roots despite the adversity faced. This return, for her, is associated with embracing the simplicity and beauty of her origins, like having a feed of swan eggs before she passes, signifying a connection to her birthplace. Her experience of displacement is steeped in historical injustices, as evidenced by her bitter recollection of being forced to leave Cummeragunja. This comment on forced relocation exposes the reality of Aboriginal displacement and the loss of home and culture.\\

These observations offer a critical perspective on the consequences of colonial policies for Aboriginal people.Nan Dear is a fiercely independent individual whose character lends gravity and authenticity to the narrative of "Rainbow's End." Her dialogue is delivered with an undercurrent of knowing defiance, revealing a woman shaped by adversity yet undiminished in spirit. She represents a generation of Aborigina lpeople who, despite systematic disenfranchisement,maintained their dignity and cultural identity. Through her words and actions, Nan Dear confronts the audience with the realities of her community's history, while also displaying an indomitable resolve to withstand and challenge these realities.\\


Two late revelations deepen Nan beyond the stern matriarch. The first is comic in form and serious in meaning: her faked coughing fit to keep Dolly from the dance, exposed by the stage direction that she ``remembers to cough again''. The manipulation is real, but so is the love underneath it -- her obstruction of Errol is later revealed to be grounded not in simple prejudice but in her own experience of violence at the hands of a white man, which recasts her control as embodied fear. Harrison compounds the irony cruelly: the assailant's surname was Fisher, yet Errol's family, originally ``Vischer'', is no relation. The secret that might have spared Dolly was kept in shame across a generation, and the play's judgement falls less on Nan than on the silence itself -- a silence colonial violence taught families to keep.\\

The second revelation is wordless. Publicly, Nan scorns the encyclopedias -- ``Ask him what it says about the Aborigines, eh?'' -- but once Gladys is out of the room, the stage directions show her taking down the latest volume to read alone. The gesture complicates any neat opposition between Western and Indigenous knowledge: curiosity is universal; what colonialism rationed was access, not desire.\\


\subsection{Errol Fisher}
Errol Fisher is a young White encyclopedia
salesman. He encounters Nan Dear, Gladys, and Dolly by
accident, after taking a wrong turn on his route selling
encyclopedias. Errol falls in love with Dolly and comes to
admire her family. Gladys demonstrates a strength that Errol
has never seen in his own mother, who is subservient to Errol’s father. The affection that Nan Dear and Gladys show Dolly is
also alien to Errol, who has grown up with a strict and unfeeling
father. Though he earnestly loves Dolly, he offends her with his
ignorance of her culture and his assumption that his way of life
must be superior to hers. He assumes she would be eager to
leave her family behind to be his wife in the city, and he
struggles to understand her insistence that she must remain
with her loved ones. Ultimately, though, Errol comes to respect
what Dolly wants for her future, and he promises to live any
kind of life she wants as long as they can do so together.\\

\subsubsection{Analysis}
Fisher
Errol Fisher emerges in Jane Harrison's "Rainbow's End" as a significant, though less central, 
character, who embodies the intersection of the Aboriginal community's life and the broader 
Australian society during the 1950s. Portrayed initially as somewhat awkward and uncertain, he is 
introduced to the audience as he nearly falls off his bike upon encountering Dolly, signaling an 
immediate spark of attraction that complicates his interactions with the family.\\

This initial interaction with Dolly sets the stage for Errol's character development. The exchange 
with Gladys further illustrates his social position and perspective. When he tentatively approaches 
Gladys as she is chopping wood, mistaking her for a man, he immediately demonstrates cultural 
ignorance and nervousness. Apologising and stating his name, "Sorry er ma’am. My name is Errol 
Fisher," he provides a contrast to the women's more assertive presence. This contrast is emblematic 
of the interplay between white Australians and Aboriginal people at the time—a dynamic of presumed 
authority and misunderstanding from the former, contrasted with an under-recognised competence and 
resilience from the latter.\\

Errol's intentions are primarily associated with business—selling encyclopedias as revealed in his 
conversation with Gladys—yet this transactional engagement serves as a backdrop to more complex 
intercultural exchanges. Gladys and Errol mention the encyclopedias simultaneously, indicating 
their previous interactions and the economic realities pressing on Gladys, who states, "We… I… 
won’t be needing them any more. All the shillings go into the meter box now," highlighting not only 
financial struggles but also the prioritisation of practical
necessities over educational luxuries.\\

His interactions with Dolly in the play are coloured with romantic interest, albeit fraught with 
the implications of such a relationship in a racially divided society. He compliments her name and 
attempts to bridge their cultural gap, however awkwardly. Since Errol also represents other 
characters in the play, his role is fundamental in bringing out the perspectives and reactions of 
Nan Dear, Gladys, and Dolly and propelling their narratives forward. Errol's presence and character 
serve as a mirror that reflects and refracts the complexities of identity and interpersonal 
relationships within the stratified, mid-20th
century Australia Australian society.



\subsection{Leon (Cousin)}
Leon, often referred to as simply “Cousin,” is
Dolly’s cousin. He is aggressive and often drunk, and he
disapproves of Dolly’s relationship with Errol. As a twisted
form of punishment for that relationship, Leon rapes Dolly,
resulting in her pregnancy. Leon’s assault of Dolly speaks to the
reality of gendered violence within marginalised communities.\\

\subsection{Papa Dear}
Papa Dear is Nan Dear’s husband, Gladys’s father,
and Dolly’s grandfather. Though he is not Gladys’s biological
father, since Gladys was conceived from Nan Dear’s rape by a
White man, Nan Dear considers him Gladys’s father. Papa Dear
is an activist for Aboriginal rights, and his work takes him all
over the country, which means he is rarely there for his family.
He never appears in person in the play; he is only mentioned in
dialogue and imagined in fantasy sequences. His absence
leaves Nan Dear, Gladys, and Dolly alone to create an
exclusively feminine space in their home.\\

\subsubsection{Analysis}
Dear
Papa Dear, as depicted in Jane Harrison's "Rainbow's End," is a paternal figure whose presence is 
felt throughout the narrative despite being physically absent for much of the play. His character 
is crafted through the recollections and expectations of the other characters, particularly Nan 
Dear and Gladys. Papa Dear's role is largely communicated through anecdotes and his influence on 
the family's structure and values.\\

Mentioned by Nan Dear as "busy doing good work. God’s work and hard work," Papa Dear is portrayed 
as a figure devoted to the betterment of the Aboriginal community. His dedication to activism and 
community work is further highlighted when Gladys recalls that he is "in Western Australia. Touring 
the communities there," and has been featured "in the newspaper and all," suggesting his actions 
and efforts are recognised and appreciated beyond their immediate surroundings.\\


Papa Dear's dream visitations reinforce his endearing place within the family, despite his absence. 
In an imagined scene, he is seen wearing "an old-fashioned hat and coat," hinting at a sense of 
nostalgia and timelessness associated with him. His brief entrance to kiss Gladys on the top of the 
head showcases a symbolic and loving connection with his family, despite his lack of direct 
intervention in their current struggles.\\

The character of Papa Dear is particularly important to Dolly during a vulnerable moment when she 
reflects on her family tree, considering naming her unborn child Reg or Regina, "after Papa Dear." 
This indicates an enduring legacy and respect for him within the family, especially for his 
commitment to their people.\\

While Papa Dear does not have a speaking role or much direct action, his presence resonates with 
the family and the audience, capturing the spirit of a generation of Aboriginal individuals who 
strived to improve conditions for their communities. Through the memories shared and the respect 
expressed by the characters, Papa Dear stands as a symbol of resilience, hope, and the ongoing 
struggle for activism and recognition within the
Aboriginal community of the time.



\subsection{The Inspector}
The Inspector is a White man employed by the
Australian government to inspect the living conditions of
Aboriginal Australians, which includes the family’s humpy. As a
representative of the Aboriginal Protection Board, the
Inspector regulates and controls the lives of Aboriginal
Australians, and that power extends to taking Aboriginal
children away from their family. The Board took Gladys away
from her family when she was a child, and she and Nan Dear
are terrified that the Inspector will take Dolly away from them.\\

\subsubsection{Analysis}
In the play "Rainbow's End," the inspector is a minor but telling character that presents a 
snapshot of the authority and surveillance that Aboriginal people endured in 1950s Australia. He 
appears in a scene where Nan and Gladys are inside the humpy, anxiously straightening everything, 
which sets a tone of tension and unease. The inspector serves as a symbol of the government 
oversight and control that loomed over the lives of the characters.\\

Paying a visit to evaluate their living conditions, the inspector enters with a superficially 
pleasant demeanor that thinly veils the oppressive power dynamics at play. He greets the efforts of 
Nan and Gladys with an offhand compliment, "I say crocheted pillow shams. Such beautiful work!" 
which Gladys attributes to her mother's handiwork. His interaction is polite yet distant, making 
notes while seemingly ignoring the weight of his presence and authority in the humpy.\\

As the scene unfolds, it becomes clear that his visit is not a casual or friendly one, but an 
inspection that could have severe implications for the family. The anxiety Nan and Gladys display 
is indicative of their vulnerability to the whims of the inspector and, by extension, the colonial 
power structure. When the inspector inquires after a Mr. Banks, his question underscores the 
traditional views of family structure and the expectation of a male head of household, hinting at 
patriarchal values as well.\\

The inspector's questions and his mode of taking notes encapsulate the intrusion of government into 
personal spaces and lives. This scrutiny implies potential consequences if living conditions or 
behaviors do not meet certain standards, though these criteria are not made explicit. This scene 
heightens the awareness of the audience to the systematic and normalised regulations that govern 
the lives of Aboriginal characters in the play, bringing the broader context of racial disparity 
and control into focus through the actions and attitudes of the inspector.



The Inspector's language deserves close reading, because Harrison gives him the play's most chilling sentence and delivers it in the mildest of tones: ``The Aborigine needs to be absorbed into the community. But how can he be absorbed until he learns to live like us?'' The verb ``absorbed'' does the ideological work -- assimilation imagined as dissolution, help defined as the disappearance of a people. His compliment ``Your whites are so white!'' measures human worth in laundry, and his marvelling ``I don't know how you do it'' converts the women's obvious skill into a miracle, which is its own condescension. He is no snarling villain; he is a courteous functionary, and the courtesy is the point -- Harrison locates racism not in individual malice but in the machinery a polite man can operate with a clear conscience.\\

The visit also fractures the family in front of him: asked about Dolly's future, Nan, Gladys and Dolly give three different answers -- the laundry, a bookkeeper, the hospital -- a small dramatic irony that maps three different accommodations with a world that limits them.\\

\subsection{Mr. Coody}
Mr. Coody is the rent collector at Rumbalara. He
imposes strict and invasive rules on his tenants, which includes
threatening to evict Dolly and her family due to Dolly’s
pregnancy. His intrusion into the women’s lives reflects that
although Rumbalara is an Aboriginal housing district, it is not a
safe haven or a stable place of community because it is run by
White men.\\

\subsection{Bank Manager}
The Bank Manager works at the bank where
Gladys wants Dolly to find a job. He resists the idea of hiring an
Aboriginal worker, making excuses as to why Dolly would not
be qualified for the job, but Gladys finally stands up to him and
he offers Dolly an opportunity to apply.\\

\subsubsection{Analysis}
Harrison uses the The Bank Manager to showcase the systemic barriers faced by the Aboriginal 
community in terms of employment and societal inclusion. His interaction with Gladys Banks in her 
attempt to secure a job for her daughter, Dolly, is telling of the racially prejudiced attitudes of 
the time. Gladys, with determination, engages the Bank Manager and puts forward Dolly's 
qualifications for a teller's position: "But she’s just completed her Leaving Certificate—the first 
in the family." Despite the qualifications presented, the Bank Manager dismisses the notion, 
indicating a lack of openness to hiring an Aboriginal person.\\

The Bank Manager focuses on Dolly's appearance rather than her capabilities when he says, "Yes. 
Very pretty face," after examining a photograph of her. This superficial assessment underscores the 
racial objectification and segregation of Aboriginal people. His concern regarding how Dolly would 
"fit in" at the bank delineates the invisible boundaries placed by a society reluctant to accept 
Aboriginal individuals into professional roles. "In a job such as this reliability is important," 
he states, questioning how Dolly would commute to work, despite evidence of her punctuality and 
capability presented by Gladys.\\

The role of the Bank Manager, while not overly complex, provides a window into both the 
institutionalised racism of the period and the socioeconomic status of Aboriginal families. His 
patronising attitude when offering a tin money bank to Gladys, likening it to "jam tins," reflects 
a condescension towards her and her family's means of saving money, which is deeply ingrained in 
the Bank Manager’s and wider society’s view of Aboriginal people.\\

The character subtly represents the gatekeeping that denies Aboriginal people equal opportunities, 
serving as a juxtaposition to Gladys’s and Dolly's aspirations for advancement. Through  this 
interaction, Harrison critiques the wider Australian society's views on race and the illusory concept of a fair, meritocratic
system, thereby reinforcing the themes of racism and
discrimination that are central to the play's message.\\



The bank scene achieves its discrimination almost entirely through manners. The stage business does the work: the manager pours tea for himself and offers Gladys none -- in an era when sharing tea ritualised social equality, the omission quietly declares that no equality is on offer. His words stay cordial while his meaning does not: Dolly's excellent results are praised, then dissolved into doubts about how she would ``fit in''. And when Gladys cannot write, his smoothing courtesy -- ``Fountain pens are a little tricky, aren't they?'' -- performs kindness while reclassifying the results of educational exclusion as a small mechanical mishap. Harrison's point is exact: this is discrimination without a single hostile word, prejudice wearing the costume of politeness, which is precisely what makes it so difficult to fight.\\

\subsection{Ester}
Ester is a neighbour of the family who represents the
most vulnerable members of the Aboriginal community. She is regularly abused by her White husband, and her sons are taken
away by the Inspector.\\

\subsection{Regina}
Regina is Dolly’s infant daughter. Her birth marks the
beginning of the fourth generation of women in the family.\\



\newpage

\newpage

\section{Character and Relationships}

\subsection{Dolly Banks}
\textbf{Key Quote:} ``And as far as what you’re offering … no thank you. This is my place. I’m staying right here with my mum and my nan.'' (to Errol, p.172)\\

Dolly’s non-reaction to her ruined magazines indicates a practised stoicism in the face of misfortune. She is studious, completing her Leaving Certificate with all ‘As and Bs’ and coming ‘top of her class in algebra’ (p.163). Remarkably ‘pretty’, she is also ‘clever’, ‘kind’, ‘modest’, and ‘straightforward’ (p.190). Dolly sometimes allows herself to dream, like Gladys winning \textit{Pick-A-Box} (p.130) or winning the Miss Mooroopna-Shepparton title at the ball (p.166). She dreams fancifully of becoming an actress or model, yet practically of a career in nursing (p.142).\\

She is unhappily aware that she might become a fruit picker like her mother and grandmother (p.141). Dolly’s exposure to racial discrimination — with ‘stones’ thrown at her and ‘snide remarks’ (p.162) overheard in the street — makes her pessimistic about her future. She has a close and loving relationship with Nan, whose overprotectiveness can sometimes be tedious. Yet their closeness and trust alleviate Dolly’s distress after she becomes pregnant, when Nan confides that Gladys was conceived in similarly traumatic circumstances. The sharing of this dark ‘secret’ brings them both comfort.\\

\subsubsection{Analysis}
Dolly Banks embodies the tensions between youthful hope and social limitation. Harrison portrays her as the bridge between aspiration and resignation — a young woman educated enough to imagine a different future but aware of the barriers imposed by race and class. Her determination to remain with her family, rather than abandon them for the city life Errol proposes, symbolises her rejection of assimilation and her embrace of identity. Dolly’s relationship with Nan and Gladys represents the generational continuum of Aboriginal womanhood: endurance, wisdom, and resistance to erasure.\\

\subsection{Nan Dear}
\textbf{Key Quote:} ``Ask him what it says in that there encyclops about the Aboriginies, eh?'' (to Gladys, p.135)\\

Nan Dear is often the cynical voice of experience, reminding Gladys that she is a lowly bean picker unworthy of the Queen’s notice (p.124). Unafraid to speak out, she angrily confronts the Rent Collector for criticising their supposed ‘flouting’ of rules, throwing eggs at him (pp.185–6). Capable and energetic, she works as a fruit picker for the canning factory, tackles physical labour, and can make a stew out of nettles when money is short.\\

She is a loving and protective grandmother who knows the danger of being ‘[taken] advantage of’ by ‘a whitefella’ (p.187). Set in her ways, she mistrusts outsiders — seeing Bob Dyer as ‘one of them flash Yanks’ (p.131) and regarding Errol as a ‘snake-oil salesman’ (p.136). Initially doubtful, she later recognises Errol’s goodness and advises Dolly to marry him ‘before someone else does’ (p.198). Nan also clashes with Gladys over Dolly’s upbringing, criticising her monarchist sympathies and sceptical of education’s value. Yet when she hears Gladys speak up on the radio, she dances with pride around the room.\\

\subsubsection{Analysis}
Nan Dear represents endurance, tradition, and the persistence of Aboriginal matriarchal strength. Her scepticism toward authority and colonial structures masks deep compassion and wisdom. Through her character, Harrison celebrates the resilience of older generations who, despite trauma and displacement, preserve dignity and family unity. Nan’s humour, practicality, and protectiveness make her both a guardian of the past and an emblem of cultural survival.\\

\subsection{Gladys Banks}
\textbf{Key Quote:} ``We demand proper schooling. And not just for us. The white people too – they need to be educated about us, and our ways.'' (p.197)\\

Initially a fervent royalist, Gladys goes to extraordinary lengths to glimpse the Queen, only to be blocked by hessian screens hiding Indigenous homes from sight (p.127). Clever and capable, she answers all questions correctly on the quiz show \textit{Pick-A-Box} but knows a ‘black’ contestant would be unwelcome (pp.129–130). She manages her household with efficiency and intelligence — chopping wood, skinning rabbits (p.129), and managing coins in jam tins — making the Bank Manager’s condescension (p.164) absurd.\\

Having been sent away as a child (p.184) and denied an education, Gladys feels the sting of illiteracy and wants her daughter to succeed where she could not. Her daydreams of Dolly in academic regalia or being hugged by the Queen (pp.126, 134) express both hope and escapism. When Dolly challenges her to ‘fix’ her own life (p.179), Gladys redirects her energy toward activism — storming a council meeting and demanding fair treatment for Aboriginal families (p.181).\\

\subsubsection{Analysis}
Gladys is a figure of transformation — from dreamer to activist. Her evolution from naive royalist to outspoken advocate mirrors the awakening of Indigenous political consciousness. She bridges the domestic and political spheres, proving that even an ordinary woman’s voice can challenge entrenched injustice. Harrison uses her courage to illustrate the theme of self-determination and the belief that change begins within the individual.\\

\subsection{Errol Fisher}
\textbf{Key Quote:} ``This is my first presentation. Well, the first one I got all the way through, anyhow.'' (p.137)\\

Errol is awkward, well-meaning, and naïve — a White encyclopedia salesman who stumbles into the Flats by mistake. His clumsy sales pitch and fearful reaction to a harmless carpet snake (p.145) reveal both humour and ignorance. Despite his flaws, his sincerity sets him apart from other male characters. His relationship with Dolly follows the familiar rhythm of romance fiction — a high at the ball, when he rescues her from Leon (p.167), and a low when he offends her by suggesting she abandon her family (p.171). Misunderstandings later see him wrongly blamed for Dolly’s assault, yet his determination to return and embrace her world reflects genuine love.\\

\subsubsection{Analysis}
Errol’s character illustrates the possibility of cross-cultural understanding, albeit imperfect and halting. Through him, Harrison explores the limits of good intentions within systems of inequality. Errol’s willingness to change, to “move to Mooroopna,” and adopt Indigenous expressions (‘snotty googles’, p.191) suggests that reconciliation begins not with assimilation but humility and empathy.\\

\subsection{Leon (Cousin)}
\textbf{Key Quote:} ``She’s no lady, she’s just a little –'' (to Errol, referring to Dolly, p.167)\\

Leon associates with the rough ‘cork-tree lads’ who treat women as objects. His violence toward Dolly reflects both personal corruption and societal neglect. Nan’s comment that the boys have ‘had it hard’ (p.139) hints at systemic roots — poor education, limited work, and widespread alcoholism. Mentioned as a ‘motherless child’ (p.131), Leon represents a generation damaged by dispossession and lack of support.\\

\subsubsection{Analysis}
Leon embodies the tragic consequences of internalised oppression and the failures of state policy. While his assault on Dolly is unforgivable, his existence indicts the colonial system that left many young Indigenous men alienated, impoverished, and angry. He is both victimiser and victim — a figure who underscores the intergenerational damage caused by racism and inequality.\\

\subsection{Papa Dear}
\textbf{Key Quote:} ``That’s how important [the Queen] is – getting a meeting with the busiest Aboriginal around!'' (Nan, referring to Papa Dear, p.128)\\

Papa Dear never appears on stage but looms large through others’ stories. A pastor who does ‘good work’ (p.143) and meets the Queen, he has gained status in white society, yet his activism yields little improvement for his people. His absence during family crises — floods, pregnancy, and poverty — highlights the emotional toll of political duty and the limits of symbolic representation.\\

\subsubsection{Analysis}
Papa Dear symbolises a generation of Aboriginal leaders who sought reform through institutional acceptance rather than confrontation. While his work is well-intentioned, Harrison implies that activism divorced from community risks perpetuating absence rather than alleviating suffering. He is a figure of both hope and disillusionment.\\

\subsection{The Inspector}
The Inspector is a smug bureaucrat who prides himself on “bringing change” (p.152), blind to the harm of assimilation. His hollow encouragement that the women should “take leadership” reveals a profound ignorance of their reality. His officiousness exemplifies paternalism — well-dressed oppression masquerading as progress.\\

\subsubsection{Analysis}
The Inspector represents state surveillance and control — the polite face of systemic violence. Harrison uses his character to expose how government policy framed Aboriginal suffering as a problem to be managed rather than a wound to be healed.\\

\subsection{Mr. Coody (The Rent Collector)}
\textbf{Key Quote:} ``Your arrangements will need to be re-evaluated, with the impending new arrival.'' (p.185)\\

Mr. Coody’s condescension and moral judgment expose the hypocrisy of colonial authority. He condemns Dolly’s pregnancy and questions Nan’s decency, despite representing a system that denies Aboriginal families adequate housing. His “re-evaluation” of their tenancy embodies the state’s power to intrude upon private life.\\

\subsubsection{Analysis}
Through Mr. Coody, Harrison critiques bureaucratic arrogance and the pretence of moral superiority that underpinned racist policies. Nan’s defiant act — throwing eggs at him — becomes symbolic rebellion, an assertion of dignity in the face of systemic degradation.\\

\subsection{The Bank Manager}
Arrogant and self-satisfied, the Bank Manager embodies white patriarchal privilege. He refuses Gladys’s request for her daughter’s employment, patronises her with a vocabulary lesson, and reinforces racist stereotypes about Aboriginal unreliability (p.164).\\

\subsubsection{Analysis}
He stands as a caricature of institutional prejudice — polite yet poisonous. Harrison uses him to illustrate how everyday racism operates behind the veneer of civility, undermining merit and silencing ambition.\\

\subsection{Nancy Woolthorpe}
\textbf{Key Quote:} ``Love your dress, Dolly. Love the fabric. My mother quite liked it too. When it was our sunroom curtains.'' (p.166)\\

Nancy is the archetypal ‘mean girl’ — wealthy, pretty, and cruel. She belittles Dolly at the ball, symbolising class snobbery and racial arrogance. Yet her hollow luxury contrasts sharply with the warmth and humour of Dolly’s family, suggesting that material comfort cannot replace love or community.\\

\subsubsection{Analysis}
Nancy Woolthorpe serves as a foil to Dolly — shallow, privileged, and spiritually impoverished. Harrison’s satire of her behaviour underscores how racism is perpetuated not only through policy but through everyday social exclusion.\\

\subsection{Ester and Her Boys}
Ester is a young Aboriginal mother living at Rumbalara with her three mischievous sons, and she is pregnant again. Her ‘whitefella husband’ (p.132) has given her a black eye, highlighting the cycles of violence and instability experienced by many Indigenous families of the time. When the Inspector visits, Nan fears that Ester’s children will be taken away (p.153) — and indeed, this fear becomes reality. Ester’s boys are removed, ostensibly for “assimilation”, but in truth to be institutionalised or adopted, separated from their family, community, and culture.\\

\subsubsection{Analysis}
Ester’s brief appearance is a powerful allusion to the Stolen Generations. Her character encapsulates the vulnerability of Aboriginal mothers under the Assimilation Policy, whose children were taken “for their own good.” Harrison’s inclusion of Ester and her boys exposes the cruelty and hypocrisy of a government that dismantled families under the guise of progress. The boys’ removal functions as both a narrative tragedy and a historical indictment, reminding the audience that these policies perpetuated trauma across generations.\\


\subsection{The Policeman}
The local policeman (the “Jungi”) appears following the devastating flood that destroys the women’s home and possessions. He exhibits little empathy for Dolly’s distress, and his irritation when she disobeys him reveals his insensitivity. Although he promises to protect their belongings, placing the encyclopaedias “high and dry,” he instead leaves them on the flooded ground (p.174), symbolising negligence and institutional indifference.\\

\subsubsection{Analysis}
The policeman’s character demonstrates the absence of meaningful protection for Aboriginal people from authorities. His behaviour reflects the wider neglect of Indigenous welfare by institutions meant to safeguard all citizens. Harrison uses his actions as a metaphor for systemic disregard — the literal and figurative abandonment of Aboriginal people during times of crisis.\\


\subsection{The Shire Councillors}
The Shire Councillors appear through a radio broadcast, discussing the “shanties” on Crown Land and claiming that Aboriginal families pose a sanitary risk to the town (p.180). Their rhetoric is filled with euphemisms and coded prejudice: they describe the Indigenous community as an “itinerant minority” while ignoring that these families are long-established residents. When Gladys interrupts their meeting to protest, the councillors order her ejection, branding her an “interloper” (p.181).\\

\subsubsection{Analysis}
The councillors epitomise bureaucratic racism and the erasure of Aboriginal voices from public discourse. Their paternalistic tone and false concern for “public health” disguise the real motive of displacement and segregation. Gladys’s interruption becomes an act of defiance against a political structure that silences Indigenous women. Harrison’s portrayal of this scene transforms a moment of exclusion into one of empowerment, as Gladys reclaims her right to be heard.\\




\newpage
\section{Scene Summary and Analysis}

\subsection{Prologue: Aftermath}
\subsubsection{Summary}
In 1950s Australia, Nan Dear and Gladys rebuild their humpy
(another word for a bush hut) after a flood. Gladys is Nan
Dear’s middle-aged daughter and the mother of Dolly, a
teenager. The women are Aboriginal Australians. Nan Dear and
Gladys tear up Dolly’s magazines to repair the walls. Dolly
returns home from school and almost protests when she sees
her ruined magazines, but ultimately she says nothing,
resigned. Nan Dear and Gladys hug Dolly.\\

\subsubsection{Analysis}
The play introduces readers and viewers to its main characters as
they rebuild their home. This establishes the home as something
that must be tended to and maintained, and it also highlights the
often invisible domestic labor that goes into maintaining a home.
That labor requires sacrifices, like Dolly’s magazines, but the women
are willing to put their personal wants aside for the sake of the
family and their home.\\

\subsection{Act 1, Scene 1 (A)}
\subsubsection{Summary}
Gladys listens to Queen Elizabeth on the radio and prepares to
travel to town to see Queen Elizabeth, the first reigning
monarch to visit Australia. Nan Dear dislikes the monarchy and
pokes fun at Gladys for being a “loyal subject.”\\

Dolly asks for help on a family tree she is making for school.
Nan Dear is bewildered that Dolly is only supposed to list
direct relatives and exclude her cousins. She evades Dolly’s
questions about her great-grandparents and reminisces about
her home by Murray River, which she was forced to leave.
Gladys continues to make herself presentable for the Queen,
and she fantasises about presenting a bouquet of flowers to
Elizabeth before finding herself holding a handful of weeds.\\

\subsubsection{Analysis}
Nan Dear disdains the British monarchy as the ultimate colonisers,
who are responsible for the oppression of Aboriginal Australians on
their own land. Gladys, on the other hand, is more assimilated into
White society and thus shares White Australians’ excitement to see
their queen. Nan Dear criticises her for this loyalty to a violently
racist government.\\

The school’s assumption that only direct relatives are significant
enough to include on a family tree demonstrates how conceptions
of family depend on culture. The school system is built around
White Australian culture, which does not place the same value on
extended family as Aboriginal culture. Meanwhile, Gladys loses
herself in the fantasy of meeting Queen Elizabeth, and the sharp
contrast between her imagined flowers and the reality of weeds
suggests that Gladys indulges in fantasies as a form of escapism.\\

\subsection{Act 1, Scene 1 (B)}
\subsubsection{Summary}
Gladys returns home after the Queen’s visit, exhausted. She
explains to Nan Dear and Dolly that her taxi never arrived,
leaving her to walk along the road hoping for a glimpse of the
Queen. Gladys failed to see the Queen, though, because the
roads have been lined with burlap to prevent Elizabeth from
seeing the humpies along the riverbank.\\

Gladys complains about the state of their housing, but Nan
Dear insists that life is better than it was in the last days at
Cummeragunja by the Murray River. Nan Dear tells Gladys that
Papa Dear had a meeting with the Queen to discuss improving
life for Aboriginal people.\\

\subsubsection{Analysis}
The burlap-lined streets are a deliberate effort to keep Queen
Elizabeth from seeing Aboriginal Australians and vice versa, and it
indicates that the Australian government does not regard Aboriginal
Australians as true Australian citizens. Through this racist mindset,
Aboriginal Australians and their humpies are seen as an
embarrassing reminder of the past on an otherwise modern and
civilised land.\\

Papa Dear’s role as a travelling activist introduces a gendered
element to the racism that the family faces. Gladys, as a woman, is
tied to her duties at home, so she can’t travel the country to fight for
her people’s rights the way her father can.\\

\subsection{Act 1, Scene 2 (A)}
\subsubsection{Summary}
Gladys listens to a radio quiz show and answers the questions
along with the contestants, getting every one of them correct
even as the contestants make mistakes. Dolly suggests Gladys
try to become a contestant on the show, but Gladys laughs off
the idea of an Aboriginal contestant.\\

Nan Dear asks Dolly about a boy she was talking to. Dolly
frustratedly explains that she won’t be marrying him because
he is her cousin, as are most boys her age in the community.
Gladys insists that Dolly won’t marry anyone until she finishes
her studies, which prompts Dolly to ask how old Nan Dear was
when she married. Nan Dear doesn’t answer. Dolly leaves to
study, and Nan Dear questions Gladys for “put[ting] ideas” into
Dolly’s head about the possibilities for her life.\\

Nan Dear leaves to tend to a neighbour, Ester, who is pregnant
and injured after an altercation with her abusive White
husband. Gladys confesses she didn’t know about Ester’s
predicament, and Nan Dear recommends Gladys pay less
attention to the radio and more attention to the world around
her.\\

\subsubsection{Analysis}
Though Gladys fantasies about a better life, she is also well aware
of the limitations she faces. She is familiar enough with racism to
know that the radio quiz show would never accept her, despite her
intelligence.\\

Nan Dear believes optimism is dangerous, since she is convinced
that the world will never improve enough for Dolly to have the kind
of life Gladys dreams for her. At the same time, Nan Dear’s inquiry
into the boys Dolly spends time with highlights that she is
concerned about Dolly’s future in her own way.\\

Ester’s situation demonstrates how gendered and racial dynamics
play out within the home. As an Aboriginal woman, Ester has
almost no power against the White man she is married to, with her
only option being to turn to the members of her community for help.\\

\subsection{Act 1 Scene 2 (B)}
\subsubsection{Summary}
A young White man, Errol Fisher, gets lost and ends up at the
humpy. He encounters Dolly, and though the two don’t speak
to each other, mutual attraction sparks between them. Errol
then approaches Gladys and tries to sell her the encyclopedias
he is peddling. Gladys turns him away at first, but she takes pity
on Errol and allows him to present his rehearsed pitch. As Errol
presents the encyclopedias, Gladys becomes interested in
them, especially when Errol notes that a previous radio quiz
show winner succeeded by studying an encyclopedia. He adds
that encyclopedias are good for school-aged children, and
Gladys has a fantasy of Dolly graduating in a cap and gown.\\

Nan Dear enters and is skeptical of Errol and his encyclopedias.
She asks what the books say about Aboriginal people, and Errol
doesn’t know the answer. Gladys considers buying a set of
encyclopedias, but Nan Dear stops her as Dolly comes in and
helps Errol sort out his map. She explains that he should go to
the White neighbourhood to sell his encyclopedias, and Gladys
and Nan Dear agree. Errol departs, but he leaves his calling
card, and Gladys is charmed by his manners and his
descriptions of the encyclopedia.\\

As Errol prepares to leave, Dolly asks him about his home in
Melbourne and teases him for coming out to her family’s
humpy in The Flats. She is confused and uncertain whether he
returns her attraction to him. She helps him retrieve his bicycle
from some local little boys, then returns to her family. Nan
Dear advises Dolly to be careful what kind of men she spends
time with. Gladys considers Errol’s card.\\

\subsubsection{Analysis}
Errol is out of place in the humpies among The Flats, but he
nervously recites the same rehearsed pitch he presents to potential
White buyers. His sales pitch appeals to Gladys’s curiosity and
thirst for knowledge as well as her desire for a better life for Dolly,
and she once again loses herself in a fantasy sequence. Her
imagined future for Dolly as a graduate seems as impossible as a
dream, but Gladys believes in it.\\

Nan Dear’s implication that the encyclopedias likely contain untrue
and offensive information about Aboriginal people highlights how
structures of knowledge and learning are not separate from the
racism that informs other cultural institutions. Nan Dear, Dolly, and
Gladys all agree that White people are the intended audience for
encyclopedias, but Gladys is still tempted by the books’ promise of
knowledge.\\

Neither Dolly nor Errol are fully comfortable making conversation
with someone of a different race, but Dolly finds herself enjoying
Errol’s company and helps him when local boys see him as an easy
target for theft. Her attraction to Errol leaves Dolly confused, and
Nan Dear and Gladys’s differing reactions to Errol don’t help that
confusion. Gladys is intrigued by Errol, while Nan Dear sees him as
a threat.\\

\subsection{Act 1, Scene 3}
\subsubsection{Summary}
Dolly rummages for linoleum at the town dump, where she
fantasies about buying linoleum from an attractive salesman
as a posh customer. Her fantasy fades and she returns home,
passing a drunk cousin of hers (Leon) who invites her to a party.
She turns him down and continues to the humpy, where Nan
Dear scolds her for taking an unsafe route home. She scolds
her further when Dolly refers derisively to her cousins who live
along the route, insisting that “they’ve had it hard” but refusing
to explain further.\\

Errol returns to the humpy to see if Gladys wants to buy an
encyclopedia set. Gladys is at work, so Errol invites himself to
wait as an excuse to spend time with Dolly. She won’t let him
inside the humpy, but they make conversation outside. After he
unthinkingly makes a racist remark, he apologies and asks if
Dolly wants to ride his bike with him. She gets on the bike by
herself and rides around him, humming, and the two sing
together.\\

Dolly and Errol talk about their lives. Errol enjoys the freedom
of being a travelling salesman, while Dolly has no sense of that
freedom; she expects she will pick fruit or work at the cannery
like everyone else in her community. When Errol asks what
Dolly dreams of doing, she muses that she would like to be a
model or a nurse, but she knows her mother has other dreams
for her.\\

Gladys arrives with a bookshelf from the dump, and, not
noticing Dolly’s embarrassment, explains to Errol that the
landfill plot is a desirable piece of land and should have been
designated for Aboriginal housing but was instead made into
the dump. Dolly leaves, humiliated, and Errol tries to salvage
the situation by assessing that the bookshelf just needs a coat
of paint. Gladys agrees, and the two laugh.\\

\subsubsection{Analysis}
Like her mother, Dolly escapes into fantasy sequences to imagine a
world where more possibilities are open to her. Nan Dear, on the
other hand, remains unwaveringly practical. She both sympathies
with Dolly’s cousins and recognises some of them (such as Leon) as
a threat to Dolly’s safety. Nan Dear believes in the importance of
family, and that demands empathy, even for relatives she dislikes.\\

Errol likes Dolly, but he was born and raised in a culture that
fundamentally disrespects her and her culture. As such, he carries
racist biases that he must consciously try to unlearn. Dolly, for her
part, forgives Errol’s casual racism and opens herself up to enjoying
her time with him, seemingly in part because he tries to shed his
racist biases.\\

The freedom that Errol enjoys, and Dolly’s lack of that same
freedom, comes from both racism and sexism. The racism of
Australian society prevents Dolly and her community from holding
any kind of job besides manual labour, while sexist expectations of a
woman’s role in the home keep Dolly tied to her house and unable
to travel like Errol can.\\

Dolly is embarrassed that Errol knows her family gets furniture from
the dump. She is already aware that, according to the commonplace
racism in their society, she and her family are of a lower status than
Errol, and their secondhand furniture is a symbol of that to her.
However, reclaiming abandoned furniture from a spot that should
have been made into Aboriginal housing is actually a small act of
resistance, as the family finds a way to use that land to make a
home despite it being zoned as a landfill.\\

\subsection{Act 1 Scene 4}
\subsubsection{Summary}
Dolly sets the table as Gladys serves dinner. Nan Dear
complains that they haven’t set a place for Papa Dear. Dolly and
Gladys point out that Papa Dear has been preaching across
Western Australia for three months, but Dolly sets a place for
him anyway. In a fantasy sequence, Papa Dear returns home
and kisses Gladys on the top of the head before disappearing.\\

Dolly tells Gladys and Nan Dear about a job as a cashier that
she hopes to get, but Nan Dear dismisses this and declares that
Dolly will continue to work with them at the cannery. She sends
Dolly out of the room and tells Gladys that a store in town will
never hire an Aboriginal girl from The Flats.\\

\subsubsection{Analysis}
Gladys’s fantasy sequence about Papa Dear makes clear that she
shares Nan Dear’s belief in the importance of family. Gladys wants
her family to be reunited, and she wants her father’s affection. This
emphasises the family dynamics present in the play; Gladys is a
mother, but she is also a daughter, and she misses her father.\\

Gladys and Nan Dear continue to butt heads over Dolly’s future.
Nan Dear sending Dolly out of the room reveals that even though
Nan Dear previously chastised Gladys for giving Dolly false hope for
the future, Nan Dear doesn’t want Dolly to hear how limited Nan
Dear believes her prospects really are.\\

\subsection{Act 1, Scene 5}
\subsubsection{Summary}
Errol is in the humpy, much to Dolly’s embarrassment, as he
sells Gladys an encyclopedia set to be paid for in instalments.
She recounts earning money by picking fruit and remembers
Dolly sleeping peacefully as a baby among the carpet snakes in
the tree paddocks. Errol asks if the snakes hurt young Dolly,
and the women are confused by the question since carpet
snakes are harmless. Errol gives Gladys Volume A of the
encyclopedia, and she displays it proudly on her new bookcase.\\

Dolly talks to Errol outside the humpy, and he asks why the
family lives in The Flats. As Dolly replies that Nan Dear and
Gladys like to live among their community, Gladys steps outside
and gestures encouragingly to Errol. He stammers nervously
until Gladys tells Dolly that Errol wants to invite her to a dance
in town. Dolly refuses, knowing that the townspeople won’t
accept Errol bringing an Aboriginal girl from The Flats as a
date, but Errol insists he doesn’t care. Dolly cautiously agrees
to the date before Nan Dear comes outside and sternly tells
Errol to leave.\\

\subsubsection{Analysis}
Errol’s misunderstanding of the story about the carpet snake
emphasises how removed he is from the cultural context that
informs the way Gladys is telling the story. Errol’s cultural context is
expressed in the encyclopedias rather than through oral stories or
practical knowledge, and it is positioned as the default culture that
Aboriginal Australians must assimilate to.\\

The family’s decision to remain in The Flats despite the constant
flooding highlights how home, for these women, is not only a matter
of physical shelter but also of community. Gladys continues to
extend her approval of Errol by encouraging the relationship
between him and Dolly, while Nan Dear continues to make clear
that she does not want to include Errol in their community. The
abuse the neighbour Ester experiences at her White husband’s hands
perhaps informs Nan Dear’s opinion.\\

\subsection{Act 1, Scene 6}
\subsubsection{Summary}
A White Inspector surveys the living conditions in the humpy,
making strained conversation with an anxious Gladys and Nan
Dear. He is a representative of the government of Victoria,
making a report on how Aboriginal people live, and he
patronisingly tells Gladys and Nan Dear that they must
assimilate to White culture. Dolly enters and the Inspector
sees her before Gladys and Nan Dear can shoo her away. After
he leaves, Gladys and Nan Dear discuss the possibility of the
government taking Dolly away.\\

\subsubsection{Analysis}
Gladys and Nan Dear’s fear of the Inspector comes from his status
as a representative of the Aborigines Protection Board, which held
absolute authority over the Aboriginal population of New South
Wales. It had, among other privileges, the right to displace
Aboriginal people from their homes and the right to remove
Aboriginal children from their families. Gladys and Nan Dear are
keenly aware that the Inspector has the power to break their family
apart, which is why they try to shoo Dolly away before he can see
her.\\


\subsection{Act 1, Scene 7}
\subsubsection{Summary}
Gladys gets dressed for a fundraiser. Gladys’s uncle intends to
raise money to negotiate for Daish’s, the dump, in order to
build Aboriginal housing. He plans to call the housing
Rumbalara, which means “end of the rainbow.” Meanwhile,
Dolly gets dressed for her date with Errol, and Nan Dear and
Gladys are astonished at how beautiful she looks. When Nan
Dear learns that Dolly is going out with Errol, however, she
pretends to have a coughing fit, and Dolly cancels her plans to
tend to Nan Dear.\\

\subsubsection{Analysis}
Though the local government denied Daish’s to the Aboriginal
community, the community continues to fight for it and for access
to stable housing. Meanwhile, Nan Dear makes clear her distrust of
Errol and her dislike of his burgeoning relationship with Dolly. Nan
Dear believes that a relationship with a White boy will endanger
Dolly, so she keeps her from the dance in an attempt to protect her
granddaughter.\\

\subsection{Act 1, Scene 8}
\subsubsection{Summary}
Dolly returns from school to find Nan Dear doing the laundry.
Dolly reports that Ester’s sons were not at school, and that
everyone suspects it’s because of the Inspector. Gladys gets
home and tells Dolly about a trainee program she heard about,
and Dolly is frustrated that her mother wants her to act like a
White person. Gladys asserts that she wants for Dolly what any
mother wants.\\

Dolly asks Nan Dear why they can’t have “a normal life,” to
which Nan Dear responds that their life is normal and there is
no point in trying to change it. They discuss the upcoming Miss
Marroopna-Shepparton Ball, and Nan Dear promises that she
will make Dolly a new dress for the occasion.\\

\subsubsection{Analysis}
Ester once again represents the level of powerlessness that
Aboriginal women, as doubly marginalised people, must contend
with. The government takes away her children, robbing her of her
status as a mother. Gladys tries to be the mother Dolly needs, but
Dolly is frustrated that her mother is teaching her to get by in White
society rather than how to thrive in Aboriginal society.\\

Dolly faces the constant expectation to assimilate into White
culture, which makes her resentful of her own Aboriginal identity.
Nan Dear encourages her to embrace her identity and traditions,
but she does this by denying that any other kind of life is possible.\\

\subsection{Act 1, Scene 9}
\subsubsection{Summary}
Gladys talks to Dolly about the latest book of the encyclopedia
to arrive. She tells Dolly that Errol asked about her and
encourages Dolly to go to the ball with him. The conversation
gets tense as Gladys urges Dolly “not to let them shame you.”
Dolly counters that Gladys should take her own advice. Nan
Dear enters, and Gladys asserts to her that Gladys is Dolly’s
mother, and she wants Errol to walk Dolly home from the ball.\\

\subsubsection{Analysis}
Gladys and Dolly both struggle with the conflicting impulses to
assimilate into White society and to honor their Aboriginal
identities. This tension is complicated by the mother-daughter
relationships at play. In this scene, Dolly stands up to Gladys while
Gladys does the same to Nan Dear.\\

\subsection{Act 1, Scene 10}
\subsubsection{Summary}
Gladys speaks to a bank manager and asks him to hire Dolly as
a teller. Elsewhere, Dolly picks fruit and sings to herself. The
bank manager questions Dolly’s qualifications and continues to
come up with reasons to deny her the job even as Gladys
rebuts each of his points. Gladys remembers Dolly asking if she
lets White people shame her, and Gladys insists that the bank
manager offer Dolly a job. He agrees to interview Dolly.\\

\subsubsection{Analysis}
Gladys has high hopes for Dolly’s future, and she sets out to make
those dreams a reality by visiting the bank and confronting the
manager. Her desire to do right by Dolly inspires Gladys to voice her
frustration with the racist systems around her, and she is further
empowered by her belief in her and her family’s capability to make a
better world for themselves.\\

\subsection{Act 1, Scene 11}
\subsubsection{Summary}
Dolly dances with Errol at the ball. She is uncomfortably aware
that the White partygoers are staring at them, but Errol
comforts her, and she imagines winning the title of Miss
Mooroopna-Shepparton Ball. Instead, the winner is a White girl
who mocks Dolly for wearing a dress made of fabric scrounged
from the dump. Dolly, humiliated, flees the dance and is
confronted by her drunk cousin Leon. To punish Dolly for
acting as if she is too good for her community, Leon attacks her
and Errol, tearing Dolly’s dress and leaving Errol with a black
eye.\\

\subsubsection{Analysis}
Dolly is uncomfortable and conscious of the danger around her, but
she lets herself dream about a world in which racism wouldn’t
prevent her from winning Miss Mooroopna-Shepparton Ball.
Though the play usually depicts this kind of daydreaming as an act
of protective optimism, in this case, Dolly’s hope for a better world is
cruelly punished by humiliation and violence.\\

\subsection{Act 1, Scene 12}
\subsubsection{Summary}
At the humpy, Gladys thinks happily of how much fun Dolly
must be having. However, Nan Dear senses that something is
wrong.\\

\subsubsection{Analysis}
Just as the violence that ensued after Dolly’s fantasy sequence
shows the hazards of losing oneself in imagination, Nan Dear now
demonstrates that her cynicism and wariness are sometimes more
realistic than Gladys’s optimism.\\

\subsection{Act 1, Scene 13}
\subsubsection{Summary}
Errol finds Dolly crying at the river, and she helps treat his
black eye. Though her dress is ruined, Errol reassures Dolly
that she is still beautiful. They tell each other about their
families, and Dolly explains that she grew up with three older
brothers and 27 male cousins, which taught her how to fight.
Errol marvels at her stories of her family. His father is strict and
traditional, and he insists that Errol’s mother is “too delicate” to
get a job.\\

Errol and Dolly get closer, but Dolly pulls away, unwilling to
subject either of them to the trials of an interracial relationship.
Errol counters that the two of them could move to the city
together for a better world. But Dolly argues that Errol’s
version of a better world is different from hers; her priority is
staying with her family, since that is what home means to her.
They fight, but she definitively turns him down and at last Errol
accepts this. Dolly starts to walk home, but she is intercepted
by Leon.\\

\subsubsection{Analysis}
Errol and Dolly have grown up with vastly different conceptions of
family, both of which are shaped by their race and culture. Errol’s
version of family is the strict nuclear family that is prioritized and
rewarded by White Australian culture. He was born into his family
but has no real attachment to it. Dolly, on the other hand, has
meaningful connections with her entire extended family. She is also
used to manual labor being expected of everyone, including women,
while Errol’s family embodies a kind of White middle-class gender
roles that forbid women from working.\\

Errol’s unthinking and casual racism has not yet hindered his
relationship with Dolly because she knows he is well-meaning.
However, it proves the undoing of their relationship as Errol can’t
wrap his head around the fact that Dolly doesn’t want to leave her
community to live an assimilated life with him in the city. She is also
more aware of the violence interracial couples face, since she is
generally more aware of racism than Errol.\\

\subsection{Act 1, Scene 14}
\subsubsection{Summary}
Nan Dear and Gladys pack up the humpy, preparing to
evacuate for a flood. They worry about Dolly, who is still not
home. A jungi (policeman) arrives to escort them to the dump,
which is the evacuation site since it is above the flood line. They
are about to leave when Dolly arrives, clearly traumatised. She
“half screams” when the jungi tries to touch her, and she
refuses to answer when Gladys and Nan Dear ask her what
happened.\\

\subsubsection{Analysis}
The flood mirrors Dolly’s implied rape by her cousin Leon. The flood
will destroy the humpy, emphasising the instability of Dolly’s home.
Meanwhile, Leon’s assault destabilises not only Dolly herself but
also her connections to her community. Her trauma is obvious to
those around her, but she refuses to share the truth with Gladys and
Nan Dear, with whom she is usually so open.\\

\subsection{Act 2, Scene 1}
\subsubsection{Summary}
The women return to the humpy after the flood. Dolly, still
“shell-shocked,” continues to wear the dress she wore to the
ball, which is now badly torn and muddied. Nan Dear and
Gladys clean up the wreckage, and Gladys finds that all her
encyclopedias are ruined. She insists she doesn’t mind. Nan
Dear muses that moving to the Rumbalara housing
development might be a good idea for Dolly. Nan Dear and
Gladys decide not to tell Papa Dear what happened to Dolly,
since it’s “women’s business.”\\

Errol arrives to apologise to Dolly, but Nan Dear and Gladys
violently drive him away, assuming that he is the one who
assaulted Dolly. Errol is devastated that they could think him
capable of such a thing, and after he leaves, Dolly finally speaks,
confirming that Errol was not the one who raped her.\\

\subsubsection{Analysis}
Dolly in unable to move past the horror of being assaulted, and her
refusal to change her dress reflects how stuck she is in the moment
of her trauma. Nan Dear and Gladys, meanwhile, try to rebuild both
Dolly and the humpy. The family lives in The Flats to be with their
family and community, but they are prepared to move to Rumbalara
now that Nan Dear believes Rumbalara might help Dolly. Nan Dear
and Gladys agree that the attack on Dolly, and the work of helping
her recover, are “women’s business,” which reflects that Papa Dear
might be an important activist, but he is not privy to the equally
important work that goes on in his own home.\\

Gladys’s affection for Errol is gone now that she suspects him of
hurting her daughter, which emphasises how committed she is to
protecting Dolly. Dolly, on the other hand, retains some loyalty to
Errol, breaking her silence to make sure her family knows that he is
innocent.\\

\subsection{Act 2, Scene 2}
\subsubsection{Summary}
A bulldozer destroys the family’s humpy, and they move into a
sterile, small, and unwelcoming concrete house at Rumbalara.
Nan Dear develops a cough, but she promises to make curtains
for the new house.\\

\subsubsection{Analysis}
The family leaves behind the humpy, which was unstable but
familiar and made by the women themselves, for the distinctly unhomelike
Rumbalara housing. The contrast between the humpy and
Rumbalara illuminates the injustice of needing to choose between
safe housing and emotionally fulfilling housing. Nan Dear’s promise
to liven up the house with curtains also underscores the importance
of domestic labour in making a house a home.\\

\subsection{Act 2, Scene 3}
\subsubsection{Summary}
Gladys brings Dolly a letter from the bank, but Dolly is
uninterested in the job as teller. She tells her mother that this
job is Gladys’s dream, not Dolly’s, and angrily suggests Gladys
“fix [her] own house.” Gladys storms out.\\

Nan Dear pays the rent to the rent collector Mr. Coody and
helps Dolly prepare a meal. They listen to the radio and are
surprised to hear Gladys arguing with city councilors about
building Aboriginal housing at Daish’s. Gladys is removed from
the council meeting, but Dolly and Nan Dear are still proud of
her.\\


\subsubsection{Analysis}
Though the family has left the humpy, the image of rebuilding a
house returns as Dolly uses a person’s house as a symbol for their
life. This metaphor emphasises the importance Dolly and her family
place on home. It also signifies how the sterile Rumbalara housing is
causing fractures within the home, which was previously a space
where the family is united.\\

Mr. Coody’s arrival to collect rent makes clear that the family’s
home is not their own anymore; it is owned and controlled by Mr.
Coody. In light of this lack of control, Gladys asserts her dignity and
a level of control by refusing to let the city councillors debate her life
without her input.\\

\subsection{Act 2, Scene 4}
\subsubsection{Summary}
Errol finds Gladys cutting wood outside the Rumbalara
housing. They discuss the encyclopedias. Gladys doesn’t want
to buy any more sets of encyclopedias, and even though she
isn’t supposed to be able to cancel her contract, Errol offers to
help her do so. Gladys reveals she can’t read or write, so Errol
helps her compose the required letter. Gladys apologies for
assuming that Errol raped Dolly, and Errol asks if he can see
Dolly. Gladys gently but firmly tells him that Dolly won’t want
to see him.\\

Gladys asks Errol to teach her to read and write. She explains
that Nan Dear attended school when the family lived at
Cummeragunja, but the mission managers were cruel, and
Gladys was sent off at a young age to do domestic work for a
large wealthy family. After that, she married Dolly’s father and
was too busy working and raising children to learn. Errol agrees
to teach Gladys to read. He still hasn’t given up hope of
reuniting with Dolly, and he hopes to prove that he “can be
worthy of her.”\\

\subsubsection{Analysis}
Gladys has always liked Errol as a boy who might be good for Dolly,
but now their relationship begins to develop on its own terms. She
priorities Dolly’s emotional well being over Errol’s desire to see her,
but she nonetheless treats Errol with gentle respect. Gladys is the
kind of strong and caring mother that Errol never had, and she
allows herself to be vulnerable with him when she reveals she can’t
read or write.\\

The reasons behind Gladys’s illiteracy are a combination of racism
and sexism. She didn’t learn to read as a child because the
Aborigines Protection Board removed her from her family, who were
already being mistreated at Cummeragunja. By the time she was old
enough to escape the most controlling parts of the Aborigines
Protection Board, Gladys found herself where she is now––bound
by the domestic responsibilities expected of her as a woman.\\


\subsection{Act 2, Scene 5}
\subsubsection{Summary}
The rent collector Mr. Coody comes by for the rent and finds
Nan Dear cooking eggs for a heavily pregnant Dolly. He notes
that he will re-evaluate their lease in light of Dolly’s pregnancy,
but Nan Dear takes him to task for consistently invading his
tenants’ privacy and imposing too many rules. She throws eggs
at him, and he flees.\\

Dolly and Nan Dear look at Dolly’s old family tree assignment,
and Dolly shares that she plans to name her baby after Papa
Dear. She confesses that she is worried she won’t love the
baby. Nan Dear says that Dolly is finally a woman and thus
ready to hear the truth: Gladys’s biological father is a White
man who raped Nan Dear when she was a teenager. This is why
Nan Dear has been so opposed to Dolly’s relationship with
Errol––especially because Errol has the same surname as Nan
Dear’s attacker. When Dolly shares that she broke things off
with Errol because he wants to take her away from her family,
Nan Dear encourages her to consider leaving. At the same
time, she insists that Dolly mustn’t resume her relationship
with Errol, and Dolly agrees.\\

\subsubsection{Analysis}
Mr. Coody represents how White authority infringes on every
moment of Aboriginal life by denying Aboriginal people privacy and
self-determination within their own homes. Nan Dear resists this
facet of White supremacy, asserting her home as a space she and
her family have the right to control.\\

While Dolly’s belief in the importance of family led to tension with
Errol, it connects her to Nan Dear, who believes the same thing.
Dolly’s decision to name the baby after Papa Dear reaffirms her
devotion to her family and her desire to honor her relative’s
traditions and legacy. Nan Dear values tradition as well, but she also
wants Dolly to make a meaningful future for herself, which is why
she encourages Dolly to leave. Nan Dear’s dislike of Errol is also
explained as not simply a general protectiveness of her
granddaughter but a reaction to trauma that she now chooses to
share with Dolly in a moment of honesty and vulnerability.\\

\subsection{Act 2, Scene 6}
\subsubsection{Summary}
Errol is teaching Gladys to read, and they practice on the
Inspector’s report of Aboriginal living conditions. Gladys
encourages Errol to call her “Aunty” and asks about his feelings
for Dolly. He admits that he still has feelings for her, but he
doesn’t know what she wants. Gladys invites Errol to a public
meeting in Melbourne the following week, where Papa Dear
will be speaking on Aboriginal housing. He accepts the
invitation, and Gladys tells Errol that even if he and Dolly don’t
end up together, she likes Errol.\\

\subsubsection{Analysis}
Errol loves Dolly, but that love has been hindered by a lack of
understanding of where she comes from. Now Errol is developing a
love for Dolly’s family and an appreciation for her community that
mirror Dolly’s own loyalty to her loved ones. Gladys’s request that
Errol call her “Aunty” is an invitation to join the family, which
demonstrates that her understanding of family is inclusive and
based on personal bonds. This contrasts Errol’s previous
understanding of family as something constraining and dependent
only on blood ties.\\

\subsection{Act 2, Scene 7}
\subsubsection{Summary}
Nan Dear, Gladys, and Dolly sit in the audience at the public
meeting. With them is Dolly’s baby, Regina. Dolly says that
Gladys should be giving a speech, but Nan Dear scoffs that
speechmaking is “men’s business.” Dolly gave up on the bank
interview, but she has been awarded a scholarship to a nursing
program in Melbourne, and Gladys is thrilled for her.\\

Errol arrives and sits with the women. Gladys pulls Nan Dear
away to go find Papa Dear. Errol is startled to see Regina, and
he finally realises what happened the night of the flood. He
apologies to Dolly “for everything,” and she apologies in turn.
He offers to move to the country so they can be together and
stay near Dolly’s family, but she tells him that she is going to
Melbourne and not ready for a romantic relationship. He
respects her decision, and they agree to remain friends.\\

Nan Dear returns, feeling ill, and Dolly hurries off to get her
medicine. Errol offers to drive her to the hospital, and when she
refuses, he offers to drive her home. She is touched. Nan Dear
asks about Errol’s surname, and he explains that his parents
adopted it after immigrating from Germany. Nan Dear is
relieved that Errol is not related to the Fisher who once
assaulted her. She kisses Errol on the cheek. Dolly sees this and
is stunned.\\

Papa Dear fails to arrive to present the petition, so Gladys does
it in his stead. She makes a speech petitioning Queen Elizabeth
to give Aboriginal Australians decent housing and to finally
treat them like Australian citizens. She begins to read from the
petition and includes her own opinions, demanding equal
education and job opportunities for Aboriginal people. She
concludes with a point that is not on the petition: she wants
White people to treat Aboriginal people with respect. The
audience applauds.\\

In a fantasy sequence, Nan Dear imagines Dolly and Errol, with
Regina between them, getting married. She returns to reality
and tells Dolly to marry Errol. She tells Errol to call her “Nan,”
and he thanks her for her blessing. Papa Dear arrives.\\

\subsubsection{Analysis}
Dolly has grown up from a girl to a woman. She has chosen a future
for herself and will be leaving the protection of her mother and
grandmother, while she herself is tasked with raising the family’s
next generation. Her family has faith in her ability to pursue a
career, but Nan Dear still believes that woman and men’s work are
divided categories, and that public activism belongs to men.\\

Errol makes clear his devotion to Dolly in this scene. He
unequivocally apologies, recognising his own wrongdoing, and
promises to prioritise Dolly’s needs over his own wants. His desire
for a relationship is not hindered at all by Dolly’s baby, which
emphasises that he is ready to commit not only to Dolly but to her
family. When Dolly turns him down, he respects her ability to
choose what she wants and is just glad to stay in her life as a friend.\\

Errol proves to Nan Dear that despite his occasional moments of
unthinking racism, he is fundamentally kind and respectful.
However, she cannot let herself like him until she knows for certain
that he isn’t related to her rapist. When she learns that Errol has no
connection to her past trauma, Nan Dear embraces him as a good
choice for Dolly.\\

Papa Dear’s work is important, as his speech provides Gladys the
foundation to begin her public activism. However, Gladys soon steps
out of Papa Dear’s shadow, asserting that politics and activism are
just as much a place for women as they are for men. Her concerns
about housing are informed by the domestic labour she has spent her
life performing, but she can also see beyond that. Her demand that
White people respect their Aboriginal countrymen indicates a
holistic understanding of how microaggressions, and seemingly
small moments of interpersonal disrespect, are the building blocks
of White supremacist ideologies that lead to racist government
policies.\\

Tellingly, the final moment of the show is centred around the
intimate workings of the family rather than the broader societal
issues the play also explores. Nan Dear finally breaks free of her
cynicism and indulges in a fantasy sequence, like her daughter and
granddaughter often do, and that fantasy inspires her to invite Errol
to join the family. As Errol, Dolly, and Regina come together to form
the next generation of parents and children, Papa Dear arrives as a
reminder of the family’s lineage.\\

\newpage

\section{Literary and Dramatic Device Tracker}



ZZPDFEMBED Tracker.pdf ZZEND



\newpage

\section{Symbols and Motifs}

\subsection{Symbols}
\renewcommand{\arraystretch}{1}

    
\begin{longtable}{|p{2cm}|p{4cm}|p{7cm}|}
\hline
\textbf{Symbol} & \textbf{Where is it seen? Scenes \& Quotes} & \textbf{What does it represent?} \\
\hline
\endfirsthead

\hline
\textbf{Symbol} & \textbf{Where is it seen? Scenes \& Quotes} & \textbf{What does it represent?} \\
\hline
\endhead

Encyclopedias &
``The encyclopedia Britannica. Your entree to the world of learning, a world of discovery, a world of interesting facts…’’ Pg 133 \newline
``Encyclops boy and he knows nothing’’ Pg 146 \newline
``You mean you’ve got to pay for those things?’’ \newline
``Could you give us a hand with the encyclopedia set?’’ \newline
``They’re only possessions. And what do they matter? People are all that matter’’ Pg 176 &
For Gladys- The embodiment of the family’s aspirations; they provide a sense of hope, particularly for the future. \newline
For Nan- Wary of the way governmental institutions have commercialized and chipped away at the knowledge of the Aboriginal people. \newline
Errol- has little need for the books, highlighting the fact that he has privilege as a white male, even though he doesn’t have the same vocabulary and knowledge. \newline
The cost of the encyclopedia collection alludes to the cost of Westernised knowledge at the expense of Indigenous knowledge. Even though the encyclopedias are destroyed, the Dear’s continue on with they life, emphasising that Western knowledge is not essential. \\
\hline

‘Que Sera, Sera’ &
``Que Sera Sera, Whatever will be will be, the futures not ours to see. Que Sera, Sera.’’ Pg 124, pg 199 \newline
``I asked my sweetheart what lies ahead’’ Pg 175 &
Que Sera Sera describes an attitude of cheerful fatalism. The lyrics of the song suggest that events are predetermined in advance and humans are powerless in determining their future. Each verse of the song, represents the character being taken through the journey of their life and as one ages and grows up. The story is bookended at the beginning and end of the story, the song has the role of opening and closing the play. At the end of the play, the song plays again, this time the song is suggesting that events are not predetermined if people take control of their own circumstances. The song explains that fate is a detriment of human destiny. \\
\hline

The Cork Trees &
``Dolly walks past the cork trees. She sees someone in the shadows’’ Pg 139 \newline
``That can be our special meeting spot’’ [cork tree] Pg 149 \newline
``Well hello Dolly… thought I’d catch up with you’’ Pg 172 &
The cork trees represent a place of separation between the innocence of Dolly/her youth and the outside world. While the cork trees are associated with good memories such as her time in the Flats and her special spot with Errol, this location is now sinister and traumatic for Dolly. At the cork trees Goomeys try to coerce Dolly to drink. The cork trees are also where she is sexually assaulted by a cousin. Pivotal good and bad moments in Dolly’s life are associated with the cork trees. The cork trees are also a danger for Aunty Ester who begins to find solace in sex and alcohol after her children are taken away. \\
\hline

Ajax &
``Ajax contains bleach, you’ll stop paying elbow tax’’ Pg 155 \newline
``Honestly their mad about whiteness’’ Pg 155 &
The Ajax bleach solution is used to explain an obsession with “whiteness” and a European society. Ajax is heavily advertised on the radio. Gladys eventually comes to reject this ideology by the middle of the play, switching off the radio. “mad about whiteness’’ \\
\hline

Radio &
``The radio is heard featuring a description of the Queens; 1954 Royal Tour of Australia’’ Pg 127 \newline
``Bob Dyers quiz show, Pick-a-Box on the radio’’ Pg 128 \newline
``Dolly out of sight, is singing to the radio’’ Pg 154 \newline
``We resume our broadcast of the Rodney Shire Council Meeting’’ Pg 180 &
The radio represents the connection the women have with the world outside the humpy and what is happening around the local shire. It offers an insight into a world they are barred from by Europeans. The radio station is one of the few things indigenous Australians and European settlers share. The radio always appears to be ominously playing in the Dear’s household. The radio is the external pressure which promotes discriminatory information. \newline
The radio serves as a parallel narrative to that of the lives of the Indigenous women. Gladys eventually has her dream of being on the radio, but not in the context of appeasing to a British radio show Pick a Box. Instead, Glady’s appearance on the radio is her fighting for her voice an indigenous Australian woman. \\
\hline

Swan Eggs &
``I’m gonna go back and die there’’ [Murray River] Pg 125 \newline
``And I’ll have a feed of swan eggs’’ Pg 125 \newline
``Swan eggs? You’re okay aren’t you Nan?’’ Pg 157 &
The swan eggs are a representation of Nans dream, to go back to her country on the Murray River and feast on swan eggs before she dies. The swan eggs show Nans longing for a desire to return home and her connection to country. Swan eggs becomes imminent when Nan’s passing is discussed \\
\hline

The humpy &
``rebuilding after a flood has devastated their home’’ Pg 123 \newline
``Dolly is really embarrassed that Gladys has allowed Errol into the humpy’’ Pg 144 \newline
``Gladys is going around the humpy emptying the mouse traps of dead mice.’’ &
The humpy is the shelter Nan and Gladys have turned into a “home” for Dolly. Nan and Gladys go to immense lengths to provide a comfortable and loving home for Dolly this is shown by them laying linoleum on the dirt floor and repairing the house when needed. The humpy symbolises the resilience of the women. Despite countless floods, they make the humble shelter a “home” \\
\hline

The house at Rumbalara- ‘Rainbow’s End’ &
``From riverbank humpy to whitehouse is quite a step’’ Pg 177 \newline
``This is the most vigorous attempt yet to solve Aboriginal housing’’ Pg 178 \newline
``Concrete, small, white and featureless? Any thing but loveable’’ Pg 178 \newline
``Not quite the ‘new deal’’’ Pg 178 &
Rumbalara is the new housing that was provided to the Indigenous Australians following the devastation of the floods. The name “Rumbalara” means at the “end of the rainbow”. The name suggests that this housing is adequate for Indigenous Australians. Unlike a rainbow, the concrete house is bare and white. Once again Gladys and Nan must try to make Rumbalara a stable home for Dolly and her child. \\
\hline

The pram &
``the lights come up on Dolly pushing a pram’’ Pg 132 \newline
``Gladys appear with a crampy old bookcase balanced on the pram’’ Pg 142 \newline
``Triumphantly, she pulls an obviously used and bashed can from the pram.’’ Pg 143 \newline
``They both look over at the pram and smile.’’ Pg 192 \newline
``He sits between her and Dolly and peeks in the pram’’ Pg 192 \newline
``Dolly and Errol, pram in the middle gets hitched.’’ Pg 198 &
Dolly is seen carrying the pram at numerous stages during the text. Dolly is seen pushing the pram when she first meets Errol, foreshadowing the end of the text when she is with him and the same pram holds her child. This pram was initially used to carry scraps and materials from Daich’s Paddock, the pram now carry’s the effects of Dolly’s assault, a child \\
\hline

The ball &
``There’s a dance on in Shepparton next Saturday’’ Pg 148 \newline
``Not just any ball. The Miss Mooroopna-Shepparton Ball’’ Pg 159 \newline
``Re this ball, why are you saying if she can step out? She’s my daughter, I have brought up 3 others before her.’’ Pg 162 \newline
``She’s going to the ball. And Errol Fisher is walking her home’’ Pg 162 &
The ball is an event Dolly longs to attend with Errol. Dolly associated the ball with finally being recognised by the European society. The ball is a way for Dolly to feel like she is not some ordinary girl from the Flats. Attending the ball signifies the one time in the text where Dolly is not protected by Gladys or Nan. Unfortunately, her attendance causes a chain of events resulting in her ending the relationship with Errol and being sexually assaulted by Leon. The ball is no longer a memory of happiness by one of pain and regret. \\
\hline

The Family Tree &
``I need a list of all our family members… our parents, their parent and so on… but not cousins’’ Pg 124 \newline
``You need to know who your cousins are so you don’t marry ‘em’’ Pg 125 \newline
``Reginald Harold Dears parents are?’’ Pg 126 &
The family trees shows the importance and the interconnectedness of the community. Nan challenges Dolly’s perception and family when she is reluctant to include cousins in the family lineage. However, the family tree is also a reminder of Nan’s past and the secret she holds onto. Nan cannot continue the family tree on Gladys paternal side due to her sexual assault and her lack of knowledge about Gladys’ real father. The incomplete and complicated family tree is the consequence of the Stolen Generations. \\
\hline

Dolly’s dress &
``Can you make me a new dress Nan? A gorgeous one with lots of fabric- pretty fabric.’’ Pg 160 \newline
``A big skirt. And a peplum’’ Pg 160 \newline
``It’s not house of Biba’’ Pg 160 \newline
``Dolly is swirling and floating around the dance floor in her full-length skirted dress’’ Pg 165 \newline
``courtesy of the town bloody tip’’ Pg 166 \newline
``sobbing her dress is slightly ripped’’ Pg 167 &
Dolly’s dress represents her desire to emulate white, European ideals and mimic those like Nancy Woolfthorpe. The resourcefulness of the family is shown with their ability to construct a beautiful dress from the remains of curtains. The dress later becomes associated with a terrible and traumatic night. At the dress rips and damages further, it signifies how Dolly is torn into by Nancy, Leon and Errol. \\
\hline

The hessian barrier &
Gladys walks into Shepparton hoping to see the Queen during the 1954 Royal Tour, only to find hessian screens erected along the route so that the visiting party will not see the humpies -- or the people who live in them. &
A physical embodiment of racial exclusion: the authorities do not fix the conditions on the Flats, they hide them. Gladys is denied even the role of spectator in the national celebration, exposing assimilation-era Australia as concerned with appearances rather than justice. The screen that blocks her view also blocks the Queen's -- concealment works in both directions, keeping white Australia comfortably ignorant. \\
\hline

The hidden shoes &
Nan hides Dolly's shoes and stages a coughing fit to stop her going out with Errol; the stage directions reveal the pretence when Nan ``remembers to cough again''. &
The shoes stand for Dolly's freedom to step out into the wider world, and Nan's concealment of them dramatises protection expressed as control. Because Harrison reveals the ruse through stage directions rather than dialogue, the audience becomes a knowing witness: we see that Nan's fear is rooted in her own past, making her neither villain nor martyr. \\
\hline

\end{longtable}

\subsection{Symbols with Links to Themes}


ZZPDFEMBED SymbolTable.pdf ZZEND


\newpage

\subsection{Motifs}


\renewcommand{\arraystretch}{1}

\begin{longtable}{|p{2cm}|p{4cm}|p{7cm}|}
\hline
\textbf{Motif} & \textbf{Where is it seen? Scenes \& Quotes} & \textbf{What does it represent?} \\
\hline
\endfirsthead

\hline
\textbf{Motif} & \textbf{Where is it seen? Scenes \& Quotes} & \textbf{What does it represent?} \\
\hline
\endhead

Music & 
``Que Sera Sera’’ \newline
``Catch a falling star’’ \newline
``Somewhere over the rainbow’’ &
In Que Sera, Sera. Each verse takes the audience though a key stage in the characters lives in which they are at the mercy of outside forces. Catch a falling star represents the very minimal luck that the family has, and that Dolly is unlikely to get the traineeship. Rather that catching stars, she is destined to be catching fruit like her mother and grandmother. Somewhere over the rainbow alludes to First Nations yearning to return home after having lost their land and being relocated to Rumbalara, which is “over the rainbow.” This song is seen in the Wizard of Oz when Dorothy longs to return home, however, she has a happy ending, suggesting such outcomes are only fiction of Hollywood. \\
\hline

Water & 
``Nan Dear and Gladys are rebuilding after a flood has devastated their home’’ Pg 123 \newline
``A tap appears from nowhere and from it flows blue jewels in an approximation of water.’’ Pg 155 \newline
``I’d have to leave the river…’’ Pg 171 \newline
``Rain, thunder, darkness’’ Pg 175 \newline
``The water was dried away but devastation has been wrought’’ Pg 176 &
Water is associated with both destruction and change. While water had the capability to devastate the Dears, it resulted in the move to Rumbalara, with water flowing from a tap. Furthermore, the Yorta Yorta people, prevalent in the region Rainbow’s End is set in, have a strong connection to waters and rivers. \\
\hline

Dreams & 
Dream sequence of Gladys meeting the Queen (Gladys) Pg 126 \newline
Gladys being a guest on Pick-A-Box (Gladys) Pg 130 \newline
Dolly parading with a robe and clapboard hat (Gladys) Pg 134 \newline
Dolly buying luxurious goods in a designer store (Dolly) Pg 139 \newline
Dream sequence of Papa Dear returning home (Gladys) Pg 143 \newline
Water flowing from the housing and Rumbalara (Gladys) Pg 155 \newline
Dolly winning Miss Mooroopna- Shepparton Ball (Dolly) Pg 166 \newline
Dolly getting married to Errol and them looking after the baby (Nan) Pg 198 &
The dream sequences reveal the unspoken desires and hopes of the characters that they often do not admit to themselves. The vivid nature of the dream sequences suggests how greatly the characters long for these outcomes. \newline
The characters frequently has dream sequences, suggesting the most optimistic and hopeful. \newline
Nan having a dream sequence towards the end of the text suggests that she too is now hopeful for Dolly’s future with Errol and her unexpected child. Dream sequences become far less frequent following Dolly’s assault, indirectly conveying the impact of this event on the characters. They become less hopeful for Dolly and her outcomes. \\
\hline

Beat & 
Pg 183, 194, 186, 187, 188, 190 &
The beat is a moment which propels the story forward. It aims to compel the audience to take stock of what may happen next while propelling the story forward. The beat emphasises what has been said and what is going to be said next. The beat helps to recapture the audience’s attention thought the play. \\
\hline

House Décor & 
``Gladys appears with a crappy old bookcase balanced on a pram’’ Pg 142 \newline
``Emptying the mice traps of dead mice and setting new ones’’ Pg 160 \newline
``Dolly hoist the bolt iron over her shoulder as the salesman fades away’’ Pg 139 \newline
``Dolly passes Gladys the jam tins. One at a time’’ Pg 144 \newline
``Just boiled up in a Kero tin, with velvet soap and a blue bag, same as everyone.’’ Pg 152 \newline
``We don’t go much for presents’’ Pg 170 &
The Dears are frequently seen collecting scraps and discarded material from Daiches paddock. Despite it’s condition, the women are still able to transform this décor into elements that make their humpy a home. Although the furniture is not elaborate and expensive, it is functional. \\
\hline

\end{longtable}


\newpage

\section{Written Exemplars}
\subsection{Sample 1 - High Scoring \underline{Rainbow's End} Essay}


ZZEXP Rainbow's End I ZZEXPEND



Topic:
How does \underline{Rainbow’s End} suggest that everything will be alright?\\


In the 2005 Indigenous play ‘Rainbows End’, playwright Jane Harrison suggests everything will be
alright for the Dear family, despite the challenges they combat. This is shown through the portrayal
of self-determination, coupled with the continued hope and optimism of the Dears. However, the
heavily damaging colonial attitudes to the Indigenous community act as a constant threat looming
over the community and their futures.\\

Throughout the play, Harrison highlights the unwavering self-determination of the Indigenous
community, suggesting everything will be alright due to the strong resolve of the Dear family. This is
depicted most vividly through the growth of Gladys, initially reserved and acquiescent, but steadily
becoming more outspoken and steadfast in her beliefs for a better future. This culminates in Gladys
taking a bold stand for the indigenous community at the town meeting, stating “I’m not an
interloper – I belong here! - this is my land!”. Gladys’ open defiance in the meeting, challenging
cultural norms and expectations, acts as a pivotal switch from Gladys’ belief everything will be
alright, to taking direct action in pursuit of a better future for her family. Similarly, Dolly’s stoicism
and resilience in the face of adversity and her assault serve as an emphatic representation of selfdetermination.
After her assault, Dolly’s honest and true nature in showing her baby to Errol, despite
severe shame, depicts her growing as an individual through hardship and pain through these
qualities of self-determination. Stating “Yes, I’m okay” when questioned by Errol about her baby,
Dolly embraces her struggles, remaining determined and continuing to seek a better future for
herself and her baby. Overall, Harrison suggests everything will be alright for the Dear family, despite
great hardship, through the continued determination of both Gladys and Dolly.\\

Harrison further suggests that everything will be alright due to the stalwart hope of the Dear family,
embodied through Papa Dear, dream sequences and the prospect of education. Throughout, Papa
Dear acts as a symbol of a positive future and hope for the family and broader Indigenous
community, travelling across Australia fighting for their rights. This is shown through Nan, who states
“its Papa Dears mission to make things better for the Aboriginal people”, suggesting that Papa Dear
acts as a figurehead and an unseen force for good in the Indigenous community. Similarly, Papa Dear
exists as a projection of hope through dream sequences with Gladys, who observes “Papa Dear, in
old fashioned hat and coat, dances in, throws his hat on the hatstand, kisses Gladys on the top of the
head, and dances out”. This suggests an omnipresent quality of Papa Dear as a symbol of hope for
the members of the Dear family; as a constant reminder that everything will be alright. Further, the
prospect of education offers hope for the future for Dolly, and is first presented in the form of
encyclopaedia’s from Errol, to which Gladys exclaims “my daughter, a graduate”. This is further
propagated by Dolly’s ambitions of becoming a nurse in Melbourne, as a catapult to a better and
more prosperous future. Ultimately, Harrison suggests that the unwavering hope of the Dear family,
through Papa Dear and the prospect of education, are a reminder that everything will be alright.\\

However, Harrison also indicates the highly destructive nature of colonial attitudes, and their ability
to prevent the Dear family from a better future, as a direct challenge to the sentiment that
everything will be alright. Harrison portrays this primarily through the housing inspector,
representing a hostile Australian government policy and its corrosive effect on the lives and futures
of the Indigenous communities. This is demonstrated through the removal of Aunty Esthers children
at the behest of the inspector, and further the repeated visits he makes to the Dear family home.
Stating “your whites are so white, with river water no less”, the inspector highlights to close scrutiny
of every detail of the Dear’s lives, and the constant threat to their future looming over them.
Similarly, Harrison suggests the continued resentment of the broader community actively works to
dissuade the Dears and Indigenous community from seeking a better future. This is displayed
through the town council, whilst attempting to silence Gladys challenging societal norms, stating
“the Chair does not recognise this…. Eject this interloper!”. Further, Harrison displays the negative
affects of such attitudes through Dolly, exclaiming “it’ll never change…. Some things never do” after
being scrutinised by Nancy at the Shepparton ball. Overall, Harrison suggests that colonial attitudes
can directly challenge and threaten the hopes and self-determination of the Dear family and broader
indigenous community, and further the prospects of everything being alright.\\

On balance, Harrison displays the potential for everything to be alright for the Dear family and
broader Indigenous community through the qualities of self-determination and hope, but also the
constant threat of colonial attitudes to the future.\\

\subsection{Sample 2 - High Scoring \underline{Rainbow's End} Essay}


ZZEXP Rainbow's End II ZZEXPEND



Topic: 'Harrison’s play celebrates characters who show strength through adversity.' Discuss.\\

Set in Victorian town, Shepparton, Jane Harrison’s 2007 play, Rainbow’s End, explores the
lives of three generations of Koori women as they navigate the setbacks and challenges
faced by First Nations people in 1950s Victoria. At a time where the government regulates
the policies of Assimilation of the Indigenous community, these women are vulnerable to the
impacts of racism, colonisation and suppression of knowledge. But despite facing
discrimination at every turn, Harrison commemorates their resilience in the face of adversity.
The women are able to become symbols of hope as they refuse to allow their struggles to
define them.\\

Throughout Rainbow’s End, the struggle for the Indigenous community to pursue their
dreams due to racial prejudice is highlighted in order to spotlight their ability to overcome it.
This is demonstrated in scene ten in the bank managers office. Gladys tries to apply Dolly
for a job as the bank teller, assuring the bank manager that she is fit for the job by
advocating for her brains and wits. But despite Dolly being a good candidate, the manager is
hesitant about letting her apply. He argues that, “reliability is important,” implying that Dolly is
unreliable despite never having met her. He makes this assumption due to his personal
prejudice to people of Dolly’s heritage. This demonstrates the racial discrimination that limits
opportunities of First Nation’s people. Regardless of the bank manager's doubts, Gladys
pushes for the job as she is reminded of Dolly’s words, “have you learnt not to be shamed by
them?”. Through this, Harrison highlights how the women’s power to overcome setbacks
comes from within themselves. She champions these women for their ability to stand up for
themselves in spite of the racially biased attitudes that are meant to limit them. This theme of
hope is prevalent throughout the play and intends to inspire strength from within to tackle
problems.\\

Moreover, Harrison explores the intricacies of cultural identities and belonging, highlighting
the strong connection to culture despite the colonisation that the characters face. Nan Dear
is a compelling example of this. Nan’s frequent allusions to spirits or warnings from her
“ancestors” demonstrate the strong indigenous traditions that survive and carry on even after
the effects of colonisation. Intergenerational wisdom is significant in indigenous culture and
Nan Dear’s frequent references to guidance from her ancestors attempts to pass that down
through her children, who may continue to succumb to assimilation but nevertheless is still
deeply rooted in their cultural identities. The reminders of tradition allows Harrison to assert
that cultural identity cannot be erased and is represented through characters like Nan Dear
who have been through cultural erasure but continues to be a strong image of indigenous
culture.\\

Furthermore, resilience is present in overcoming adversity from the suppression of
knowledge that characters throughout the play go through. Gladys is a representation of this
as it is revealed to the reader that Gladys is illiterate. This was due to the assimilation policy
that separated Gladys from her family and forced her into working at a young age rather than
getting an education. Harrison demonstrates Gladys overcoming this in the last scene of the
play where Gladys delivers a speech at the public meeting. Not only does she deliver it
flawlessly, she delivers it with confidence, which Gladys has mentioned that she “gets
nervous in front of people.” Harrison acknowledges the optimistic and determined attitude
that Gladys uses as a tool to fight for her ambitions despite the systemic racism that has denied her basic human rights. Through Gladys Harrison aims to write about the ‘unsung’
heroes who were not expected to be a hero but was able to become one in spite of the
adversity that she faces.\\

In conclusion, “Rainbows End” champions the characters for their ability to overcome
adversity due to their unwavering determination, offering a powerful testament to the power
of human spirit when up against hardship. Through the use of quotes and direct references
to the characters, Harrison tells a compelling story that aims to bring hope to the audience,
reminding them that the difficulties that humanity faces is incomparable to the resilience that
lies within each of us.\\


\subsection{Sample 3 - High Scoring \underline{Rainbow's End} Essay}


ZZEXP Rainbow's End III ZZEXPEND



Topic: ‘Least here we do things our way- no one breathin’ down our necks’\\
Home is central to identity in Rainbow’s End.
To what extent do you agree?\\

Amongst the racism and oppression, it is clear that ‘home’ is of the utmost importance to the
identity of Harrison’s Indigenous characters. Set against a backdrop of white Australian
superiority in the 1950s, ‘Rainbow’s End’ by Jane Harrison explores the central theme of
‘home’ in the context of identity. For the three generational main characters, their ‘humpy’
alongside the riverbank of the Mooroopna flats is integral to the make up of who they are
and essential for Harrison to showcase the strength and resilience needed for First Nations
to fight their marginalised status. Beyond the physical ‘home’, identity is shaped by the family
and culture surrounding the characters through the lens of white and Indigenous Australians.\\

Throughout the first part of the play, ‘home’ is physically represented as the ‘humpy’, a hut
with a dirt floor and hessian-covered walls. Despite the poverty of the main character’s living
conditions, it isn’t the place where they sleep that is important to them as much as the
culture surrounding them. For the character of Nan, "the Murray…Is your place isn’t it” due to
the immense sense of belonging she feels after her oppressive experience with the
Cummeragunja mission and the close proximity of her Aboriginal ‘family.’ In this way
Harrison highlights the role culture and a physical ‘home’ plays in an individual's identity. In
addition, the pride taken in the appearance and homeliness of the ‘humpy’ is a further aspect
of identity illustrated by Harrison. When in the play it comes time for the inspector to visit to
the ‘humpy’, Harrison has Nan and Gladys ‘nervously’ straightening the cushions and draws
attention to the ‘beautiful work’ of the crocheted pillows, emphasising the importance of
‘home’ to the characters. This is brought to the attention of the audience again during the
move to the ‘white and featureless’ Rumberala housing, and the characters immediately
think of ‘making curtains’ as a way to reassociate their identity with the new housing. The
contrast between the ‘home’ of the white Australian characters such as Errol and Nancy, and
the housing in Rumbalara and the flats is reflective of the identities of the characters and
how they feel entitled to act. For example, Nancy is portrayed as wealthy and privileged and
therefore acts in a privileged and entitled way, whilst the main characters are substantially
underprivileged and act humble, with kindness despite adversities.\\

The playwright suggests ‘home’ is more than a physical establishment but extends into a
‘home’ found within family and culture. The notion of family and culture as a ‘home’ is a
central theme throughout the play, explicitly explored between the misalignment of Errol and
Dolly’s definition of a ‘home’. For the character of Errol, he believes a ‘home’ is a place to
live, where only the wealth and geographical opportunities are relevant, however, Harrison
utilises his conversation with Dolly to pinpoint to ideal that “a real home is where people are
looking out for each other,” regardless of how many ‘presents’ are under the tree.
Furthermore, Errol reflects on the differences of their families, expressing to Dolly that his
Dad “expects me to call him sir” and ultimately concluding that despite the poverty of the
flat’s he bets ‘you have more fun than we do’. This conversation was pivotal to underlining
the impact of a ‘home’ with a supportive family and rich culture in creating a luster for life and
powerful sense of identity.\\

For the characters in the play, ‘home’ is closely intertwined with colonial suppression, the
white Australian’s control over Aboriginal housing had a heavy influence on the identity of
the characters. Throughout the play, Harrison makes reference to a difference in the
perception of Aboriginals from the flat’s and ‘townie types of cross-over Aborigines’ depicting
those who lived in town as slightly better off, with better opportunities. This differentiation
underpins the impact of location on identity as the ‘white fellas’ feel the Indigenous
characters need to ‘be absorbed into the community’ in order to be ‘assimilated’ and forced
to act in alignment with white culture. The control that the white characters, such as the
inspector, assert over housing is detrimental to the identity of the characters, it withdraws
their control and forces them to abandon their culture in order to be accepted. Moreover,
Harrison suggests that a ‘home’ can affect social mobility, Gladys in particular is obsessed
with ‘fitting in’ and being accepted and it becomes a large part of her identity. The ‘white
gloves’ and ‘Auntys shoes' are a symbol of her need to assimilate to white culture, acting in
direct contrast to the Aboriginal culture that makes up her ‘home’ and thereby her identity.\\

Throughout the entirety of the play, Harrison brings forth the different types of ‘home’, both
physical and as a culmination of family and culture. In this way it is evident to the audience
that home is closely entwined with the identity of all the characters in the play, from Errol
who has never experienced a ‘home’ within a family, to Dolly whom does not associate a
‘real home’ with a fancy establishment but instead with a loving family and cultural traditions.
The exploration into the impacts of oppressive attempts from the ‘white fellas’ to force
‘assimilation’, displays how a ‘home’ is more than a roof over one's head, but instead an
embodiment of an individual's unique identity as formed by relationships, culture, pride and
resilience.\\

\subsection{Sample 4 - High Scoring \underline{Rainbow's End} Essay}


ZZEXP Rainbow's End IV ZZEXPEND



Topic: Harrison's play celebrates characters who show strength through
adversity. Discuss.\\

Jane Harrison’s play Rainbows End is a powerful exploration of the lives of three
generations of Indigenous women, Nan Dear, her daughter Gladys, and granddaughter
Dolly, in 1950s rural Australia. Harrison presents their daily struggles with poverty, limited
access to education, and systemic racism, not to victimise them, but to celebrate their
resilience and determination. Through her portrayal of these women, Harrison critiques the
entrenched hierarchies that marginalise Indigenous communities, while affirming the
strength required to survive and resist. Ultimately, the play celebrates the capacity of
Indigenous women to endure adversity, even as it recognises the costs of their survival.\\

Harrison portrays how racial discrimination restricts Indigenous Australians’ opportunities
and dreams. This is evident in Scene Ten, when Gladys visits the bank in the hope of
securing Dolly a position as a teller. Gladys proudly emphasises her daughter’s academic
success, noting that Dolly has completed her “leaving certificate” and is “top of her class,” in
an attempt to counteract the stereotypes imposed upon Indigenous people. However, the
bank manager dismisses Dolly’s achievements, questioning how she would “fit in” and
whether she could be “relied on” to attend work daily, despite never having met her. Through
the portrayal of the bank manager, Harrison highlights how deep-seated structural racism
embedded in Australian institutions causes harmful stereotypes that prevent Indigenous
communities from fulfilling their full potential in society, denying them equal access to
employment. By juxtaposing the women’s ambitions with the bank manager’s prejudices,
Harrison critiques the injustice of a culture that confines Indigenous families to “picking” and
other forms of low-status labour regardless of merit. Yet she also affirms Gladys’ resilience,
showing her refusal to abandon her daughter’s dreams despite repeated setbacks. In doing
so, Harrison not only exposes how systemic racism denies equality but also celebrates the
persistence of Indigenous women who continue to assert dignity and demand recognition
against such barriers.\\

Harrison exposes the entrenched poverty endured by Indigenous families to critique the
systemic marginalisation of First Nations communities. Rather than presenting poverty as a
personal failing, Harrison illustrates how social and economic exclusion force the Dear family
into degrading conditions and low-cost labour, while simultaneously highlighting their
resilience in preserving dignity and hope. Nan Dear’s ironic description of their riverbank
“humpy” as “our own little palace” re-frames deprivation as pride, transforming humiliation
into humour. Similarly, Gladys expresses frustration at Dolly’s reliance on “hand-me-downs,”
revealing her determination to secure independence for her daughter. By infusing moments
of humour and defiance into depictions of poverty, Harrison refuses to allow Indigenous
identity to be defined by shame. Instead, she asserts the capacity of women like Nan and
Gladys to maintain agency and self-respect in the face of adversity. Dolly’s dream sequence
after her visit to the town tip symbolises a longing for a future beyond dispossession,
foregrounding aspirations that colonial systems sought to deny. Through these portrayals,
Harrison critiques structural inequalities, particularly the exploitation of Indigenous people as fruit pickers and scavengers, while simultaneously celebrating the resilience and strength of
Indigenous women to endure and reimagine possibility. Thus, Harrison’s authorial intent lies
not only in exposing the material realities of Indigenous poverty but also in affirming the
power of humour, pride and aspiration as acts of survival against systemic injustice.\\

Knowledge is presented in the play as both a pathway to empowerment and a means of
representing the differences in wisdom. Gladys has a strong belief that schooling will elevate
Dolly beyond ingrained disadvantage, insisting that she “get an education” so she can “be
someone.” This determination highlights her strength as a mother, willing to challenge
generational cycles of poverty. Harrison deliberately frames this ambition with irony: while
Dolly’s academic success is undeniable, she remains judged through stereotypes and
institutional racism, suggesting that education alone cannot dismantle systemic barriers. At
the same time, Harrison complicates the notion of knowledge by staging contrasts between
Indigenous oral tradition and Western authority. When Errol dismisses Dolly’s experience
with snakes, Gladys’ clarification that it was a harmless “carpet snake” undermines the
authority of the encyclopedia and asserts that lived experience offers forms of wisdom
inaccessible to formal schooling. By juxtaposing these knowledge systems, Harrison
critiques the way white institutions privileged their own frameworks as universal while
marginalising Indigenous perspectives. Her authorial intent is doubled to expose how
education in the 1950s was arranged by race and class, yet also to celebrate the resilience
of families like the Dears, who embraced education without relinquishing pride in their
cultural knowledge. Harrison both honours the courage of Indigenous families who value
education as a tool for advancement and critiques the system of schooling planted between
classes in the 1950s.\\

Harrison’s Rainbow’s End critiques the systemic racism that restricted Indigenous
Australians while celebrating the resilience of women who refused to be silenced. By
exposing institutional prejudice, entrenched poverty and competing forms of knowledge, she
highlights both the barriers that denied equality and the strength of Indigenous families who
asserted dignity despite them. In doing so, Harrison affirms that hope, humour and cultural
pride are not just survival strategies but powerful acts of resistance against injustice.\\

\subsection{Sample 5 - High Scoring \underline{Rainbow's End} Essay}


ZZEXP Rainbow's End V ZZEXPEND



Topic: '\underline{Rainbow's End} is a bitter condemnation of colonisation.' Discuss.\\
Cemented within a world enveloped by overpowering white culture, Rainbows End seeks to
explore the lengths of condemnation expressed by Indigenous Australians, juxtaposed with
the desire to be included. Within the play, Harrison blatantly reveals consistent resistance
and dislike aimed towards white authority yet contrasts this with younger generations who
desire integration and inclusion in society. By depicting Indigenous women in this way,
Harrison can characterise older generations as resistant to white knowledge and eternally
furious at the erasure of Indigenous culture yet depicts younger characters as open to
colonisation to succeed in their corrupted society. Overall, Harrison criticises the unjust
dilution of Indigenous culture, simultaneously advocating for Indigenous Australians to
continue fighting to be heard over white dominance.\\

Harrison calculatingly embeds Indigenous resistance throughout the play, through
unwavering denial, which represents a resilience to maintain Indigenous culture, despite
attempts to silence it. Harrison evidence this though Nan Dear showing her despise for
Errol’s encyclopedias, asking him what it says, “in that there encyclops about Aboriginals”.
Harrison’s lexical choice of characterising Nan through her informal and grammatically
incorrect language of “that there” and “encyclops”, illustrates Nan’s refusal to submit to white
assistance in education, as the encyclopedias are based off Eurocentric knowledge,
maintaining her stance in continuing Aboriginal knowledge and culture no matter the state of
her academics. Nan Dear also directs this hate towards Errol, as he is the delivery boy for
them, and is white, as she sees him as a messenger for white authority, being cautious of
letting him get too close. Harrison lambasts the erasure of Indigenous knowledge, yet
cautiously condemns the anger being taken out on unrelated and innocent white individuals.
Harrison juxtaposes this condemnation of colonisation through dream sequences, when
Gladys begins to dream of the possibility of Dolly becoming “a graduate”. Gladys’ willingness
to accept and integrate into the complexities of white society in order to provide her daughter
with a proper education, contrasts that of Nan’s ideologies relating to white assistance,
demonstrating the differences in generational beliefs. As encyclopedias are based on white
knowledge, this further illustrates Gladys’ will to accept their knowledge as superior to
succeed. The dream sequence represents Gladys’ subconscious suggestion that her
daughter may never be able to be “a graduate” due to her needing to follow Nan’s beliefs as
the eldest of the family, being constrained to the limitations rightful morals to uphold
Indigenous independence. Harrison champions the pursuit of education despite society
restraining full ability yet criticises the neglect of protecting Indigenous culture and forgetting
true identity.\\

Harrison impatiently lambasts the evident erasure of Indigenous culture, through white
authorities directly and indirectly pushing it, which exposes their never-ending desire to
abolish Aboriginal knowledge and existence. Harrison represents this through the humpy
inspector asking if Gladys has heard of “the concept of assimilation”, as a promise that
“things will change.” The contradictory nature of the inspectors’ intentions and promises
represent white ignorance towards Indigenous perspectives and opinions, as the inspector
didn’t let Gladys answer if she has heard of the “concept” of assimilation, continuing to talk
about it as if she isn’t there. The inherent dominance from the inspector pushing for the
absorption of Indigenous Australians into white society conveys the collective idea of
abolishment, and then stating that “knowledge is power” continues to falsely represent the
freedom of Indigenous Australians white authorities distribute. Harrison demonises the
dominance of white imposed hierarchies that put Indigenous Australians at an unfair disadvantage, keeping them there without proper consultation. Harrison further bolsters this
through the stage direction of Gladys and Nan “exchange[ing] worried looks”, after the
inspector’s suggestion. Their “worried looks” emphasise the desire to remain where they are,
away from the “mercy of the manager” which Nan Dear recalls from her traumatic days of
living at a mission. The family can easily accept help from white authorities by submitting to
their dominance yet refuse to so they can maintain their independence and keep their
cultural heritage and identity intact, living of their own free will. Their desire to stay at their
humpy and not receive better living conditions demonstrates their appreciation for the
opportunity to live by themselves and the gratitude for where they are, no matter the
difficulties that surround these unfavourable living circumstances. Harrison bolsters the idea
of maintaining Indigenous heritage but questions the resistance towards white assistance if it
is at the expense of health and safety.\\

Harrison diametrically contradicts the stereotypical resentment towards colonisation through
Indigenous youth desiring to submit to white dominance, which raises questions about the
rightful path to follow when aiming to achieve Indigenous independence. Harrison depicts
this through Dolly’s desire to be included in white society, when Dolly is attending a dance in
Shepparton, and exclaims how “everyone’s going to be there”. Dolly being the youngest of
the family is the most exposed to whitewashing and diluted Indigenous culture, hence her
powerful ambition to attend dances and wear Eurocentric 1950s outfits and accessories, to
be like “everyone” who’s going “to be there”. Her outfit is symbolic of the tensions between
her socio-economic status and her dreams, acting as a constant reminder of her inability to
completely be accepted in white society. Harrison praises the confidence of Indigenous
youth to break free from limiting social and gender norms in order to achieve one’s dreams
yet warns against the dangers of being denied acceptance from such rebellion. Harrison
further reinforces a rebellion in the opposition to white superiority through Gladys’ dialogue
when she dares to believe that Dolly can succeed by not letting “her address stop her in life”.
Gladys’ rhetorical question towards Nan exemplifies the difference and experience of
generations, Gladys believing Indigenous women can succeed without the acceptance of
white people. She is met by Nan Dear’s cold response telling her to “get a grip”, snapping
her back to reality and reminding her of the fact that they are a minority that is
disadvantaged in a backwards society that will never favour them. Gladys daring to stray
from the wisdom of her mother represents an underlying feeling of hope for change, despite
claims and actions that proves change is unlikely to occur. Gladys and Dolly don’t harbour
severe resentment towards their society like Nan Dear, as they have less experience with
unwavering discrimination. It outlines their desires to blend into the white community at the
expense of their autonomy, in order to gain increased access and succeed in achieving a
better life than what they have. Harrison champions the unwavering objective of striving to
be better in one’s life, while cautioning of the implications of aiming too high and ignoring the
safety provided by one’s community.\\

Conclusively, Harrison continuously exhibits justified resistance to colonisation of Indigenous
culture yet remains hopeful that younger generations may have the courage to break the
discriminatory cycle. Harrison encapsulates this through continued vilification of white
knowledge and education, and the injustice of institutional authority, while postulating a
juxtaposed change in views through the acceptance of colonisation within younger
generations. Harrison ultimately condemns the unjust erasure of Indigenous culture and
identity, while pushing the notion that the amalgamation of cultures may lead to increased
opportunities for Indigenous Australians and should not be seen as an entirely demonised
occurrence.\\

\subsection{Sample 6 - High Scoring \underline{Rainbow's End} Essay with Annotations}


ZZEXP Rainbow's End VI with Annotations ZZEXPEND



Topic:
\textbf{How does Harrison explore the concept of 'home' in \underline{Rainbow's End}?}\\

\textit{Introduction}

In \underline{Rainbow's End}, Jane Harrison explores the concept of home by presenting Australia through the lens of three generations of Indigenous Australian women (1). Set in Shepparton/Mooroopna during the 1950s, against the backdrop of assimilation policies and the aftermath of the Cummeragunja walk off (2), the play sets up home as a source of conflict between cultures. For white Australia home is a house full of modern conveniences, and therefore superior to a humpy on the floodplains, that needs to be monitored by government inspectors (3). Through three generations of Indigenous women, Harrison proposes that a house containing a nuclear family is not the same thing as a true home built on harmony with nature, connection to Country and community, and that full reconciliation depends on white Australia learning that difference (4).\\

\textit{Annotations:}

(1) Naming the structural device (three generations of women) in the first sentence signals to a marker that the essay has a clear analytical architecture before a single quote appears — the body paragraphs can then be trusted to deliver on this promise.

(2) Historical specificity like this is worth far more than a vague "in the past" reference. The Cummeragunja walk-off is real context that grounds Nan Dear's later dialogue; using it here shows research beyond the play itself, which examiners reward as sophistication.

(3) This sentence defines the opposing side of the argument (the white Australian view of home) before the essay has even reached a body paragraph — setting up the contention as a genuine debate rather than a one-sided assertion.

(4) A strong contention doesn't just answer the question, it makes an arguable claim with a "because" attached — here, that reconciliation is conditional on white Australia changing, not just Indigenous families adapting. That specificity is what separates a contention from a topic restatement.\\

\textit{Body Paragraph 1}

White Australia's idea that a house makes a home is set against the Indigenous idea that home is mostly about staying connected to country and culture (1). Older Indigenous generations belong to a time before the government controlled Indigenous people's independence through assimilation policies aiming to eradicate their culture, policies that included the Stolen Generations, that were only wound back in the late 1960s (2). For Nan Dear, home is the Murray River and Cummeragunja, the mission where managers "forced us to leave... to leave Cummeragunja. Our home," she says. Through this, Harrison shows that Nan Dear's memories of Cummeragunja aren't just nostalgic but political, tied to the loss of autonomy under assimilation policies (3). Her fierce independence, evidenced by her "we sink or swim on our own," is why she'd rather live in a humpy, insisting "it's ours." The humpy, respectfully described as "clean and homely" by Harrison in her setting descriptors, is a symbol of self-determination, its ownership set against the government-controlled housing offered to Indigenous people elsewhere (4). But for Nan Dear, home means independence and agency, a connection to land and heritage, rather than four walls or new fittings. Even the floods don't shake this as she accepts that "that's nature," in contrast to Gladys and Dolly, who keep "fighting against how things are." Harrison respects a generation willing to endure struggle, if it means holding on to agency and connection to country. As such, home is more than a place for one family to live, but a broader cultural, historical and geographical concept, and Nan Dear's contentment in the humpy shows this isn't sacrifice, but freedom (5).\\

\textit{Annotations:}

(1) The topic sentence names both sides of the tension (white materialism versus Indigenous connection) rather than just describing Nan Dear — this keeps the paragraph anchored to the actual question about "home" rather than drifting into a character study.

(2) Embedding a factual clause about assimilation policy and the Stolen Generations mid-sentence, rather than as a separate statement, is an efficient way to fold context into analysis without stalling the argument.

(3) This is the paragraph's strongest analytical move: reframing a quote that could be read as simple grief ("our home") as evidence of a political argument about autonomy. Markers reward this kind of reinterpretation over surface-level quote description.

(4) Citing Harrison's own stage direction ("clean and homely") as evidence is a technique many students overlook — playwrights' descriptive language is authorial comment in a way dialogue alone isn't, and using it shows attention to the whole text, not just spoken lines.

(5) The paragraph's final sentence links back to the contention's key terms ("more than a place... but a broader... concept") rather than simply summarising the plot — this is what makes a concluding sentence do analytical work instead of just closing a paragraph.\\

\textit{Body Paragraph 2}

Harrison also asserts that change under government control has only ever meant less independence for Indigenous Australians (1). This distrust of government comes through when the Inspector announces that "things will change," regarding the humpy's deficient living conditions such as sanitation and no electric lights, a statement met with "worried looks between Nan Dear and Gladys," as Harrison's stage directions depict. This is clearest in the ironically titled "Home Sweet Home" scene (2), where the radio spruiks the move to Rumbalara as "the most vigorous attempt yet to solve Aboriginal housing," but Harrison's stage directions describe the houses as "concrete, small, white and featureless," and it's here that Nan Dear develops a cough (3). This is not the kind of house they expected, Harrison even portraying the Dears having to sleep on the roof during the hotter nights. When Gladys turns her want for decent housing into activism, over the radio she declares that she's going "to fix my own house," because "us Aboriginals are not welcome in the townships" and are left with nothing but "concrete humpies" the council dares to "call… houses." Harrison's point is that Indigenous people must fight just to get a decent home, and that fight shouldn't cost them their identity or their right to decide what happens within their own walls (4). Thus, Harrison demonstrates that Indigenous people cannot rely on the white man's government to improve living conditions or recognise their land rights, as Gladys declares unequivocally, "I belong here - this is my land!"\\

\textit{Annotations:}

(1) Opening with an authorial claim ("Harrison also asserts...") rather than a plot description keeps the analysis operating at the level of what the playwright is arguing, which is what Section A criteria are actually assessing.

(2) Naming the scene title and noting its irony in the same breath is a compact, efficient piece of structural analysis — it shows the essay is tracking the play's own use of dramatic irony, not just quoting from memory.

(3) Pairing the government's own promotional language ("the most vigorous attempt") against Harrison's flat, clinical stage direction ("concrete, small, white and featureless") is a strong juxtaposition technique — letting the text's own contradictions do the analytical work rather than over-explaining them.

(4) This sentence does double duty: it states the play's argument and pre-empts a counterpoint (that better housing might have been enough), which shows awareness of the topic's complexity without derailing into a full alternate paragraph.\\

\textit{Body Paragraph 3}

Harrison also argues that modern generations of Indigenous women can and should aspire to go to universities, even if that means leaving those you love in the rural areas (1). Dolly's role is to reflect the agency of choice for Indigenous women, not because a white man wants to marry her, but because of the educational opportunities there. When Errol offers Dolly "a better life" in Melbourne, she refuses because that would mean leaving her Nan, her mother, and her community in Shepparton and Mooroopna. With good intentions, he wants Dolly to have "a real home," not understanding that this "world" isn't hers, because for Dolly, "it don't matter if the floors are dirt." Their life, she insists, echoing Nan Dear, "isn't perfect but... it's ours" (2). Harrison shows the reconciliation between white and Indigenous Australia through this relationship, as Errol eventually comes to accept that where Dolly belongs, and her family, is important - "to you, and me," he says. By the end, Dolly and Errol symbolise a reconciliation between black and white Australia based on mutual respect, not on Dolly having to become someone else to fit into a white Australian world (3).\\

\textit{Annotations:}

(1) This paragraph deliberately introduces a new dimension of "home" (choice and generational change) rather than repeating the previous paragraph's argument — a common weakness in weaker essays is three body paragraphs that all make the same point with different quotes.

(2) Noticing that Dolly's line echoes Nan Dear's earlier "it's ours" is intertextual analysis within the one play — tracing a phrase across generations of characters demonstrates a command of the whole text rather than isolated scenes.

(3) Using the Dolly-Errol relationship as a symbol for broader reconciliation, rather than just a subplot romance, elevates the paragraph from character analysis to thematic analysis, which is what the higher bands of the criteria are looking for.\\

\textit{Conclusion}

Looking at all three generations together, Harrison makes the point that home can't just be a house on Crown land (1). White Australia's focus on modern features and material comfort ends up looking empty because it comes without family closeness, while the Indigenous characters, even with their differences, keep coming back to land rights, community and love as what really makes a home. Nan Dear holds onto that connection no matter what the government tries to take from her, Gladys fights to have it recognised on her own terms, and Dolly shows that it can survive change and distance without being lost (2). Home, for Harrison, isn't something white Australia can hand over through better housing policy. It must come with genuine respect for Indigenous identity, land and culture. That's ultimately what Dolly and Errol's relationship offers, and what the play argues Australia still should strive to achieve (3).\\

\textit{Annotations:}

(1) A strong conclusion returns to the contention's own language ("home can't just be a house") rather than introducing it for the first time here — this shows the essay has been building towards this point all along, not arriving at it by accident.

(2) Summarising each generation with one precise verb (holds, fights, shows) is an efficient way to synthesise three paragraphs of evidence into a single sentence without simply repeating quotes.

(3) Ending on what "Australia still should strive to achieve" extends the essay's implications beyond the play itself, which is a hallmark of a top-band response: the argument matters because it says something about the real world, not just about the text.\\

\cleardoublepage

\clearpage
\thispagestyle{empty}
\hfill


\cleardoublepage

\clearpage
\thispagestyle{empty}
\hfill
\clearpage
\part{Analytical Text Response Essays}
\clearpage
\clearpage
\thispagestyle{empty}
\hfill
\clearpage

\section{Introduction}
From the very beginning of your schooling career, you have been introduced to several different methods of writing essays. Although all structures have their advantages and disadvantages, it is essential that you have a fool-proof method of writing essays for your end of year examination.\\


Having a solid method of writing analytical essays and practicing these throughout the year will result in a strong foundation, enabling you to tackle any essay question with confidence. The key to success lies not only in understanding the content but also in being able to present your ideas in a clear, structured, and coherent manner. With a well-practised approach, you will be able to break down complex texts, analyse them effectively, and construct a well-supported argument that demonstrates both depth and clarity of thought.\\

In the following, we will explore a structured approach to writing an analytical text response. By familiarising yourself with this framework, you will be better equipped to organise your ideas, manage your time efficiently, and ultimately maximise your performance in the exam.\\

\newpage
\section{What is Needed in an Essay}

An essay is a structured piece of writing that presents an argument, analysis, or narrative on a particular topic. The purpose of an essay is to communicate a clear, well-reasoned point of view, often backed by evidence and examples. Essays are typically composed of several key parts: an introduction, body paragraphs, and a conclusion. Each section plays a specific role in conveying the writer’s ideas and helping the reader understand and engage with the topic at hand.

\subsection{Introduction}
The introduction serves as the starting point of the essay, setting the stage for the discussion to follow. It introduces the topic, provides necessary background information, and presents the thesis statement, which outlines the main argument or purpose of the essay. A strong introduction grabs the reader’s attention, often through a compelling hook, and establishes the importance of the topic. It should also provide a clear roadmap of what the essay will cover, guiding the reader's expectations and providing a sense of structure.

\subsection{Body Paragraphs}
The body of an essay is where the main arguments or points are developed and supported with evidence, examples, or analysis. Typically, an essay contains three body paragraphs, although this can vary depending on the length and depth of the topic. Each paragraph should focus on a single aspect of the thesis statement and contribute to the overall argument.

\subsection{Conclusion}
The conclusion is the final part of the essay, where the writer sums up the key points made in the body paragraphs and reinforces the thesis statement. It provides closure to the discussion and leaves a lasting impression on the reader. A strong conclusion restates the main argument in a fresh way, summarising the key insights presented in the essay. It may also offer suggestions for further exploration or suggest the broader implications of the topic.
\newpage


\section{Planning}
\subsection{Deconstructing the Prompt}
When first planning your essay, you must directly respond to the prompt that is presented to you, and consider the key authorial messages that could help you respond with insight. This involves three important steps:
\begin{itemize}
    \item Find the key concepts. This will most likely correlate with the key concepts you have studied throughout the unit, you should also make sure to identify synonyms of these concepts too!
\item Find any instructional verbs. These are words that instruct you to answer the prompt in a
certain manner, in English the main verbs to look out for is discuss, do you agree, to what extent or how.
\item Identify any expectations/limitations. These ore words that limit what you con include in
your essays and set on expectation on what you should be discussing. 

\end{itemize}

\textbf{Example prompt:}
\textit{`The dream sequences in \underline{Rainbow's End} are the characters' only escape from the depressing reality of life.' Do you agree?}\\

This is a very straightforward prompt—you need to demonstrate a multiple-sided response to what Harrison perceives as the role of dreams and aspirations in her characters' lives. There are no limits to what should be discussed here, but you are expected to focus your discussion on the ideas of dreams, escape, and reality as they relate to the values, identities, and choices of Nan Dear, Gladys, and Dolly. Consider the historical context of 1950s assimilationist Australia, racism, class, and the characters' evolving sense of agency throughout the play.

\subsection{Contention}
Now that you understand what the prompt is asking, we can begin to think about what our interpretation of Harrison’s message is. Let’s consider the first prompt again:\\

\textit{`The dream sequences in \underline{Rainbow's End} are the characters' only escape from the depressing reality of life.' Do you agree?}\\
We can approach this in multiple ways. Here are three simplistic interpretations (agree, disagree, somewhere in between):
\begin{itemize}
    \item Harrison suggests that dreams offer the characters comfort, but they cannot change the reality of their circumstances.

\item Harrison implies that dreams are the characters' only refuge from a society that denies them dignity and opportunity.

\item Harrison demonstrates that it is action, not dreaming, that ultimately allows the characters to escape their circumstances.
\end{itemize}

However, these interpretations are quite general and lack specificity. They don’t tell the reader exactly what you are arguing. Let’s improve them with more precise language and deeper knowledge of the text:

\begin{itemize}
    \item Harrison portrays dreams as both sustaining and limiting, suggesting that Gladys’s fantasies of meeting the Queen preserve her dignity while simultaneously revealing how deeply colonial ideals of acceptance have been internalised.

\item Harrison implies that, for a family denied housing, education, and respect, the imagination becomes a vital site of survival and self-worth in a society that offers no material escape.

\item Harrison critiques passive dreaming by charting Gladys’s transformation from daydreamer to activist, suggesting that genuine escape from injustice comes not from fantasy but from finding one’s voice and demanding change.
\end{itemize}

These improved interpretations use specific language, engage with the complexity of the characters’ motivations, and demonstrate a sophisticated understanding of the play’s ideas.

\subsection{Extending Sentences - Connectives}
To create a 'sophisticated' idea you need to avoid single ideas, all of the contentions above are singular and lack connections to other ideas. We can connect them together using three simple connectives (and their advanced synonyms!)

\begin{longtable}{|p{2.8cm}|p{6.8cm}|p{5.2cm}|}
\caption{Useful Conjunctions and Their Purposes} \label{tab:conjunctions} \\
\hline
\textbf{Conjunction} & \textbf{Purpose} & \textbf{Synonyms} \\
\hline
\endfirsthead

\hline
\textbf{Conjunction} & \textbf{Purpose} & \textbf{Synonyms} \\
\hline
\endhead

\hline
\endfoot

\textbf{Because} & This conjunction explains why the information you have stated is true, demonstrating greater understanding of the author's purpose. & As, as a result of, considering, being, for, through, seeing as. \\
\hline
\textbf{But} & This conjunction adds contradictory information, allowing for greater debate. & However, nevertheless, yet, despite, even so, nonetheless, except. \\
\hline
\textbf{So} & This conjunction allows for inferencing the cause and effect of the statement, which demonstrates greater thinking and ideas about the text. & Therefore, consequently, accordingly, leading to. \\
\hline

\end{longtable}

From here, we can take these singular idea contentions and extend them with these connectives in order to make multi-idea contentions that can create sophisticated links between ideas:

\subsection{Topic Sentences}
Now we can create topic sentences that allow to explore complex ideas linked using BECAUSE, BUT \&
so.

\subsection{Authorial Message/Intent - NOT PLOT OR CHAPTER!}

The sentence should be focused on an authorial message, not just a concept or character­ make the author active -the most common start is the \textbf{author}+ verb:\\

E.g \textbf{Harrison} explores/communicates/illustrates etc.

\subsubsection{Authorial Verbs}
\begin{longtable}{|p{3.5cm}|p{3.5cm}|p{3.5cm}|p{3.5cm}|}
\caption{Analytical Verbs for Text Analysis} \label{tab:analyticalverbs} \\
\hline

\hline
\endfirsthead

\hline

\hline
\endhead

\hline
\endfoot

Conveys     & Explores    & Demonstrates & Emphasises \\
Communicates & Depicts     & Illustrates  & Challenges \\
Identifies   & Frames      & Reveals      & Highlights \\
States       & Implies     & Suggests     & Underlines \\
\hline

\end{longtable}

\subsubsection{Prepositional Phrases}
Another way to structure this sentence is to use a \textit{prepositional phrase} at the beginning of the sentence that introduces ideas to vary your sentence structure:

\textbf{Throughout the text, Harrison} explores..\\

\begin{longtable}{|p{3.5cm}|p{3.5cm}|p{3.5cm}|p{3.5cm}|}
\caption{Textual Analysis Phrases and Sentence Starters} \label{tab:analysisphrases} \\
\hline
\hline
\endfirsthead

\hline
\hline
\endhead

\hline
\endfoot

Throughout           & By (authorial verb)     & In keeping with    & With an eye to \\
Through              & From the outset         & In the absence of  & Under the \\
By                   & Beyond the              & Within the         & To the extent \\
During               & Concerning the          & Without            &               \\
\hline

\end{longtable}

\subsubsection{Specificity and Relevancy}
Examiners punish vague or general statements, especially those that lack clear links to the topic or the author’s message. For example:\\

\textit{By emphasising the importance of dreams in the play, Harrison shows that they matter to Gladys.}\\

This is not technically incorrect—but it is far too general and superficial. It doesn’t demonstrate insightful analysis, nor does it show a clear understanding of Harrison’s authorial intent. This kind of vague sentence would not receive full marks.\\

Here’s an improved version:\\

\textit{By emphasising Gladys’s fantasies of royal recognition, Harrison portrays how dreams function not as an escape from reality, but as a fragile shield of dignity for those denied respect and opportunity by an assimilationist society.}\\

This is better. It is clear, specific, well-structured, and shows awareness of both the theme (dreams vs reality), character motivation, and authorial message. It also directly links to the contention: that dreams sustain the characters, even in circumstances that offer no material escape.\\

The issue you will face is that there is only so much you can discuss in your paragraph from here, as you have only examined one idea—dreams sustain dignity. To fix this, lets' go back to an old friend:\\

\subsubsection{Because, But, So - Make a Complex Sentence!}
Too often, students fail to complete long paragraphs because they are stuck repeating themselves or have run out of things to say about their topic they introduced. Therefore, if we have two ideas in our topic sentence, we can 'hinge' these sentences together which demonstrate complex understanding, writing and structure - all of them are key parts of the rubric you are judged on!
Here's an example:\\

\begin{quote}
    By emphasising Gladys’s daydreams of royal recognition, Harrison conveys how, for a family denied dignity and opportunity, the imagination becomes a vital refuge from an unjust reality—ultimately suggesting that dreams alone cannot deliver change, and that true escape demands the courage to speak and act.
\end{quote}

Now, you have two linked ideas:
\begin{itemize}
    \item Dreams preserve dignity and hope for the characters in an unjust society.
    \item Dreaming alone, without action, cannot deliver genuine change.
\end{itemize}

This provides a clear structure for an entire paragraph and sets you up to use a hinge sentence later to develop further layers—e.g., how this 'dreaming” can also reveal internalised colonial ideals or delay confrontation with injustice.\\

By planning your paragraph in this way, you’ll ensure:
\begin{itemize}
    \item Specific and sophisticated analysis
    \item Strong authorial intent
    \item Complex understanding of the characters’ motivations and Harrison’s critique of assimilationist policies in 1950s Australia
\end{itemize}












\newpage
\section{Analytical Text Response Structure}

\subsection{Topics}
There are three types of questions you may receive on a SAC or exam. These are the 'Discuss', 'To what extent' and 'Do you agree?' statements. These are all tackled in a slightly varying manner.\\

\begin{itemize}
    \item Discuss : This is probably the most common type of prompt that you will encounter. When dealing with such prompts, it is important that you discuss as many aspects of it as possible. One way to improve your discussion here is to include a paragraph that rebuts the prompt, like what you do in your persuasive essays. However, since it is a discuss prompt, you do not need to make the two sides of your argument extremely distinct.\\
    \item To what extent: These topics require you to make a decision. Never say 'it is to some extent' in your contention. You really need to assign a portion of something. The best way to do this is to start with subordinating conjunctions like 'whilst” or 'although”. To What Extent questions rarely want you to agree or disagree - they want you to explore a shade of grey.\\
    \item How: The 'How' topics usually focus on the construction of a text - your evidence should contain metalanguage like 'characterisation' for instance.\\
\end{itemize}

\subsection{Introduction}
The introduction is the first piece of your writing that will be examined. For this reason, it should be clear, direct and show your understanding of the topic. The introduction is there to show your examiner how you wish to respond to the set topic.\\

\textit{If your topic contains a quote, ensure that you mention it somewhere in your introduction to show you are addressing the topic.}


\subsubsection{Structure}
An introduction must include the following:\\

\paragraph{A General Line of introduction}
Too often our introductions lack any insight and mainly begin with author+text+background information which has nothing to do with the actual topic. We should attempt to be relevant with our opening statement and discuss why this concept is important.\\

\begin{longtable}{|p{4.5cm}|p{4.5cm}|p{6cm}|}
\caption{Textual Analysis Starters, Verbs, and Noun Phrases} \label{tab:analysislanguage} \\
\hline
\textbf{Phrases to Introduce} & \textbf{Analytical Verbs} & \textbf{Noun Phrases and Structures} \\
\hline
\endfirsthead

\hline
\textbf{Phrases to Introduce} & \textbf{Analytical Verbs} & \textbf{Noun Phrases and Structures} \\
\hline
\endhead

\hline
\endfoot

Throughout literature... & Conveys & the challenges of... \\
Authors are often preoccupied with... & Explores & the pressures and difficulties of... \\
Much of human thought has focused on... & Emphasises & the dangerous nature of... \\
In these ... worlds created by authors... & Directs & the devastating impact of... \\
Most stories are characterised... & Engages & the importance of... \\
The importance of... is often highlighted... & Establishes & the consequences of... \\
 & Depicts & the conflict between... \\
 & Grapples & the ... it requires... \\
 & Highlights & the way ... can... \\
 & Illustrates & the ways in which ... can... \\
 & Challenges & the way ... shapes and influences... \\
 & Composes & how ... leads to... \\
 & & how ... is able to... \\
 & & how ... creates... \\
 & & how ... causes... \\
 & & how ... challenges us to... \\
 & & how ... forces us to... \\
 & & how ... compels us... \\
 & & how not only ... but also... \\
 & & what it means to be... \\
 & & what it means for... \\
 & & what happens when... \\
 & & how a world where... \\
 & & how societies in which... \\
 & & how times of... \\
 & & how situations in which... \\
\hline

\end{longtable}

\paragraph{Introduce the Text, Author and Contention}
Now we introduce the text and author and the audience. Then state how it relates to the topic specifically. These need to be connected fluently, when writing we should be relying on connectives to make sure there is a logical flow in our expression.\\

Depending on the prompt itself, it might be a wise idea to link either/both of characters and setting to the topic of the essay prompt. We can do this as such:\\

Linking a Character to a Topic:
\begin{longtable}{|p{4.5cm}|p{4cm}|p{6cm}|}
\caption{Sentence Starters: Character Discovery and Purpose} \label{tab:discoveryphrases} \\
\hline
\textbf{Introductory Phrase} & \textbf{Verb of Discovery} & \textbf{Linking Phrase} \\
\hline
\endfirsthead

\hline
\textbf{Introductory Phrase} & \textbf{Verb of Discovery} & \textbf{Linking Phrase} \\
\hline
\endhead

\hline
\endfoot

As the main character...    & Discover/s    & in order to... \\
As the leading character... & Find/s        & so that... \\
As the hero/ine...          & Learn/s       & by... \\
As the main focus...        & Understand/s  & to become... \\
As the protagonist...       & Accepts/s     & to work towards... \\
                            & Admit/s       & since... \\
                            & Challenges/s  & it is not as important as... \\
                            & Confront/s    & isn't as necessary as... \\
                            & Seeks/s       & can never be a substitute for... \\
                            & Realise/s     & \\
\hline

\end{longtable}




\textit{N.B. that the name of the book is always underlined.
I.e. \underline{Sunset Boulevard} or \underline{Rainbow's End}}.\\

\textbf{Contention}\\
A contention is the author's point of view on the topic. Since you are the author in this scenario, you need to write what your position on the topic. A contention should be short, to the point, and engage with the topic. On this falls the crux of your essay.\\

\textit{Do not be afraid to agree/disagree with a topic. You are allowed to have a conditional contention.}\\

\newpage
\subsection{Body Paragraphs}
The body paragraphs form majority of your essay. The body paragraphs are the main reasons that the contention your provided is true. Your body paragraphs need to be able to demonstrate your understanding of the text and the fluency of your language. You should ensure that every sentence (and thus every word) brings something towards your argument.

\subsubsection{Structure}
Below is the structure for a body paragraph using 'Th.E. S.A.M.L.' Method. This method of writing body paragraphs was developed in order to maximise possible marks. Only the 'TH' and 'L' elements of this structure need to be in the order given.



ZZSAML



\subsection{Basic Example 'Th.E S.A.M'}
\begin{quote}
   'Obsession with one’s self is what ultimately leads to the downfall of characters. This can be seen with Norma exclaiming ‘who else?’ when Joe questions her about her script. The script itself is also a representation of Norma’s obsession as she has written several hundred pages where she stars in every scene. Wilder uses this to demonstrate that Norma is unable to move on from anything besides herself. The slow camera pan whenever the script is shown foreshadows Norma’s impending doom which links is inherently linked to her self-absorption. In this way, obsession with self ultimately leads to the downfall of several characters in \underline{Sunset Boulevard}.'

\end{quote}

\subsection{Writing About Ideas: Adding Adjectives to Nouns}
One of the easiest ways of improving your writing is to add adjectives to nouns, this helps us to create an 'idea phrase'. Here is a list of commonly used nouns in analytic writing, a good way is to use these as springboards to finding even more complex and advanced nouns.


\begin{longtable}{|p{2.8cm}|p{2.8cm}|p{2.8cm}|p{2.8cm}|p{2.8cm}|}
\caption{Idea Nouns for Analytical Writing} \label{tab:ideanouns} \\
\hline
\textbf{Column 1} & \textbf{Column 2} & \textbf{Column 3} & \textbf{Column 4} & \textbf{Column 5} \\
\hline
\endfirsthead

\hline
\textbf{Column 1} & \textbf{Column 2} & \textbf{Column 3} & \textbf{Column 4} & \textbf{Column 5} \\
\hline
\endhead

\hline
\endfoot

Growing up     & Survival     & Justice       & Friendship     & Loss \\
Independence   & Conflict     & Humanity      & Family         & Guilt \\
Coming of age  & Courage      & Prejudice     & Relationship   & Redemption \\
Maturity       & Strength     & Oppression    & Belonging      & Hope \\
Identity       & Bravery      & Injustice     & Loneliness     & Grief \\
Discovery      & Persistence  & Power         & Loyalty        & Sorrow \\
Revelation     & Fortitude    & Greed         & Betrayal       & Memory \\
Concealment    & Heroism      & Selfishness   & Love           & Future \\
Blossom        & Resilience   & Competence    & Hatred         & Past \\
Deteriorate    & Cowardice    & Potential     & Envy           & Melancholy \\
\hline

\end{longtable}

We then can pair these ideas with adjectives of degree - these tell us how big, small, frequent, fast or slow something is, here is a list to help build your writing vocabulary:\\

{small
\begin{longtable}{|p{2.5cm}|p{3cm}|p{3cm}|}
\caption{Adjectives of Degree for Analytical Writing} \label{tab:adjectivesdegree} \\
\hline
\textbf{Category} & \textbf{BIG (Intensifying)} & \textbf{SMALL (Minimising)} \\
\hline
\endfirsthead

\hline
\textbf{Category} & \textbf{BIG (Intensifying)} & \textbf{SMALL (Minimising)} \\
\hline
\endhead

\hline
\endfoot

\textbf{SIZE}         & Deep, Immense, Extensive, Overwhelming         & Empty, Insignificant, Little, Tiny \\
\textbf{STRENGTH}     & Extreme, Powerful, Fierce, Determined           & Weak, Uncertain, Fragile, Feeble \\
\textbf{SPEED}        & Rapid, Swift, Impulsive, Sudden                 & Slow, Reluctant, Hesitant, Glacial \\
\textbf{CONSISTENCY}  & Frequent, Constant, Incessant, Unchanging       & Unsteady, Unreliable, Erratic, Unstable \\
\hline

\end{longtable}

}

\subsection{Hinge Sentences}
In your topic sentence, you can introduce two related ideas. This creates room for a more complex, insightful paragraph. But remember: when you shift focus from the first idea to the second, you must signal this change to the reader clearly. This is where the hinge sentence comes in.\\

A hinge sentence acts like a bridge—it acknowledges what you've already explored, then transitions to the second idea you introduced in your topic sentence. Here’s how that works:
\begin{quote}
    By emphasising Gladys’s daydreams of royal recognition, Harrison conveys how, for a family denied dignity and opportunity, the imagination becomes a vital refuge from an unjust reality—ultimately suggesting that dreams alone cannot deliver change, and that true escape demands the courage to speak and act.
\end{quote}

\subsubsection{Evidence + Analysis of First Idea}
In the dream sequences where Gladys imagines being received by the Queen, Harrison frames fantasy not as idle escapism but as a shield of dignity for a woman excluded from the ``decent'' society she longs to join. Yet Nan Dear’s blunt dismissal—``And what bloody good would that do? Daydreams!'' (p.~154)—exposes the limits of this refuge, highlighting how dreaming leaves the family’s material circumstances untouched.\\

\subsubsection{Hinge Sentence (transitioning to Second idea)}
However, while Gladys’s dreams preserve her sense of self-worth, Harrison also reveals that fantasy alone cannot deliver her family from injustice.

\subsubsection{ Analysis of Second Idea}
Gladys’s transformation culminates in her address to the Rodney Shire Council, where she demands ``proper schooling … The white people too — they need to be educated about us, and our ways'' (p.~197). Harrison complicates the notion of escape by showing that it is voice and action, not fantasy, that begin to change the family’s circumstances, suggesting that dreams are only powerful when they inspire agency.\\

By using a hinge sentence, you show your reader that you're not just stacking ideas but developing an argument, deepening your interpretation, and showing the author’s nuanced message.\\
\begin{longtable}{|p{9cm}|p{6.5cm}|}
\caption{Concluding Dependant Phrases Paired with Analytic Verbs} \label{tab:concludingphrases} \\
\hline
\textbf{Concluding Dependant Phrase} & \textbf{Analytic Verb} \\
\hline
\endfirsthead

\hline
\textbf{Concluding Dependant Phrase} & \textbf{Analytic Verb} \\
\hline
\endhead

\hline
\endfoot

Even if Sophocles recognises...         & questions... \\
Whilst Sophocles endorses...            & illustrates... \\
Whether or not Sophocles implies...     & reveals... \\
In order to explore ... Sophocles ...   & demonstrates... \\
\hline

\end{longtable}






\newpage

\subsection{Conclusion}
If you think of a conclusion as less of summarising everything, and more of combining all of these
conflicting ideas and messages into one overall conclusive point. We can do this by using the following structure.

\subsubsection{Structure}
The structure for a conclusion can follow the same structure as an introduction. However, instead of introducing the text, you are concluding your discussion.\\

Here is an alternative method:\\


Topic:
\begin{quote}
    `The dream sequences in \underline{Rainbow's End} are the characters' only escape from the depressing reality of life.' Do you agree?
\end{quote}


\begin{longtable}{|p{3cm}|p{5.25cm}|}
\caption{Conclusion Structure for Analytical Essay Writing} \label{tab:conclusionstructure} \\
\hline
\textbf{Purpose in Conclusion} & \textbf{Example Sentence} \\
\hline
\endfirsthead

\hline
\textbf{Purpose in Conclusion} & \textbf{Example Sentence} \\
\hline
\endhead

\hline
\endfoot

Overview of how the whole text demonstrates the key concepts in the essay topic
&
Fundamentally, \underline{Rainbow's End} explores the tension between dreams and reality, particularly how fantasy can offer dignity, but not deliverance, to those denied opportunity by an unjust society. \\
\hline

Summary of the complications of this topic
&
Throughout the play, Harrison complicates the notion of escape by showing that while Gladys’s daydreams preserve her sense of self-worth, they leave her family’s circumstances unchanged, challenging the romanticisation of hope without action. \\
\hline

Focus on the key authorial message in conclusive language
&
Therefore, Harrison critiques a society that forces its most marginalised to retreat into fantasy, asserting that genuine change lies not in dreaming of acceptance but in finding the voice to demand it, as Gladys does before the Shire Council. \\
\hline

\end{longtable}

\paragraph{The Initial Concluding Sentence}
Here is a list of concluding and verb phrases you can use at the beginning of your conclusion.

\begin{longtable}{|p{5cm}|p{5cm}|}
\caption{Concluding Sentence Starters and Verb Phrases} \label{tab:conclusionphrases} \\
\hline
\textbf{Concluding Phrase} & \textbf{Verb Phrase} \\
\hline
\endfirsthead

\hline
\textbf{Concluding Phrase} & \textbf{Verb Phrase} \\
\hline
\endhead

\hline
\endfoot

Essentially,             & explores \\
Fundamentally,           & is a text about... \\
At its core,             & critiques \\
At its heart,            & is a novel that... \\
Ultimately,              & examines \\
In the end,              & is a science fiction tale \\
In the final analysis,   & exposes \\
Everything considered,   & that... \\
\hline

\end{longtable}

\subsection{Sentence Starters}


ZZPDFEMBED SS.pdf ZZEND




\cleardoublepage

\clearpage
\thispagestyle{empty}
\hfill
\clearpage

\section{Proofreading}

\subsection{Take a Break Before Proofreading}
After finishing your essay, your mind is still fresh with ideas, but it may be too familiar with the text to catch all mistakes. Step away from the essay for at least 10–15 minutes. This will help you approach it with fresh eyes and a clearer mindset when you return to review it.

\subsection{Read Aloud}
Reading aloud forces you to slow down and focus on each word. It also helps you hear awkward phrasing or sentences that don’t flow well. Read your essay aloud, paying attention to how each sentence sounds. If something feels off or too long-winded, make adjustments.

\subsection{Check for Structure and Coherence}
A strong essay is logically structured with clear arguments and smooth transitions. Ensure each paragraph has a clear topic sentence. Verify your argument flows logically from one point to the next. Ensure your introduction sets up your essay well and your conclusion summarises without introducing new ideas.

\subsection{Look for Clarity and Precision}
Vague or unclear language can confuse your reader, even if the argument itself is solid. Check if any sentences can be more specific or concise. Remove unnecessary filler words like 'really,' 'just,' or 'quite.' Replace vague words like 'things' or 'stuff' with more precise terms.

\subsection{Verify Your Thesis and Topic Sentences}
Your thesis statement should be the central focus of your essay, and topic sentences should align with it. Double-check that:
\begin{itemize}
    \item Your thesis is clear and present in the introduction.
    \item Each body paragraph clearly supports your thesis with an argument or evidence.
    \item Topic sentences match the content of their respective paragraphs.
\end{itemize}

\subsection{Proofread for Grammar and Syntax}
Grammar mistakes, awkward phrasing, or incorrect sentence structure can weaken the professionalism of your essay. Read through your essay for common mistakes (e.g., subject-verb agreement, incorrect punctuation, run-on sentences). Use tools like Grammarly or even Microsoft Word’s spelling/grammar checker for a quick second look.

\subsection{Check for Word Choice and Tone}
The right word choice ensures your ideas are clear and impactful. Avoid overly complex words unless they are necessary for the context. Ensure you maintain an academic tone throughout the essay. Replace informal language (e.g., 'kids' or 'stuff') with more formal, specific terms.

\subsection{Look for Repetition}
Repeating the same word or idea can weaken your essay's overall impact. Identify words or phrases that appear too often. Use synonyms or restructure sentences to avoid redundancy. Consider varying sentence structures to keep the reader engaged.


\begin{figure} [h!]
    \centering
    \includegraphics[width=1\linewidth]{Screenshot 2025-07-24 105659.png}
    \caption{\textit{Insight's Proofreading Checklist}}
    \label{fig:enter-label}
\end{figure}

\subsection{Check Punctuation and Spacing}
Simple punctuation errors or inconsistent spacing can distract the reader. Verify commas, full stops, and apostrophes are used correctly. Ensure there is proper spacing after full stops and between paragraphs.

\subsection{Ask Someone Else to Read It}
Fresh eyes can often spot mistakes you may have missed. Ask a peer, teacher, or family member to read your essay. Be open to their feedback and use it to improve your writing.

\subsection{Ensure It Meets the Task Requirements}
Your essay should always answer the question and meet the criteria specified in the prompt. Revisit the question to ensure your essay directly addresses it. Check that all key parts of the task are covered and that your argument is relevant.




\cleardoublepage

\clearpage
\thispagestyle{empty}
\hfill
\clearpage

\part{Creating Texts}
\clearpage

\section{Audience}
The audience refers to anyone who is viewing, reading or listening to a text.\\

1. Begin to identify your audience by answering the following questions:

2. What are you writing about?

3. Who would be interested in this topic?

4. Is there more than one audience who might be interested in this?

5. What other texts exist that are like my idea? Who are their audience?

\begin{figure}
    \centering
    \includegraphics[width=1\linewidth]{aud.png}
    \caption{\textit{Audience Diagram}}
    \label{fig:enter-label}
\end{figure}

\newpage

\section{Purpose}
\subsection{Writing to Express}
Use your imagination and creativity to explore and convey ideas

Examples include novels, short stories, plays, poems and films.

Usually involves imagined characters, places, situations and events.

Usually contain tension and conflict as characters seek to achieve their coals.

\subsection{Writing to Explain}
Use your imagination and creativity to explore and convey ideas

Examples include novels, short stories, plays, poems and films.

Usually involves imagined characters, places, situations and events.

Usually contain tension and conflict as characters seek to achieve their coals.

\subsection{Writing to Reflect}
Use your imagination and creativity to explore and convey ideas

Examples include novels, short stories, plays, poems and films.

Usually involves imagined characters, places, situations and events.

Usually contain tension and conflict as characters seek to achieve their coals.

\subsection{Writing to Argue}
Use your imagination and creativity to explore and convey ideas

Examples include novels, short stories, plays, poems and films.

Usually involves imagined characters, places, situations and events.

Usually contain tension and conflict as characters seek to achieve their coals.
\begin{figure}
    \centering
    \includegraphics[width=1\linewidth]{aud2.png}
    \caption{\textit{Language}}
    \label{fig:enter-label}
\end{figure}

\newpage
\section{Personal Journeys}
\textbf{From the Study Design}\\

\textit{Explorations of 'life' or biographical explorations - telling our stories, telling others' stories, the problem of telling stories, appropriation of stories, who tells the stories and our history, missing stories, marginalised and elevated stories. Students could explore personal milestones, the effects of key events on their lives, or explore these ideas through the eyes of others.}\\

\textit{Students who have migrated can explore their stories of movement and disruption. They can explore the expectations of change and the language of a new place and culture.}\\

\begin{figure}
    \centering
    \includegraphics[width=0.5\linewidth]{Screenshot 2025-07-24 105411.png}
    \caption{\textit{Insight's Different Forms of Text}}
    \label{fig:enter-label}
\end{figure}

\subsection{Writing}
\begin{itemize}
    \item Writing about personal journeys explores both the physical world and the internal world of oneself and others.
    \item Reading and crafting texts about journeys provide opportunities to consider new places, perspectives, and self-discovery.
    \item Personal journeys involve various experiences, including challenges, growth, and self-discovery.
    \item Writing about personal journeys allows individuals to share unique experiences with others.\\
\end{itemize}


\textbf{REMINDER: You have access to the exam generator and the prompt generator to practice your Section B skills.}\\

\subsection{Ideas}
\textit{A personal journey can involve:}

\begin{itemize}
    \item Travel to a physical place.
    \item Psychological or emotional growth and change.
    \item A combination of both, where physical travel leads to new insights.
    \item Interior journeys may take a long time to occur due to significant events.
    \item The most important aspect of a personal journey is change – the individual or group is fundamentally altered by it.
\end{itemize}

\begin{figure}
    \centering
    \includegraphics[width=1\linewidth]{2024ExamSectionB.png}
    \caption{\textit{2024 VCAA English Exam - Section B for Personal Journeys}}
    \label{fig:enter-label}
\end{figure}

\section{Practice Stimuli}
\subsection{Stimuli List I}


ZZPDFEMBED PJ1.pdf ZZEND


\newpage




ZZPDFEMBED PJ2.pdf ZZEND



\cleardoublepage

\newpage
\clearpage
\thispagestyle{empty}
\hfill
\clearpage

\section{Sample Creating Texts Piece (19/20)}


ZZEXP Creating Texts: 'The Path Less Taken' (to reflect / express) ZZEXPEND


A high-scoring (19/20) Section B piece written to reflect and express, responding to the title below.
\begin{quote}
\textit{Title: 'The Path Less Taken'}
\end{quote}

I was at a family dinner over Easter earlier in the year, and one of my aunts asked me a seemingly-innocuous question that made me quite uncomfortable. Usually, I am an upbeat, bubbly person, and most would tell you that I have high confidence. Unfortunately, the opposite is true: every day I question the choices I have made and the impact they have had on my relationships.

When I entered high school, I was faced with a rather uncomfortable choice: I would have to pick a school sport to play during winter, or join a band. The extracurricular lady peered at me over her bifocals, frowning, eyebrows drawn together and raised like little anthills. The only band that had any space was the orchestra, and they needed a piccolo player.

Now, anyone who is acquainted with woodwind instruments would know that the piccolo is a tiny, shrieking version of the flute, and that it should be played only by an experienced flautist with excellent tone. But for a reason unknown both to my parents and myself, I, standing in my oversized blazer with my dress almost at my ankles, boldly offered my services to the school band. The lady in charge of extracurriculars frowned a little harder, and I could almost imagine the tiny ants running in and out of those hills. Then she sat back, sighed, and put my name down, one eyebrow still raised.

Obviously, I was terrible. However, I started practising more and more. I took a certain tormented delight in knowing that while other students were tucking into dinner, or doing a little homework, I was running up and down the same five notes with deliberate focus. As painful as my practice was for all within earshot, the improvement it allowed me in the orchestra gave me surging motivation.

Sometimes in life, you are captured by an idea which is inexplicable to those who haven't felt it. The sense of mastery I felt playing the piccolo made all other subjects at school pale in comparison. What was maths, with all its useless formulae repeated over and over, when music was a space with infinite room for improvement? Where every day's practice contributed to notable growth, however small? Music was not just a hobby for me, but a pursuit which gave me the deepest sense of meaning. Over time, I came to love the piccolo more than anything else.

Unfortunately, as the years went by, my parents, who had initially been in favour of my playing an instrument, slowly started to prioritise other things. I should have been studying more maths, more science, writing more essays. I knew that they expected me to study like my friends, and to aim for medicine so I could become a doctor. As much as I tried, my enthusiasm for those subjects never matched what I felt for music. I practised more, and my grades fell further. I tried to hide reports from my parents (with little success). Finally, when I entered year 12, I committed a cardinal sin: I listed a music degree as my first university preference.

The decision had a negative impact on my parents. Our relationship turned to more fighting, and then stony, awkward dinners. Parents often have high expectations for their children. Life is short. You only buy books for a new highschool year six times with your child, and they only get the ages of 18-22 to do undergrad with people their own age. For parents, not only does time with children fly by, but in many cases the having of children cut their own professional lives short. I remember my mother constantly reminding me that perhaps she would have been so much further ahead in her workplace if not for the family, as much as she loved us. I think for her in particular, the lives of my siblings and myself were vicarious extensions of her own, through which she could live out different "maybe" paths.

You can understand, then, why I would have been so perturbed by my aunt's question: "And what will you be doing at the end of the year?"

For my part, I could never see why my parents were so upset. My three older siblings all left secondary education and diligently commenced medical school. I, on the other hand, could only see one thing: I had one life to live, and I'd rather do anything than spend it dissecting bodies and examining skin conditions. Once you realise that life is short, you come to terms with the fact that you may as well spend it doing things you enjoy. How sad it would be to die doing something you hated out of a feeling of social obligation.

And yet, this is what so many of us do. Many of us end up trapped in degrees, then internships, then grad roles, then middle-management jobs chosen for us by our high-schooler selves, with extremely little experience to guide the decision. Most people enter this life because it is the socially acceptable thing to do, and also because it is extremely difficult to find direction. Artists, inventors, and scientists have to bet against conventional wisdom to pursue their innate dreams. However, conventional wisdom is usually right (which is why we call it wisdom), and so if you don't have a good reason to ignore it, you probably shouldn't. Perhaps, given how much I know my parents care about me, I do understand their worry after all.

So when my aunt asked me this question, the mood in the room dropped somewhat. I thought carefully about my response. Should I lie, tell her that I'll be doing medicine too? Should I say I was prevented by my poor grades? It was almost true. But then I felt the weight of the truth settle upon me, and it was heavy and liberating all at once. I was different. Who cared? This was my life. "So, then? Will you be following your good older siblings?"

"You know, aunty," I twirled some spaghetti around my fork. "I've decided to walk the path less taken."

\newpage

\section{Sample Creating Texts Piece: 'Finding My Way' (10/10)}


ZZEXP Creating Texts: 'Finding My Way' (full marks) ZZEXPEND


A full-mark Personal Journeys piece responding to the title 'Finding My Way'.
\begin{quote}
\textit{Stimuli used: a pair of feet on asphalt with painted arrows pointing in many directions; 'The caterpillar transforms without thought. Its journey to a beautiful butterfly, inevitable. Why must my own metamorphosis be shaped by so many choices?'}
\end{quote}

I remember the day that I got off the plane so vividly. I remember the sun that day, a cruel god beating down upon me from its court of cloudless blue. It scorched my skin and left no prisoners. It was nothing like the benevolent sun over my homeland of Calcutta, where the gentle rays of morning would shine over the backs of a hundred thousand people, wrapped up in their idiosyncratic lifestyles since twilight.

But there was no one here. There was just me, my father, and the red earth that stretched out before me in every direction, desolate and uninviting. My father, a self-titled opportunist, yearning to abandon the stagnancy that is the defining characteristic of the Indian middle class, saw Laxmi, the Hindu goddess of prosperity, hidden in every crack of the drought-stricken land.

I always thought my father was a little bit of an oddball. His brain was rattled around by the many, many wars it had been exposed to. In the same spirit, his nervous yet paradoxically poised energy set our house alive. It would wake up with the soft, repetitive chants of his Sanskrit prayers in the morning and fall asleep to his many lectures about the ideal governmental system come nightfall. It smelled like incense, asafoetida, and cumin, instead of the burnt smoke of the great Australian outback outside our porch door. It was the only piece of home in a strange new country. I cherished it with all my heart.

He was quick to find his way. In three years' time, he was listening to The Beatles whilst fortifying himself with double helpings of Milo and Bengali egg fry, perusing an assortment of Hindi newspapers he got internationally delivered to us. The blokes at his sheep shearing job liked to alternate between calling him African and Indigenous Australian, slurs for both inclusive. He paid no heed, jovially punctuating his latest tales of workplace discrimination with, "Their stones and sticks can only rattle my bones! In Calcutta, we used gulel when we really wanted to hurt someone." I didn't have the heart to quip about his idiomatic misquoting. His spirit was indomitable. Nothing could faze him.

I, on the other hand, could not extricate myself from the diminishing gazes of my peers. Every remark on the failure of my foreign tongue to mold itself to the valleys and troughs of a new language, every question underpinned by perfunctory distaste, every snigger swirled in my head each night, consuming the liminal space between waking and sleeping. The movies I watched, the songs I liked, the food I ate, all underwent a strict metamorphosis like the caterpillar confined within the fecund dark of a chrysalis. Slowly, one choice at a time, aligning more and more with the Australian lifestyle. Slowly, one choice at a time, departing from my roots.

That was, until the diagnosis. "Hepatitis B," the doctor called it, the only bloke in town willing to look at a 'coloured' patient. "From infected needles, or maybe passed down in the womb from his mother." My father, try as he might, could not pronounce the name of the illness eating his liver alive. He jokingly called it a gift from our motherland. Over the next few months, I watched him disintegrate into a shell of himself. He could barely open his eyes most days, and when he could, he would talk. He would rasp his smooth tenor into a small tape recorder into the early hours of the morning, flitting from English to Hindi to Bengali as the disease continued its unending assault into his being.

When they complained to me, he laid his papery, wizened hand, which once, in a bygone era, would emanate vitality so potent it was palpable in his air, upon my shoulder and looked at me in the eye with a desperate so intense I felt it in my very soul. "You must find your way, beta. You must tell my story before my time is up."

When he passed, I lost the only person rooting for me, both on this side of the Indian Ocean and the one where I came from. I was lost without him. I felt stagnant, stuck. Our house suddenly felt as desolate as the golden land on my very first day. Over time, I have found my way. I honour his memory by telling his extraordinary story to anyone that will listen. I continue his tradition by hanging his portrait above my mantel and pestering my eldest daughter to hang a garland of fresh marigolds underneath it every morning, just as my father did with his father's portrait, and how his father would have done with his father before him.

Because of him, I have spent a lifetime working as a social worker, aiding those fragmented because of migration, granting me the insurmountable pleasure of listening to a thousand stories. Many of them are like my father and I's, each marked with their own anecdotes and idiosyncrasies. People, struggling, trying, grappling with making sense of a new world. I see the glint of my father's smile in the face of a stranger coming to Australia to make a better life for themselves and their families. I hear his voice in their adversities. It feels like home.

\subsubsection{What made this a 10/10}
Every scene threads the stimuli through fresh sensory imagery, recurring motifs, and a quiet, honest voice. That layering is what lifts a Section B piece from competent to top band.

\newpage

\section{Sample Creating Texts Piece: 'This is Me' (10/10)}


ZZEXP Creating Texts: 'This is Me' (full marks) ZZEXPEND


A full-mark Personal Journeys piece responding to the title 'This is Me'.
\begin{quote}
\textit{Stimuli used: 'Life isn't about finding yourself. Life is about creating yourself.'; 'My story is ready. There's enough to share. As I have made plenty. Come, sit at my table. Listen and learn. And you will see me.'}
\end{quote}

"Good evening customers. Your next service to depart from Platform 2 will be the four-fifty-six Flinders Street running express from Newport to Footscray, Footscray to North Melbourne, and then stopping all stations to Flinders Street. Always remember to touch on your Myki at the start of your journey." The train announcer's voice is a worn chime, filling the stark station like a frayed tapestry in an abandoned household. The station smells like it always does, redolent with the acrid tang of piss, sweat, and smoke. There is something self-righteous about the smell of a train station; none are exempt from its slightly carcinogenic, slightly metallic embrace.

A low mechanical groan swells beneath his feet as the train approaches, its headlights cutting through the dimness of the already-darkening winter sky with a tired reluctance. The doors slide open with a hydraulic sigh. He steps inside. The 4:56 pm train from Newport Station on Wednesday, the 20th of August, the 2025th year of the Common Era is almost deserted.

He looks down at his hands, admiring the soft webbing connecting the valleys between his fingers. His calluses are peeling, revealing the angry red of new skin underneath. It distantly reminds him of the animal documentaries he used to watch as a kid. He recalls his grandfather teaching him about the rattlesnakes as they appeared on their family's old LCD; how their cyclical shedding of skin added a new segment to their rattles. Each interlocked and keratinised bit of memory ultimately conjoined to create their characteristic rattles, allowing them to alert nearby friends or warn foe when shook. He feels a bit like that rattlesnake now, shedding another layer of him, who he was (is?), as his sole ward against something bigger than him.

He straightens, instinctively brushing his hands against the worn fabric of his jeans, as though he can rid himself of that thin, flaking skin, of the evidence that he is still stuck in the process of becoming something else.

The carriage is quieter than he expects. Subdued, like a room where the conversation has only just ended and the last spoken consonant lingers. A man in a hi-vis vest dozes with his head tilted back, mouth slightly ajar and a tenuous string of drool tracing its way down the contours of his jaw. Two teenagers sit across from one another, sharing a set of earbuds, their heads bowed together in conspiratorial symmetry. The girl in green does not look up. He takes a seat by the window.

As the train lurches forward, the station recedes into abstraction. The steel beams dissolve into streaks of grey, then into nothing at all. He watches his reflection instead. It hovers faintly against the glass, superimposed upon the passing city: a face that feels both intimately his and faintly unfamiliar. The train gathers speed.

But at the city, the carriage doors part and the world rushes in. The platform is dense with motion. He takes note of the sudden layering of language, the multiplicity of accents folding over another like tenuous silk organza of his mother's sarees. Someone coughs sharply; someone else argues into a phone; a group of schoolboys with reckless affection. The air smells of coffee, of baked bread, and of the faint tang of iron from the rails.

He lingers. He feels it again, that distance. Not exclusion, really, but a kind of peripheral existence. As though he is still watching the city through the glass when in reality, he's out in the cold, the gale lashing unforgivingly like a whip against his face. He looks down at his hands, at the peeling skin, scrutinising the rawness beneath. Shedding, again. Becoming, again.

He steps off the train and follows the crowd instinctively, letting the current of its movement take him toward Queen Victoria Market. It is not the quiet barter of daytime commerce that he expected. Tonight, the market radiates luminous against the winter chill, illuminated by strings of fairy lights and the glow of fire pits. The scents are the first to hit him: the rich spices from the steaming laksa, sweet smoke from grilled takoyaki, and the inviting swirl of mulled wine and spiced Milo in ceramic mugs held tight between gloved hands.

Luxuriating in the sensory overload, he moves deeper into the market. He spots Hungarian chimney cakes next to Malaysian satay, Latin American empanadas beside Turkish gozleme; the diversity of flavour becoming a rare map of the world in itself, condensed into the congested borders of one crowded laneway. At a food stall selling steaming xiaolongbao, an older man smiles at him and nods, the gesture small but genuine. The delicate crows feet of his smile remind him of his grandfather, the passionate twinkle in his eye as he taught the young boy before him, so many moons ago.

Suddenly trapped in the fossilised resin of memory, he remembers his grandfather's parting words to him, "Life isn't just about finding yourself, beta. You must remember to listen to the stories around you." He nods back to the older man, a little bit awkwardly, but finally laden with the understanding that had eluded him prior.

He realises, watching the ebb and flow of the throng, locals and newcomers, laughter and music, old and young, all engaged in the shared act of eating, moving, and breathing as one under the winter light, that belonging is not a singular destination. It is a mosaic of interactions, layered and overlapping, like the rattles on the snake's tail. Distinct segments forming something whole, something enduring.

He leaves the market with nothing in his hands, but a warmth settled somewhere in his throat. A loosening, however slight, of the tight, invisible vice of unfamiliarity.

The next evening, he boards the 4:56 again. The carriage is fuller now, or perhaps he is simply noticing more. A conversation drifts across the aisle, punctuated by laughter. As the train comes to a grinding halt, someone shifts to make space without being asked. He catches his reflection in the window. It is still there, sullen and superimposed against the speeding city. However, the edges have softened. The distance, though not demolished, has narrowed into something traversable. This is me, he thinks.

"Now arriving at Flinders Street Station..." The train announcer's voice is a worn chime and the station smells like it always does, like piss and sweat and smoke, but this time, he does not flinch. He listens instead, as the city moves with him, no longer a current he moves against, but something he is, slowly, learning how to inhabit.

\subsubsection{What made this a 10/10}
Every scene threads the stimuli through fresh sensory imagery, recurring motifs, and a quiet, honest voice. That layering is what lifts a Section B piece from competent to top band.

\newpage

\part{Analysing Argument}

\clearpage
\thispagestyle{empty}
\hfill
\clearpage

\section{Essay Structure}

\subsection{Introduction}

\subsubsection{What you are trying to achieve}
The purpose of an introduction is to 'introduce' both the irrefutable facts about the text you will be analysing, and your interpretation of the text. The best way to do this is to connect the text's plain details with your early analysis of what the contention, tone and intended audience.

\subsubsection{Include}
\begin{itemize}
    \item Context and contention
    \item Details (writer, publication, date and text type)
    \item Major persuasive devices
    \item Tone (consistant or shifting)
    \item Intended audience and how you know this
\end{itemize}

\subsubsection{Revise}
A common mistake students make is that they turn their introductions into lists of everything they need to include. Remember to express yourself clearly, professionally and with style. This is the introduction to your analysiss and sets your assesor's expectations for the rest of the piece.

\subsubsection{Sentence structure}
The first part of your introduction should focus on what issue is being discussed and why it is important. This will help you to create an opening to demonstrate a strong understanding of the article, for example:
\begin{quote}
    The current trend of dangerous Tik-Tok challenges has once again sparked debate over whether we should be restricting access to the wildly popular app. 
\end{quote}

If we break it down it becomes:
\begin{itemize}
    \item The current trend - \textbf{Phrase of recency}
    \item of dangerous Tik-tok challenges - \textbf{Incident}
    \item has once again sparked - \textbf{Verb phrase}
    \item debate - \textbf{Issue}
\end{itemize}

Here are some sample sentence builders:

\begin{longtable}{|p{2.5cm}|p{3cm}|p{2.25cm}|}
\caption{Framing Contemporary Issues in Writing: Phrases of Recency, Cause, and Contention} \label{tab:recencyphrases} \\
\hline
\textbf{Phrases of Recency} & \textbf{Verb Phrases} & \textbf{Issues} \\
\hline
\endfirsthead

\hline
\textbf{Phrases of Recency} & \textbf{Verb Phrases} & \textbf{Issues} \\
\hline
\endhead

\hline
\endfoot

The latest...         & ...has once again raised... & issue \\
The present...        & ...has sparked...           & debate \\
The recent...         & ...has led to...            & argument \\
The earlier...        & ...has motivated...         & discussions \\
The current...        & ...has caused...            & disputes \\
Once again...         & ...has generated...         & altercation \\
Recently...           & ...has engendered...        & discourse \\
Lately...             & ...has provoked...          & consultation \\
Over the past year... & ...has propagated...        & dialogue \\
\hline

\end{longtable}

\subsubsection{Key Information Phrase}
In this section we should include key information that is usually found in the background information section of the text, including:
\begin{itemize}
    \item Author's name
    \item Connection of author to text
    \item The form of the piece, the date and where it was published
    \item How this links back to the issue you identified in your first sentence
\end{itemize}


For example, this sentence can look like:
\begin{quote}
    Dr Marko Shellius, CEO of Psychiatric Safe Australia, has written a submission to the parliamentary inquiry into social media, in particular addressing Tik-Tok’s burgeoning influence over Australian teens. 
\end{quote}

If we break it down, it looks like this:
\begin{longtable}{|p{2cm}|p{2.5cm}|p{2cm}|p{2.5cm}|}
\caption{Authorial Context and Publication Framing for Argument Analysis} \label{tab:authorcontext} \\
\hline
\textbf{Author} & \textbf{Connection of Author to the Issue} & \textbf{Form and Publisher} & \textbf{Link Back to the Issue} \\
\hline
\endfirsthead

\hline
\textbf{Author} & \textbf{Connection of Author to the Issue} & \textbf{Form and Publisher} & \textbf{Link Back to the Issue} \\
\hline
\endhead

\hline
\endfoot

Dr Marko Shellius & CEO of Psychiatric Safe Australia & Has written a submission to the parliamentary inquiry on social media & In particular addressing TikTok’s burgeoning influence over Australian teens \\
\hline

\end{longtable}


This order can also be moved around as well:

\begin{longtable}{|p{2.75cm}|p{2cm}|p{2.75cm}|p{2.75cm}|}
\caption{Form, Author, and Contextual Framing of the Issue} \label{tab:formauthorcontext} \\
\hline
\textbf{Form and Publisher} & \textbf{Author} & \textbf{Connection of Author to the Issue} & \textbf{Link Back to the Issue} \\
\hline
\endfirsthead

\hline
\textbf{Form and Publisher} & \textbf{Author} & \textbf{Connection of Author to the Issue} & \textbf{Link Back to the Issue} \\
\hline
\endhead

\hline
\endfoot

In a written submission to the parliamentary inquiry into social media, & Dr Marko Shellius & CEO of Psychiatric Safe Australia & Addresses TikTok’s burgeoning influence over Australian teens \\
\hline

\end{longtable}

Here are some sentence starters to help build these sentences:
\begin{longtable}{|p{2.5cm}|p{3cm}|p{2.5cm}|}
\caption{Author Connection, Form of Publication, and Issue Framing Phrases} \label{tab:authorformissue} \\
\hline
\textbf{Connection of Author to Issue} & \textbf{Form and Location} & \textbf{Link to Issue} \\
\hline
\endfirsthead

\hline
\textbf{Connection of Author to Issue} & \textbf{Form and Location} & \textbf{Link to Issue} \\
\hline
\endhead

\hline
\endfoot

Chief … of…                & has written in/on…            & in response to… \\
Leader of…                 & has spoken at/in…             & addressing… \\
Instigator…                & a blog published…             & about… \\
a writer…                  & an opinion piece published…   & concerning… \\
an opponent of…            & a… newsletter                 & regarding… \\
a supporter of…            & a speech presented to…        & encircling… \\
a member of…               & a debate in front of…         & surrounding… \\
a key…                     & published a…                  & acknowledging… \\
an important…              & created a…                    & responding to… \\
one of the…                &                               & confessing… \\
\hline

\end{longtable}


\subsubsection{Tones}
Analysing argument involves understanding how authors use language to persuade their audience. This includes identifying the author's tone, which refers to the overall feeling or attitude conveyed in the text. Recognising shifts in tone is crucial as it often accompanies changes in argument strategies. Authors use tone to appeal to the reader's emotions, establish credibility, and guide the reader's interpretation of the argument.\\

This is the perfect place for you to show the breadth of your vocabulary. Choose a specific word that relates to the article. Always avoid using the terms 'Positive', 'Negative' and 'Neutral'.''

\begin{longtable}{|p{2cm}|p{2cm}|p{2cm}|p{2cm}|p{2cm}|}
\caption{Positive Tone Vocabulary Bank} \label{tab:positivetonewords} \\
\hline
\hline
\endfirsthead

\hline
\hline
\endhead

\hline
\endfoot

Amazed & Calm & Appreciative & Affectionate & Amiable \\
Astonished & Casual & Approving & Benevolent & Amused \\
Attentive & Collected & Assuring & Compassionate & Cheerful \\
Curious & Composed & Confident & Concerned & Delighted \\
Eager & Content & Determined & Considerate & Ecstatic \\
Interested & Peaceful & Encouraging & Consoling & Elated \\
Keen & Pleasant & Grateful & Empathetic & Energetic \\
Polite & Relaxed & Hopeful & Emphatic & Enthusiastic \\
Startled & Relieved & Inspiring & Friendly & Excited \\
Stunned & Serene & Optimistic & Loving & Exuberant \\
Surprised &  & Pleased & Merciful & Happy \\
 &  & Promising & Romantic & Humourous \\
 &  & Proud & Soothing & Jovial \\
 &  & Respectful & Supportive & Joyful \\
 &  & Reverent & Sympathetic & Jubilant \\
 &  & Sanguine & Thoughtful & Playful \\
 &  & Satisfied &  & Vibrant \\
 &  & Thankful &  & Vivacious \\
\hline

\end{longtable}

\begin{longtable}{|p{2.5cm}|p{2.5cm}|p{2.5cm}|}
\caption{Neutral Tone Vocabulary Bank} \label{tab:neutraltonewords} \\
\hline
\hline
\endfirsthead

\hline
\hline
\endhead

\hline
\endfoot

Admonitory & Factual & Patriotic \\
Allusive & Formal & Personal \\
Authoritative & Frank & Picturesque \\
Balance & Honest & Questioning \\
Blunt & Informal & Reflective \\
Candid & Informative & Reminiscent \\
Colloquial & Knowledgeable & Resigned \\
Contemplative & Learned & Scholarly \\
Controlled & Lyrical & Serious \\
Conversational & Naïve & Sublime \\
Discursive & Noble & Virile \\
Distinct & Nostalgic & \\
Evocative &  & \\
Expectant &  & \\
\hline

\end{longtable}

\newpage

\begin{longtable}{|p{3.5cm}|p{3.5cm}|p{3.5cm}|p{3.5cm}|}
\caption{Negative Tone Vocabulary (increasing intensity from left to right)} \label{tab:negativetones} \\
\hline
\multicolumn{4}{|c|}{\textbf{Negative Tones (increasing intensity from left to right →)}} \\
\hline
\textbf{Negative and passive} & \textbf{Negative thoughts} & \textbf{Negative and uncontrolled} & \textbf{Negative and forceful} \\
\hline
\endfirsthead

\hline
\textbf{Negative and passive} & \textbf{Negative thoughts} & \textbf{Negative and uncontrolled} & \textbf{Negative and forceful} \\
\hline
\endhead

\hline
\endfoot

Apathetic & Apologetic & Agitated & Accusing \\
Bored & Critical & Alarmed & Aggravated \\
Cold & Doubtful & Anxious & Angry \\
Dejected & Envious & Apprehensive & Annoyed \\
Depressed & Foreboding & Defenceless & Belligerent \\
Despaired & Frustrated & Distressed & Calculating \\
Disappointed & Gloomy & Disturbed & Condemnatory \\
Discontented & Guilty & Embarrassed & Condescending \\
Disinterested & Judgmental & Fearful & Contempt \\
Dispirited & Pessimistic & Helpless & Disgusted \\
Gloomy & Regretful & Humiliated & Factious \\
Hopeless & Remorseful & Mortified & Furious \\
Hurt & Shameful & Nervous & Harsh \\
Melancholy & Solemn & Powerless & Hateful \\
Miserable & Somber & Shocked & Insulting \\
Regretful & Suspicious & Stressed & Irritated \\
Sad &  & Tensed & Manipulative \\
Upset &  & Troubled & Outraged \\
 &  & Uneasy & Quarrelsome \\
 &  & Vulnerable & Sarcastic \\
 &  & Worried & Sardonic \\
 &  &  & Vexed \\
\hline

\end{longtable}


\subsubsection{Audience}
The audience refers to anyone who is viewing, reading or listening to a text.\\

When writing an analysing argument essay, it isn’t enough to just say your audience is ‘everyone’ or ‘all Australians’. You need to refine your idea of who you are writing for!\\

Three questions to ask yourself
\begin{itemize}
    \item What is the author writing about?
    \item Who would be interested in this topic?
    \item Is there more than one audience who might be interested in this?
\end{itemize}

\begin{figure} [h!]
    \centering
    \includegraphics[width=0.75\linewidth]{image.png}
    \caption{Refining Audience}
    \label{fig:enter-label}
\end{figure}

\textit{The key take-away is that you need to be as specific as possible when stating the intended audience.}


\newpage
\subsection{Body Paragraphs}
\subsubsection{What you are trying to achieve}
Each body paragraph should explore, in sequence, the major arguments that the author uses to support their contention. Remember you are being assessed on your ability to analyse the arguments the author uses and how the language they have chosen supports these.

\subsubsection{Step 1: What?}
Outline what the author is arguing in your own words. Make a comment on structure; if this is the first or final argument then what is the reason the author introduced this at this time.
\subsubsection{Step 2: How?}
Begin by identifying the persuasive language features that the author uses to substantiate their argument, include a quote illustrating your point (but no more than 5 words). But this is not enough, then you need to analyse the effectiveness of this approach and discuss how it positions the audience. If there is an image associated with this argument, you should talk about it here - treat it like a persuasive device. Depending on the argument, you may want to talk about how a group of persuasive devices work together.\\

\textit{Note: Try to avoid repetition by constantly stating 'The author uses X which makes the audience feel Y.} This can come across tacky and quite robotic. To see examples of how this can be done in a successful way, look at the three high-scoring essays below.\\

\subsubsection{Step 3: Why?}
Finally, you need to infer why the author chose this approach, what is their persuasive intent? How would it appeal to their particular audience? How does it impact the tone of the article? How does this support their contention?
\subsubsection{Step 4: Examine (Optional)}
Does the tone change for this argument? Does the structure change? How does this position the audience?

\newpage

\subsection{Conclusion}
\subsubsection{What you are trying to achieve}
In your conclusion, you are exploring both how the text concludes and wrapping up your analysis at the same time. It is an opportunity to discuss the text as a whole, and emphasises your analytical skills.

\subsubsection{Include}
\begin{itemize}
    \item Comment on the structure of the article, including whether there is a consistent or changing tone of writing.
    \item Focus on the last paragraph of the article, evaluating how it supports the writer's contention.
    \item Evaluate whether the article developed more trust in the readers and explain why.
\end{itemize}

\subsubsection{Note}
Some students find they don't have time for a conclusion. This isn't ideal, but it isn't the end of the world, as most of this information should have been explicit throughout. If you are running short on time, a sentence summarising how each argument leads to the  contention is sufficient.
\newpage


\subsection{Complex Sentence and Paragraph Building}

When analysing persuasive texts, it’s not enough to simply spot techniques—you need to explain what is happening, why it’s there, and how it affects the audience.\\

Students often start with short, seperate sentences:\\
\textit{'The writer uses emotive language. This makes the audience feel sympathy'}\\

A stronger version combines ideas into one complex sentence:\\
\textit{‘Describing the victims as “cruelly abandoned” evokes a deep sense of sympathy in the audience, prompting them to see the government as neglectful and morally failing in its duty of care.’}\\

\subsubsection{Techniques for Stronger Sentences}
\begin{itemize}
\item Subordinating conjunctions – because, although, while, since, if
\item Embedded quotes – weave quotes into the sentence naturally
\item Relative clauses – who, which, that
\item Appositives – a technique aimed at…
\end{itemize}

By applying these sentence-building tools, your analysis will become more precise, connected, and persuasive—mirroring the sophistication expected in high-level responses.\\

To help you identify and discuss these elements in a clear, structured way, we will now look at a step-by-step framework that breaks persuasive writing into its key components.\\


\subsection{W.A.S.T.E.: F.A.T.!}
This acronym is designed to assist you in crafting complex and sophisticated sentences that strengthen your analytical writing. In higher-scoring essays, writers deliberately weave these elements into their work to add depth, clarity, and precision. While it is not necessary to include every element in every individual sentence, you should aim to integrate all of them across the body of your essay. This means that, by the end of each body paragraph, the reader should be able to identify how you have employed these elements to develop your argument, enhance your expression, and engage the audience effectively. Consistent and purposeful use of these techniques will not only improve sentence complexity but also demonstrate a mature and well-structured writing style.\\

\begin{itemize}
    \item \textbf{W}ords – Focus on the most persuasive words or phrases of the piece.
    \item \textbf{A}rgument – What is the supporting argument of each section?
    \item \textbf{S}tructure – Where is this located in the piece and why?
    \item \textbf{T}one – How does the author apply tone to be more persuasive in each section?
    \item \textbf{E}ffect upon the audience – How does it position the audience to:
    \begin{itemize}
        \item \textbf{F}eel?
        \item \textbf{A}ct?
        \item \textbf{T}hink?
    \end{itemize}
\end{itemize}

\subsubsection{W.A.S.T.E.: F.A.T. Medium Example}


ZZEXP W.A.S.T.E.: F.A.T. Medium Scoring Example ZZEXPEND


\textit{The editorial uses emotionally charged words such as “climate catastrophe”, “tipping point”, and “irreversible damage” to stress the seriousness of the issue. The supporting argument is that urgent government action is necessary to prevent catastrophic environmental collapse. These phrases appear in the introduction so that the reader is immediately confronted with the scale of the problem. The author adopts an alarmist and urgent tone to heighten concern and prompt immediate attention. This positions the audience to feel fear and anxiety about the future, to act by pressuring politicians or signing petitions, and to think that if society does not act now, it will be too late.}

\subsubsection{W.A.S.T.E.: F.A.T. High Example}


ZZEXP W.A.S.T.E.: F.A.T. High Scoring Example ZZEXPEND


\textit{The editorial uses emotive phrases such as “climate catastrophe”, “tipping point”, and “irreversible damage” in its introduction to immediately confront the reader with the scale of the problem. By presenting the argument that urgent government action is needed to prevent environmental collapse, the writer establishes the issue as both pressing and unavoidable. Through an alarmist and urgent tone, the piece positions the audience to feel fear, act by pressuring politicians or signing petitions, and think that immediate intervention is essential.}

\subsection{Fill-in-the-blanks Analysing Argument Template for an Introduction}

Your introduction should discuss the writer’s contention and intention, the specific audience, the
form, and the publication details, as well as the context of the issue. You might like to make a
comment about the tone of the piece in the introduction.\\

Here is a sample introduction structure:\\
The issue of ….. (issue) was addressed in a ….. (text type) for ….. (publication) on ….. (date). …..
(Author) contends, in a ….. (tonal word) and at times ….. (tonal word) voice that ….. (contention).
The piece is predominantly ….. (words about overall structure or style) in style, and seeks to
position its ….. (audience) to accept ….. (intention).\\

Another sample introduction structure could look like this:\\
Following ….. (context), debate resurfaced regarding ….. (issue). In a ….. (text type) for …..
(publication) on ….. (date), ….. (author) argues in a ….. (tonal word) and ….. (tonal word) fashion
that ….. (contention), and seeks to ….. (intention). The piece, targeted at ….. (audience), ….. .\\

For example:\\
In the wake of the recent sentencing of murderer John Smith, Sarah Farah’s letter to the editor of
The Daily Telegraph (9 June 2014), raises the issue of perceived injustice in the treatment of
Australian citizens. In an aggressive tone, she contends that the penalty imposed on this man was
too harsh and that such punishment represents a form of anti Australian discrimination. Thus, she
hopes to prompt the Australian Government to take action to protect its citizens.\\



\section{Analysing Written Language}
\subsection{Terms and Strategies}

\subsubsection{Specific Word Choice (word connotation) }
Writers choose particular words to provoke an emotional reaction in the reader or emphasise an idea.\\

For example, if writing about abortion, the author may choose;
"I read with disgust that the government is to ...
kill, slay, slaughter, murder, butcher, exterminate, destroy, sacrifice, execute ...
Emotive language is language that is not logical or rational - it reflects and expresses an emotional or subjective reaction. Such words help create mood. Many words have connotations or associations that the writer uses with deliberate intent to create a desired effect in the reader. Words may have positive or negative connotations. They may evoke intense imagery. You will be expected to discuss some targeted words and their relative connotations in your language analysis piece.\\


\subsubsection{Particular language structure and features}
Inclusive language (we, our) may be used to give the audience a sense of being part of a wider
community of thought/action.\\

Modifiers and qualifiers are words used to denote a level of impact. There is a difference
between a writer choosing to use 'some' instead of 'all' or 'most' instead of 'many'. Phrases
such as 'very good' or 'basically good' also suggest different meaning to the reader. Students
analysing language need to consider why specific lexical strings have been chosen by the writer
as a form of persuasion.\\

\subsubsection{Euphemisms}
Euphemisms are words sometimes used to soften a more accurate but seemingly more offensive
word or phrase. For example, euphemisms for 'die' include; passed away, lose one's life, kick
the bucket, croak. A company may say they are 'downsizing' instead offiring workers from the
job.

\subsubsection{Hypothetical Scenarios}
Hypothetical scenarios are used to prompt the reader to imagine possible outcomes, often beginning with phrases like \textit{“Imagine if...”}. These speculative examples draw attention to the consequences of actions or inactions, even if they are not real or likely. Writers use this technique to provoke reflection and emotional investment from the reader.\\

For example: \textit{“Imagine if the world was ruled by tyrants who dominated our every thought daily...”}\\

\subsubsection{Repetition}
Repetition involves deliberately repeating words or phrases for emphasis and to make an idea more memorable.\\

\textbf{Anaphora:} repetition at the start of successive clauses.\\
\textit{e.g. School is distracting. School is boring. School sucks.}\\
\textbf{Epistrophe:} repetition at the end of clauses.\\
\textit{e.g. Follow the rules. Obey the rules. These are the rules.}\\

\subsubsection{Puns}
Puns are humorous wordplays based on similar sounds or multiple meanings. They are often used to lighten tone or ridicule opposing views. Some puns also carry double meanings, especially in persuasive or satirical writing.\\

\textit{e.g. “Time flies like an arrow; fruit flies like a banana.”}\\

\subsubsection{Irony}
Irony occurs when the intended meaning of a statement contrasts with its literal meaning. Writers often use irony to highlight absurdities or criticise ideas indirectly.\\

\textit{e.g. “Lovely day,” said during a storm.}\\

\subsubsection{Sarcasm}
Sarcasm is a biting, mocking form of expression usually intended to hurt, ridicule, or convey contempt. It often overlaps with irony but has a more cutting tone.\\

\textit{e.g. “Nice face. Shame about the rest.”}\\

\subsubsection{Exaggeration / Hyperbole}
This technique involves overstating a point for dramatic effect. It can be used to emphasise importance or provoke an emotional response.\\

\textit{e.g. “I’d rather cut off my arm than go back to work under those conditions!”}\\

\subsubsection{Humour}
Humour engages the reader by entertaining them. It can also make controversial ideas more palatable or draw attention to flaws in opposing arguments.\\

It is often used as a non-threatening way to critique.\\

\subsubsection{Satire}
Satire combines humour and criticism to ridicule individuals, institutions, or ideas. It often exposes hypocrisy or absurdity, making serious points in a humorous way.

\textit{e.g. A TV show like \textit{Frontline} that uses mock-journalism to critique the media.}\\

\subsubsection{Oxymoron}
An oxymoron combines two contradictory terms to highlight contrast or provoke thought. It is often used to emphasise complexity or irony.\\

\textit{e.g. deafening silence, wise fool, bitter sweet}\\

\subsubsection{Assonance}
Assonance is the repetition of vowel sounds within words or phrases to create rhythm, mood, or emphasis. Often used poetically or in slogans.\\

\textit{e.g. “cool, calm and collected”}\\

\subsubsection{Simile}
A simile compares one thing with another using the words \textit{‘like’} or \textit{‘as’}. It helps readers understand unfamiliar concepts by likening them to familiar ones, often to create vivid imagery.\\

\textit{e.g. My bedroom is like a pigsty. She is as skinny as a rake.}\\

\subsubsection{Metaphor}
A metaphor is an implied comparison that states one thing is another, thereby transferring qualities from one to the other. Metaphors can make writing more persuasive by enhancing imagery and emotional appeal.\\

\textit{e.g. I stepped into my pigsty, shut the door and climbed into bed. He is a pillar of strength.}\\

\subsubsection{Rhetorical Questions}
These are questions posed for effect rather than to elicit an answer. They imply that the answer is self-evident, guiding the audience to agree with the writer’s point.\\

\textit{e.g. Do you want to see your wife and children out on the streets in rags? Of course you don’t!}\\

\subsubsection{Scapegoats}
A scapegoat is a person or group blamed for something they are not truly responsible for. This technique simplifies complex issues and focuses the reader’s anger on a convenient target.\\

\textit{e.g. All our crime and unemployment problems are because of the government.}\\

\subsubsection{Assumptions}
Assumptions involve treating beliefs or opinions as facts without proof. These are often subtle and can influence readers by appealing to what is generally believed to be true.\\

\textit{e.g. He must be smart because he is on the computer every day.}\\

\subsubsection{Presumptions}
Presumptions occur when statements are accepted as truths without critical examination. They appeal to common beliefs or cultural norms.\\

\textit{e.g. If it’s natural, it’s good for you.}\\

\subsubsection{Analogy}
An analogy compares two things that are alike in some respects to clarify a point or persuade readers that what applies to one must apply to the other.\\

\textit{e.g. If birds can fly because they have wings, then maybe if we develop mechanical wings we can fly too.}\\
\textit{e.g. The human brain is like a computer.}\\

\subsubsection{Expert Opinion}
This technique involves citing someone perceived to have authority or expertise on a topic. It lends credibility to an argument and can be persuasive even if the “expert” is unnamed.\\

\textit{e.g. In the thirty-five years I was Head Surgeon at...}\\

\subsubsection{Case Studies and Anecdotes}
Case studies and anecdotes are used to personalise or humanise an argument. Anecdotes are usually emotive and personal; case studies more formal and factual. Both aim to build credibility and emotional appeal.\\

\textit{e.g. Pool fencing should be mandatory because my neighbour’s child drowned in an unfenced pool.}\\

\subsubsection{Empty Truths}
Also known as tautologies, these are statements that repeat the same idea unnecessarily, often sounding meaningful but lacking substance.\\

\textit{e.g. Boys will be boys. Enough is enough. Business is business.}\\

\subsubsection{Cause and Effect}
This technique suggests that one event directly causes another. It is persuasive because it offers logical explanation, though causality may not always be valid.\\

\textit{e.g. Peter crashed his car. Peter died.}\\

\subsubsection{Ends and Means / Means to an End}
This approach justifies an action based on the outcome, even if the action itself is questionable.\\
\textit{e.g. A criminal donates to a playground; the donation is good, but the source of funds is illegal.}\\

\subsubsection{Consequence of a Fact}
This presents a pragmatic response to an existing situation by arguing that something should be done simply because a condition exists.\\

\textit{e.g. Nuclear power exists, so we might as well use it.}\\

\subsubsection{Judging an Action by its Consequences}
This technique justifies or condemns an action based on its results. It focuses on outcomes rather than intent.\\

\textit{e.g. Stealing from the rich to give to the poor is acceptable because the poor are hungry.}\\

\subsubsection{Sweeping Statements / Generalisations}
These make broad claims that apply to a group without considering individual differences. Generalisations can sound convincing but often oversimplify.\\

\textit{e.g. All teenagers are irresponsible.}\\

\subsubsection{Cliché}
Clichés are overused phrases that offer familiar shortcuts to meaning. Though persuasive due to their recognisability, they can also weaken writing if not used purposefully.\\

\textit{e.g. What goes around comes around.}\\

\subsubsection{Co-existence}
This assumes an automatic relationship between ideas or traits, such as assuming someone is good simply because they attend church.\\

\textit{e.g. She goes to the local church, so she must be a good person.}\\

\subsubsection{Guilt by Association}
This technique links people, ideas or groups with undesirable individuals or arguments in order to discredit them.\\

\textit{e.g. Premier Kennett has behaved like Hitler since coming to power.}\\

\subsubsection{Questioning Motive}
This involves undermining an argument by questioning the motives of the speaker, diverting attention away from the argument itself.\\

\textit{e.g. The Health Minister’s decision to close the hospital is just an excuse to build a monument to himself.}\\

\subsubsection{Figures and Statistics}
Numbers and data are used to give the illusion of objectivity and reliability. However, the reader must critically assess source, sample size, relevance, and interpretation.\\

\textit{e.g. 300 people died on our roads last year...}\\

\subsubsection{Discredit Another Argument}
This technique aims to dismantle opposing viewpoints by exposing flaws or inconsistencies in them. It draws attention away from the original argument and weakens the opposing side’s credibility.\\

\subsubsection{Attack}
This technique involves criticising or mocking opposing arguments or individuals. It is designed to weaken the opposition's credibility and make their position less persuasive. Writers may use name-calling, ridicule, or other personal attacks to influence the audience’s perception.\\
One form of attack is labelling the opposition in an attempt to belittle them.

\subsubsection{Flattery and Praise}
This technique involves listing individuals or groups who agree with the writer’s position in a positive light. By using terms that convey admiration or respect, the writer bolsters the authority of their own argument.\\

\textit{e.g. “Leading scientists support this research.”}

\subsubsection{Use of Punctuation and Printing Style}
Writers use punctuation and formatting deliberately to add emphasis, create subtext, or manipulate tone and flow.\\
Common strategies include:
\begin{itemize}
  \item Inverted commas – suggest doubt or irony, even when quoting.
  \item Exclamation marks – show emotional emphasis or urgency.
  \item Dashes – indicate pauses, hesitation or parenthetical ideas.
  \item Long sentences – may signal sophistication or rambling.
  \item Short sentences – create punch, pace and immediacy.
  \item Italics, CAPITALISATION, \textbf{bold}, \underline{underlining} – all used for emphasis.
\end{itemize}

\subsubsection{Vocabulary Use}
The choice of vocabulary significantly affects tone and credibility. Writers use sophisticated or precise word choice to sound more authoritative or to target specific reader emotions. Students should be aware of:
\begin{itemize}
  \item Emotive connotations (e.g. \textit{slaughter} vs \textit{kill})
  \item Word strength and subtlety (e.g. \textit{sporadic} vs \textit{occasional})
  \item Replacing vague or weak terms with more accurate or academic language.
\end{itemize}
Improving vocabulary can enhance clarity and persuasiveness. For example:
\begin{itemize}
  \item Replace “straight away” with “immediately”
  \item Replace “in the next paragraph” with “in the following section”
  \item Replace “always” with “consistently”
\end{itemize}

\newpage
\subsection{Table of Common Persuasive Techniques}

\begin{center}
\textbf{\LARGE Persuasive Techniques Reference Table}
\end{center}

\begin{longtable}{|p{2cm}|p{3.25cm}|p{3.25cm}|}
\hline
\textbf{Persuasive Technique} & \textbf{Definition} & \textbf{How It Persuades} \\
\hline
\endfirsthead

\hline
\textbf{Persuasive Technique} & \textbf{Definition} & \textbf{How It Persuades} \\
\hline
\endhead

Adjective & Describing words usually connected to a noun (naming word). \newline e.g. \textit{The bright red house} (bright and red = adjectives; house = noun) & Adds emphasis. \\
\hline
Alliteration & Repeating and playing upon the same letter. \newline e.g. \textit{A back-breaking job...} & Adds emphasis and reinforces meaning, especially where an idea is repeated. Often used to create a highly emotive image. \\
\hline
Anecdote & A short account or story of an entertaining or interesting incident. \newline e.g. \textit{In my experience...} & Usually makes reader sympathetic and receptive to the point; can set up a character, then position reader to accept or reject that individual character. \\
\hline
Associations & Words that create meanings beyond the literal meaning of the word. \newline e.g. \textit{attractive, alluring, pleasing, charming} & Suggests or influences reader response in certain ways, often subconsciously. Implied meanings need to be identified. \\
\hline
Attacks/praise & The writer attacks or praises an opponent or idea. This is really an assertion of the writer's viewpoint. \newline e.g. \textit{denigrating or humiliating the person, discrediting opponents as unreliable, dishonest, suggesting unsatisfactory associations} & Often effective in placing the reader in a position of agreement with the writer/speaker through forceful language and insistent claims, even without evidence. \\
\hline
Bias & One-sidedness in presentation of view/opinion. \newline e.g. \textit{Eminem’s tour promoter obviously sees no danger in what Eminem says.} & Can subjectively influence the reader by intentionally only presenting one side of the argument. \\
\hline
Clichés & Worn out, overused expressions. \newline e.g. \textit{Fit as a fiddle, Turn over a new leaf, Pushing the envelope} & Are familiar, often colloquial, so can offer a shortcut to convey meaning. \\
\hline
Colourful words and descriptive language & Words that are heightened, vivid, lively, and full of interest. \newline e.g. \textit{Hot = blistering, sultry, muggy, suffocating...} & Produce a picture and/or induce an emotion. Engage reader by gaining attention and often put a new slant on familiar events and issues. \\
\hline
Contention & The main line of argument in a text. A statement of the main idea being argued or debated. Focuses on WHAT is being argued. & Pinpoints the issue and the writer’s point of view. Helps the reader identify the issue and evaluate the techniques used to persuade them to agree. \\
\hline
Emotional Appeals & Play on people's emotions such as fears, insecurities, hopes, desires, and values. \newline Can target: \newline • A sense of justice or injustice \newline • Family values \newline • Tradition or custom \newline • Patriotism \newline • Moral values & Frequently appeal to our values and attitudes, sometimes subtly. Can manipulate reader to take notice of issues by triggering emotional responses. \\
\hline
Emotive Language & The deliberate use of strong emotive words to play on reader’s feelings. \newline e.g. \textit{sleazy, slimy, vicious, disgusting, outrageous} & Evokes strong emotional response to coerce or force agreement. Used to stimulate strong emotional reactions and manipulate responses. \\
\hline
Loaded Words (Labelling) & Words that are loaded with associations. \newline e.g. \textit{capitalist, communist, salesman} & Shorthand way of discrediting or praising by suggesting associations. Exploits emotional baggage of such words to sway opinion. \\
\hline
Repetition & Repeated words, phrases, sentence patterns, ideas. & Gives emphasis and prominence to a point or idea; reinforces and makes reader remember the point. \\
\hline
Rhetorical Questions & Questions that have the answer embedded in them; often use irony. & Manipulate the reader to agree because the answer is assumed. Disagreement would mean dismissing a point that commands agreement. \\
\hline
Tone & Refers to the voice of the writer; the writer's attitude to the subject and reader. Tone can be emotive. \newline e.g. \textit{angry, sympathetic, sarcastic} & Reflects writer's attitude, which can influence the reader to agree or reject. Tone can signal a new direction or perspective in the writing. \\
\hline

\end{longtable}
\newpage

\subsection{Table of Verbs for Authorial Intent}

\begin{center}
\textbf{\LARGE Verbs Indicating Authorial Intent}
\end{center}

\begin{longtable}{|p{3cm}|p{3cm}|p{3cm}|p{3cm}|}
\hline
\textbf{Positive Verbs} & \textbf{Neutral Verbs (``say'', ``suggest'')} & \textbf{Subtle Verbs (``imply'', ``elicit'')} & \textbf{Negative Verbs} \\
\hline
\endfirsthead

\hline
\textbf{Positive Verbs} & \textbf{Neutral Verbs (``say'', ``suggest'')} & \textbf{Subtle Verbs (``imply'', ``elicit'')} & \textbf{Negative Verbs} \\
\hline
\endhead

Encourage & Suggest & Insinuate & Condemn \\
Endorse & Foreground & Reveal & Criticise \\
Support & Emphasise & Illuminate & Critique \\
Laud & Highlight & Hint & Attack \\
Celebrate & Illustrate & Infer & Undermine \\
Extol & Delineate & Intimate & Contest \\
Inspire & Showcase & Allude & Refute \\
Praise & Underscore & Elucidate & Provoke \\
Applaud & Explore & Evoke & Galvanise \\
Commend & Express & Engender & Vilify \\
Venerate & Assert & Invite & Disparage \\
Glorify & Demonstrate & Induce & Belittle \\
Elevate & Propose & Explicate & Scorn \\
Bolster & Announce & Interpret & Censure \\
Exalt & Declare & Clarify & Castigate \\
Glamorise & Portray & Exemplify & Reprimand \\
Reinforce & Depict & Characterise & Harangue \\
-- & Convey & Resolve & Denigrate \\
\hline

\end{longtable}

\newpage
\subsection{Table of Adjectives Describing Tone}

\begin{center}
\textbf{\LARGE Tone Reference Table}
\end{center}

\begin{longtable}{|p{2cm}|p{2cm}|p{2cm}|p{2cm}|p{2cm}|}
\hline
\textbf{Happy Tones} & \textbf{Sad Tones} & \textbf{Angry Tones} & \textbf{Humorous Tones} & \textbf{Neutral Tones} \\
\hline
\endfirsthead

\hline
\textbf{Happy Tones} & \textbf{Sad Tones} & \textbf{Angry Tones} & \textbf{Humorous Tones} & \textbf{Neutral Tones} \\
\hline
\endhead

Light-hearted & Disappointed & Accusatory & Satirical & Formal \\
Hopeful & Shameful & Condescending & Cynical & Matter-of-fact \\
Elated & Resigned & Aggressive & Taunting & Informative \\
Enthusiastic & Sympathetic & Disgusted & Facetious & Reflective \\
Optimistic & Bleak & Incredulous & Amused & Contemplative \\
Encouraging & Sombre & Demeaning & Animated & Apprehensive \\
\hline

\end{longtable}
\newpage
\subsection{Table of Linking Words and Phrases}

\begin{center}
\textbf{\LARGE Linking Words and Phrases Reference Table}
\end{center}

\begin{longtable}{|p{4cm}|p{4cm}|p{6cm}|}
\hline
\textbf{Points of Comparison} & \textbf{Points of Difference} & \textbf{Points of Causality} \\
\hline
\endfirsthead

\hline
\textbf{Points of Comparison} & \textbf{Points of Difference} & \textbf{Points of Causality} \\
\hline
\endhead

Also & Conflictingly & Consequently \\
Likewise & In contrast & As a result \\
Equally & Contrastingly & Because of this \\
Synonymously & Contrarily & Thus \\
Analogously & Unlike & Hence \\
Akin & Nonetheless & Therefore \\
Just as & A point of disparity & Ultimately \\
So too & Conversely & Inevitably \\
Similarly & Counter to & Accordingly \\
Furthermore & Contradictorily to & Ergo \\
Moreover & Opposed to & As a consequence \\
In the same sense & Inversely & This leads to \\
-- & Alternatively & In turn \\
-- & However & In light of this \\
-- & Although & Given that \\
-- & On the other hand & Since \\
-- & Despite this & This culminates \\
Yet & -- & This foreshadows \\
Whereas & -- & For that reason \\
Incongruously & -- & This is a repercussion of \\
\hline

\end{longtable}


\newpage

\section{Analysing Visual Language}
Many students find the analysis of images an accessible and enjoyable part of Year 12 English. It is important, however, to prepare for a variety of visual elements that may appear on your Unit 3 SAC task on analysing argument, as well as on Section C of the end-of-year examination. Below are some elements to keep an eye out for.

\subsection{Context and Position}
In order to fully understand the meaning of the visual elements in a text, you need to understand the context. Be sure to read any background information accompanying the task so that you know what the main issue is. Both the background information and the written material will help give you a sense of the visual’s purpose.\\

Remember that the focus of Area of Study 2 in the Unit 3 SAC and in Section C of the exam is on the analysis of argument. Noticing where and why a visual element is placed in the written piece can help you to tackle the idea of how the argument is structured.\\

\begin{itemize}
\item Is the visual at the start of the article to elicit empathy or shock?
\item Is it in the middle to reinforce a particular piece of evidence or reason?
\item Is it at the end to consolidate the argument as a whole?
\item Is it a frame or border to the written piece that contributes to the overall approach and tone?
\item Alternatively, if the visual is a completely separate text from the written piece, does it present the same or a different point of view on the issue?
\end{itemize}

\textbf{Ask yourself:} where is the visual positioned and why does it matter?

\subsection{Colour and Design Elements}
Although the task material for Section C of the exam will be printed in black and white, your Unit 3 SAC might well contain colour images. If this is the case, you can certainly explore the mood that is created by the use of colour, as well as any relevant connotations of the colour scheme. In a black-and-white image it is possible to discuss degrees of shading and how these affect meaning.\\

It is important to remember that visual language can pertain to design elements in the text you are analysing. For example, in the task material for Section C of the VCAA sample exam the speaker’s first slide has illustrations incorporated into the year ‘2010’. The embedded images of fish, a flamingo, waves, a tree, and an adult and child holding hands help to promote the speaker’s arguments about the beauty of biodiversity and humanity’s central involvement in its protection.\\

Also be alert to the design elements of internet text types such as web pages and blogs. You can, for example, discuss the use of icons, how they position the viewer and the symbolism embedded in them.\\

\textbf{Ask yourself:} how do design elements, colour and/or shading of the visual contribute to its meaning?

\begin{figure} [h!]
    \centering
    \includegraphics[width=1\linewidth]{Screenshot 2025-07-24 105039.png}
    \caption{\textit{Insight Example Image Analysis}}
    \label{fig:enter-label}
\end{figure}

\subsection{Cartoons}
Cartoons often accompany an opinion piece; they can also be stand-alone texts offering a viewpoint on a current issue – especially one that involves politicians or high-profile figures. To understand a cartoon it can be important to have an understanding of the issues circulating in the media in any given week. As I write, there are two quite different cartoons in The Age and The Australian that contain an image of a whiteboard, among other elements. To make sense of both cartoons, a reader would need to know that earlier in the week security guards in Parliament House used a whiteboard to shield a politician from photographers and journalists. Without this knowledge it would be difficult, if not impossible, to interpret the cartoons. Try to stay informed about current affairs as you practise analysing media texts, so that you can grasp their full meanings.\\

Cartoons tend to exaggerate the features of those they represent and use minimal – often quite pithy – text. Be mindful of where this text is: does it appear in a speech bubble or balloon and therefore represent a person’s thoughts or speech? Or is the text elsewhere – on a sign, say, or framed as a caption – providing additional context or comment?\\

The tone of a cartoon is an important feature to note. Political cartoons in particular are often sarcastic and position the reader to be critical of politicians and/or their views and policies.\\

\textbf{Ask yourself:} how does the cartoon make me feel? Your response will give you a guide to the tone of this visual.

\subsection{Photographs}
Photographs can provide a wealth of material for analysis. Look for a range of elements, including:

\begin{itemize}
\item whether the shot is low-angle (looking up at the subject), high-angle (looking down on the subject) or straight-on – for example, a low-angle shot of a child might give the impression that they have been able to overcome adversity and achieve success, whereas a high-angle shot might suggest that a person is vulnerable or remorseful
\item whether the subject is in the foreground (making them very prominent) or background (emphasising their context)
\item whether elements of the photograph are blurred or in soft focus
\item whether a filter (e.g. warm or cool) is used
\item where the horizon is placed in a landscape shot (e.g. far away can signal isolation or opportunity)
\item what is actually in the frame (mise en scène) and what might have been left out
\item whether a text overlay has been used and how that frames your reading of the photograph.
\end{itemize}

\textbf{Ask yourself:}What would the impact be if I changed an element of the photograph (e.g. camera angle or focus)? Answering this question will give you better insight into what the photograph is intending to achieve.


\subsection{Charts and Diagrams}
Charts and diagrams containing numerical information can be used to present evidence. Although you are not in a position to question the data, you can explore how the information and its visual presentation are used to strengthen and develop an argument. It is good to watch for:
\begin{itemize}
\item what the visual elements suggest and how they connect to each other
\item whether a background image can be seen beneath a written overlay (e.g. an image of a gun beneath statistics about the number of deaths caused by shooting)
\item whether data is attributed to an authority (such as the Australian Bureau of Statistics).
Sometimes charts and diagrams can look dry and intimidating.
\end{itemize}
\textbf{Ask yourself:}how does the information relate to the written text or texts? Does it support or contradict a written text and, if so, how?

\subsection{Concluding Remarks}
For the Exam, it is imperative that you compare the written and visual material supplied. It is also wise to compare the written and visual texts in the examination. Comparing written and visual material in terms of the arguments they present and the ways in which they position the audience will give your analysis more depth and cohesion, and demonstrate a stronger understanding of each.\\
\begin{figure}[h!]
    \centering
    
\includegraphics[width=1\linewidth]{Screenshot 2025-07-23 161319.png}
    \caption{\textit{Section C image from 2024 VCAA English Exam}}
    \label{fig:enter-label}
\end{figure}

\newpage
\section{Practice Articles for Analysis}

\subsection{This Book Week, let’s celebrate the power of reading (not just the costumes)}

\textit{By Tracey Cheung — August 21, 2023, 11.25am}\\

It’s officially that time of year when parents scour the shops for the perfect Book Week costume. For some, making sure their little one looks just like their chosen storybook character involves face-painting, finding the right glasses, hat, bow or cape — sometimes tail and ears.\\

Among all this, it can be easy to lose the true meaning behind Book Week. Wendy Rapee, chair of Children’s Book Council of Australia, says the event is an opportunity to draw attention to the joy of reading and for advocating for stories in children’s lives.\\

Rapee explains this is the reason behind this year’s CBCA Children’s Book Week theme: \textbf{Read, Inspire, Grow}.\\

\begin{quote}
``Like dressing up, reading books presents the opportunity to walk in someone else’s shoes, to experience adventures, feel things they have never felt and hopefully learn something about the world they live in,'' says Rapee.\\
\end{quote}

Over the years, I have found that reading stories to my daughter has helped me teach her helpful life skills and ways of thinking. And I’m not the only one with this experience.\\

Sue Cahill, a student wellbeing leader at St Charles Borromeo Primary School in Melbourne, says she finds storybooks extremely powerful in teaching her school children how to deal with various issues, such as school refusal, friendships and bullying. She used the book \textit{I Am Jack} by Susanne Gervay to help two boys in a bullying situation.\\



In the book, Jack is bullied. Cahill says the story helped the boy being bullied to see that, just like Jack, he could also find his inner strength, reach out to others, be heard and change the situation. It also helped the bully see how it made the other child feel.\\

Cahill explains that children see themselves in book characters, which helps make the message feel real.\\

\begin{quote}
``After reading to them, I would ask, ‘How would you feel? What would you have done differently?’'' says Cahill.
\end{quote}

\begin{quote}
``They feel safe when they’re not talking about themselves, but about the third person. This gets them to see it from another’s perspective and builds empathy. You can start to unpack the values and learnings, using the story in the book as a trigger for conversation.''\\
\end{quote}
\begin{figure}
    \centering
    \includegraphics[width=0.75\linewidth]{Screenshot 2025-07-25 085051.png}
    \caption{Student wellbeing leader Sue Cahill with her granddaughters Ella and Audrey.}
\end{figure}


\subsubsection*{Learning valuable relationship skills}

Venita Dimos is a lawyer, mediator and the author of children’s book series \textit{Mini and Milo}. Each book in the series has a different theme, including mindset, communication and listening.\\

As a lawyer and mediator, Dimos has witnessed the profound impact of unresolved conflict in families and workplaces. Her series was driven by a heartfelt need to teach children key relationship skills early on in life.\\

\begin{quote}
``Managing conflict is not a skill you are born with – it’s a skill you need to learn,'' she says. ``Unfortunately, most people don’t learn this in early childhood, so grow up learning on the job with their relationships at stake.''\\
\end{quote}

Recently, I read Dimos’ \textit{Mini and Milo: The Teeny-Tiny Voice}. At the time, my daughter had picked up a swear word from somewhere. When I told her to stop saying it, she said to me, ``I want to stop but the voice inside my head keeps telling me to say it.'' The book taught her that we have greater power over our thoughts than we think, which eventually helped her to stop swearing.\\

\subsubsection*{Discussion of sensitive topics}

Books have also helped encourage conversations with my daughter about issues I have otherwise felt unsure about how to open myself. For example, while I found it awkward to start conversations about the topic of body safety, reading \textit{My Body! What I Say Goes!} by Jayneen Sanders helped break the ice.\\

Lucy Kippist, 44, a community manager of Mums \& Co and mother of two from Sydney, found books helpful when helping her children get through the grief of losing their grandparents.\\

She says \textit{How Can I Help Papa? A Child’s Journey Through Loss and Healing} by Elissa Al-Chokhachy helped give her children new words to describe their feelings around grief. While reading, they learned that part of grieving someone you loved is to remember the happy times.\\

\begin{quote}
``It elicited some happy memories with their grandfather. We talked about the more fun things, like how he used to do Lego with them,'' says Kippist.\\
\end{quote}

Alyson Simpson, a professor of English and literacy education at the University of Sydney, says parents don’t need to expect that a reading session will be a literacy lesson.\\

\begin{quote}
``Allow children to explore the books,'' she says. ``My grandson finds more things in the pictures than I do sometimes.''\\
\end{quote}

\subsubsection*{Fostering friendships}

Elisha Casagrande, 35, found that her son was able to find a message in a book without any prompts from her. The Brisbane-based health coach recently started home-schooling her son, who was having trouble being apart from his friends, so she decided to read \textit{The Giving Tree} by Shel Silverstein to him.\\

In the book, a boy makes brief visits to a 'living” tree several times over the years, until he becomes an old man.\\

\begin{quote}
``My son picked up that relationships can still last, provided there’s some connection over time, like the man who kept connecting with that tree,'' says Casagrande.\\
\end{quote}

When I asked my 5-year-old daughter who she was going to dress up as for her first Book Week, it was an easy choice: Cinderella. Although I know she is drawn to the story because of the pretty princess, after reading the book together, I also know she understands the morals behind the story, such as the importance of kindness.\\

\newpage
\subsection{I’m Melbourne Grammar’s school captain and I’m gay. The Presbyterian Church would have me sacked}

\textit{By Daniel Cash — Melbourne Grammar School Captain\\April 18, 2023 — 7.30pm}

The Presbyterian Church of Australia has declared its right to refuse school leadership opportunities to students based on sexuality or premarital sexual activity. This comes in response to the Australian Law Reform Commission’s review of discrimination laws.\\

The Church has reasoned that, in the case of discriminating against students based on sexual preference, gay students \textit{``would not be able to give appropriate Christian leadership in a Christian school which requires modelling Christian living.''}\\

I write this as captain of Melbourne Grammar School. I write this as someone who is gay, a fact about myself which I cannot change. I also write this as someone confused — surely 'modelling Christian living” is not well exemplified in the practice of exclusion or discrimination?\\

The Christian faith is built on the example Christ set in his life. His revolutionary compassion, so radical then as it still is now — care for the outcasts, the lepers, the unclean — inform the values of Christianity and direct the moral constitution of the Church. Jesus shocked the established powers of his time, his love the instrument of a legacy which left a world faith developing in its wake.\\

His example leads me to wonder what Christ would make of this situation. If he was confronted with a body of young people — all, to remind the Presbyterian leaders, made in the image of God, in the \textit{imago dei} — would he divide them and cast inequality among their ranks?\\

Would he tell a portion of them they could not lead their friends, simply because those they love are of the same sex? Well, Jesus did not condemn the adulterers, did not avoid the pariahs, and surely would not deem some more capable of Christian living than others because of the details of their sexuality.\\

I should not have to pen this opinion piece. Not in modern Australia. But here I am, a student, writing an open letter to a group of adults who have lived a good deal longer than I have, and my intention here (at the risk of sounding quite bold) is to educate. I ask the Presbyterian Church of Australia to withdraw their response which outlines the right to discriminate among children and adolescents based on sexuality or sexual activity.\\

Firstly, it is a violation of the privacy of students and a rather clumsy attempt at preserving a worldview now thankfully frowned upon. Secondly, it is the most remarkable antithesis to the true Christian values (the central commandment \textit{'love thy neighbour as thyself”} should be the directing maxim here).\\

\begin{figure}
    \centering
    
\includegraphics[width=0.75\linewidth]{Screenshot 2025-07-25 085238.png}
\end{figure}

And, thirdly, this desire to discriminate among young people is plainly unsettling. The sexuality of students is not something which should be the subject of so much scrutiny — as Stephen Fry once pointed out, the stricter branches of the Church often like to brand homosexual people as exhibiting unnatural sexual attitudes, though it is not the gay community which displays an unhealthy obsession with the regulation, repression and restriction of sex; rather, the shoe is on the other foot there.\\

I am very lucky — I am secure in myself; I attend a school which values all its students equally. Yet so, so many children do not yet have this same confidence. The influence of the Presbyterian Church of Australia in declaring the right to discriminate against students based on sexuality is therefore mightily dangerous to some of the most vulnerable in our society.\\

Those people behind this submission to the ALRC need to exhibit the compassion which their faith is built on, as do the other religious groups which have submitted similar requests.\\

Organised religion is a powerful force for positivity. Those groups which tarnish the name of spirituality by using it to justify prejudice do a disservice to all people of faith, and they endanger all those who fall under their care.\\

\begin{flushright}
\textit{— Daniel Cash, Melbourne Grammar School Captain}
\end{flushright}

\newpage
\subsection{Confession of a Knox Grammar old boy: my son won’t lose sleep as I have}

\textit{By Malcolm Knox — Journalist, author and columnist\\September 2, 2022 — 10.15am}\\

An old schoolmate invited me to his home for dinner. We hadn’t seen each other in the 30 years since we had left Knox Grammar, the school that hit the headlines this week after boys were suspended or left after posting what their headmaster, Scott James, called ``inappropriate images and \ldots\ offensive commentary in an online private chatroom''. Much of the content was so stomach-turning it could not be published by \textit{The Daily Telegraph}, which broke the story.\\

Back to my schoolmate. Towards the end of a pleasant evening, he revealed he had never got over his bitterness at the horrendous abuse he had received at Knox (from other students, not the teachers who were secretly abusing boys sexually in that same era).\\

The damage he revealed, even as an adult with a good life deriving from his mighty intelligence, his lovely young family and the educational headstart he’d had in life, was chilling. I thought: \textit{Uh-oh.} After a silence, I said, ``I want to know – was I responsible for any of that?'' He said with a sharpness that showed how hurt he still was, ``I wouldn’t have you in my home if you were one of them.'' The awful truth for me was, I honestly wouldn’t have known if I had or hadn’t been ``one of them''.\\

In the discussion of the most recent Knox Grammar disgrace (like many private schools, it can count scandal among its many traditions), much focus is on online malignancy, in this case in an app the boys used called Discord. Attention has been drawn to the \$10 million or so that taxpayers provide the school, including inadvertent state overfunding last year of some \$1.8 million, on top of more than \$80 million poured in by parents and donors. If you want to use the word ``offensive'', annual school fees are up to \$35,130. If they had kept pace with inflation since I was there, those fees would now be about \$9000. The resultant shift towards economic elitism speaks for itself.\\

When the blame is not placed on entitlement, toxic masculinity or social media, it is shifted back to those individuals and their parents. Poor examples at home, absence of morals, no limits. These incidents recur year after year, or the tiny sliver that is uncovered and reported, and responsibility becomes so diffuse that the remedy shrinks back to the simple punishment of the children involved. And so the preconditions continue.\\

I have no idea whether adolescents today behave worse or better than in the past, nor whether teenage boys are any different from teenage girls. Many of us would like to think that the narrowing of economic entitlement produced by elevated school fees has lowered morality to a kind of local Eton-Harrow gutter. The online environment generates its own sewer mentality, which is also soothing to those who already believe electronic devices are the source of all evil, but I’m not so sure any of those factors is more to blame than any other.\\

What I do know, from my schoolmate’s ongoing anger and my phantom guilt, is that the worst bullying from teenaged boys and girls stems from strength in numbers, and it is essentially a performance in front of a group. It is simplistic to think these boys are actually glorifying Adolf Hitler. What they are doing arises from an intoxicated over-exuberance that renders the perpetrators too stupid and too full of themselves to have the first idea of what they are doing.\\

These particular teenagers are not solitary wolves brooding alone over why they are victims in this world (those are the really dangerous ones who, in countries where they have easy access to guns, go out and use them). The bully groups believe they are among the strong, whose group identity is consolidated in attacking the weak. They do know right from wrong; the problem is that ``wrong'' is just another boundary that they can compete with each other in crossing. You want to cause shock. You escalate it to anti-Semitism, gross misogyny and homophobia precisely because you know they are forbidden places; you want your friends to gasp and giggle at your audacity. Drunk on group approval, enclosed in your little group, you think your crime is victimless. This is not solo violence; it is one big circle-jerk.\\

I recently stumbled across my HSC copy of Jane Austen’s \textit{Pride and Prejudice}, which I had annotated in those days when our group’s competitive over-exuberance might have unwittingly – witlessly – harmed schoolmates. Aside from getting top marks for stating the bleeding obvious, Young Malcolm’s notes showed a peculiar viciousness towards Mrs Bennet. Austen makes her a comical figure, of course, but her interfering and alarmist ways evidently got right under Young Malcolm’s skin. He felt sorry for Elizabeth but sorriest for poor old Mr Darcy, the uber-wealthy landowner victimised by the hysterically anxious mother-of-five. Young Malcolm’s notes escalate into misogynistic territory, egged on by extra comments by friends, towards what we all knew were outrageous slurs against a type of woman. So much for the civilising power of Austen. We knew it was wrong; being wrong was what made it funny. This was two decades before social media, two decades before the private school fee explosion. It wasn’t as hateful or propagated as widely as what Knox boys in groups do in 2022, but it was on the way.\\

So, although all the usual suspects feed into today’s brand of excess, I really believe the enduring toxicity goes back to competition, group dynamics and boys performing for each other. Unfortunately, it lies in a continuum with high-status performances in sport, on the stage, in the classroom, and in social life. Excess, good and bad, comes out of competitiveness. If we are going to deal with the dark side of boys trying to outdo each other, we also have to come to a more mature and rational reckoning with all sides of the competitive encouragements of late-stage capitalism. Boys conceal all sorts of emotions, not least their fear of the future and fear of who they really are, under the cloak of competitiveness that they are constantly told is a good thing.\\

I do harbour hope. My son would never have dreamt of behaving like those Knox boys, even though he was on his device a lot and had a privileged education. I don’t put it down to good parenting. I put it down to his innate morals, his empathy after having been bullied himself, and his not being part of an alpha male group inebriated on their own performances.\\

He’s a better person than I was at the same age, despite – or really, because of – not having as many friends, not being as full of himself and not being over-rewarded for competitive success. It also helped that he went to a school that did not see competitive success as the be-all and end-all, or ingrain the fear in young people that life would eventually sort them into winners and losers. I know for sure that in 30 years, he won’t be kept awake wondering if he unknowingly made one of his schoolmates’ life a misery.\\

\newpage
\subsection{Don’t ask me to give the Queen a minute’s silence, ask me for the truth about British colonialism}

\textit{By Lidia Thorpe\\Wednesday, 14 September 2022 — 12.07 AEST}

\textit{'We’re not grieving a singular human life, we’re reeling from the violence that is the legacy of the monarchy.”} — Photograph: Mick Tsikas/AAP

News of the Queen’s death broke on the same day as my cousin’s funeral. My cousin, like more than 500 First Nations people in the last 30 years alone, died in custody. I was notified on Monday of another death in custody.

The institutions that British colonisation brought here — from the education that erases us to the prisons that kill us — are designed to destroy the oldest living culture in the world. That’s the legacy of the crown in this country.

The ``British empire'' declared a war on these shores, against this country’s First Nations peoples. This led to massacres. And you want a minute’s silence from me?

Their war continues and is still felt today — on our children, our men, our land, our water, the air we breathe. Yet we’re meant to kneel to the colonising force with our hands on our hearts?

It is insulting that the prime minister, Anthony Albanese, has called the 22 September public holiday a \textit{``National Day of Mourning for Her Majesty The Queen''}, when First Nations people have called for 26 January to be acknowledged as a Day of Mourning since 1938.

We called for a Day of Mourning so that this country could understand how we’re still affected by colonisation today. \textit{We’re not grieving a singular human life, we’re reeling from the violence that is the legacy of the monarchy.}

Who gave permission for our flag to be lowered to half-mast? That power has been taken away from us, again.

Our young people are turning to suicide, because they don’t see a future for themselves in the systems the monarchy created in this country.

I’m writing this in solidarity with First Nations people who are triggered by the glorification of our oppressor, to the mothers who will have to listen to people grieving the Queen while living in fear of their babies being stolen, and our kids who will have to pay their respects at school despite the British empire’s attempts to wipe us out.

Colonial governments are complicit in continuing the crimes against First Nations people because that is the agenda of the coloniser against First Nations people around the world. To think that people can stand and celebrate that is ignorant and painful.

First Nations people are proud, sovereign people. The British empire tried to wipe us out, but they failed. We are still here and we have an opportunity to do things differently in this country.

The ``demise of the crown'' comes with the opportunity to listen to grassroots First Nations people, who are telling the truth about the monarchy in this country.

Breaking the cycle of colonial oppression starts with truth-telling, and truth-listening. Understanding that this is our shared heritage will help us move forward and heal this country.

Truth-telling holds up a mirror to our society and asks: who are we? Where are we coming from, and where do we want to be?

Do we want to be a nation who can’t control our own affairs? Where our head of state is whoever was born into the right family? Do we want to champion stolen wealth and privilege, or are we a country that cares about equality and democracy?

We’re one of the few Commonwealth countries that doesn’t have a treaty with First Nations people.

\textbf{Treaty is an end to the war.} It’s when we come together to negotiate how we can live alongside each other, peacefully. Treaty will create a new national identity — one we can all feel proud of and a part of.

If we want to mature as a nation, the only mechanism to truly unite us is a treaty.

\textbf{Join us in creating a nation that we can all celebrate.}

\newpage
\section{Section C from Previous Exams}
\subsection{Strapleton Community Voice: Public Post by Daphne Lower}

\textbf{Background information}

\textit{The website \textbf{Strapleton Community Voice} has a social media page for the exclusive use of residents of the town of Strapleton. Members use \textbf{Strapleton Community Voice} to share information and concerns about the local area. Daphne Lower is a member of the group and created the post on the following page.}

\medskip

\textbf{Strapleton Community Voice} \\
Public group \\
11,727 followers\\

\textbf{Daphne Lower} \\
\textit{2 October at 11:35}\\


I want to share with you a dilemma I encountered this morning as I was taking my usual stroll around our picturesque park and lovely lake area. I noticed a new sign near the water saying: \textbf{DO NOT FEED THE BIRDS.} Metres away from the new sign stood a man and two children busily feeding a flock of wild ducks, totally ignoring the very clear message not to do so.\\

What should I do, if anything? This is what went through my mind …\\

When I was growing up, we often fed wild birds. My mum had a little wooden bird feeder hanging from a tree in our front yard. Every morning, she’d lower the rope attached to the bird feeder, place a sprinkling of seed inside and carefully hoist it back up. That bird feeder attracted all sorts of beloved parrots and we would watch them, in awe of their majesty. On family trips to the beach, we’d eat hot chips and throw some to the swarm of seagulls begging at our feet. Even on the few times we ventured overseas, we saw lots of people donating scraps of their lunch to pigeons or scattering bread for the ducks at the local park. Everybody, everywhere, fed wild birds back then. Today, millions of Australians continue to feed wild birds. Why do we do it?\\

It’s easy to see why we do it. Feeding birds feels like an act of kindness and seems to be mutually beneficial. We help the birds survive and, in return, we get enjoyment, feeling as though we are giving back to our community and connecting with the world around us. Taking time to appreciate nature is a great way to improve psychological wellbeing. In our increasingly technology-saturated world, it’s healthy and happy, even uplifting, to unplug and recharge. The duck-feeding dad probably worked pretty hard to coax his kids away from their screens in order to do something outdoors for a change.\\

Unfortunately, despite the immense health benefits that humans gain from feeding birds, the same can’t be said for the birds! Yes, in some parts of the world, certain species of birds would not survive the long winter without humans lending a helping hand. But I wish more people understood that supplementary feeding, even though it might feel like an act of love and nurturing, is often the opposite.\\

Giving wildlife the wrong food in the wrong amount isn’t only unhelpful, it can be downright dangerous. This is the reason I am posting today.\\

I have read that wildlife experts agree that supplementary feeding creates an imbalance in nature that can have drastic consequences. It increases the numbers of some species and this, in turn, decreases those of other species. The rainbow lorikeet, for example, has taken great advantage of its good looks, so its numbers have soared due to supplementary feeding. Sadly, the rainbow lorikeet’s population boom has caused numbers of the much plainer scaly-breasted lorikeet to plummet.\\

Even the most well-intentioned bird lover can upset Mother Nature without meaning to do so. Left uneaten in damp conditions, seeds and bread rot, grow mould and introduce disease to the birds. The uneaten food also attracts pests such as rats, mice and cockroaches. The more food these pests get, the more they breed.\\
\begin{figure} [h!]
    \centering
    \includegraphics[width=1\linewidth]{Screenshot 2025-07-25 085732.png}
\end{figure}

Feeding birds is unnecessary. Native birds will be healthy and happy if left to eat only their normal diet of insects, nectar, seed or fruit. Most people don’t realise that bread — the most popular choice of food to give to wild birds — is actually one of the worst things to feed to wild birds because it contains huge quantities of sugar and salt. But should we put Strapleton’s ducks on a diet or ban bird feeding altogether? Perhaps the real problem is not so much that we are feeding wild birds, but rather that people don’t understand how often, how much and what to feed them. Seeds, duck pellets, a sprinkle of oats or a small handful of peas or corn are acceptable options. If you must feed the birds in our neighbourhood, \textbf{PLEASE} try to do so as an occasional treat, not as an everyday occurrence. Think of it as giving them a healthy snack, not a five-course meal!\\

As I watched those kids and their dad endlessly lob carbs into our beautiful, clean lake, the wild ducks scrambled over one another to fight for every doughy morsel that they could stuff their beaks with. I wrestled with my conscience. I didn’t want to start a confrontation or accuse the dad and his kids of committing a crime against nature, but I desperately wanted to take them under my wing and politely explain a few choices that we can make to better care for our feathered friends. Hence my dilemma …\\

\textbf{What action would you have taken?}\\

\newpage
\subsection{Our Local Gym: Emergency Meeting Speech by Sandy Gestalt}

\textbf{Background information}

\textit{Our Local Gym is a fitness centre that caters to a local community. Until recently, the facility was supported by a council subsidy. This funding has been withdrawn and reallocated to other council projects. The Our Local Gym committee has called an emergency meeting for members of the gym to vote on a proposal to keep the facilities operating.}\\

\textit{The chair of the Our Local Gym committee, Sandy Gestalt, presented the following speech with accompanying slides.}\\

\medskip

Thank you for attending this emergency meeting to vote on the future of Our Local Gym. We find ourselves in an uncomfortable position, but as a group of people who have worked together – worked out together – for years, we know that together we can save Our Local Gym, and ensure its future.\\

Before we move on to possible solutions to our situation, I want to acknowledge the history and the value of Our Local Gym to both you, our members, and to the wider community. Our Local Gym has been in operation for 47 years. We built it from scratch with council support, local fundraisers and fees from you, our members. We have, for years, been a safe space for all people, and for the young and young-at-heart. This space is where you have built your fitness, your strength and your confidence. While larger commercial gyms have flourished, they cater for a different clientele – for people who perhaps are looking for a more competitive place to achieve their fitness goals. People who flock like moths to the flame to shiny new equipment, giant flatscreen TVs, obnoxiously loud music, matching athleisure wear – and who love posting gym selfies on their socials. But Our Local Gym has never been that; we just provide an opportunity for anyone to experience success at any level of fitness. At Our Local, we value acceptance for all bodies and all fitness journeys. We train for life.\\

However, we have recently lost the annual council grant we have relied on for decades to keep Our Local open. As you know, we tried to appeal to the council to reverse its decision but, unfortunately, we missed the deadline and the funding has already been reallocated. The council grant was important to Our Local for two reasons: it supported the maintenance of our building and it enabled us to keep our fees low, in line with the needs and lived experience of our community.\\

In seeking alternative funding, we have considered a series of bleak options: raising admission costs beyond what our community (including young families, refugees and seniors) can afford; privatising the fitness centre; or closing the facility entirely and watching Our Local be demolished and turned into a carpark.\\

But we have found another option. A lifeline.\\

We have been approached by a company about the neglected shopfront on the south side of our building. This group sees real potential in the vintage, grassroots community vibe of our building.\\

The proposal includes an agreement that they will upgrade the shopfront at their expense, and they will pay us rent, which will cover what we have lost in the council grant. This will guarantee Our Local Gym’s low membership and admission fees, and may even leave extra funds for upgrades.\\
\begin{figure} [h!]
    \centering
    
\includegraphics[width=0.5\linewidth]{Screenshot 2025-07-23 161319.png}
\end{figure}
The group is Long Live Lollies.\\

Long Live Lollies is a popular confectionery company. You might be familiar with this brand – they sell lollies at reasonable prices in a retro-themed shop. Recently, their own lease has been terminated. This offers us a chance for a win–win situation.\\

Clearly, while this proposal solves our immediate financial problems, it does present challenges. Our Local Gym and Long Live Lollies do not, at first glance, share much in terms of values. Our Local offers education about nutrition and the best fuel for our bodies and supports members of the community for whom fitness can feel out of reach, while Long Live Lollies sells foods with limited nutritional benefit. There is no doubt that bringing these two businesses together in one building presents a serious clash.\\

Taking into account this clash of purposes – we are for community and they are for profit – I want to make it clear why we support this proposed partnership. The first is – unsurprisingly – money. Long Live Lollies have been nothing but a sweet delight to work with. They understand our dire financial position, offering us up-front rent payments that will keep the lights on and the doors open. They are keen to see our gym thrive and be part of the Our Local Gym community. So they will also pay for gym memberships for their employees – a concrete gesture that shows their support. The second: while eating lollies excessively isn’t great, the occasional treat, or sugar hit, can be a joy. At Our Local, we know that fun and laughter are as important to our health as regular exercise.\\

This partnership can bring colour and deliciousness to our doorstep. And so, we hope that a new community will be enticed to Our Local. Perhaps parents who don’t know we exist will take their children to buy treats from Long Live Lollies, discover us – and join up! We don’t have the advertising budget of large commercial gyms. This is an opportunity to broaden our reach by tapping in to the customer base and marketing expertise of Long Live Lollies.\\

Before we vote, remember that we are not just voting for ourselves. We are voting for the many who cannot attend tonight and for those we hope to reach in the future. This is the only viable solution, and we as a community understand the need for compromise. To save Our Local Gym – this wonderful place where we all get so much for our bodies and souls – we will need to bend and stretch some uncomfortable and perhaps stiff muscles. It will not be perfect but, if we can make it work, it will allow our imperfect community to flourish.\\
\newpage

\subsection{Lighting the way: a new year, a new celeberation}
by Jack Adult, 2025

\textbf{Background information}

\textit{Every year the regional town of Timberoona hosts a fireworks display to usher in the new year. This annual tradition has been an important event on the community calendar for more than 50 years. Over the past few years, however, some residents have been concerned about the impact of the event on the town and have been seeking alternative ways to celebrate New Year's Eve.\\
Local resident JackAdut submitted the following article with accompanying images to the community newsletter Timberoona: Talk of the town.}\\


Like most people in our community, I consider New Year's Eve to be a highlight on my social 
calendar. I look forward to participating in the many festivities held in the beautiful setting of 
Timberoona Park. The family picnics, roaming entertainment, food vans and showcasing of local musical talent in the rotunda create the perfect backdrop for the main event: the annual fireworks display. The spectacular pyrotechnics can be viewed from vantage points across the park and 
beyond, with veterans of the event knowing exactly where to set up their picnic rugs and folding 
chairs to secure the best view. It is with great anticipation that I await the collective countdown to midnight and the first big bang that reverberates through the air! 
As you can see, I'm all for the opportunity to welcome the new year with a big celebration, but 
it's time to re-think the use of fireworks. It wasn't until I saw firsthand the impact that our annual 
celebration has on the environment, some members of the community and their pets (not to 
mention the huge expense!) that I came up with a solution -a dynamic projected lightshow! 
Every year the fireworks display has an environmental effect on our local area. Exploding nine 
tonnes of fireworks generates waste and debris that is blown into the air and lands in waterways and nature reserves, posing a threat to our native flora and fauna. The smoke created by fireworks also pollutes the air and the explosions contribute to noise pollution. Imagine how much the carbon footprint of Timberoona could be reduced if we just replace fireworks with beams of light that are non-toxic, silent and, most importantly, a cleaner option. Add to this the fact that fireworks pose a serious safety and fire risk, especially in December at the height of summer. A single mistake, a 
single wayward spark, can have devastating consequences. The change to a projected lightshow is a no-brainer.\\


\begin{figure} [h!]
    \centering
    
\includegraphics[width=1\linewidth]{dog1.png}
\end{figure}


I know how upsetting the loud explosions and high-pitched noises of traditional fireworks can be 
for little Fido or Fluffy. I understand the yearly dilemma: do I stay home with my beloved pet? Or do I go out and enjoy myself with my family and friends and hope my pet will be okay? Lights silently illuminating the buildings and landscapes in the park are the perfect solution. No more panicked pets cowering under the couch or climbing the curtains. No more scared and lost pets roaming the streets in dangerous traffic. No more frantic owners desperately searching the Timberoona Animal Shelter in the early hours of New Year's Day. How great would that be? \\


A projected lightshow could also offer the chance for some residents to attend the New Year's Eve celebrations for the first time ever. I know of many locals who experience sensitivity to loud sounds and are unable to attend the spectacle. Over the years I have seen some children covering their 
ears and crying in distress because of the intensity and volume of the noise. It would be wonderful to think that these children could come along to enjoy the event without the fear of being so 
panicked and upset. \\
A projected lightshow is a cheaper and smarter investment. The cost of New Year's celebrations will drop from \$1.6 million annually to \$750 000. The money saved could fund more of the 
programs that are vital to the community, such as our meals for the homeless and our breakfast 
club programs. Money that would be much better spent on programs we need, rather than a single environmentally and financially disastrous event that lasts for 10 minutes. \\
I've heard criticisms of the proposal to change to a light extravaganza because of the traditional 
appeal of fireworks. There is much nostalgia about maintaining the old-fashioned tradition of 
singing Auld Lang Syne'1 under a sky illuminated by colour and light with the thunderous sounds of the fireworks as backing music. Not to mention the tradition of having a romantic kiss with a loved one while feeling as if no-one else in the world is watching! Are these New Year's Eve traditions 
really dependent on fireworks? Or, would this all be more spectacular without the deafening din 
and suffocating smoke? A lightshow offers the same ambience with so many advantages. We could have actual music that is choreographed to the movement of the lights! 
Many of you attended the lightshow held in Timberoona Park earlier this year, pictured below, 
to celebrate Community Week, and I think you'll agree it was incredible!\\

\begin{figure} [h!]
    \centering
    
\includegraphics[width=1\linewidth]{dog2.png}
\end{figure}



Crowds came from far and wide to listen to local musicians and DJs, and interact with the 
mind-bending illusions that were created with lights and projections. This event was a stunning 
success, proving that we can create magic in the night sky without a single explosion.\\
Ultimately, this decision is about balance. It's about finding a way to celebrate that keeps the 
wonder of New Year's Eve alive while also being mindful of the needs of our community members, our environment and our future.\\
Together, we have an opportunity to lead the way with a celebration that is spectacular, safe and 
sustainable.\\



\newpage
\subsection{An Open Letter from the Student Music Leadership Group, Narrow Valley}

\textbf{Background information}

\textit{The Narrow Valley region comprises small rural towns that border a large regional city. A group of young musicians who live in Narrow Valley have decided they want more opportunities to perform in their area. The group has written an open letter to the Narrow Valley community, which has been published on social media and in the local newspaper. In addition, copies have been displayed in shops and businesses. The following text is this open letter with images added by the young musicians.}\\

\medskip

\textbf{An open letter to the Narrow Valley community, young and old, who wish to see our region rejuvenated \ldots}\\

Please allow us to introduce ourselves!\\

We are a growing group of high-school music students, teachers and mentors from the five Narrow Valley secondary colleges. We are eager to seize opportunities to bring dynamic and long-overdue changes to our region through an inclusive, youth-driven, energetic and shiny new music festival. Our festival will provide a space for all budding musicians, regardless of age and abilities, tastes and styles to come together and enjoy music.\\

Our passion, talents and all-encompassing abilities enable us to experiment with, explore and master all types of music, including rock, pop, opera, hip-hop (to name just a few), which are all played by our various orchestras, ensembles and solo artists. We celebrate diversity in our local region and already have thousands of followers on social media, who are sharing our endeavours every day.\\

\textit{‘What?!’,} we hear you say. \textit{‘An unrelenting cacophony disrupting the peace and quiet of our streets? Corrupting the impressionable? Shattering windows? Driving the dogs, cats and budgerigars crazy? What \textbf{type} of music did you say? \textbf{Another} music festival? The \textbf{cost?} And \ldots what about our Gorgeous Gardens Festival?’}\\

Yes. What about our garden festival? Who doesn’t love a garden? Like music, gardens never fail to nourish the soul. However, our garden festival is tired. The number of visitors is dwindling – fast. Council data shows that the visitor numbers have dropped by 50 per cent over the past five years and it can only be expected that this trend will continue if we don’t plant some new ideas. With a decline of those proportions, the income from the festival is dwindling dangerously, too. It’s time for a change in tempo and style!\\

For generations, the same Gorgeous Gardens organising committee has done a remarkable job of planning and hosting this event – planting and replanting – but we now need some companion planting. Consider our proposed music festival as ‘festival fertiliser’ that will re-germinate interest and energy in our regional towns, otherwise stagnation will fully set in and we will end up in the compost bin.\\

Few music festivals enjoy the geographical advantages that our region offers. We have a natural amphitheatre overlooking the valley, which could provide not only perfect acoustics but also a stunning backdrop for performances, particularly at twilight. Imagine this cathedral-esque space with the surrounding trees lit up for an evening performance, casting light and shadow on our home-grown superstars. Imagine the sounds of melodic folk tunes wafting across the valley, of hard rock and metal exploding into the night, or the sounds of jazz sizzling under the summer sunset.\\

Like all semi-rural areas, we have buildings that have been standing silent and empty for a long time, including old town halls that could provide brilliant spaces for lessons, workshops and jamming. Our towns can proudly boast that our spaces are accessible to people of all abilities, so musicians and audiences alike can fully enjoy the festival offerings wherever they are staged. The district’s gardens could be opened for a different purpose – not just for admiration but also for musical events. Rock with your roses or mix some Mozart among the marigolds. Dare to put some tubas with your tubers! Melodious offshoots will enrich your world!\\

We know that our Narrow Valley region has seen some tough economic times lately and it might seem that a music festival is a frivolous waste of money. However, we see it as an opportunity to establish innovative partnerships with regional businesses – including fresh food producers, cafes and delicatessens, local artisans and artists – as one solution to overcome financial concerns. Our student bodies don’t just include musicians who can be involved but also media students, foodies, technology gurus and writers. We are confident that we can entice visitors, tourists and music lovers from far and wide to stay, eat and explore. We can see only \textit{increased} prosperity for our region.\\

\begin{figure} [h!]
    \centering
    
\includegraphics[width=1\linewidth]{Screenshot 2025-07-25 085941.png}
\end{figure}

Imagine, too, how a visitor to our festival will see new musicians not in the cutthroat world of television talent shows but against the wholesome natural backdrops and the quaint historic buildings that abound in our regional towns. Visitors will eat fresh produce made into delicious meals and buy local. Visitors, residents and local businesses will all be much richer, and not just financially!\\

We are the youth of our region. Younger voices need to be heard so that our communities can find new energy, new purpose and new perspectives. It is time for generational change. We have talents that deserve to be seen and appreciated and nurtured – and not just on social media. Let’s share our triumphs. Let’s share our aspirations. Let’s provide opportunities for \textit{all}.\\

We invite you, the residents of the Narrow Valley region, to our initial meeting, next Tuesday, in the Narrow Valley Botanical Gardens. The meeting will be followed by a performance by the Narrow Valley North College band. Please bring your ideas, your energy and even any concerns, so that we can begin our shared journey towards a new, prosperous future through our state-of-the-art music festival.\\

\vspace{0.5em}
\textbf{Yours, in harmony,} \\
\textbf{The Student Music Leadership Group, Narrow Valley}

\newpage
\subsection{\textit{Hapsey Happenings} – Transcript of podcast, episode 32}

\textbf{Background information}

\textit{\textit{Hapsey Happenings} is a podcast hosted by lifelong Hapsey resident Janelle Tanley. She discusses local events and issues, and promotes the conservation of the town’s local environment. In the most recent episode of her podcast, Tanley discusses a proposed development plan for the Hapsey Wetlands by Nature+Adventure Group, who run commercial excursions into unspoilt areas. The following text is a transcript of this episode, with images provided by Tanley.}

\medskip

Hello and welcome to \textit{Hapsey Happenings.} I’m your host, Janelle. I’m sure all of you listening have heard about the recent development plan that has been proposed for the Hapsey Wetlands, one of the very last protected sanctuaries in this local region. Nature+Adventure Group’s plan includes bike paths, boat ramps, kiosks and cafes. The local newspaper is being flooded with letters and emails from people for and against the plan, and my socials have also been overflowing.\\

Let’s consider our wetlands. They are quite a distance from the town, so they’re mainly visited by nature lovers – people who don’t mind a walk from where they can leave their car. The paths are narrow and unpaved. Walkers have to be careful but this means the native flora is taken care of. My friends and I frequently meet in the wetlands to relax but we talk softly without disturbing the tranquil surroundings, so the native fauna can live in peace.\\

But now, our wetlands seem to be just another area of Hapsey under threat of destruction. In my lifetime, I have often seen people bring proposals for developments to council meetings, which are disguised as plans to protect the environment, claiming that they will provide access for people who would otherwise be unable to enjoy the natural beauty of the area. The suggestion that this proposed ‘revitalisation’ will protect the environment and preserve the magnificence of this natural space for future generations sounds like something that is said to placate those of us who want this habitat to be kept as it is. In fact, the destruction of the natural home of over 40 different species of grasses to put up money-making ventures – cafes and so on – serves the interests of humans and businesses rather than the precious environment.\\

As locals, we are the guardians of this area. Many of us have worked tirelessly to protect these native plants and animals from all kinds of danger. We know from experience what will happen if the development plan is approved. First, Nature+Adventure Group will send in digging machinery and desecrate the home of fragile species, which will be damaged beyond repair. Then, they’ll build footpaths and bike paths. Inevitably, one careless tourist will go stomping on a delicate fern to look at a frog and will take out an entire species that is unique to our wetlands. How are we going to protect these plants, flowers, insects and other animals from human intrusion?\\

This current development plan will only be the start. Attracting walkers, cyclists and café patrons will only give impetus to further development. With increased traffic, Nature+Adventure Group will build car parks (and charge for parking) and facilities such as toilets and gift shops. If we let this proposal through, it will set in motion a chain of proposals that will erode the protection that we have provided for the wetlands all these years. Suddenly, this area will have been concreted over, and the Hapsey Wetlands will have become the Hapsey Wetlands Shopping Centre or the Hapsey Amusement Park!\\

If we’re looking at increasing the opportunities for outdoor adventures in the Hapsey area, there are other equally scenic and tranquil places to the south of the town, where installing walking and cycling paths would not disturb the delicately balanced ecology of the natural wetlands.\\
\begin{figure} [h!]
    \centering
    \includegraphics[width=1\linewidth]{Screenshot 2025-07-25 090126.png}
\end{figure}
These would-be developers claim in their proposal that they have an ‘excellent record in terms of environmental responsibility’. They say that they see cycling and boating as leaving a ‘small footprint on the environment’ and as being environmentally responsible activities because these activities use no fossil fuels. They acknowledge that transporting equipment, such as bikes and boats, will leave a bit of a carbon footprint but that they are working towards reducing the impact of transporting their bikes and boats. Believe that if you like!\\

I realise that these people feel that they are offering an excellent opportunity to the Hapsey community. They claim that they are very protective of our precious resource and that they share the goal of preserving this special place. But will an increase in the number of visitors promote the preservation of the wetlands? Will all these new visitors, grasping disposable coffee cups, share the values of those protecting the native inhabitants of the wetlands? Is the wetlands’ survival helped by people falling off bikes and out of boats? Aren’t these arguments just a selfish attempt to paint a commercial venture in a favourable light? Can we trust Nature+Adventure Group’s mission statement: ‘Above all, we acknowledge that environmental and ecological concerns are our number one priority in the expansion and development of our business’?\\

We must consider what is best for our wonderful wetlands. I would love to hear from you, listeners of \textit{Hapsey Happenings} – how do you feel about this? Leave your comments on my website. Don’t forget to follow the podcast to receive updates about the proposal for the wetlands.\\


\newpage
\section{Sample Analysing Argument Essays}

Below is the article of which these responses are based.

\subsection{Article}
Below is the article of which these responses are based.

\begin{center}
    

\textbf{The Guardian}\\
\textbf{I started using heroin decades ago. The rejection of a Melbourne CBD injecting room feels like a rejection of me}\\
\end{center}
\textit{Anonymous}\\
The government seems to want teams to go scoop up drug users and put them in treatment – whether they’re ready for it or not
Fri 26 Apr 2024

\vspace{1.5cm}

I’ve used heroin off and on for decades now. I still do.\\

Throughout my 20s, I had used drugs in ways that were often destructive to me. Back then, I worked 40 hours a week in a paid job a lot of the time, but at other times I was homeless, seemingly always in withdrawal, and desperate for dollars or drugs.\\

Looking back, I was incredibly destructive. I ripped off shops, crashed cars and got in trouble with the cops. I had hepatitis C (of which I have been cured), have overdosed multiple times and developed infections from injecting – often in an alleyway or an abandoned house or a public toilet or a train.\\

While people have issues with alcohol and other acceptable drugs that make their lives difficult, they don’t go to jail just for buying or consuming their drugs. That fact – and the drive to avoid the police, along with the awful reality of opioid withdrawal – has ensured my decision-making is sometimes focused on immediate relief or safety, rather than long-term health choices.\\

That is why harm reduction and injecting rooms – such as the one in North Richmond – are so crucial for keeping people alive and healthy.\\

When I heard the government in Victoria was not going ahead with the CBD injecting room, I felt both sad and angry. If you are using drugs and are dependent on them, you need somewhere safe to use. 'Treatment”, 'care coordination”, 'stigma-free care” and the rest of those nice words about getting people off drugs mean nothing when you are sick and hanging out and just want to get well. Injecting rooms work because they provide safety at this moment of vulnerability.\\

\begin{figure} [h!]
    \centering
    \includegraphics[width=0.75\linewidth]{Screenshot 2025-07-25 075742.png}
    \caption{A person demonstrates using a syringe during a tour of a medically supervised injecting room in North Richmond}
\end{figure}


This immediate safety is something I might seek out if I was using in the city again. Then once that has been sorted, I maybe would seek other support. I have been let down and discriminated against so many times by so many health services I find it hard to ever trust them.\\

‘Will cost lives’: addiction experts condemn decision to scrap proposed safe injecting room in Melbourne’s CBD\\

For all those in the media and the CBD hoping to get rid of the people like me: I hope you know that people will continue to use drugs in your doorways or overdose in your back alleys, and know that people only care up until the point it inconveniences them. We are real people – who care, hurt, love and feel. But we are only ever defined by drug use.\\

It is so depressing that when I look at the list of things in the government’s 'action plan” released this week, there seems to be a focus on treatment. It looks, to a cynic like me, as though they are funding a massive team of people to go out and scoop up people and get them into treatment to be helped – whether they want it or not. Sprinkled in there are some things that may immediately save lives, like more naloxone and overdose assistance, but it really is a cold comfort next to the loss of the injecting room.\\

As a mate of mine said this week: 'Fuck, it’s so typical ... it was going good with methadone changes [that made it cheaper on the PBS] ... I had some money in my pocket and felt like a person again, but now they do this and make us feel like shit again.”\\

This feels like a betrayal at a fundamental level for someone like me, who has voted for this government and party for years.\\

It is not just about the injecting room. It is because it is so clear that there are votes in words like 'treatment” and 'care” – and listening to shop-owners – but not in doing difficult things for parts of the community who aren’t squeaky clean.\\

The rejection of the second injecting room feels like a rejection of me as a person, partly because the first injecting room here in Melbourne felt like, maybe, they really do see people like me as humans and part of society.\\

It feels as if I was very naive to believe that.\\

Even if all the other things they said they were going to do work out well, that rejection will hurt for a long time.\\

\begin{quote}
    The author lives in Melbourne and has recieved treatement for his heroin use.
\end{quote}

\newpage



\subsection{Sample Essay I}


ZZEXP Analysing Argument Essay I ZZEXPEND


In their article for The Guardian Australia on the 26th of April 2024, titled ‘I started using heroin decades ago. The rejection of the Melbourne CBD injecting room feels like a rejection of me’, an anonymous author with a history of drug use describes how the decision not to open a second safe injecting room in Melbourne is an abandonment of drug users across the state. To an audience of mainly residents of Melbourne, but also to the wider Australian public, the author uses an upset and exasperated tone to outline how the Victorian Government’s decision to scrap plans for an injecting room in Melbourne’s CBD is essentially letting down those who use drugs and feel that safe injecting rooms are a critical part of the solution.\\

The author initially argues why drug users can be so hesitant to ask and accept help in regard to their drug use. Through the use of anecdotal evidence, they discuss some of the experiences they have been through while using drugs. By saying 'Looking back, I was incredibly destructive”, the author highlights to the audience how dire a situation can get for those using drugs and do not have the necessary support networks around them. As such, the audience get an insight into the mindset of a drug user and start to see things slightly more from their point of view. Thus, the audience are more inclined to develop a sense of empathy for drug users, as they can somewhat see what drug users go through without any help. Additionally, the author also appeals to the audience’s sense of fear by recounting some of the more unsavoury incidents that have occurred whilst they were on drugs. The author describes occasions of 'having hepatitis C, overdosing and developing infections”, designed to gently scare the audience into understanding some of the grim outcomes that exist with taking drugs unsupervised. The audience start to feel that this something that they do not want to see on the streets of the city, and consequently start to believe that actions need to be taken to reduce these sorts of incidents from happening out in the open. Furthermore, the author questions how users of alcohol and non-illicit drugs do not end up in jail, whilst people like them who use heroin are so quick to be arrested by police. By suggesting that users like them have a 'drive to avoid police”, the author uses reasoning and logic to demonstrate to the audience that the preconceptions by the police and the general public towards users is a major reason why they do not seek help. As such, the audience are more likely to see a different side to drug users, one that presents them as real humans. This causes the audience to feel uncomfortable and somewhat ashamed the drug users like the author cannot trust those who are wanting to help. This demonstrates how the author feels that drug users cannot ask for help without getting into trouble.\\

The author moves to point out how safe injecting rooms are actually a place of sanctuary and respite to drug users. They start by playing on the audience’s view of common sense by arguing that safe injecting rooms are 'so crucial” in keeping users safe. This challenges the audience’s previous ideas about the role of injecting rooms, as they may have previously seen them as facilities that only enhance addiction and increase drug use. As such, the audience are more likely to see the role of injecting rooms from a user’s perspective and start to build an understanding of why they should exist. The author also outlines some of the key phrases made by the Victorian government following their decision to not open an injecting room in the CBD. By describing these statements as 'nice words”, the author sarcastically calls out the government for their decision. The use of sarcasm helps the audience see how much drug users value injecting rooms, through their anger of the facility not opening. As a result, the audience are more likely to see the benefits of injecting rooms to users and potentially start to support them being opened in more locations across Melbourne. The author also highlights how injecting rooms provide 'immediate safety”, aimed to appeal to the audience’s sense of safety. This causes the audience to feel more sympathetic towards drug users as they now see how crucial they are when drug users are at their most vulnerable. As such, the audience are more likely to understand why drug users are so desperate for these facilities to remain in the city. This represents how crucial safe injecting rooms are viewed as by drug users, and how much they are needed by them.\\

The author also attacks the decisions made by the Victorian government regarding injecting rooms, suggesting that it does not properly address the major issue surrounding drug use. The tone of the article shifts to being more cynical towards the issue, demonstrated by the author describing the government’s action plan as 'depressing”. The audience are more inclined to feel more uncomfortable about the situation that drug users find themselves in and as such start to believe that the government’s policy regarding drug users may not be working in handling the wider issue. The author also includes an attack on the media, who may have sensationalised the risks that drug users present to the community. They argue that media only care about the issue until 'it inconveniences them”, causing the audience to question their trust of these media outlets and their descriptions of drug users. As such, the audience are more likely to sympathise with the author as they now have an insight into how all the perceptions and stigmas that are created by the media have a significant impact on the author, as well as other drug users. This results in the audience to further consider the type of media coverage that they might follow in regard to drug users. Additionally, the author also calls the alternative actions to the injecting rooms that government proposed as a 'cold comfort” to drug users. The combined use of a metaphor and alliteration summarises the author’s feelings regarding the abandonment of an injecting room in the CBD to the audience, resulting in them having a somewhat changed view on injecting rooms and how necessary they are to the issue of drug use. This causes the audience to not only have increased levels of empathy towards drug users and the issues that they face daily, but possibly also view the government’s work regarding drug intake as ineffective in handling this major issue. This could also result in the audience maybe losing a level of trust with the government to be effective in dealing with other issues. As such, the author is able to demonstrate how the Victorian Government has not made the appropriate decisions with handling drug use.\\

The controversial topic of safe injecting rooms is explored in a Guardian article by an individual with experiences of taking drugs such as heroin. Through the lived experiences and appeals at government, they argue how the decision not to open an injecting room in Melbourne’s CBD is not the right approach to take in responding to drug use. They strongly believe that injecting rooms keep users safe and healthy and should be a key part in the long-term solution to drug use.\\

\newpage

\subsection{Sample Essay II}


ZZEXP Analysing Argument Essay II ZZEXPEND


The Victorian Labor government plan to abandon the installation of an injecting room in Melbourne CBD, in favour for a health hub that connects to rehabilitation services, has sparked the issue of whether the Labor government actually cares for people struggling with druge use. The Guardian recently published an article on Friday 26th of April 2024 including personal account of a drug user, who contemprarly addresses thhe Labor governemnet's rejection of the injecting room in the CBD for favour of the local residents. The former drug user also known as the anonymous author of the article, assertively attacked both the general public and the Labor government for outright discriminating drug users.\\

The anonymous author initiates his argument by first explaining his personal anecdote as a drug user wherein he repeatedly claims how ‘destructive’ and how’ awful reality’ was the severe ‘withdrawals’ of drug use. Highlighting his attempt to connect with his audience the residents of CBD in the Labor government, in a deeper level by giving them the chance to sympathise with his predicament, thus showing that drug users are genuine humans who are victims of the environment and only need ‘immediate relief’ for safety in which the injecting rooms wholeheartedly provide.\\

Upon establishing the benefits that the injecting room provides, the author then forwards how rehabilitation centres do not adequately solve the problem logically claiming the empty words of ‘treatment’, ‘care coordination’, and ‘stigma free care’ will not do anything when people just ‘want to get well.’ Resulting in rehabilitation programmes to be stained in a negative light, and in turn point the Labor government who provided these statements is ignorant in solving the problems of drug use. Purposefully presenting them as ‘votes’ driven mongrels that do not support care for the ‘discriminated’. Furthermore the author then compares the injecting rooms with ‘many health services’ in a dogmatic way, presenting the injecting rooms as vastly superior as it provides ‘safety’ in the ‘moment of vulnerability’. Attempting to again strengthen the need for an injecting room in the CBD for their sake.\\

Another key argument the author provides is that the residents of the CBD are just as ignorant as the Labor government. He achieves this by assertively attacking them for defining people for their ‘drug use’ and not viewing them as ‘real people’ and only caring about the problems of drug use when it ‘inconveniences’ them. This consequently shames the general populace for their lack of humanity and compassion to other human beings and presenting the mas no better than the vote hungry Labor government. Additionally, the author then depressive Lee used for the Labor government ‘betrayal’ for favouring the ‘cold comfort’ of ‘treatment’ outright honestly claiming his ‘naïve’ nature for trusting them ‘for years’. Consequently serving to guilt trip the Labor government for betraying the left-wing nature, and the ‘community who are not squeaky clean’, and rejecting the author as a ‘person’ or ‘part of society’.\\

Overall, the guardian and the anonymous author effectively displayed their views of the need of an injecting room in the CBD by first establishing themselves and their poses in equal footing as humans by sharing their perspectives as a drug user and then painting their opponents as ignorant fools that must comprehend the perspectives of drug users and the need of injecting rooms, and health services, that will help them regardless if they are ready for it or not.\\

\newpage
\subsection{Sample Essay III}


ZZEXP Analysing Argument Essay III ZZEXPEND


The personal essay and opinion piece written for the guardian by an anonymous heroin user provides a blunt breakdown of the author's reaction following the recent rejection of a Melbourne CBD safe injecting room. The author outlines the importance of such injecting rooms, and criticises the biases held by many Melburnians opposed the preposition, contending that the decision serves as a betrayal of Melbourne’s most vulnerable.\\

The author opens his piece with a series of brutally honest depictions of his ‘often-destructive’ experience with heroin use. As well as being an engaging introduction, the description served a former personal connection from author to ruler. The use of first person and conversational confession positions a ruler to emphasise with his experience and feel less alien from the issues most concerned stakeholders. His descriptions of using in ‘an allegory’, ‘abandoned house’,  ‘public toilet or a train’, both makes for contrasting imagery but serves A logical purpose. By highlighting these public and often unseen locations, the author implies the very need for safe injecting facilities. The author convinces to address a notion of hypocrisy that heroin uses face more dire pressures than ‘acceptable’ drug users, addressed almost as the reader themselves. This effect is quite confronted confrontational and positions the opposed ruler to reassess his prejudice, while reinforcing a sense of moral standing in the grateful stop here the authors sends to assert the importance of injecting rooms as ‘crucial to ‘keeping people alive.’ The sentiment is forthright and sees to revise the readers empathy.\\

Following the more fundamental assertions of the authors protest, he elates to address the opposition directly. His tone becomes more experienced, with long flowing structure leading to the crux of the argument. No longer a plea for empathy, this section is a seething criticism of the ‘people hoping to get rid of the people like me.’ The author undermines common rebuttal as more ‘inconvenience’, employing a sense of moral authority. All of this served not to It's a sense of empathy, but guilt; A guilt only exacerbated by the prior emotional framing of drug users as ‘vulnerable.’ This component essence of utilising first empathy and then guilts creates a structure that is particularly Difficult to read but. The author particularly the conservative one to which this section is addressed, is caught in a ‘catch twenty-two’ wherein this cannot take a position of moral integrity, without conceding the value of a safe injecting facility.\\

The author then transitions to a sensitive voice, criticising the governments’s ‘action plan’ for its prioritisation of treatment over relief. The author is again controversial, and involves a quote from his ‘mate’, asserting the’ betrayal’. This strengthens his case as not just the testimony of one, the betrayal of an entire community of stakeholders. The reader is then positioned to perceive the authors view as one of relevance and authority. This further results by the political lens next taking, wherein the author elects to admire to the reader that he has voted for his government and party for years. This again is most effective on a reader whom does not share his view, but may share his vote. By asserting a common ground, readers positional to assess their own biases, and furthers the personal connection constructed at the opinion of the peaceful stop the author then flips the politicisation on its head, criticising where votes are sound From a moral view. This again leverages a sense of moral guilt, that the government would prioritise ‘shop owners’ over doing ‘things for the community’. These statements escalate on the guilt restored previously frankly the issue as a single moral dichotomy in the readers eye.\\

The author asserts the efficacy of safe injecting rooms and criticises of rejection as a betrayal of vulnerable people. Believing tone shift, from gathering empathy to explicit criticism the author seeks to frame the issue in a moral and emotional light; Hoping to persuade the less vehement opposition amongst Melburnians.\\

\cleardoublepage

\clearpage
\thispagestyle{empty}
\hfill
\clearpage
\cleardoublepage
\part{The Exam}
\clearpage
\thispagestyle{empty}
\hfill
\clearpage


\section{Exam Format}
The VCE English exam is a three-hour exam, plus 15 minutes of reading time. It has three sections, and you are required to complete all three sections. Each section is worth the same number of marks, so you should spend the same amount of time on each. 

\subsection{Exam Specifications}
\subsubsection{Overall Conditions}

\begin{itemize}
    \item The examination will be sat at a time and date to be set annually by the Victorian Curriculum and Assessment Authority (VCAA). VCAA examination rules will apply. 
    \item There will be 15 minutes of reading time and 3 hours of writing time.
\item The examination will be assessed by a panel appointed by the VCAA.
\item The examination will contribute 50 per cent to the study score.

\end{itemize}


\subsubsection{Format}
The Task Book will be sixteen pages in length. The front page contains the general cover notes and instructions, including rules about what you can and can’t take into the exam.\\

\textbf{An English and/or bilingual printed dictionary is permitted.}\\
\textbf{Mobile Phones and/pr any unauthrosied electronic devices are not permitted.}\\

\paragraph{Specifications from VCAA:}
\begin{itemize}
    \item The examination will be in the form of a Task Book.
\item The examination will consist of three sections.
\item Section A will consist of an analytical response to a text and will be worth a total of 20 marks.
\item Students will be required to write an analytical response to one of the two texts studied from List 1 of the VCE English and English as an Additional Language (EAL) Text List for the year in question. There will be two topics for each of the 20 texts on List 1. Students will be required to write a response to one topic on one text.
\item Student responses should demonstrate knowledge and understanding of the text, its structure and the ideas, concerns and values raised by the topic, and support their response with close reference to the text. If the selected text is a collection of poetry or short stories, students may write on several poems or stories, or on at least two in close detail. 
\item The response will be assessed according to the assessment criteria specified on the following page.
\item Section B will consist of creating a text and will be worth a total of 20 marks. 
\item Students will be required to create one written text. For the purposes of this task, a ‘written text’ does not include song, poetry or verse. In creating their text, students must use the title provided, relevant idea(s) from at least one unseen stimulus and be connected to one of the Frameworks of Ideas. The text must be developed with a clear purpose(s) – to explain, to express, to argue, to reflect – and should be an effective and cohesive piece of writing. Students should demonstrate fluency through the selection of vocabulary, text structure(s) and language features, and employ an appropriate voice in their text. 
\item The place of mentor texts is to provide access to ideas and inform the teaching of writing, but students will not be required to make explicit reference to them in the examination.
\item The text will be assessed according to the assessment criteria specified below. Section B will be worth a total of 20 marks.
\item Section C will consist of an analysis of argument and language and will be worth a total of 20 marks. 
\item Students will be required to analyse the use of argument(s) and language to persuade an intended audience to share the point of view expressed in an unseen persuasive text. For the purposes of this task, ‘language’ refers to written and spoken language, and ‘visuals’ refers to images and graphics. Students will not be required to evaluate the relative merits or successes of argument(s).
\item The response will be assessed according to the assessment criteria specified below.
\item Section C will be worth a total of 20 marks.
\item The total marks for the examination will be 60.
\item Answers to Section A, B and C are to be recorded in the Answer Book.

\end{itemize}

\section{Strategies for Addressing the Requirements}
\subsection{Section A}
In Section A you will write an essay that responds to one of the two given topics on a text you have studied. You can choose to write on either your Unit 1 text or your Unit 4 text. In your answer book you will need to identify the text on which you are writing and the topic (i or ii) you are responding to by shading the relevant box.\\

Spend a few minutes unpacking the topic you feel most comfortable responding to. Look carefully at each term in the topic and make sure you understand its meaning and significance. Use the strategies for identifying topic words, direction words and limiting words outlined on pages 48–50 in English Year 12.\\

Spend about 5 minutes writing a brief plan that includes your contention, three or four topic sentences and one or two pieces of textual evidence for each topic sentence.\\

The following strategies will help you to address the Section A requirements and criteria.\\

\begin{itemize}
    \item Address the topic in full. Strong responses take account of every word in the topic, including the limiting words, rather than writing about one key term in a general way.

\item Use your dictionary. Make sure you understand the meanings of all the words in the topic. Your dictionary will also give you some synonyms for key terms in the topic, and using these will give your writing more variety and precision.

\item Consider a range of textual features. Don’t limit your reasons and evidence to just what the characters say and do. While authors do use characters to explore ideas, they also use features such as narrative voice, structure, setting and imagery. Film techniques should be referred to when analysing a film, and an analysis of poetry will be enhanced by a consideration of poetic form and techniques. Discussing these features shows a deeper understanding of the text and a more sophisticated response to the topic.

\item Integrate quotes and explain their significance. It is useful to include some direct quotes, as this is part of how you show your ‘knowledge and understanding of the text’ (criterion 2 for Section A). However, make sure you incorporate quotes into your discussion and thoroughly explain how they support the point you are making. If you use a quote and then mention it on its own, or place it at the end of a sentence without further explanation, it will be unclear to your assessor how the quote helps to advance your argument.

\item Write precisely and expressively. Develop a wide vocabulary during the year by creating word banks (see the examples on pages 58–63 of English Year 12) and use a variety of words in your practice essays, so by the time you sit the exam you can make effective vocabulary choices in all your responses.
\end{itemize}


\subsection{Section B}
In Section B you will write a text that responds to a given title and at least one piece of stimulus material. In the answer book, you will need to shade the box of your chosen Framework, and the box (or boxes) corresponding to the stimulus (or stimuli) you have responded to.\\

The following strategies will help you to address the Section B criteria and weave the stimuli and title into your creative piece.\\

\begin{itemize}
    \item Ensure your writing is clearly relevant to your chosen Framework. This involves integrating concepts and perspectives that you’ve explored during the year. When you are revising in preparation for the exam, refer to your notes and written pieces to stay aligned with the core ideas of the Framework.

\item Integrate the title and at least one stimulus into your writing. You must use at least one of the stimuli. The stimulus should not just be referenced superficially; it should be woven into the fabric of your piece. This could mean, for example, using a text stimulus as the basis for a central idea; using a text stimulus as a direct quote at a key point in your text; or drawing inspiration from a visual stimulus to set a scene. Similarly, the title provided for your Framework should provide an idea, a character arc, a central image or another key element of your writing.

\item Give your piece cohesion and structure. Your piece should have a clear beginning, middle and end, regardless of whether it’s a narrative, an argument or a reflective piece. Each part should seamlessly transition into the next, maintaining a logical flow of ideas.

\item Make appropriate use of language, style and voice. The language and style of your writing should suit the purpose and audience. This includes choosing the most appropriate vocabulary, tone and sentence structures. In your Unit 3 drafts and assessment pieces, you will have worked hard to develop your own unique voice. You likely have a preferred form of writing (fiction or nonfiction, story, article etc.) and you should lean into that during the exam.

\item Write with fluency and expression. Your writing should read smoothly, with a natural flow of ideas and language. This can be achieved by varying your sentence structures, giving your writing a subtle rhythm and clearly articulating your ideas. Transitions between paragraphs or sections should be logical and easy for the reader to follow.
\end{itemize}


\subsection{Section C}
In Section C you will analyse a text in which the writer seeks to persuade a particular audience to agree with their point of view on an issue. The text will be primarily written but will also contain one or more visual elements.\\

This task has been part of the VCE English exam for many years and as part of your preparation you may use past exam papers as sources of practice material. In some cases you will find that the task material consists of more than one written piece, but the exam specifications for 2024–2027 indicate that there will only be a single ‘unseen persuasive text’ in the English exams during this period.\\

The following strategies will help you to address the Section C requirements and criteria.\\

\begin{itemize}
    \item Understand the contention. What, essentially, is the writer saying? An accurate understanding of the writer’s central argument underpins your analysis. Clearly state the writer’s contention in your introduction.

\item Read the background information carefully. This is where you will find out the intended audience and the context for the text. Your analysis will be strengthened by a consideration of how the writer’s choices of argument and language are shaped by their particular audience and context (see the final bullet point below).

\item Examine the sequencing of the argument. Writers will arrange their material and their points of argument in a particular order so as to engage their intended audience and influence their thoughts and feelings. Look at the overall shape of the argument, then drill down into the specific points and language choices. Shifts in tone are important aspects of the structure, too.

\item Examine language choices. Writers make deliberate choices of words to add impact to their arguments and to evoke the audience’s emotions. Identify words with strong positive or negative connotations and associations and spend some time in your analysis ‘unpacking’ their likely effects on the audience. You will not have time to analyse every use of language in the text, but discuss a representative sample of language choices.

\item Use effective verbs. In your analysis of a persuasive text, verbs enable you to describe what the writer is doing and how they are doing it. Repeating the verbs ‘positions’ and ‘persuades’ doesn’t show your understanding of how argument and language are being used. Verbs such as ‘reinforces’, ‘emphasises’, ‘instils’, ‘elicits’ and ‘undermines’ can be used to show your understanding of how a writer is attempting to persuade. See the verbs in the activity on page 210 of English Year 12 for some other examples.

\item Analyse how argument, language and visuals are working together. Don’t treat these as separate elements of the text, since an audience’s response will be a result of how all these elements are working in combination. Look at how language is used to present an argument, and how an image complements and reinforces particular reasons or pieces of evidence.
\end{itemize}

\clearpage
\thispagestyle{empty}
\hfill
\clearpage
\cleardoublepage
\part{Practice Exams}
\clearpage
\thispagestyle{empty}
\hfill
\clearpage

This section gathers everything you need for timed practice: links to official VCAA past examinations, externally hosted practice papers, the school's purchased trial exams, and a generator that assembles a fresh exam on demand.

\section{English Exam Generator}
The exam generator now runs directly on this site. It assembles a complete three-section practice paper from the question banks and downloads it as a Word document in the authentic task-book format.

\href{ZZFILE::__EXAMGEN__}{Open the Exam Generator}

\clearpage
\thispagestyle{empty}
\hfill
\clearpage
\cleardoublepage
\part{Exam Assessment Criteria}
\clearpage
\thispagestyle{empty}
\hfill
\clearpage

\begin{center}
\textbf{From VCAA's documentation.}\\
\end{center}


Section A will be assessed against the following criteria:
\begin{itemize}
    \item knowledge and understanding of the text, its structure, and the ideas, concerns and values it explores
\item development of a coherent analysis in response to the topic
\item use of evidence from the text to support the analysis
\item use of fluent expression through appropriate use of vocabulary and conventions of Standard Australian English
\end{itemize}
\vspace{1cm}
Section B will be assessed against the following criteria:
\begin{itemize}
    \item use of relevant idea(s) drawn from one Framework of Ideas, the title provided and at least one piece of stimulus material
 \item creation of a cohesive text that connects to a clear purpose(s) and incorporates an appropriate voice
\item use of suitable text structure(s) and language features to create a text
\item use of fluent expression, including the appropriate use of vocabulary 
\end{itemize}
\vspace{1cm}
Section C will be assessed against the following criteria:
\begin{itemize}
    \item understanding of contention, argument(s), and point of view
\item analysis of the ways in which written and spoken language and visuals are used to present an argument(s) and to persuade an intended audience
\item use of evidence from the text to support the analysis
\item use of fluent expression through appropriate use of vocabulary and conventions of Standard Australian English

\end{itemize}
\vspace{1cm}

Assessors mark holistically, relating student performance to the published criteria and ranking students over the full range of marks available. Determination of the mark is assisted by the descriptors in ‘Expected qualities for the mark range’. These descriptors have been written to reflect the level of achievement expected at a particular mark or mark range. Both the criteria and the descriptors are fully explored and directly related to the range of student responses during the assessor training process.\\

For Section B, the Framework of Ideas, textual forms, audience and purpose are explained in the section ‘Framework of Ideas’ on pages 22 and 23 of the VCE English and English as an Additional Language Study Design 2024–2027. These concepts underpin the treatment of key knowledge and key skills, and they are an explicit criterion in the expected qualities for Section B.\\

\newpage

\section{Section A Mark Allocation}

\begin{longtable}{|c|p{13cm}|}
\caption{Section A Exam Rubric} \label{tab:sectiona}\\
\hline
\textbf{Mark} & \textbf{Expected qualities} \\
\hline
\textbf{9--10} & 
\begin{itemize}
  \item Demonstrates a close and perceptive reading of the text, considering complexities of its ideas, concerns and values explored through the structure of the text and its language
  \item Demonstrates a clear understanding of the implications of the topic, using an appropriate strategy for dealing with it, and exploring its complexity from the basis of the text
  \item Develops a cogent, controlled and well-substantiated discussion using precise and expressive language
\end{itemize} \\
\hline
\textbf{8} & 
\begin{itemize}
  \item Demonstrates a thoughtful reading of the text, exploring its ideas, concerns and values through a discussion of the structure of the text and its language 
  \item Demonstrates an understanding of the implications of the topic, exploring it from the basis of the text
  \item Develops a detailed, substantiated and coherent discussion using language fluently and confidently
\end{itemize} \\
\hline
\textbf{7} & 
\begin{itemize}
  \item Demonstrates detailed knowledge of the text throughout, including some of its ideas, concerns and values by appropriately acknowledging aspects of the structure of the text and its language 
  \item Understands the topic, developing a sustained and well-supported response
  \item Develops an organised piece of writing using language accurately and appropriately
\end{itemize} \\
\hline
\textbf{6} & 
\begin{itemize}
  \item Demonstrates a clear knowledge of the text, including some elements of its ideas, concerns and values and some awareness of structure and language
  \item Develops a response to the topic, supported by appropriate evidence from the text
  \item Presents a generally organised piece of writing using language that is mostly accurate and appropriate
\end{itemize} \\
\hline
\textbf{5} & 
\begin{itemize}
  \item Demonstrates an adequate knowledge of the text and makes some reference to its ideas, concerns and values
  \item Presents a response that shows an understanding of the topic referring to evidence from the text
  \item Presents a piece of writing that communicates adequately, displaying some elements of organisation
\end{itemize} \\
\hline
\textbf{4} & 
\begin{itemize}
  \item Demonstrates basic knowledge of the text with limited reference to ideas, concerns and values
  \item Presents a response that shows some understanding of the topic 
  \item Shows adequate expression and language control
\end{itemize} \\
\hline
\textbf{3} & 
\begin{itemize}
  \item Shows a familiarity with the text
  \item Shows a limited awareness of the topic 
  \item Shows basic expression and language control
\end{itemize} \\
\hline
\textbf{1--2} & 
\begin{itemize}
  \item Shows limited familiarity with the text
  \item Shows a very limited awareness of the topic
  \item Uses language that is not always clear
\end{itemize} \\
\hline
\textbf{0} & 
\begin{itemize}
  \item Shows no knowledge of the text and/or no attempt to engage with the topic and/or only minimal control of language
\end{itemize} \\
\hline
\end{longtable}

\newpage
\section{Section B Mark Allocation}

\begin{longtable}{|c|p{13cm}|}
\caption{Section B Exam Rubric} \label{tab:sectionb}\\
\hline
\textbf{Mark(s)} & \textbf{Expected Qualities} \\
\hline
\textbf{9--10} & 
\begin{itemize}
  \item Demonstrates insightful consideration of the ideas raised by the title and stimulus in connection with a Framework of Ideas
  \item Creates a cohesive text with an explicit purpose(s) and an appropriate voice
  \item Demonstrates sophisticated control of language and text structure(s), that makes rich use of vocabulary and language features
\end{itemize} \\
\hline
\textbf{8} & 
\begin{itemize}
  \item Demonstrates an astute exploration of the ideas raised by the title and stimulus in connection with a Framework of Ideas
  \item Creates a coherent text with an explicit purpose(s) and an appropriate voice
  \item Demonstrates confident control of language and text structure(s), that makes thoughtful use of vocabulary and language features
\end{itemize} \\
\hline
\textbf{7} & 
\begin{itemize}
  \item Demonstrates a detailed connection to the ideas raised by the title and stimulus related to a Framework of Ideas
  \item Creates a coordinated text with a clear purpose(s) and an appropriate voice
  \item Demonstrates sound control of language and text structure(s), that makes clear use of vocabulary and language features
\end{itemize} \\
\hline
\textbf{6} & 
\begin{itemize}
  \item Develops a clear connection to the ideas raised by the title and stimulus with reference to a Framework of Ideas
  \item Creates a connected text with a clear purpose(s) and an appropriate voice
  \item Demonstrates clear control of language and text structure(s), makes effective use of vocabulary and language features
\end{itemize} \\
\hline
\textbf{5} & 
\begin{itemize}
  \item Develops an adequate connection to the ideas raised by the title and stimulus with reference to a Framework of Ideas
  \item Creates an organised text that is linked to a purpose(s) and has an appropriate voice
  \item Demonstrates adequate control of language and text structure(s), that makes some use of vocabulary and language features
\end{itemize} \\
\hline
\textbf{4} & 
\begin{itemize}
  \item Shows a basic connection to the ideas raised by the title and stimulus
  \item Shows some connection to a purpose(s) and some awareness of voice
  \item Shows basic control of language and text structure(s), that makes simple use of vocabulary and language features
\end{itemize} \\
\hline
\textbf{3} & 
\begin{itemize}
  \item Shows limited connection to the ideas raised by the title and stimulus
  \item Shows limited awareness of purpose or appropriate voice
  \item Shows limited control of language and structure
\end{itemize} \\
\hline
\textbf{1--2} & 
\begin{itemize}
  \item Shows little or no apparent connection to ideas raised by the title and stimulus
  \item Shows little awareness of the purpose(s) of their text
  \item Uses language that is not always clear
\end{itemize} \\
\hline
\textbf{0} & 
\begin{itemize}
  \item Shows no knowledge of the task and/or no attempt to engage with the task
\end{itemize} \\
\hline
\end{longtable}

\newpage
\section{Section C Mark Allocation}
\begin{longtable}{|c|p{13cm}|}
\caption{Section C Exam Rubric} \label{tab:sectionc}\\
\hline
\textbf{Mark(s)} & \textbf{Expected Qualities} \\
\hline
\textbf{9--10} & 
\begin{itemize}
  \item Demonstrates a perceptive understanding of contention, the development of argument(s) as constructed in the text and the point of view expressed
  \item Demonstrates sophisticated insight into the ways in which written and spoken language and visuals complement the argument(s) and are used to persuade the intended audience
  \item Uses sophisticated and precise language
\end{itemize} \\
\hline
\textbf{8} & 
\begin{itemize}
  \item Demonstrates a thoughtful understanding of contention, the development of argument(s) as constructed in the text and the point of view expressed
  \item Demonstrates sound insight into the ways in which the written and spoken language and visuals work together to build the argument(s) and persuade the intended audience
  \item Uses language confidently
\end{itemize} \\
\hline
\textbf{7} & 
\begin{itemize}
  \item Demonstrates detailed understanding of contention, argument(s) presented in the text and the point of view expressed
  \item Demonstrates insight into the ways in which the written and spoken language and visuals in the text are used to persuade the intended audience
  \item Uses fluent expression
\end{itemize} \\
\hline
\textbf{6} & 
\begin{itemize}
  \item Demonstrates a clear understanding of contention, the argument(s) presented in the text and the point of view expressed
  \item Demonstrates some awareness of the ways in which the written and spoken language and visuals in the text are used to persuade the intended audience
  \item Uses competent expression
\end{itemize} \\
\hline
\textbf{5} & 
\begin{itemize}
  \item Demonstrates adequate understanding of contention, the argument(s) presented and the point of view expressed
  \item Demonstrates basic awareness of the ways in which the written and spoken language and visuals in the text are used to persuade the intended audience
  \item Uses language adequately
\end{itemize} \\
\hline
\textbf{4} & 
\begin{itemize}
  \item Demonstrates a basic understanding of the contention, the argument(s) and point of view expressed
  \item Describes the ways in which written and spoken language and visuals are used to persuade an audience
  \item Uses basic language
\end{itemize} \\
\hline
\textbf{3} & 
\begin{itemize}
  \item Shows a limited knowledge of the argument(s) or point of view
  \item Shows a limited knowledge of the ways in which written and spoken language and visuals are used
  \item Attempts basic language
\end{itemize} \\
\hline
\textbf{1--2} & 
\begin{itemize}
  \item Shows little understanding of the material
  \item Shows only minimal knowledge of the task
  \item Uses language that is not always clear
\end{itemize} \\
\hline
\textbf{0} & 
\begin{itemize}
  \item Shows no understanding of the requirements of the task
\end{itemize} \\
\hline
\end{longtable}


\clearpage
\thispagestyle{empty}
\hfill
\clearpage
\cleardoublepage

\section{How the Marker Reads Your Essay}
\addcontentsline{toc}{section}{How the Marker Reads Your Essay}

Examiners do not read your essay the way a friendly teacher does -- they read it in passes, testing it against the criteria. Knowing the passes lets you write for them.\\

\subsection{The Six Passes}
First, the scope audit: the topic is broken into its explicit and implicit parts, including modifying words (``unrelenting cruelty'', not just ``cruelty'') and both halves of a double-barrelled prompt. A dimension you never address keeps you out of the 8--10 range. Second, the summary filter: every sentence is silently sorted into retelling, explaining, or analysing how the author constructs meaning -- and if retelling dominates, the essay cannot reach the upper band no matter how fluent it is. Third, the topic-sentence thread: conceptual topic sentences that build a progressive argument separate 9--10 essays from 7--8 essays that give three parallel examples of one idea. Fourth, the authorial hand: markers look for sentences whose subject is the author (``Harrison constructs\ldots'', ``Wilder juxtaposes\ldots'') -- writing about characters as though they are real people is a mid-band signature. Fifth, evidence gravity: quotes that stand alone as their own sentences read mid-band; short quotes woven into your own grammar, drawn from across the whole text, read upper-band. Sixth, the gap analysis: the difference between your mark and the next band is always something specific and nameable -- which is why good feedback tells you what was missing, not just what was there.\\

\subsection{What the Bands Feel Like}
A 5 is relevant but basic: simple structure, simple vocabulary, evidence that supports plot points. A 7 is clear and detailed, uses evidence effectively, but stays at the level of what the text shows rather than how and why it was built that way. A 9 or 10 is sophisticated and perceptive: it reframes the topic conceptually, treats the text as a construction, uses precise metalanguage correctly, and integrates evidence so smoothly the quotation marks could be removed without breaking the sentence.\\

\subsection{Section-Specific Traps}
In Section A, a contention that simply agrees or disagrees without a bigger conceptual claim caps the essay in the middle band, and ignoring the text's social and historical values keeps the top band out of reach. In Section B, a competent story that never engages the Framework's ideas sits mid-band, as do clich\'ed premises and unearned endings -- while adapted echoes of your mentor texts signal studied craft. In Section C, ``the readers'' is a mid-band audience; the upper band names who the writer targets and what they are being positioned to think, feel and do. Naming a technique without explaining its function for that specific audience is label-picking, and it caps the analysis no matter how many techniques you spot. Track the argument's arc and tonal shifts; treat any visual as part of the argument, not a separate exhibit.\\

\part{Exam Checklists}
\clearpage
\thispagestyle{empty}
\hfill
\clearpage
\section*{General Checklist for Section A Response}

\subsection*{Textual Understanding and Interpretation}
\begin{itemize}
  \item Have I demonstrated a clear and accurate understanding of the text?
  \item Have I explored the ideas, concerns, and values present in the text?
  \item Have I considered the structure and language used in the text to convey meaning?
  \item Have I shown awareness of the complexities or subtleties in the text’s ideas?
\end{itemize}

\subsection*{Response to the Topic}
\begin{itemize}
  \item Have I clearly addressed the essay topic throughout my response?
  \item Have I interpreted the implications of the topic (not just the literal meaning)?
  \item Have I developed a relevant and thoughtful argument or discussion?
  \item Have I used an appropriate strategy to explore the topic (e.g., thematic, character-based, structural)?
\end{itemize}

\subsection*{Use of Evidence}
\begin{itemize}
  \item Have I supported my ideas with specific and relevant evidence from the text?
  \item Have I explained how my evidence supports my argument?
  \item Have I used quotations effectively and accurately?
\end{itemize}

\subsection*{Organisation and Coherence}
\begin{itemize}
  \item Is my essay well-structured, with a clear introduction, body, and conclusion?
  \item Are my ideas logically sequenced and connected?
  \item Have I written clear topic sentences and developed each paragraph?
\end{itemize}

\subsection*{Language Use}
\begin{itemize}
  \item Is my language accurate, appropriate, and mostly free of errors?
  \item Have I used precise and expressive vocabulary to convey meaning?
  \item Is my writing fluent and confident in tone?
\end{itemize}


\section*{General Checklist for Section B Response}

\subsection*{Engagement with Ideas and Stimulus}
\begin{itemize}
  \item Have I considered the ideas raised by the title and stimulus in a thoughtful and meaningful way?
  \item Have I connected these ideas to a clear and relevant Framework of Ideas?
  \item Have I explored the implications or complexities of the ideas, rather than treating them superficially?
\end{itemize}

\subsection*{Purpose and Voice}
\begin{itemize}
  \item Have I clearly established the purpose(s) of my text?
  \item Have I maintained an appropriate and consistent voice throughout the piece?
  \item Is the form of the piece suited to the intended audience and purpose (e.g., speech, monologue, letter)?
\end{itemize}

\subsection*{Control of Language}
\begin{itemize}
  \item Have I used language that is clear, expressive, and appropriate for the form?
  \item Have I used a rich and precise vocabulary to strengthen my writing?
  \item Have I varied sentence structure and language features to create impact and cohesion?
\end{itemize}

\subsection*{Control of Text Structure}
\begin{itemize}
  \item Is my writing well-organised and easy to follow?
  \item Have I structured the text purposefully, with a clear beginning, middle, and end?
  \item Have I used paragraphs and transitions effectively to guide the reader?
\end{itemize}

\newpage
\section*{General Checklist for Section C Response}

\subsection*{Understanding of Argument and Contention}
\begin{itemize}
  \item Have I accurately identified the contention of the text?
  \item Have I understood how the argument(s) are structured and developed across the text?
  \item Have I recognised the overall point of view being expressed by the author/speaker?
\end{itemize}

\subsection*{Analysis of Persuasive Language and Visuals}
\begin{itemize}
  \item Have I analysed how language features (written or spoken) are used to persuade the intended audience?
  \item Have I explained how visuals (if any) support or enhance the written or spoken argument?
  \item Have I discussed how these elements work together to construct the overall persuasive strategy?
\end{itemize}

\subsection*{Expression and Language Use}
\begin{itemize}
  \item Have I used language that is clear, precise, and appropriate to an analytical task?
  \item Have I written fluently and used analytical vocabulary confidently?
  \item Have I maintained a consistent and formal tone suitable for academic analysis?
\end{itemize}


\clearpage
\thispagestyle{empty}
\hfill
\clearpage


\part{Key Takeaways from the 2024 Assessment Report}
\clearpage
\thispagestyle{empty}
\hfill
\clearpage

\clearpage

The 2024 VCE English exam, the first under the revised study design, demonstrated impressive student engagement and strong outcomes. According to the VCAA report, students were generally well prepared and confident. However, the report also identified key areas for improvement. This guide breaks down practical strategies for success in each section of the exam, grounded in the report’s findings.

\section{General Exam Advice}
\subsection{What Worked}
\begin{itemize}
    \item Students clearly understood the structure and expectations of the exam.
\item Many responses demonstrated well-developed writing and analytical skills.
\end{itemize}

\subsection{What Needs Improvement}
\begin{itemize}
    \item Some students forgot to clearly label their chosen text (Section A) or Framework of Ideas (Section B).
    \item Some responses misinterpreted the prompt or overlooked key words in the question.
\end{itemize}

\subsection{Practical Tips}
\begin{itemize}
    \item Label your work clearly. Always begin with:\\
    Section A: 'Text: Macbeth | Prompt: ‘Ambition leads to downfall. Do you agree?’”\\
    Section B: 'Title: Resist and Persist | Framework: Protest”

\item Build your vocabulary. Familiarise yourself with words like inescapable, perceptions, central concern, implication, etc.

\item Select the best-fit prompt. Choose the one that allows you to demonstrate your strongest understanding, not just the one that sounds easiest.
\end{itemize}

\newpage
\section{Section A: Analytical Response to a Text}
Your task: Write a structured essay in response to one of two prompts for your studied text.

\subsection{What Worked}
\begin{itemize}
    \item High-scoring responses interpreted the topic with nuance and clarity.

\item Students demonstrated a strong understanding of authorial intent.

\item Responses were well structured and supported with relevant textual evidence.
\end{itemize}

\subsection{What Needs Improvement}
\begin{itemize}
\item Some students only responded to part of the question (e.g. discussing 'youth” but not 'infatuation with youth”).

\item Others summarised the plot or characters instead of analysing the text’s ideas and values.
\end{itemize}

\subsection{Practical Tips}

\begin{table}[ht]
\centering
\begin{tabular}{|l|l|}
\hline
\textbf{Effective Approach}                      & \textbf{Ineffective Approach}                                  \\ \hline
Interpreting all keywords in the topic           & Ignoring modifiers (e.g. ``unrelenting cruelty'' vs ``cruelty'') \\ \hline
Using a clear ``what/how/why'' structure         & Listing unrelated examples                                      \\ \hline
Connecting ideas across paragraphs             & Treating each paragraph as a standalone point                   \\ \hline
Focusing on how the author constructs meaning    & Just retelling the storyline                                    \\ \hline
\end{tabular}
\caption{Section A - Effective vs Ineffective Approaches}
\end{table}
'
\textbf{Tip:} Write practice introductions where you break down the prompt. Ask:
\begin{itemize}
    \item What is the main idea or concept?

\item Are there two ideas in relationship?

\item What are the implications or assumptions?
\end{itemize}

\newpage
\section{Section B - Creating a Text (Framework of Ideas)}
Your task: Create an original piece of writing that responds to a title and stimulus, grounded in one of the four Frameworks of Ideas.

\subsection{What Worked}
\begin{itemize}
    \item High-scoring pieces treated the title and stimulus as central, not secondary.

\item Strong voice, clear purpose, and tailored language made the writing stand out.

\item Responses showed thoughtful exploration of ideas relevant to the Framework.
\end{itemize}

\subsection{What Needs Improvement}
\begin{itemize}
    \item Some students ignored or misunderstood the Framework (e.g. writing about sports under the Protest framework).

\item Others forced pre-prepared stories to fit the title.

\item Lack of planning resulted in unstructured or unfocused pieces.
\end{itemize}

\subsection{Practical Strategies}
\subsubsection{Planning Method}
\begin{itemize}
    \item Framework: Personal Journeys

\item Title: Finding My Way

\item Stimulus: Butterfly emerging from a cocoon

\item Purpose: Reflective

\item Audience: Teenagers navigating cultural identity

\item Text Type: Memoir
\end{itemize}

\subsubsection{Tips for Strong Writing}
\begin{itemize}
    \item Use recurring imagery or motifs to strengthen cohesion.

\item Match tone to purpose: persuasive, reflective, emotive, or narrative.

\item End with intention: resolution, reflection, or call to action.
\end{itemize}

\newpage
\section{Section C - Analysing Argument}
Your task: Analyse how a written and visual text positions its audience.

\subsection{What Worked}
\begin{itemize}
    \item Identifying the author’s contention and target audience.

\item Analysing tone, persuasive language, structure, and visuals.

\item Exploring how language is used to shape reader response.
\end{itemize}

\subsection{What Needs Improvement}
\begin{itemize}
    \item Avoid retelling the article.

\item Go beyond spotting techniques; explain their effect and purpose.
\end{itemize}

\subsection{Practical Strategies}
\begin{table}[ht]
\centering
\begin{tabular}{|l|l|}
\hline
\textbf{Strategy}                    & \textbf{Example}                                                                                      \\ \hline
Identify contention and purpose       & 'The writer argues that vaping is a growing public health crisis and seeks to alarm parents.”          \\ \hline
Link language to effect               & 'By using the phrase ‘a ticking time bomb’, the writer creates a sense of urgency and panic.”          \\ \hline
Analyse visuals                       & 'The image of a distressed child reinforces the emotional appeal used in the article’s opening.”       \\ \hline
Use precise metalanguage              & Appeal to fear, inclusive language, repetition, expert opinion, etc.                                   \\ \hline
\end{tabular}
\caption{Section B: Effective Strategies and Examples}
\end{table}


\newpage
\section{Final Exam Checklist}
\begin{itemize}
    \item Have I practised responding to a range of prompts and titles?

\item Do I clearly understand my Frameworks of Ideas?

\item Can I write for a clear audience and purpose?

\item Have I reviewed unfamiliar vocabulary and metalanguage?

\item Can I plan and write all three sections under timed conditions?
\end{itemize}

\section{Final Advice}
The 2026 exam will reward students who think deeply, write purposefully, and tailor their writing to the demands of the task. Whether you’re analysing a text, creating your own piece, or evaluating argument and language, success lies in showing insight, clarity, and control.\\

Strong writing does not come from memorisation—it comes from understanding. Prepare thoughtfully, practise deliberately, and revise reflectively.
\cleardoublepage

\clearpage
\thispagestyle{empty}
\hfill
\clearpage

\cleardoublepage
\part{Effectively Studying For Exams}
\clearpage
\thispagestyle{empty}
\hfill
\clearpage

\clearpage

\section{Assess your current study habits}
\subsection{Identifying strengths and weaknesses}
Take a few minutes to jot down what's working well in your current study routine.
What areas could use some improvement?
Consider factors like time management, focus, and understanding of different subjects.

\subsection{Setting realistic goals}
Based on your assessment, you could start by setting 2-3 specific, achievable goals for improving your study habits.
Make these goals SMART: Specific, Measurable, Achievable, Relevant, and Time-bound.
Write them down and place them somewhere visible as a daily reminder.\\

\section{Create a personalised study schedule}
Now that you know where you stand, it's time to map out a path to exam success. Think of your study schedule as a roadmap to help you navigate your way through VCE studies.

\subsection{Time management tips}
List all your subjects and the topics you need to cover for each.
Assess how much time you have until your exams.
Consider prioritising subjects based on difficulty and your current performance.
Using a study schedule template
Remember: Be realistic with your schedule. If you're not a morning person, a sunrise study session may not suit you. Try to work with your natural rhythms, and not against them.

\subsection{Find a study schedule template that works for you – digital or physical, whichever works best.}
Block out fixed commitments like school, extracurriculars, and family time.
Allocate study blocks for each subject, ensuring a mix of subjects each day.
Don't forget to schedule breaks and 'me time' to avoid feelings of burnout.\\

\section{Organise your study materials}
A cluttered space can contribute to a cluttered mind. Here are some strategies to effectively organise your study materials.

\subsection{Gathering resources}
Collect all your textbooks, notes, and past papers for each subject.
Organise digital resources into clearly labelled folders on your computer.
Consider creating a 'study toolkit' with essentials like highlighters, sticky notes, and index cards.
\subsection{Creating a study-friendly environment}
Find a quiet, well-lit space where you can focus without distractions.
Ensure you have a comfortable chair and desk at the right height.
Keep your study area clean and organised – a tidy space can do wonders for your concentration.\\


\section{Implement active recall and spaced repetition}
Below are some tried and true study techniques you might like to explore.

\subsection{The power of self-testing}
After studying a topic, close your books and try to recall as much as you can.
Create flashcards for key concepts and test yourself regularly.
Explain concepts out loud as if you're teaching someone else.
\subsection{Spacing out your study sessions}
Avoid cramming by spreading out your study sessions over time.
Review material at increasing intervals: e.g. 1 day later, 3 days later, 1 week later, etc.\\


\section{Utilise practice exams and past papers}
Another way to effectively prepare for exams is by practising past exams. Here are some techniques you might find helpful.

\subsection{Simulating exam conditions}
Find past papers or practice exams for your subjects.
Set up a quiet space and time yourself as you would in a real exam.
Avoid using your notes or textbooks during these practice sessions.
\subsection{Analysing your performance}
Mark your practice exams using the relevant marking schemes, if available.
Identify areas where you lost marks and focus on improving these in your study sessions.
Keep a log of common mistakes to avoid repeating them in the real exam.\\

\section{Develop a healthy study-life balance}
It's important to look after yourself and practice self-care during this time. Find what works best for you, or try our ideas below.

\subsection{Incorporating breaks and relaxation}
Use techniques like the Pomodoro method: 25 minutes of focused study followed by a 5-minute break.
Schedule longer breaks for activities you enjoy – maybe a quick game of basketball or a chat with friends.
Practice mindfulness or meditation to help manage stress and improve focus.
\subsection{Maintaining physical and mental well-being}
Aim for 7-9 hours of sleep per night – your brain consolidates memories and learning during sleep.
Incorporate regular exercise to improve your mood and reduce stress.
Eat a balanced diet with brain-boosting foods like blueberries, nuts, and fatty fish.\\


\section{Collaborate with study groups}
Two heads are better than one, and a study group can be a powerhouse of knowledge, inspiration and motivation.

\subsection{The benefits of peer learning}
Explaining concepts to others can help solidify your own understanding.
Gain new perspectives and insights from your peers.
Stay motivated and accountable with group study sessions.
\subsection{Organising effective study sessions}
Set clear objectives for each study session.
Use tools like shared Google Docs or virtual whiteboards for collaborative note-taking.\\

\section{Seek help when needed}
Remember, asking for help isn't a sign of weakness.

\subsection{Utilising teacher support}
Don't hesitate to approach your teachers with specific questions.
Attend any extra help sessions offered by your school or community.
Ask for feedback on your work.

\subsection{Exploring additional resources}
Explore online tutorials or video explanations for challenging topics.
Consider getting a tutor for subjects where you're really struggling.
Join online forums or study groups for your specific subjects.\\

\section{Fine-tune your exam day strategy}
When your exams are approaching – make sure you're feeling prepared.

\subsection{Pre-exam preparation}
Get a good night's sleep before the exam.
Eat a nutritious breakfast to fuel your brain.
Arrive at the exam venue early to avoid last-minute stress.

\subsection{During the exam tactics}
Read through the entire paper before starting to answer questions.
Manage your time wisely – allocate minutes to each section based on mark values.
If you get stuck on a question, move on and come back to it later.



\cleardoublepage

\clearpage
\thispagestyle{empty}
\hfill
\clearpage

\cleardoublepage

\clearpage
\thispagestyle{empty}
\hfill
\clearpage

\end{document}

